% Framework S3 finite-response exact comparison integration v3.035: source-scoped companion certificates and ownership firewall.
% Framework v3.027: exact 600-cell/Hurwitz return channels, Cartan rank, native clutch, Clifford constraints and conditional physical gates.
% Framework v3.025: exact character-ring and isolated heat-face reference computations, physical action conditional.
% Paper-series v2.382: audit-identified standalone repair integration; mathematical scopes and open physical owners preserved.
% Paper-series v2.381: complete algebraic exceptional torsion atlas and all-cocycle twist-stability; original claims retained.
% Paper-series v2.380: unified bibliography, exact exceptional double/length/gauge mathematics in I-B; previous theorem labels retained.
% Paper-series v2.379: full equivariant Dirac/Hodge/eta/torsion atlas in Paper I-B; paper-root references synchronized.
% Paper-series v2.378: 2026 exceptional literature consolidation and reference synchronization. The mathematical additions are located in Paper I-B.
% Paper-series v2.377: reference-only rebind to current Paper I-B v1.1; original mathematics retained except new Paper I-B comparison chapter.
% Framework v2.370: derived exceptional-root/Hurwitz comparison in I-A/I-C; other papers reference-only synchronized.
% Framework v2.368: preserve prior II-B convention/unital fixes and verified II-D corrections; current references synchronized.
% Framework v2.365: scoped Hurwitz/600-cell hadron-source integration; physical APBI/decay frontier remains open.
% Framework v2.358: current 30-root source reference synchronization; selectively updated III-E, III-F, VI-B, VI-C.
% Framework v2.357: normal-jet/Clifford/hadron-shell observable proof additions at scoped mathematical, conditional, or QCD-source levels.
% Framework v2.356: bibliography/source-registry synchronization; mathematical body unchanged.
% Framework v2.350: repaired Paper VII mathematics; exact finite identities separated from analytic/physical certificates.
% Framework v2.345: F40 native three-readout and source-visibility integration.
% Framework v2.344: Twisted Hurwitz C3/Petersen bridge, native baryon source and same-action curvature firewall.
% Framework v2.323: Petersen finite strong-source coherence, exact multi-core mirror spectra, QCD action-incidence and continuum firewalls. Existing source theorems preserved.
% Framework v2.275: exact positive dense-twin word cubature, finite repeatability and eighth-moment classification, nested sampling frames, joint action and POVM-jet exports; preserves v2.274.
% Framework v2.274: two-scalar positive convolution certificates, first-leaked exceptional coset detection, branch-symmetric twin mixing, and finite endpoint moments; all v2.273 branches preserved.
% Framework v2.273: restored missing Problem 3 theorems from v2.268; exact exceptional frames, twin geometry, convolution repeatability and projector/occurrence firewalls; merged on v2.272 without overwriting independent branches.
% Framework v2.272: merged v2.270 covariant native-readout mathematics with v2.271 Problem 4 finite-action Hopf hypercharge/holonomy corrections; no universal bulk-gauge uniqueness claim.
% Independent branches: v2.270 readout/guard exports and v2.270 K16 action-consumed charge closure; all prior scoped assumptions retained.
% Framework v2.267: exact Paper II-A converse/observable certificates, III-F action tests, V-B right orbit and V-C native gauge scope; prior v2.266 preserved.
% Framework v2.266: exact Paper II-C outcome-count/readout converse results and scoped downstream imports; 30-root references synchronized.
% Framework v2.265: II-C coordinate-stage constitution and exact first-symbol host inheritance; prior mathematical domains retained.
% Framework v2.264: II-B reverse and action-bridge integration; preserves prior mathematical domains; active 30-root references synchronized.
% Framework v2.263: Paper I-D adds bridge-controlled root-carry parity and source-relative old-cap cancellation; other root mathematical bodies retained; active 30-root references synchronized.
% Framework v2.262: merges Paper II-F three-gap closure from v2.260 into v2.261, preserving newer II-B/II-D repairs; 30-root references synchronized.
% Framework v2.261: merged Paper II-B GP2-level/geometric repairs; preserves v2.260 II-D repairs; active references synchronized.
% Framework v2.260: Paper II-D signed RS, finite-symbol typing, and conditional curvature/Jordan repairs; references synchronized.
% Framework v2.259: conditional native SU(2) atlas and explicit later discharge; v2.258 mathematical bodies retained elsewhere.
% Framework v2.258: preserves parallel II-A/dynamics update and adds complete readout/frequency closures.
% Framework v2.257: exact minimal complete position readout and arbitrary-preparation frequency law.
% Framework v2.256: exact native inertia and sharp bounded ordered-source closure; preserves concurrent v2.255 work.
% Framework v2.255: preserves parallel detector integration and adds exact finite readout closures.
% Framework v2.254: shared bibliography/reference sync; substantive Pasterski detector-Ward additions only IV-B/IV-C/VII-C.
% Framework v2.253: typed raw coordinate defect in I-A and synchronized VI-A import; active references synchronized.
% Framework v2.252: optional shell-frame source-return calculation in III-G; active references synchronized.
% Framework v2.251: forward common-action execution, joint source jets and scoped coupled control.
% Framework v2.251: exact full-heat variations, fermion source jets and explicit physical completion frontier.
% Framework v2.251: certified interacting source jets, pole stability and scoped completion proofs.
% Framework v2.251: interacting finite-source evaluation and coupled local vacuum classification.
% Framework v2.251: exact finite source execution and native causal common-parent bridge.
% Framework v2.251: constructive parent and scoped source-incidence audit; see release report.
% Framework v2.251: common-action two-jet, corrected source/response closure and synchronized references.
% Framework v2.244: scoped rooted source-order Schur diagnostics and synchronized paper references.
% Framework v2.242: signed rooted phase lift, complete transverse return and exact family cubic.
% Framework v2.241: proved rooted conormal geometry, conditional cap realization and scoped action tests.
% Framework v2.239: coherent 30-paper reference rebinding; selective new proofs on established owners.
\documentclass[12pt,a4paper]{article}
\PassOptionsToPackage{hidelinks}{hyperref}
\usepackage{microtype}
\usepackage{etoolbox}
\usepackage{caption}
\usepackage{makecell}
% Macros-v2_184: active Paper III-E v1.200/III-F v1.267 binding for framework v3.031; mathematical macros unchanged.
% Macros-v2_183: updated shared cross-reference bindings for Paper1B/III-E/III-F v3.030; mathematical macros unchanged.
% Macros-v2_182: active Paper II-F/V-A cross-references for framework v3.029; definitions unchanged.
% Paper-series v2.381: new I-B torsion/modular results, bibliography and current source roots.
% Paper-series v2.380: new exceptional bibliography and synchronized paper references.
% Paper-series v2.379 current 31-paper macro/aux registry.
% Paper-series v2.378: current 31-paper AUX/bibliography registry updated for Paper I-B v1.2 literature additions.
% Framework v2.371: current-source reference registry refreshed; old theorem labels retained.
% Framework v2.370: current 30-paper source/AUX root rebinding; native mathematical macros unchanged.
% Framework v2.368: synchronized merged 30-root registry; mathematical notation unchanged.
% Framework v2.358: active 30-paper source registry synchronized; existing mathematical macros unchanged.
% Framework v2.357: selective new source-owned certificates; active reference roots updated for five revised papers.
% Framework v2.356: active 30-root registry synchronized to the 30 papers in this release.
% Framework v2.275: exact positive dense-twin word cubature, finite repeatability and eighth-moment classification, nested sampling frames, joint action and POVM-jet exports; preserves v2.274.
% Framework v2.274: filename and root-source synchronization.
% Framework v2.273: restored missing Problem 3 theorems from v2.268; exact exceptional frames, twin geometry, convolution repeatability and projector/occurrence firewalls; merged on v2.272 without overwriting independent branches.
% Framework v2.272: merged v2.270 covariant native-readout mathematics with v2.271 Problem 4 finite-action Hopf hypercharge/holonomy corrections; no universal bulk-gauge uniqueness claim.
% Independent branches: v2.270 readout/guard exports and v2.270 K16 action-consumed charge closure; all prior scoped assumptions retained.
% Framework v2.269: synchronized 30 active AUX roots; mathematics and bibliography preserved from Macros-v2_161.
% Framework v2.266: Paper II-C converse-derived outcome resources; II-E, III-C, IV-A/B, V-A and VI-B downstream mathematical closures; 30-root registry synchronized.
% Framework v2.265: II-C coordinate-stage constitution and exact first-symbol host inheritance; prior mathematical domains retained.
% Framework v2.264: paper II-B reverse/audit and IV-A observable-minimal exchange; synchronized external references to 30 active roots.
% Framework v2.263: Paper I-C adds bridge-controlled root-carry parity and source-relative old-cap cancellation; other root mathematical bodies retained; active 30-root references synchronized.
% Framework v2.262: integrates Paper II-F radius, Casimir and readout results; active 30-root references synchronized.
% Framework v2.261: merged Paper II-B GP2-level/geometric repairs; preserves v2.260 II-D repairs; active references synchronized.
% Framework v2.260: Paper II-D signed RS, finite-symbol typing, and conditional curvature/Jordan repairs; references synchronized.
% Framework v2.259: conditional native SU(2) atlas and explicit later discharge; v2.258 mathematical bodies retained elsewhere.
% Framework v2.258: preserves parallel II-A/dynamics update and adds complete readout/frequency closures.
% Framework v2.257: exact minimal complete position readout and arbitrary-preparation frequency law.
% Framework v2.256: exact native inertia and sharp bounded ordered-source closure; preserves concurrent v2.255 work.
% Framework v2.255: preserves parallel detector integration and adds exact finite readout closures.
% Framework v2.254: Pasterski bibliography and synchronized 30-paper external-reference registry.
% Framework v2.253: typed raw coordinate defect in I-A and synchronized VI-A import; active references synchronized.
% Framework v2.252: synchronized source registry; optional III-G shell-frame calculation.
% Framework v2.251: forward common-action execution, joint source jets and scoped coupled control.
% Framework v2.251: exact full-heat variations, fermion source jets and explicit physical completion frontier.
% Framework v2.251: certified interacting source jets, pole stability and scoped completion proofs.
% Framework v2.251: synchronized active source registry; mathematical notation unchanged.
% Framework v2.251: synchronized active source registry; mathematical notation unchanged.
% Framework v2.251: synchronized active source registry; mathematical notation unchanged.
% Framework v2.251: synchronized source registry; no mathematical notation changes.
% Framework v2.244: synchronized 30-paper AUX registry and rooted phase-source action-order tests.
% Framework v2.242: signed rooted phase lift, complete transverse return and exact family cubic.
% Framework v2.241: proved rooted conormal geometry, conditional cap realization and scoped action tests.
% Framework v2.239: all 30 active source roots synchronized; no mathematical macro changes.
% Framework v2.238: Paper IV A/B/C split and 30-root reference registry.
% Framework v2.237: Paper VII-D technical split and 28-root reference registry.
% Framework v2.236: Paper VI-C technical-volume reference added; all current root bindings advanced.
% Framework v2.235: 26-root AUX and bibliography bindings synchronized; mathematical macros unchanged.
% BEGIN INLINE SOURCE: Paper00-TeXBuildCompatibility-v1_0.tex
% Rendering-only repairs for inherited long appendices and bitmap text fonts.
% Integer counters through 26 retain the original labels; later values use AA,
% AB, ... . This file changes no action coefficient, state or theorem premise.
\RequirePackage{alphalph}
\makeatletter
\let\BHPrimitiveAlph\@Alph
\newalphalph{\BHExtendedAlph}{\BHPrimitiveAlph}{26}
\let\@Alph\BHExtendedAlph
\AtBeginDocument{%
  \@ifpackageloaded{microtype}{\microtypesetup{expansion=false}}{}%
}
\makeatother

% END INLINE SOURCE: Paper00-TeXBuildCompatibility-v1_0.tex

% Framework v2.164: current reference/bibliography profile; mathematical definitions retained.
% Framework v2.162: finite metric, functional Hilbert and scoped Ward integration.
% Framework v2.161: inherits Macros-v2_57.tex exactly except the current external-reference binding.
% Macros v2.57 / framework v2.156: current reference registry only; mathematical definitions unchanged.
% Framework v2.143 / Macros v2.55: current external-reference bindings; mathematical notation unchanged.
% Framework v2.110 / Macros v2.54: current external-reference registry and bibliography; mathematical notation inherited.
% Framework v2.67 / Macros-v2_53: adds Born-action-cell and joint-EH readout symbols; absolute unit realization is metrological; MATCH_{J/scal}^{EH} remains the sole G24 normalization bridge.
% Framework v2.35 / Macros-v2_52: grouped paper numbering, active AUX registry, and Born-v6_20 bibliography binding synchronized; notation unchanged.
% Framework v2.26: active-source AUX and bibliography bindings synchronized; notation unchanged.
% Framework v2.09: reviewed early-reciprocity integration; scoped physical imports.
% Framework v2.08: Born--Euler mathematics and conditional complete-action potential criterion.
% Framework v2.07: active AUX registry includes Paper I-B v1.2; shared bibliography Born-v6_18.
% Macros-v2_46 / framework v2.04: mathematical atlas and downstream realization correction.
% Macros-v2_46 / framework v2.03: scoped structure integration and current dependency bindings.
% Macros-v2_46 / framework v2.01: shared source bindings; prior revision comments are historical.
% v2.42 / framework v1.99: Synchronized RG/AS owner and bibliography bindings.
% v2.41 / framework v1.99: active manuscript AUX bindings and geometric-closure bibliography synchronized.
% Release-v1.80 centralized source synchronization for scalar-circle/comparison-band integration.
% Release v1.79: Hurwitz hierarchy and kinetic-normalized action response integration; physical mass matching remains conditional.
% Release v1.78: approved M8 closure-batch integration; physical and analytic domains retained.
% Current v1.78 registry: all active TeX roots are synchronized below.
% Earlier release comments are historical scope records, not current status.
% Release-v1.77 synchronization: native same-owner Noether closure, generator--flow/history typing, analytic-certification firewalls, scalaron/K16, LF/QCD, quantum-Ward, and RAR boundary-owner transport are registered; shared mathematical notation is unchanged.
% Release-v1.76 synchronization: registry points to the v1.76 active paper roots.
% Release-v1.76 synchronization: registers the local A+Phi+g Noether/constraint closure and dynamics-frontier propagation; shared mathematical notation is unchanged.
% Release-v1.75 synchronization: registry points to the v1.75 active paper roots.
% Release-v1.75 synchronization: integrates the joint-Dirac Einstein bulk real form and the local odd-scalar heat-coefficient action on explicit analytic first-jet/interior domains; v1.74 Yang--Mills closure is retained. Fermionic reality, pure-metric curvature-squared/exact-nonlocal terms, noncanonical boundaries, physical vacuum selection, and actual solution-history tokens remain open.
% Release-v1.74 synchronization: registry points to the v1.74 active paper roots and retains the v1.73 theorem labels.
% Release-v1.74 synchronization: canonical Macros-v2_34.tex; generated local nonlinear Yang--Mills action, constrained connection-field phase space, source-free Gauss momentum map, and Legendre closure are integrated on their proved domains. The remaining non-gauge ACT_cl sectors and actual solution-history selection remain open.
% Release-v1.71 synchronization: closes the former JOINT_x/JOINT_5
% connection-occurrence gap by canonical structure-group extension of the
% nonflat Born connection, faithful associated representation, and the local
% cotangent--Clifford operator constructor.  Actual fluctuation members,
% actions, vacua, solutions, calibrations, and analytic QFT domains remain
% separately typed.  The active-AUX registry is advanced once.
% Release-v1.71 canonical macro and active-AUX registry.
% ======================================================================
% Macros-v2_39.tex -- canonical publication macro file
% Release v1.71 retains the v1.70 null-ray/QCD boundaries and additionally
% removes the duplicate structural connection actuality premise.  It does not
% claim global bundle topology, select a gauge/Higgs member, or promote the
% finite construction to an interacting continuum theory.
% Release v1.68 constructs the finite light-front total-fraction operators and
% parton number/momentum operator measures, proves their simplex and Hausdorff-limit
% theorems, and separates positive cut measures from scheme-renormalized
% light-ray distributions.  It also separates finite real Schur poles, stable
% norm-resolvent/Riesz pole-residue limits, and continued resonance
% pole/residue limits.  The Paper-VII-C scattering package is claim-indexed:
% local QFT and spectral claims import no S-matrix or infrared premise, while
% asymptotic amplitudes, LSZ, S-matrix poles, completeness, and cross sections
% import only their named massive or massless/long-range clauses.
% Release v1.67 factorizes BND48_H into a same-action Hessian/Schur
% descendant and still-open typed boundary-trace, critical-lift, and
% refinement-limit conditions.  It proves the finite-collar Dirichlet-to-
% Neumann law sqrt(y)tanh(L sqrt(y)), its nonuniform critical limit, and the
% finite-return ultra-infrared turnover.  The Paper-I-A Hardy constant mode is
% an exact geometric near-hit, not yet the physical metric boundary mode.
% The current local f(R) Hessian remains massive and therefore still refutes
% a RAR prediction on the minimal action branch.
% Release v1.66 derives the transverse depth-six normalizer bridge from the
% already generated C6 Reynolds image and Paper-I-A binary-dihedral normalizer.
% The bridge is the unique least 2D3-stable scalar-containing realization of
% the faithful quotient 2D3/C6=C2, has trace defect six and hence K=48, but is
% transverse to the nested active carrier.  Exact K=44 and K=52 contraction,
% semigroup, and energy-midpoint identities are registered as structural
% readouts rather than projector events.  The active E48 carrier gap remains
% unchanged.  A same-action boundary-symbol/Calderon intertwiner is isolated
% as the necessary and still unproved owner bridge to the static DtN sort;
% representation equivalence alone neither derives RAR nor creates a K=48
% action junction, pole, occupation, or cosmological transition.
% Release v1.65 derives the native E_48-centered carrier mass gap from the
% GP1--GP3-generated active spectrum, distinguishes its carrier-free interval
% from a particle located at E_48, and withdraws the unsupported promotion to
% a positive continuum Yang--Mills mass gap.  The latter remains a valid
% conditional transfer theorem on an enlarged continuum branch, but it is not
% an active framework prediction or an unfinished native-gap proof.
% Release v1.64 synchronizes Derivation Batches 14--22.  It registers the
% GP1--GP3 heat-profile orbit, same-Hessian electroweak response and pole
% ownership, K16/K48 action-response corrections, perturbation state and
% non-Gaussianity typing, finite-vacuum/gap/QCD limit repairs, factorized
% position-area ownership, the actual static metric response, and the ordered
% Wilson/string-tension limits.  Strict continuum convergence, positive
% physical gap, surviving higher cumulants, and positive tension remain open.
% Release v1.61 closes TIMEWIDTH16 as an independent direct-production datum.
% JUNC_{2|7} continues to own the finite intrinsic microscopic profile/width,
% while the actual same-action background/production history determines its
% physical-time pullback, support, and duration.  The current Paper-I-A--VII-C AUX
% registry is advanced once; shared mathematical notation is unchanged.
% Release v1.60 derives the final-octet primitive-remnant depletion, canonical
% Hopf-current representative, exact Hopf descent to dOmega/(4pi), and
% closed-FLRW a^{-3} neutral-current shape.  It retypes K=48 as the onset
% of the final one-unit geometric transfer and records cDM/RAR as distinct
% physical projections of that common neutral parent while preserving
% occurrence, Hilbert-owner, stress-matching, and static-kernel premises.
% Release v1.59 transports the v1.58 exceptional-shell/K16 closure through
% Papers III-F and VI-B--VII-C.  It replaces Paper VI-B's obsolete mixed-Hessian
% output typing by the third-variation pair vertex and Bogoliubov occupation
% support, derives the post-seed Higgs-pair feedback vertex and its exact
% alpha--delta kernel factors, propagates the unique K16/K48 action-junction
% split, identifies the QCD delta count with the upstream 2O shell atom only
% as a common-scalar numerical compatibility, and narrows THERM16 to dynamical
% approach/beta selection because the finite Gibbs/KMS target is already
% unique on a supplied simple observable algebra once its Hamiltonian and
% inverse temperature are supplied; direct-sum central weights remain typed.
% Release v1.58 registers the exceptional nearest-shell closure, the derived
% transverse K=16 action event, the unique interior action junction, the
% type-correct ordered pair vertex, the thin-layer squeeze law, and the
% radiation-like/matter-like kinematic split.  REG_0, transfer magnitude,
% physical layer width, and thermalization remain explicitly conditional.
% Release v1.57 registers the active-depth-six obstruction, the native
% K=48 carrier desert, the typed 7|9 owner reduction, and the dynamical
% reheating handoff across Papers III-B--VI-B.
% Release v1.56 registers the least-support n_2 inflation dressing of the
% persistent n_7 Laurent response, the analytic global-convexity and positive
% endpoint-measure proofs, the derived inflation core and K=16 coefficient
% freeze, and the no-particle dark-sector split into static RAR plus the
% conditional homogeneous effective-curvature response at reflected K=48.
% There is no n_6 action, scalaron-DM mass, or second Einstein--Hilbert term.
% Release v1.55 registers the protected-octet resolved/residual reflection
% q_sp+Z_W=8, its j<->8-j and K<->64-K scale covariance, and the derived
% rank-one-exit quadratic grade C_1^2/U=C_2 in Papers III-C, III-E, and VIII.
% The reflection is a typed trace/scale theorem, not a time reversal or a
% physical event-occurrence theorem; reheating transfer remains conditional.
% Release v1.54 registers the exact positive-part Laurent indices, the
% active-domain/terminal-endpoint action split, the no-second-EH firewall, and
% the q_DM=6/K_DM=48 scalaron readout typing across Papers III-D/E/F, VI,
% VIII, and IX.
% Release v1.53 registers the common nonnegative Laurent obstacle projection
% for n_1^+ and n_7^+, discharges the three former exit/freeze assumptions,
% and advances the Paper-III-D/E/F, VIII, and IX consumers.  Shared notation
% and bibliography ownership are unchanged.
% Release v1.51 registers the PQTI derivation batch and its updated
% Paper-III-A/C/D/E/F, VIII, and IX consumers.  Shared notation and
% bibliography ownership are unchanged.
% Release v1.50 advances the Paper-III-B AUX target to v1.40 after the
% rigorous audit repair; shared notation and bibliography ownership are unchanged.
% Release v1.49 advances the Paper-II-E/F AUX target to v3.14 after the
% minimal audit repair; shared notation and bibliography ownership are unchanged.
% Release v1.46 synchronizes Paper II-A's minimal frozen constructor-origin
% theorem and its five dependent AUX targets.  Shared notation is unchanged.
% Release v1.45 updates the Paper-II-B AUX target to v2.23 after the
% GP3-rigidity and appendix-placement repair; shared notation is unchanged.
% Release v1.44 advances the centralized Paper-I-A--VII-C AUX registry after the
% Paper-I-A v5.16 protected-completion scope repair.  Every current
% paper imports this one file; no paper-level opt-in or \externaldocument list
% is permitted.  The bibliography remains centralized here through Born-v6_13.
% Shared bibliography, theorem infrastructure, colors, notation alignment,
% and current Paper-I-A--VII-C external-reference registration.
% This file centralizes common imports, bibliography sorting, colors, theorem
% infrastructure, and the notation shared by the complete series.
% v2.9 (release v1.43) made the Paper-I-A--VII-C AUX registry automatic and
% centralized the bibliography and all external-reference ownership.
% v2.8 release-registry refresh (v1.33) registers the rooted generated-parent
% cascade and the POL/REG/free-law/11+4/Born-group closure revisions.  No
% shared mathematical notation or Paper-II-A fundamental principle is changed.
% v2.8 release-registry refresh (v1.32) registers the universal generated-
% continuation revisions in Papers II-A, III-A, III-C, IV, and VIII.  No
% shared mathematical notation or Paper-II-A fundamental principle is changed.
% v2.8 release-registry refresh (v1.30) registers the CRP4/projective-state
% source revisions and their direct registry-only consumers.  No shared
% mathematical notation or Paper-II-A axiom is changed.
% v2.7 release-registry refresh (v1.27) updates only the current external
% Paper-I-A--VII-C basenames after the placement/assumption synchronization; no
% shared mathematical notation or Paper-II-A axiom is changed.
% v2.7 synchronizes the geometric defect-repair revision: Paper-II-A v3.42,
% Paper-III-C v1.26, Paper-III-D v1.27, Paper-III-E v1.32, and Paper-III-F
% v1.34.  It registers the common trace-defect/spectral-filtration theorem,
% its owner-typed response pullback, and the reduced GRED and TERM interfaces.
% No shared mathematical notation is changed.
% v2.6 synchronizes Paper-III-E v1.31's cone-type repair, minimal-response
% classification, and explicit classical action-role boundary, together with
% Paper-III-F v1.32's downstream response-atlas synchronization.  No shared
% mathematical notation is changed.
% v2.5 synchronizes the Paper-III-D v1.26 faithful-trace-duality proof,
% exact ordered-owner joint locus, and typed alternative-presentation
% bookkeeping repair, together with its Paper-III-E and closure-frontier
% references.  No shared mathematical notation is changed.
% v2.4 synchronizes the Paper-III-C v1.25 orbit/process/owner repair, its
% downstream Papers III-B--F, V, VI, VIII, and IX revisions, and Born-v6_13.
% v2.3 synchronizes the Paper-II-A v3.41 witness/support/principal inheritance
% layer, the Paper-III-A v1.24 joint-realizability/readout firewall, all dependent
% host-rate revisions, and the current Paper-I-A--VII-B source registry.
% v2.2 synchronizes the generated active-realization, determinant-depth,
% dimension-weighted hierarchy, and typed cosmological-event interfaces.
% v2.1 synchronizes the exceptional-support interface revisions: Paper I-A
% v5.9, Paper II-A v3.30, Paper II-E/F v3.7, and Paper IV v4.16.
% v2.0 introduces the stage-safe G0 occurrence notation, synchronizes the
% active Paper-I-A--IV source registry (current rigorous Paper-III closure:
% III-A v1.17, III-B v1.22, III-C--E v1.17, III-F v1.14, IV v4.14), and makes
% external-reference self-skip depend on logical paper role rather than a
% versioned job name.
% ======================================================================

% ======================================================================
% macros.tex — Shared macros for the Born Geometry series
% ======================================================================

% ======================================================================
% 1. PACKAGES, BIBLIOGRAPHY, COLORS, AND UNIT DEFINITIONS
% ======================================================================
% Core packages shared across the complete Paper-I-A--VII-C series.  Loading them here makes
% the macro file a genuine series preamble while remaining compatible with existing
% paper preambles that load the same packages before \input.
\usepackage{amsmath}
\usepackage{amssymb}
\usepackage{amsthm}
\usepackage{mathrsfs}
\usepackage{array}
\usepackage{booktabs}
\usepackage{xcolor}
\makeatletter
\@ifpackageloaded{geometry}{}{\usepackage{geometry}\geometry{margin=1in}}
\makeatother
\usepackage{mdframed}

% Shared support packages required by the macros below and by most papers.
\usepackage{ragged2e}
\usepackage{longtable}
\usepackage{tabularx}
\usepackage{enumitem}
\usepackage[most]{tcolorbox}

% Series-wide numeric bibliography convention using the portable BibTeX path.
% All Papers I-A--VII-C use the ordinary \cite interface.
\usepackage[numbers,sort&compress]{natbib}
\providecommand{\printbibliography}{%
  \bibliographystyle{plainnat}%
  \bibliography{Born-v6_59}%
}

% Shared colors.  \providecolor avoids clobbering a paper-specific override.
\providecolor{postulatecolor}{RGB}{230,242,255}
\providecolor{compare}{RGB}{230,245,230}
\providecolor{definitioncolor}{RGB}{230,242,255}

\providecolor{openriskcolor}{RGB}{255,235,235}
\providecolor{bridgecolor}{RGB}{240,240,255}
\providecolor{gold}{RGB}{255,215,0}

% ======================================================================
% 2. CUSTOM ENVIRONMENTS
% ======================================================================
% Common table column types used in all papers.  Guarded so that a paper may
% define its own versions before loading this file.
\makeatletter
\@ifundefined{NC@find@L}{%
  \newcolumntype{L}[1]{>{\RaggedRight\arraybackslash}p{#1}}%
}{}
\@ifundefined{NC@find@C}{%
  \newcolumntype{C}[1]{>{\Centering\arraybackslash}p{#1}}%
}{}
\@ifundefined{NC@find@Y}{%
  \newcolumntype{Y}{>{\RaggedRight\arraybackslash}X}%
}{}
\makeatother

% ======================================================================
% 3. THEOREM STYLES AND ENVIRONMENTS
% ======================================================================
% Custom underlined theorem style (useful for remarks or notes)
\newtheoremstyle{underlined}
  {3pt}   % Space above
  {3pt}   % Space below
  {}      % Body font (empty means normal/roman text, typical for remarks)
  {}      % Indent amount
  {}      % Theorem head font (can be \itshape, \bfseries, etc.)
  {.}     % Punctuation after theorem head
  { }     % Space after theorem head
  {\underline{\thmname{#1}\thmnumber{ #2}}\thmnote{ (#3)}} % Theorem head formatting

% Standard plain theorem style environments.  Each environment is guarded so
% the shared macro file can be loaded after a paper-specific theorem setup.
\theoremstyle{plain}
\makeatletter
\@ifundefined{theorem}{\newtheorem{theorem}{Theorem}[section]}{}
\@ifundefined{lemma}{\newtheorem{lemma}{Lemma}[section]}{}
\@ifundefined{proposition}{\newtheorem{proposition}{Proposition}[section]}{}
\@ifundefined{corollary}{\newtheorem{corollary}{Corollary}[section]}{}
\makeatother

% Definition style environments (typically non-italicized body)
\theoremstyle{definition}
\makeatletter
\@ifundefined{definition}{\newtheorem{definition}{Definition}[section]}{}
\@ifundefined{postulate}{\newtheorem{postulate}{Postulate}[section]}{}
\@ifundefined{closure}{\newtheorem{closure}{Closure}[section]}{}
\makeatother

% Applying the custom underlined style
\theoremstyle{underlined}
\makeatletter
\@ifundefined{remark}{\newtheorem{remark}{Remark}[section]}{}
\@ifundefined{convention}{\newtheorem{convention}{Convention}[section]}{}
\@ifundefined{principle}{\newtheorem{principle}{Principle}[section]}{}
\makeatother

% Framed proof environment.
% Every \begin{proof} ... \end{proof} is automatically framed/shaded.
% This avoids mdframed/trivlist nesting, which can drop the top border after
% theorem-like environments.
\makeatletter
\renewenvironment{proof}[1][\proofname]{%
  \par\addvspace{6pt}%
  \noindent
  \begin{tcolorbox}[%
    enhanced,
    breakable,
    colback=black!4,
    colframe=black!35,
    boxrule=0.35pt,
    arc=1pt,
    left=6pt,
    right=6pt,
    top=4pt,
    bottom=4pt,
    boxsep=0pt,
    before skip=0pt,
    after skip=6pt,
    parbox=false
  ]%
  \pushQED{\qed}%
  \normalfont\itshape #1\@addpunct{.}\enspace
  \normalfont\ignorespaces
}{%
  \unskip\nobreak\hfill\popQED
  \end{tcolorbox}%
  \par\addvspace{2pt}%
}
\makeatother

\mdfdefinestyle{postulatebox}{
  backgroundcolor=yellow!10,
  linecolor=black,
  linewidth=0.4pt,
  skipabove=\baselineskip,
  skipbelow=\baselineskip,
  nobreak=true
}

% ======================================================================
% PROOF-READABILITY HELPERS
% ======================================================================

\newcommand{\ProofPlan}[1]{%
  \par\smallskip
  \noindent\textbf{Proof plan. }#1%
  \par\smallskip
}

\newlist{proofsteps}{enumerate}{2}
\setlist[proofsteps,1]{
  label=\textbf{(\alph*)},
  ref=\alph*,
  leftmargin=2.2em,
  itemsep=0.45\baselineskip,
  topsep=0.35\baselineskip
}
\setlist[proofsteps,2]{
  label=\textbf{(\roman*)},
  ref=\roman*,
  leftmargin=2.0em,
  itemsep=0.25\baselineskip,
  topsep=0.2\baselineskip
}


% ======================================================================
% 4. NUMBER SYSTEMS AND ALPHABETS
% ======================================================================
\providecommand{\RR}{\mathbb{R}} % Real numbers
\providecommand{\PP}{\mathbb{P}} % Projective space / Probability
\providecommand{\CC}{\mathbb{C}} % Complex numbers
\providecommand{\ZZ}{\mathbb{Z}} % Integers
\providecommand{\NN}{\mathbb{N}} % Natural numbers
\providecommand{\HH}{\mathbb{H}} % Quaternions / Hyperbolic space
\providecommand{\LFRayGeom}{\mathsf{LF\!\mbox{-}RAY\!\mbox{-}GEOM}}
\providecommand{\LFRep}{\mathsf{LF\!\mbox{-}REP}}
\providecommand{\LFState}{\mathsf{LF\!\mbox{-}STATE}}
\providecommand{\LFLim}{\mathsf{LF\!\mbox{-}LIM}}
\providecommand{\LFScat}{\mathsf{LF\!\mbox{-}SCAT}}
% Number systems / projective spaces
\providecommand{\CP}{\mathbb{CP}}
\providecommand{\RP}{\mathbb{RP}}

% ======================================================================
% 5. GROUPS AND LIE THEORY
% ======================================================================
\providecommand{\SO}{\ensuremath{\mathrm{SO}}}     % Special Orthogonal group
\providecommand{\SU}{\ensuremath{\mathrm{SU}}}     % Special Unitary group
\providecommand{\Spin}{\ensuremath{\mathrm{Spin}}} % Spin group
\providecommand{\CS}{\ensuremath{\mathrm{CS}}}     % Chern-Simons (capitalized)
\providecommand{\cs}{\ensuremath{\mathrm{cs}}}     % Chern-Simons (lowercase)

% ======================================================================
% 6. COMMON MANIFOLDS AND TOPOLOGY
% ======================================================================
\providecommand{\Sone}{S^1}                 % Standard circle

% ======================================================================
% 7. MATH OPERATORS (Declared via \DeclareMathOperator)
% ======================================================================
\DeclareMathOperator{\End}{End}   % Endomorphisms
\DeclareMathOperator{\Aut}{Aut}   % Automorphisms
\DeclareMathOperator{\spec}{spec} % Lowercase spectrum
\DeclareMathOperator{\tr}{tr}     % Lowercase trace (if you already have \Tr)
\DeclareMathOperator{\supp}{supp} % Support
\DeclareMathOperator{\TrA}{Tr_{A}} % Partial trace over subsystem A
\DeclareMathOperator{\TrB}{Tr_{B}} % Partial trace over subsystem B
% Lie/algebra operators
\DeclareMathOperator{\Ad}{Ad}
\DeclareMathOperator{\Sym}{Sym}
\DeclareMathOperator{\su}{su}
% ======================================================================
% 8. STANDARD OPERATOR-LIKE SYMBOLS (Safe definitions via \providecommand)
% ======================================================================
\providecommand{\vol}{\operatorname{vol}}   % Volume form
\providecommand{\Vol}{\operatorname{Vol}}   % Total volume
\providecommand{\Ric}{\operatorname{Ric}}   % Ricci tensor
\providecommand{\Rm}{\operatorname{Rm}}     % Riemann tensor
\providecommand{\Tr}{\operatorname{Tr}}     % Uppercase trace
\providecommand{\Spec}{\operatorname{Spec}} % Uppercase spectrum
\providecommand{\diag}{\operatorname{diag}} % Diagonal matrix
\providecommand{\rank}{\operatorname{rank}} % Rank
\providecommand{\scal}{\mathcal R} % Scalar curvature
\providecommand{\Cl}{\operatorname{Cl}}     % Clifford algebra
\providecommand{\diam}{\operatorname{diam}} % Diameter

% ======================================================================
% 9. FUNCTIONS, TRIGONOMETRY, AND CALCULUS
% ======================================================================

% Codifferential naming conventions
% (Avoid \ddag: commonly defined by physics.sty and symbol packages)
\providecommand{\codiff}{\mathrm{d}^{\dagger}} % Codifferential operator

% Quantum mechanics bra-ket notation (needed for basis tables)
\providecommand{\ket}[1]{\lvert #1 \rangle}
\providecommand{\bra}[1]{\langle #1 \rvert}
\providecommand{\braket}[2]{\langle #1 \mid #2 \rangle}

% ======================================================================
% 10. PHYSICAL SCALES AND CONSTANTS (Paper 1c / Born Geometry)
% ======================================================================
\providecommand{\Rp}{\ensuremath{R_{\mathrm P}}}     % matched common EH curvature scale (not defined from observed G)
\providecommand{\hBorn}{\ensuremath{h_{\mathrm B}}} % physical representative 2\pi\hbar of one primitive Born scalar action cell
\providecommand{\UEHjoint}{\ensuremath{U_{J}}} % joint-Dirac EH curvature-coordinate scale 8 kappa_J
\providecommand{\Gread}{\ensuremath{G_{\rm scal}}} % scalar-action coordinate c_op^3/(h_B U); physical matching explicit
\providecommand{\GreadJoint}{\ensuremath{G_{J,\mathrm{raw}}}} % raw ordinary-trace coordinate = G_scal/96
\providecommand{\Rds}{R_{\mathrm{dS}}}    % de Sitter Curvature
\providecommand{\Mpl}{M_{\mathrm{Pl}}}    % Planck mass
\providecommand{\lp}{\ell_{\mathrm{Pl}}}  % Planck length
\providecommand{\Pmax}{P_{\mathrm{max}}}  % Maximum momentum
\providecommand{\Mplbar}{\bar{M}_{\mathrm{Pl}}} % Reduced Planck mass
\providecommand{\eff}{\mathrm{eff}}       % Subscript for effective quantities
\providecommand{\Cosmo}{\ensuremath{\Lambda}}                % Cosmological constant (use \Cosmo in source)
\providecommand{\LCDM}{\ensuremath{\Cosmo\mathrm{CDM}}}  % \LambdaCDM (typesets with \Cosmo)
\providecommand{\Lcut}{\ensuremath{\Lambda_{\mathrm{cut}}}}  % Cutoff scale (brace when used as script)
\providecommand{\Rcs}{\ensuremath{R_{\mathrm{cs}}}}      % Configuration-space radius
\providecommand{\CosmoEff}{\ensuremath{\Cosmo_{\eff}}}   % Effective cosmological constant

% Disambiguation: \Rmom is the momentum-space (S^3) radius. 
% Reserve plain R for the 4D curvature scalar in f(R).
\providecommand{\Rmom}{\ensuremath{R_{\mathrm{mom}}}} 
\providecommand{\Punit}{P_{0}} % Momentum unit set by \Rmom (= \hbar/\Rmom)
\providecommand{\Rvec}{\ensuremath{R_{\mathrm{vec}}}} % invariant vector-space radius in the reduced round sector
% Specific numeric values and expressions
\providecommand{\Vlogcap}{L_{\mathrm{Hopf}}^{(1)}}
\providecommand{\Hopf}{\operatorname{Hopf}}
% Scalar phase (S^1) sector + universal cover notation (Paper I-A)
\providecommand{\tper}{T_\star}      % Time period/scale
\providecommand{\Eper}{E_\star}      % Energy period/scale
\providecommand{\ttilde}{\tilde t}   % Universal cover time
\providecommand{\Etilde}{\tilde E}   % Universal cover energy

% ======================================================================
% 10A. SERIES-WIDE LOGICAL / NOTATIONAL DISAMBIGUATIONS
% ======================================================================
% Use these symbols to avoid collisions among: (i) the Paper--III
% no-insertion/geometric-closure class, (ii) Paper--I-A trimer diagonals,
% (iii) Paper--V-A/V-C family-space diagonals, (iv) scalar-index running,
% and (v) protected shell-clock labels.
\providecommand{\Rprim}{\ensuremath{\mathsf R_{\mathrm{prim}}^{\Sigma}}} % primitive finite structural readouts
\providecommand{\Rloop}{\ensuremath{\mathsf R_{\circ}^{\Sigma}}} % closed-loop readout role
\providecommand{\Rbdry}{\ensuremath{\mathsf R_{\partial}^{\Sigma}}} % boundary-conditioned readout role
\providecommand{\Rgen}{\ensuremath{\mathsf R_{\mathrm{fin}}^{\Sigma}}} % finite generated closure of primitive roles
\providecommand{\PhysOcc}{\operatorname{PhysOcc}}
\providecommand{\Actual}{\operatorname{Actual}}
\providecommand{\GeomLic}{\operatorname{GeomLic}}
\providecommand{\InvDist}{\operatorname{InvDist}}
\providecommand{\Dcl}{\ensuremath{\mathsf{D}_1}}      % unmarked positive Dirac-even closed spectral-readout form class
\providecommand{\Dheat}{\ensuremath{\mathsf D_{\mathrm{heat}}}} % one-parameter heat-semigroup subclass
\providecommand{\DtriI}{\ensuremath{D^{\mathrm{tri}}_1}} % Paper--I-A cyclic-trimer diagonal
\providecommand{\DtriII}{\ensuremath{D^{\mathrm{tri}}_2}}
\providecommand{\DfamI}{\ensuremath{D^{\mathrm{fam}}_1}} % Family-space diagonal selector
\providecommand{\DfamII}{\ensuremath{D^{\mathrm{fam}}_2}}
\providecommand{\alpharun}{\ensuremath{\alpha_{\mathrm{run}}}} % scalar spectral-index running
\providecommand{\Ksh}{\ensuremath{K_{\mathrm{sh}}}} % protected shell-clock coordinate
\providecommand{\ddef}{\ensuremath{\delta_{\mathrm{def}}}} % universal defect scalar after Paper--III realization

% ======================================================================
% 10B. SERIES-WIDE LOGICAL STATUS TAGS
% ======================================================================
% Upright small-caps wrapper for audit/status labels.  It prevents LaTeX from
% trying to use unavailable small-caps italic font shapes inside theorem bodies
% such as TU/lmr/m/scit.
\providecommand{\AuditLabel}[1]{}
\renewcommand{\AuditLabel}[1]{{\normalfont\scshape #1}}

% These tags are bookkeeping labels for consistency audits. They do not
% create new mathematical hypotheses; they record how a datum is being used
% in the proof chain.
\providecommand{\StatusPrimitive}{}
\renewcommand{\StatusPrimitive}{\AuditLabel{Primitive}}
\providecommand{\StatusDerived}{}
\renewcommand{\StatusDerived}{\AuditLabel{Derived}}
\providecommand{\StatusImported}{}
\renewcommand{\StatusImported}{\AuditLabel{Imported}}
\providecommand{\StatusConditional}{}
\renewcommand{\StatusConditional}{\AuditLabel{Conditional}}
\providecommand{\StatusBridge}{}
\renewcommand{\StatusBridge}{\AuditLabel{Bridge-derived}}
\providecommand{\StatusComparison}{}
\renewcommand{\StatusComparison}{\AuditLabel{Comparison-only}}

% ======================================================================
% 10C. SERIES-WIDE CLAIM-WORDING AUDIT LABELS
% ======================================================================
% These are textual audit labels used to keep modal language consistent
% across the series.  They introduce no new mathematical assumptions.
\providecommand{\ClaimForced}{}
\renewcommand{\ClaimForced}{\AuditLabel{forced}}
\providecommand{\ClaimUnique}{}
\renewcommand{\ClaimUnique}{\AuditLabel{unique}}
\providecommand{\ClaimCanonical}{}
\renewcommand{\ClaimCanonical}{\AuditLabel{canonical}}
\providecommand{\ClaimMinimal}{}
\renewcommand{\ClaimMinimal}{\AuditLabel{minimal}}
\providecommand{\ClaimClosed}{}
\renewcommand{\ClaimClosed}{\AuditLabel{closed}}
\providecommand{\ClaimPrediction}{}
\renewcommand{\ClaimPrediction}{\AuditLabel{prediction}}

% ======================================================================
% 10D. SERIES-WIDE THEOREM-CHAIN AUDIT LABELS
% ======================================================================
% These labels mark the direction of load-bearing theorem chains. They are
% audit annotations only; they introduce no new mathematical assumptions.
\providecommand{\ChainInput}{}
\renewcommand{\ChainInput}{\AuditLabel{Chain input}}
\providecommand{\ChainOutput}{}
\renewcommand{\ChainOutput}{\AuditLabel{Chain output}}
\providecommand{\ChainGuard}{}
\renewcommand{\ChainGuard}{\AuditLabel{Chain guard}}
\providecommand{\ChainNoBackflow}{}
\renewcommand{\ChainNoBackflow}{\AuditLabel{No backflow}}

% ======================================================================
% 10E. SERIES-WIDE LOCAL THEOREM-AUDIT LABELS
% ======================================================================
% These labels mark a statement-level dependency audit around individual
% theorem/definition nodes. They are bookkeeping labels only.
\providecommand{\LocalAudit}{}
\renewcommand{\LocalAudit}{\AuditLabel{Local audit}}
\providecommand{\UsesOnly}{}
\renewcommand{\UsesOnly}{\AuditLabel{Uses only}}
\providecommand{\Exports}{}
\renewcommand{\Exports}{\AuditLabel{Exports}}
\providecommand{\DoesNotUse}{}
\renewcommand{\DoesNotUse}{\AuditLabel{Does not use}}

% ======================================================================
% 10F. SERIES-WIDE NUMERICAL-ROW AUDIT LABELS
% ======================================================================
% These labels classify individual rows in numerical, benchmark, status,
% and comparison tables.  They are audit annotations only and do not change
% any theorem, value, formula, or comparison number.
\providecommand{\RowAudit}{}
\renewcommand{\RowAudit}{\AuditLabel{Numerical-row audit}}
\providecommand{\RowTheorem}{}
\renewcommand{\RowTheorem}{\AuditLabel{Theorem-fixed row}}
\providecommand{\RowBridge}{}
\renewcommand{\RowBridge}{\AuditLabel{Bridge-fixed row}}
\providecommand{\RowRepresentative}{}
\renewcommand{\RowRepresentative}{\AuditLabel{Representative row}}
\providecommand{\RowComparison}{}
\renewcommand{\RowComparison}{\AuditLabel{Comparison row}}
\providecommand{\RowCalibrated}{}
\renewcommand{\RowCalibrated}{\AuditLabel{Calibrated row}}
\providecommand{\RowEstimate}{}
\renewcommand{\RowEstimate}{\AuditLabel{Estimate row}}
\providecommand{\RowOpen}{}
\renewcommand{\RowOpen}{\AuditLabel{Open row}}
\providecommand{\RowNoSelector}{}
\renewcommand{\RowNoSelector}{\AuditLabel{No selector}}

% ======================================================================
% 10G. TABLE HELPERS
% ======================================================================
% Table helpers used by Paper V-C.  The unused readout-label preset has been
% removed; only active table helpers remain.
\providecommand{\gpsection}[1]{%
  \addlinespace[0.55em]
  \multicolumn{3}{@{}l}{\textsc{#1}}\\[-0.25em]
  \midrule
}
\providecommand{\mathcell}[1]{%
  \(\begin{gathered}#1\end{gathered}\)%
}

\newcommand{\DerBy}[1]{%
  \begin{tabular}[t]{@{}l@{}}#1\end{tabular}%
}

\newcommand{\TableNotes}[1]{%
  \par\medskip
  \begin{minipage}{0.94\linewidth}
  \footnotesize\emph{Notes.} #1
  \end{minipage}\par
}
% ======================================================================
% 10H. PAPER VI-A / QCD NOTATION
% ======================================================================
% Shared symbols used by Paper VI-A and downstream QCD/Yang--Mills handoffs.
\providecommand{\Pqcd}{\mathfrak P_{\rm VIII}^{\rm QCD}}
\providecommand{\Iqcd}{\mathfrak I_{\rm QCD}^{(7)}}
\providecommand{\muQCD}{\mu_{\rm QCD}^{\rm int}}
\providecommand{\LambdaQCD}{\Lambda_{\rm QCD}}
\providecommand{\LambdaBH}{\Lambda_{\rm BH}^{\rm QCD}}
\providecommand{\HQCD}{\mathcal H_{\rm QCD}^{\rm BH}}
\providecommand{\HQCDfin}{H_{\rm QCD}^{\rm fin}}
\providecommand{\Hraw}{\mathcal H_{\rm raw}}
\providecommand{\HphysQCD}{\mathcal H_{\rm phys}^{\rm QCD}}
\providecommand{\Vcol}{V_C}
\providecommand{\Ps}{P_{\rm sing}}
\providecommand{\PiOct}{\Pi_8}
\providecommand{\MSbar}{\overline{\rm MS}}
\providecommand{\YM}{\mathrm{YM}}
\providecommand{\Ktri}{K_{\triangle}}
\providecommand{\KYM}{K_{\YM}^{\rm BH}}
\providecommand{\aBH}{a_{\rm BH}}
\providecommand{\epsBH}{\varepsilon_{\rm BH}}
\newcommand{\RowConstructive}{\textnormal{\textsc{Constructive}}}
\newcommand{\RowObligation}{\textnormal{\textsc{Bridge obligation}}}
\newcommand{\RowTarget}{\textnormal{\textsc{Target row}}}
\newcommand{\RowValidation}{\textnormal{\textsc{Validation row}}}
\providecommand{\slashed}[1]{\not{#1}}

% ======================================================================
% 10I. PAPER VI-B / CONTINUOUS-QFT NOTATION
% ======================================================================
% Shared symbols used by Paper VI-B and any downstream continuous-QFT or
% Yang--Mills handoff paper.  All definitions are guarded so individual
% papers may override them locally before loading this file if needed.
\providecommand{\A}{\mathfrak A}
\providecommand{\Hphys}{\mathcal H_{\rm phys}}
\providecommand{\QCD}{\mathrm{QCD}}
\providecommand{\BH}{\mathrm{BH}}
\providecommand{\loc}{\mathrm{loc}}

\providecommand{\RBH}{\mathcal R_{\rm BH}}
\providecommand{\DG}{\mathfrak D_G}
\providecommand{\PW}{\mathcal{P\!W}} % Peter--Weyl spectral-symbol subspace
\providecommand{\Toe}{\operatorname{Toe}} % Toeplitz compression/representative
\providecommand{\Hol}{\operatorname{Hol}} % Holonomy

% ======================================================================
% 11. LATE-LOADED PACKAGES (Must be at the end of the preamble)
% ======================================================================
% Load late packages only if a paper has not already loaded them, avoiding
% option clashes in paper-specific preambles.
\makeatletter
\@ifpackageloaded{hyperref}{}{\usepackage{hyperref}}
\makeatother
% --- Series-wide disambiguations (Pass 3) ---
\providecommand{\diracdelta}{\delta_{\rm D}} % Dirac delta distribution (delta reserved for hierarchy)
\providecommand{\Hg}{H_{\mathrm g}} % Hodge harmonic projector (keep H for Hubble)
\providecommand{\Gr}{G_{\mathrm r}} % Green operator (keep G for Newton)
\providecommand{\FD}{\mathrm{D}} % functional/variational derivative symbol (delta reserved)
\providecommand{\Id}{\mathbf{1}} % identity/Kronecker tensor (avoid \delta_{ij})


% BEGIN INLINE SOURCE: Paper00-ActiveExternalReferences-v2_130.tex
% Framework v2.191: current reference binding; existing definitions retained.
% Framework v2.190: current22-root registry with native clock-evolution proofs.
% Framework v2.189: exact current22-root external registry, including retained v2.188 hardening owners.
% Framework v2.187: current registry synchronization only; source release v2.186 preserved.
% Framework v2.186: authoritative current 22-paper registry.
% Framework v2.182: all current 22 logical paper AUX bindings.
% Framework v2.170: current I-B v1.43 external AUX binding.
% Framework v2.167: exact current 22-root registry after CP--BL2 integration.
% Framework v2.166: exact current 22-root compact-release AUX registry.
% Framework v2.165: authoritative current root bindings.
% Framework v2.164: exact current 22 manuscript AUX bindings.
% Framework v2.162: finite metric, functional Hilbert and scoped Ward integration.
% Framework v2.161: current 22 manuscript AUX bindings; source only.
% Framework v2.156 / registry v2.113: synchronized current 22 manuscript roots; source-only registry.
% Framework v2.143 / registry v2.111: active manuscript AUX basenames; no generated AUX or PDF included.
% ======================================================================
% Central external-reference imports, numbers only, no PDF jumping
% ======================================================================

% xr-hyper is loaded centrally so that papers may retain their existing
% \usepackage{xr-hyper} lines without redefining \externaldocument after this
% registry is installed.  If plain xr was loaded earlier, retain it.
\makeatletter
\@ifpackageloaded{xr-hyper}{}{%
  \@ifpackageloaded{xr}{}{\usepackage{xr-hyper}}%
}
\makeatother

% Make external-document registration idempotent.  The registry identifies a
% logical paper by its Paper<number><optional letter> prefix.  Prefix-qualified
% imports remain independent.  Current paper roots contain no local
% \externaldocument calls; this guard also makes accidental duplicates harmless.
\ExplSyntaxOn
\prop_new:N \g_bh_external_refs_prop
\cs_new_eq:NN \BHOriginalExternalDocument \externaldocument

\cs_new_protected:Npn \bh_external_document:nnnn #1#2#3#4
 {
  \tl_set:Nn \l_tmpa_tl {#3}
  \regex_replace_once:nnN
   { \A (?: xref/ )? (Paper[0-9]+[A-F]?) .* \Z }
   { \1 }
   \l_tmpa_tl
  \tl_set:Nx \l_tmpb_tl { #1 | \tl_use:N \l_tmpa_tl }
  \prop_if_in:NVTF \g_bh_external_refs_prop \l_tmpb_tl
   {
    \PackageInfo{BornHopf}
     {Skipping~duplicate~external~refs~for~#3}
   }
   {
    \prop_gput:NVn \g_bh_external_refs_prop \l_tmpb_tl {#3}
    \IfNoValueTF{#4}
     { \BHOriginalExternalDocument[#1][#2]{#3} }
     { \BHOriginalExternalDocument[#1][#2]{#3}[#4] }
   }
 }

\RenewDocumentCommand{\externaldocument}{O{} O{} m o}
 { \bh_external_document:nnnn {#1}{#2}{#3}{#4} }

\NewDocumentCommand{\BHExternalRefs}{m}
 {
  % Do not import the aux file of the logical paper currently being compiled,
  % even when the job and registry use different version suffixes.
  \tl_set:Nx \l_tmpa_tl { \c_sys_jobname_str }
  \regex_replace_once:nnN
   { \A (Paper[0-9]+[A-F]?) .* \Z }
   { \1 }
   \l_tmpa_tl
  \tl_set:Nn \l_tmpb_tl { #1 }
  \regex_replace_once:nnN
   { \A (Paper[0-9]+[A-F]?) .* \Z }
   { \1 }
   \l_tmpb_tl
  \str_if_eq:eeTF { \tl_use:N \l_tmpa_tl } { \tl_use:N \l_tmpb_tl }
   {
    \PackageInfo{BornHopf}
     {Skipping~external~refs~for~#1~because~this~is~the~current~document}
   }
   {
    \IfFileExists{xref/#1.aux}
     { \externaldocument[][nocite]{xref/#1}[#1.pdf] }
     {
      \IfFileExists{#1.aux}
       { \externaldocument[][nocite]{#1} }
       {
        \PackageInfo{BornHopf}
         {External~aux~file~#1.aux~not~found;~no~labels~imported}
       }
     }
   }
 }
\ExplSyntaxOff

% External references are automatic.  This file is the sole owner of both the
% bibliography and the current Paper-I-A--VII-D AUX registry (30 roots including I-C).  Missing auxiliaries
% are informational on the first pass and are populated by the source-only
% two-pass series build.  Basenames deliberately omit temporary download/copy
% suffixes such as ``(1)'' or ``-1''.
  % Paper I-A: intrinsic S3 and exceptional group introduction.
  \BHExternalRefs{Paper1A-Geometry-v5_178}

  % Paper I-B: complete exceptional geometry and finite spectral structures.
  \BHExternalRefs{Paper1B-ExceptionalGeometry-v1_17}
  % Paper I-C: reciprocal product M8 geometry.
  \BHExternalRefs{Paper1D-TechnicalMaterial-v1_67}
  % Paper I-D: terminal technical proofs and calculations.
  \BHExternalRefs{Paper1C-M8Geometry-v1_142}

\ifdefined\BHMathOnly
  % Papers I-A and I-B import only the mathematical atlas labels.
\else
  % Paper II: five physical-mathematical main papers plus one technical proof volume.
  \BHExternalRefs{Paper2A-GeometricConstitution-v3_224}
  \BHExternalRefs{Paper2B-MomentumCarrier-v2_145}
  \BHExternalRefs{Paper2C-ReciprocalLocalization-v2_170}
  \BHExternalRefs{Paper2D-SymmetrySpectra-v2_143}
  % The former II-F export labels remain in II-E; current II-F is the technical volume.
  \BHExternalRefs{Paper2E-KinematicExports-v3_138}
  % Paper II-F: complete technical proof and calculation owner for II-A--II-E.
  \BHExternalRefs{Paper2F-TechnicalMaterial-v1_68}

  % Paper III: Dynamics and Intrinsic Geometric Constitution (six components).
  \BHExternalRefs{Paper3A-GeneratedDynamics-v1_177}
  \BHExternalRefs{Paper3B-SpectralAction-v1_189}
  \BHExternalRefs{Paper3C-ProtectedHierarchy-v1_191}
  \BHExternalRefs{Paper3D-RunningResponse-v1_166}
  \BHExternalRefs{Paper3E-GravityAction-v1_200}
  \BHExternalRefs{Paper3F-CompleteDynamics-v1_268}
  % Paper III-G: complete preserved technical proof owner.
  \BHExternalRefs{Paper3G-TechnicalMaterial-v1_54}

  % Paper IV-A: generated finite quantum mechanics.
  \BHExternalRefs{Paper4A-FiniteQuantumMechanics-v4_215}
  % Paper IV-B: complete finite quantum processes and physical response.
  \BHExternalRefs{Paper4B-QuantumResponse-v1_53}
  % Paper IV-C: extended proofs and finite technical calculations.
  \BHExternalRefs{Paper4C-TechnicalMaterial-v1_49}

  % Paper V-A: The Standard Model.
  \BHExternalRefs{Paper5A-FiniteMatterGeometry-v3_213}

  % Paper V-B: Finite Spectrum and Flavor Closure.
  \BHExternalRefs{Paper5B-TerminalSpectrum-v5_201}

  % Paper V-C: Finite Interaction Boundaries of the Standard Model.
  \BHExternalRefs{Paper5C-InteractionResponse-v2_180}
  % Paper V-D: preserved technical proof and calculation owner.
  \BHExternalRefs{Paper5D-TechnicalMaterial-v1_55}

  % Paper VI-A: Gravity and Closed Cosmology.
  \BHExternalRefs{Paper6A-GravityCosmology-v2_263}

  % Paper VI-B: Geometric Dark Sector and Cosmological Readouts.
  \BHExternalRefs{Paper6B-CosmologyReadouts-v1_210}

  % Paper VI-C: complete technical proof and comparison owner.
  \BHExternalRefs{Paper6C-TechnicalMaterial-v1_51}

  % Paper VII-A: Finite QCD Kernel and Nonperturbative Closure.
\BHExternalRefs{Paper7A-QCDKernel-v1_177}

  % Paper VII-B: Hadronic, Thermal, and Nuclear QCD Readouts.
  \BHExternalRefs{Paper7B-QCDReadouts-v1_164}

  % Paper VII-C: Continuous Quantum Field Theory.
  \BHExternalRefs{Paper7C-ContinuousQFT-v2_206}

  % Paper VII-D: complete technical material for VII-A through VII-C.
  \BHExternalRefs{Paper7D-TechnicalMaterial-v1_58}

\fi

% END INLINE SOURCE: Paper00-ActiveExternalReferences-v2_130.tex


\providecommand{\LambdacomparisonBH}{\Lambda_{\rm BH}^{\rm comp}}
\providecommand{\deltageom}{\delta_{\rm geom}}
\providecommand{\alphasseed}{\alpha_s^{\rm seed}}
\providecommand{\Hseedzero}{\mathsf H_{\rm seed}^{0}}
\setlength{\tabcolsep}{2pt}
\AtBeginEnvironment{longtable}{\footnotesize}
\setlength{\emergencystretch}{3em}
\title{\textbf{Born--Hopf Geometry in a Closed Universe VII-D: Technical Material}\\[0.4em]
\large Complete finite-QCD, observable and bounded-continuum calculations}
\author{Peter Kugelmann\thanks{\texttt{pkugelmann@proton.me}}}
\date{Version 1.60 --- 19 November 2026 \\ S$^3$ finite-response integration v3.035}
\setcounter{tocdepth}{2}
\BHExternalRefs{Paper7A-QCDKernel-v1_177}
\begin{document}
% Publication typesetting: controlled prose line-breaking, without scaling mathematics.
\emergencystretch=2em
\maketitle
\begin{abstract}
This companion contains the inherited full auxiliary mathematical results
supporting Papers VII-A--VII-C: exact Hurwitz and color matrix calculations,
gauge heat/spectral bounds, chiral and source comparisons, hadron/thermal and
optional finite LFRep proofs, fixed-ultraviolet analytic and RG certificates,
finite subsystem classification, and conditional numerical targets. In
particular the predecessor's $E_{36}$/$E_{40}$ Hurwitz split-grammar numerical
comparisons are retained in full as phenomenological targets whose physical
operator derivation is not yet established. No theorem is promoted from a
conditional comparison into an unconditional QCD prediction by extraction.
\end{abstract}
\tableofcontents
\listoftables
\clearpage

\section{Finite Hurwitz, Haar projection, and auxiliary carrier calculations}
\label{p13:sec:v237-d1}
% Source Paper VII-A, original line 7845; inherited proof bodies unchanged.
\subsection{Supplementary ordinary overlap-insertion Gram}
\label{p10:app:exroot-overlap-gram}

\begin{proposition}[Finite ordinary operator Gram of the mass insertion]
\label{p10:prop:exroot-overlap-gram}
On the finite overlap chart of
Corollary~\ref{p10:cor:v165-overlap-scalar-source-derivative},
let \(\epsilon_H\) and \(\Gamma_5\) be the Hermitian involutions,
\(V=\Gamma_5\epsilon_H\), and \(d\) the complex matrix dimension.
Write \(O_m=\mathcal O_m^{\rm BH}=(I-V)/2\) for the existing mass
insertion. It satisfies, in the ordinary
complex Hilbert--Schmidt metric,
\[
 \boxed{\|O_m\|_{\rm HS}^2
 =\frac{d-\operatorname{Re}\operatorname{Tr}V}{2}
 =\sum_{j=1}^d\sin^2(\theta_j/2),}
\]
where \(e^{i\theta_j}\) are the eigenvalues of the unitary \(V\).
Consequently this ordinary Gram equals \(240\) precisely when
\(\operatorname{Re}\operatorname{Tr}V=d-480\).
\end{proposition}
\begin{proof}
The product of the involutions is unitary, so
\[
 O_m^\dagger O_m=(2I-V-V^\dagger)/4.
\]
Take the ordinary trace, then use the spectral theorem for \(V\).
\end{proof}

This finite insertion identity supplies a diagnostic for the actual
operator already present in the overlap constructor. It does not
identify the ordinary Hilbert--Schmidt metric with a stationary CAR
source pairing, a renormalized interacting scalar vertex, or the
named-scheme coefficient transported by Paper~VII-C.
In particular the condition on \(\operatorname{Tr}V\) has not been
evaluated on the actual charged physical member in the exceptional-root
calculation. The existing G16 source normalization and pole transport
retain their current domains; no \(E_{32}/\sqrt{240}\) unity-crossing
claim follows.

% END INLINE SOURCE: Paper7A-OverlapSourceGram-Appendix-v1_0.tex

% END EXCEPTIONAL-ROOT-INTEGRATION v2.173

% BEGIN XFE-EXTENSIVE-FERMION-DIAGNOSTICS v2.180

% BEGIN INLINE SOURCE: Paper7A-ExtensiveFermionConstructorDiagnostics-Appendix-v1_0.tex
% Paper7A-ExtensiveFermionConstructorDiagnostics-Appendix-v1_0.tex
% Framework v2.180.  Appendix-only diagnostics under the Born-Series placement rule.

% Source Paper VII-A, original line 7894; inherited proof bodies unchanged.
\subsection{extensive-fermion constructor diagnostics}
\label{p10:app:v2180-xfe-diagnostics}

This appendix records negative and optional results that delimit the live
extensive-fermion constructor obligation.  None of the candidate constructions
below is promoted to a constituted physical quark regulator.

\begin{proposition}[Sharp-band compression does not represent arbitrary local color multiplication]
\label{p10:prop:v2180-compressed-local-gauge-no-go}
Let $P_m$ be a nontrivial finite sharp scalar-band projector and let $M_g$ be
pointwise multiplication by a local color transformation $g$.  Put
$T_m(g)=P_mM_gP_m$.  Then
\[
 \boxed{T_m(g)T_m(h)-T_m(gh)
 =-P_mM_g(1-P_m)M_hP_m.}
\]
Hence $T_m$ is a representation of the full local multiplication group only
when the band is invariant under all such multiplications, which a nonzero
finite harmonic band is not.
\end{proposition}
\begin{proof}
Insert $I=P_m+(1-P_m)$ between $M_g$ and $M_h$ and use
$M_gM_h=M_{gh}$.  For a nonconstant smooth multiplier, multiplication of a
top-shell harmonic has a component outside the finite band, so the leakage
term is generically nonzero.
\end{proof}

\begin{proposition}[Finite retained holonomies do not determine the exact local Dirac compression]
\label{p10:prop:v2180-finite-holonomy-dirac-no-go}
Let $\Gamma_L$ be a finite retained path graph in a three-dimensional smooth
spatial manifold.  In general there is no functional of only the retained
edge holonomies that equals the exact Galerkin connection term $P_Lc(A)P_L$
for every smooth connection $A$.
\end{proposition}
\begin{proof}
The finite union of retained one-dimensional paths leaves an open ball $B$
disjoint from all retained edges.  Choose a nonzero smooth Lie-algebra-valued
one-form $\delta A$ supported in $B$.  Every retained edge holonomy is
unchanged under $A\mapsto A+\delta A$, whereas a matrix element of
$P_Lc(\delta A)P_L$ can be chosen nonzero.  Thus the retained holonomy list
cannot determine the exact compressed local connection term.
\end{proof}

\begin{lemma}[Positive massive determinant from finite chiral anti-Hermiticity]
\label{p10:lem:v2180-positive-mass-chiral-antihermitian}
If a finite massless operator satisfies
\[
 D^\dagger=-D,\qquad \{D,\Gamma_5\}=0,
\]
then its nonzero eigenvalues pair as $\pm i\lambda_j$ and, for every real
$m>0$,
\[
 \boxed{\det(D+m)=m^{n_0}\prod_j(m^2+\lambda_j^2)>0.}
\]
Thus massive vectorlike positivity alone does not require transporting the
existing overlap regulator to a future extensive carrier.
\end{lemma}
\begin{proof}
Anti-Hermiticity puts the spectrum on the imaginary axis and the chiral
anticommutator maps an $i\lambda$ eigenvector to a $-i\lambda$ eigenvector.
Multiplying each conjugate pair gives $m^2+\lambda^2$, while every zero mode
contributes $m$.
\end{proof}

\begin{remark}[Covariant graph-Laplacian band is an extension candidate only]
\label{p10:rem:v2180-graph-laplacian-extension-candidate}
Given independently supplied unitary covariant shifts $T_i(U)$ and a physical
bundle frontier $\Lambda$, the finite operators
\[
 L_U=\sum_i\frac{2-T_i-T_i^\dagger}{h^2},\qquad
 \nabla_{i,U}=\frac{T_i-T_i^\dagger}{2h},
\]
\[
 P_\Lambda(U)=\mathbf1_{[0,\Lambda^2]}(L_U),\qquad
 D_{\rm phys}(U)=P_\Lambda(U)\sum_i\gamma^i\nabla_{i,U}P_\Lambda(U)
\]
are exactly gauge covariant, anti-Hermitian and chirally odd when the gamma
matrices commute with the spin-even projector.  For the free periodic graph,
every nonzero naive doubler has $L\ge4/h^2$, so $\Lambda<2/h$ excludes it from
this proposed low band.  These algebraic facts do not constitute the
construction in Paper~VII-A: the active framework does not generate the needed
fermion/bundle frontier $\Lambda$ or the cofinal extensive covariant shifts.
\end{remark}

\begin{remark}[Finite matrix gauge calculus is likewise unconstituted]
\label{p10:rem:v2180-fuzzy-gauge-extension-candidate}
A finite matrix calculus can make
$D_A=\gamma^i(\Pi_i+A_i)$ exactly covariant under
$A_i\mapsto UA_iU^\dagger-[\Pi_i,U]U^\dagger$, but the active framework does
not identify the full unitary group of the auxiliary finite matrix algebra
with physical local color.  This is therefore an optional extension candidate,
not a current physical constructor.
\end{remark}

% END INLINE SOURCE: Paper7A-ExtensiveFermionConstructorDiagnostics-Appendix-v1_0.tex


% BEGIN INLINE SOURCE: Paper7A-CovariantSpectralBand-Appendix-v1_0.tex
% Framework v2.186: conditional mathematical comparison, not physical constitution.


\section{Gauge heat coefficients, spectral, Schur and coercivity verification}
\label{p13:sec:v237-d2}
% Source Paper VII-A, original line 6912; inherited proof bodies unchanged.
\subsection{APBI action-owner diagnostics: homothetic area, cotangent direction, and conditional electric heat}
\label{p10:app:v2159-apbi-action-owner-diagnostics}

\begin{proposition}[Fixed-combinatorics homothetic Hurwitz faces cannot supply the APBI area bound]
\label{p10:prop:v2159-fixed-hurwitz-area-nogo}
The unit-radius Hurwitz triangular face has spherical side $\pi/3$ and area
\[
 a_H=3\arccos(1/3)-\pi=0.5512855984\ldots .
\]
Consequently a homothetic copy on $S^3(R_x)$ has
\[
 A_H(R_x)=a_HR_x^2.
\]
On the homogeneous reciprocal-position branch of
Corollary~\ref{p4:cor:exact-reciprocal-position-resolution},
$R_x/\ell_F=m_x+O(1)$, so $A_H(R_x)\to\infty$ as $m_x\to\infty$.
Therefore a reciprocal-position comparison presentation obtained only by a
fixed-combinatorics homothetic transport of one Hurwitz $24$-cell cannot obey
a level-independent APBI bound $A_f\le A_+<\infty$.
\end{proposition}
\begin{proof}
Paper~I-A, Theorem~\ref{p1:thm:s3one-structural-invariants}, gives the
Hurwitz edge length $\pi/3$ and, by the spherical cosine law,
$\cos\alpha=1/3$ for the equilateral face angle.  Spherical excess gives
$a_H=3\alpha-\pi$.  Rescaling the round radius multiplies two-area by
$R_x^2$.  The quoted reciprocal-position resolution theorem gives
$R_x/\ell_F=m_x+O(1)$.
\end{proof}

\begin{proposition}[Holonomy alone does not define a forward cotangent symplectic descent]
\label{p10:prop:v2159-holonomy-cotangent-firewall}
Let \(E\) be a finite oriented edge set in a smooth spatial region and let
\[
 f:=\operatorname{Hol}_E:\mathcal A_\Sigma\longrightarrow G^E
\]
be the finite holonomy map. The canonical cotangent construction for a
general map \(f:X\to Y\) is the \emph{fiber-product pullback}
\[
 \boxed{X\times_{Y}T^*Y\longrightarrow T^*X,\qquad
 (x,\eta_{f(x)})\longmapsto (df_x)^*\eta_{f(x)}.}
\]
It is not, without a section/choice of \(x\), a map
\(T^*Y\to T^*X\). For holonomy this means
\(\mathcal A_\Sigma\times_{G^E}T^*G^E\to T^*\mathcal A_\Sigma\),
with the last target interpreted in a cotangent class allowing the
edge-supported distributional covectors produced by the differential.
In left trivialization, finite link momenta
\(J_e\in\mathfrak g^*\) give the line-integral variation covector at
\emph{each specified connection} \(A\); a generic smooth Yang--Mills
momentum density \(\Pi\) is not of this edge-supported form. Hence
holonomy alone yields neither a canonical forward finite flux map
\(T^*\mathcal A_\Sigma\to T^*G^E\) nor a full-phase-space
symplectomorphism. A forward flux map needs additional generated data,
such as a dual-cell smearing/projection or same-owner cylindrical
construction.
\end{proposition}
\begin{proof}
For an edge \(e\) and a specified connection \(A\), the holonomy
variation is, in a consistent left/right convention,
\[
 (\delta U_e)U_e^{-1}
 =\int_e\operatorname{Ad}_{U_{e,s}}\!\bigl(\delta A(\dot e_s)\bigr)\,ds.
\]
Thus \((df_A)^*\) sends a link cotangent vector to the resulting
line-supported functional \emph{at} \(A\). Without specifying \(A\)
over the link configuration \(f(A)\), its target point and even its
linear differential are undetermined; this is exactly why the fiber
product is required. The continuum canonical one-form
\(\Theta_{\mathcal A}=\int_\Sigma\langle\Pi,\delta A\rangle\)
allows generic smooth momentum densities, which cannot all be represented
by finite sums of edge-supported functionals. No inverse or symplectic
descent from holonomy alone follows.
\end{proof}

\begin{theorem}[Conditional finite electric-kinetic heat realization]
\label{p10:thm:v2159-conditional-electric-heat}
Suppose a same-owner finite cylindrical realization of the generated electric
Yang--Mills kinetic term has already been supplied on
$\mathcal C_L=SU(3)^{E_L}$ in the form
\[
 H_{E,L}=\frac12\sum_{e\in E_L}\alpha_{e,L}\,
 q_C^{-1}(J_e,J_e),
 \qquad \alpha_{e,L}>0,
\]
with the canonical link cotangent brackets.  Quantize in the normalized Haar
representation by $J_e^a\mapsto-i\hbar X_{e,a}$, where the $X_{e,a}$ are
left-invariant vector fields for the Casimir metric $q_C$.  Then
\[
 \boxed{\widehat H_{E,L}
 =-\frac{\hbar^2}{2}\sum_e\alpha_{e,L}\Delta_e}
\]
and, for Euclidean transfer time $\tau>0$,
\[
 e^{-\tau\widehat H_{E,L}/\hbar}
 =\bigotimes_e\exp\!\left(\frac{t_{e,L}}2\Delta_e\right),
 \qquad
 \boxed{t_{e,L}=\tau\hbar\alpha_{e,L}}.
\]
Thus its integral kernel is the product of compact-group heat kernels
$\prod_eK_{t_{e,L}}(U_eU_e'^{-1})$.  This theorem is conditional on the
same-owner finite cylindrical kinetic realization; it neither constructs that
realization nor selects the nonlinear magnetic/spatial face action or equates
$t_{e,L}$ with the face heat time of
Corollary~\ref{p10:cor:v185-common-heat-face-stiffness}.
\end{theorem}
\begin{proof}
The quadratic cotangent Casimir quantizes as
$q_C^{-1}(J_e,J_e)\mapsto-\hbar^2\Delta_e$.  Distinct edge operators
commute.  On an irreducible Peter--Weyl block of Casimir $C_2(\lambda)$ the
kinetic transfer eigenvalue is
$\exp[-\tau\hbar\alpha_{e,L}C_2(\lambda)/2]$, exactly the heat-kernel
coefficient $\exp[-t_{e,L}C_2(\lambda)/2]$.
\end{proof}


% Source Paper VII-A, original line 7016; inherited proof bodies unchanged.
\subsection{native topology/measure boundary after G20 retirement}
\label{p10:sec:v161-g20-retirement}
The former G20 wrapper is no longer an independent physical-owner gap.
The finite overlap/GW owner, native integer charge, open unit-charge sectors,
positive primitive pure-heat sector magnitude, and the declared positive
massive fixed-carrier vectorlike measure are derived on their stated
domains. The quantitative massive bounds and exact unsaturated massless
endpoint are proved in Section~\ref{p10:sec:v2189-native-weighted-topology}.  The strong sector character and
quark determinant orientation are fixed upstream by the generated action and
rooted source geometry.

Two stronger tasks remain when explicitly claimed: (i) unrestricted
\(\nu_{\rm ov}=\nu_{\rm clutch}\) comparison over a larger admitted family,
and (ii) global determinant-line flatness/holonomy after introducing a
configuration-dependent chiral/Weyl subspace.  These are
continuum/chiral-regulator certification tasks and are exported to Paper~VII-C,
not prerequisites for native finite topology.


% Source Paper VII-A, original line 7034; inherited proof bodies unchanged.
\subsection{Supplementary unit-clutching regulator check}
\label{p10:app:v190-unit-clutching-check}
The three-layer member is strengthened by the residual certificate and clock-cover theorem in the main text.  The following retained stress test also records the earlier four- and five-layer checks.
For the canonical clutch $g_1(q)=\operatorname{diag}(q,1)$, a finite
interpolation using principal edge logarithms and a closing temporal ribbon
was assembled on $N_\tau=3,4,5$ layers at one representative point
$r_H=r_\tau=\rho=1$.  Direct diagonalization gives
\[
 \nu_{\rm ov}(g_1)=+1,\qquad
 \nu_{\rm ov}(g_1^{-1})=-1,\qquad
 \nu_{\rm ov}(1)=0,
\]
with nonzero minimum gaps on all tested representatives.  Flipping the flat
temporal spin sign leaves the tested index unchanged.  This is a finite
stress test and generator normalization for the stated family, not the
universal overlap/clutching comparison theorem; the latter is retained only
for the corresponding continuum/geometric export claim.


% Source Paper VII-A, original line 7052; inherited proof bodies unchanged.
\subsection{Exact integer certificate for the native index centers}
\label{p10:app:v256-exact-inertia-certificate}
This appendix states a \emph{conditional exact-certificate route} for the native
kernel, with a specified integer matrix constructor and the inequalities a
dyadic verifier must check. Numerical diagonalization can propose dyadic
witness arrays, but cannot replace their exact replay. The named sixteen
witness arrays, executable exact stencil/checker and rational results cited
in the historical \texttt{Audit/} directory are \emph{not present} in the
PapersMinimal v3.033, DAG Atlas v1.183, ProofChain Audits v1.75 or S3
Structure Network v1.36 archives available for this source integration.
Consequently their particular numerical outcomes are not independently
replayed here. The following algebraic deductions require the exact
inequalities as explicit, \emph{as-yet-undischarged certificate premises}.

\begin{remark}[Historical phase-05 payload claim and replay boundary]
\label{p10:rem:v258-phase05-payload}
The historical v2.256 phase-05 file was reported as empty, and a later
v2.257 archive was reported to replace it.
The named replacement verifier
\path{Paper00-NativeExactInertiaChecks-v1_1.py}, its manifest,
and \path{Paper00-NativePhase05ExactInertia-v1_1.json}
are referenced by name but are \emph{not included} in the current four
supplied framework archives listed above. Their claimed integer replay and
inertia $(433,431)$ therefore remain historical assertions pending the
actual byte-level witness and independent execution. No gap/number has
been recalculated or negated by the present source correction.
\end{remark}

\begin{theorem}[Conditional exact native mesh inertia and gap]
\label{p10:thm:v256-exact-native-inertia}
Use precisely the native three-layer parameters and Clifford/link
conventions of Theorem~\ref{p10:thm:v199-native-unit-component} and the
closing phase family of Theorem~\ref{p10:thm:v199-native-clock-covers}.
For each $\phi_j=j\pi/8$, $0\le j<16$, assume that the integer
matrix stencil below is instantiated and an independently replayable dyadic
pair $(Q_j,D_j)$ is provided whose exact rational Gram, congruence,
separation and inertia inequalities are precisely
\eqref{p10:eq:v256-exact-certificate-inequalities}.
\emph{Under this finite witness premise}, the resulting
$864\times864$ Hermitian kernel satisfies
\begin{equation}
 (n_-,n_+)=(433,431),\qquad
 \operatorname{gap}H(\phi_j)>
 \frac{2576}{10001}>\frac{32}{125}=0.256.
 \label{p10:eq:v256-exact-mesh-gap}
\end{equation}
Once the exact rational witness inequalities have been checked, these
conditional conclusions do not need an IEEE arithmetic or BLAS error model.
Then the existing phase Lipschitz and clock Fourier arguments give
\begin{equation}
 \operatorname{gap}H(\phi)>\frac{417}{7000},\qquad
 \nu_{\rm ov}(H(\phi))=1\quad(\phi\in\mathbb R),
 \qquad
 \nu_{\rm ov}(H_{3k})=k\quad(k\ge1)
 \label{p10:eq:v256-exact-cover-index}
\end{equation}
on the same native clock-cover and spin branches as before.
\end{theorem}
\begin{proof}
We first identify the matrix being certified, before using any witness.
Write each Hurwitz vertex as $q=p/2$, with
$p\in\{\pm2e_\mu\}\cup\{(\pm1,\pm1,\pm1,\pm1)\}$.
Adjacency is the exact integer test $p\cdot p'=2$.
Let $g_2(q)=q_0I-iq_a\sigma_a$ and
$R_2=2g_2(q)g_2(q')^*$. For each adjacent pair put
$\sqrt3\,n_{qq'}^a=\epsilon_a\in\{\pm1\}$ and
$G=\sum_a\epsilon_a\Gamma_a$.
The unit-parameter native stencil has diagonal $3\Gamma_5\otimes I_3$,
spatial edge block
\[
 \frac38\Gamma_5(\Gamma\cdot n_{qq'}-I_4)
       \otimes U_{qq'}(s_j),
 \qquad (s_0,s_1,s_2)=(0,1/2,1),
\]
and forward temporal block
$\frac12\Gamma_5(\Gamma_4-I_4)\otimes V_{j,q}$, with backward
blocks their adjoints. The temporal transporters are $I,I,g(q)$,
with the closing scalar phase on the last one.

On the color doublet the three spatial transports are
$I$, $(I+R_2/2)/\sqrt3$, and $R_2/2$; the third color is fixed.
Set $E=\operatorname{diag}(0,0,1)$,
$C=\operatorname{diag}(2I_2+R_2,0)$ and
$\widetilde R=\operatorname{diag}(R_2,2)$.
Before the common left factor $\Gamma_5\otimes I_3$, the spatial
coefficients of $16H=A+\sqrt3B$ are, respectively,
\[
\begin{array}{c|cc}
 s&A&B\\ \hline
 0&-6I_4\otimes I_3&2G\otimes I_3\\
 1/2&G\otimes C-6I_4\otimes E&-I_4\otimes C+2G\otimes E\\
 1&-3I_4\otimes\widetilde R&G\otimes\widetilde R
\end{array}
\]
The diagonal coefficient is $48\Gamma_5\otimes I_3$.
The temporal coefficient is $4\Gamma_5(\Gamma_4-I_4)\otimes(2V)$.
Thus every coefficient is a Gaussian integer; this is exactly the
previous native stencil, not a substituted comparison operator.
For a closing phase, separate its real and imaginary parts to write
\[
 16H(\phi)=A_0+\sqrt3B+\cos\phi\,C_0+\sin\phi\,D_0,
\]
where the four Gaussian-integer coefficient matrices are constructed
and checked Hermitian using integer arithmetic.

For each sample a certificate \emph{would consist of} a dyadic pair
$Q=(Q_R+iQ_I)/2^{24}$ and
$D=\operatorname{diag}(d_a)/2^{24}$, with integer arrays.
A corresponding exact verifier must compute $Q^*Q-I$ and the congruences of all four
coefficient matrices exactly. It encloses $\sqrt3$ and the mesh's
$\sqrt2$, $\sqrt{2+\sqrt2}$ and $\sqrt{2-\sqrt2}$ by rational
intervals whose endpoints are certified by integer square inequalities.
This gives exact rational upper bounds $\varepsilon_G,\varepsilon_C$
for
\[
 \|Q^*Q-I\|_{\rm op}\le\varepsilon_G,\qquad
 \|Q^*H(\phi_j)Q-D\|_{\rm op}\le\varepsilon_C.
\]
Here a Hermitian matrix's spectral norm is at most its absolute row-sum
norm, and $|a+ib|\le|a|+|b|$ supplies a rational entrywise majorant.
Before every fixed-width integer operation the verifier bounds every
product and partial sum strictly below $2^{63}-1$; final interval sums
use arbitrary-precision integers. Replay rejects optimized Python modes
that would disable assertions. The requisite exact inequalities, to be checked for every sample, are
\begin{equation}
 \begin{gathered}
 \varepsilon_G<\frac1{10000},\quad
 \varepsilon_C<\frac1{10000},\\
 \min_a|D_{aa}|>\frac{2577}{10000},\quad
 (n_-(D),n_+(D))=(433,431).
 \end{gathered}
 \label{p10:eq:v256-exact-certificate-inequalities}
\end{equation}
In particular $Q$ is invertible. The Hermitian perturbation
$Q^*HQ-D$ cannot move an eigenvalue of $D$ through zero, so Sylvester
inertia gives the displayed inertia of $H$.
Also, writing $K=Q^*HQ$, one has
$s_{\min}(K)\ge\min|D|-\varepsilon_C$ and
$H^{-1}=QK^{-1}Q^*$, whence
\[
 \operatorname{gap}H\ge
 \frac{\min|D|-\varepsilon_C}{1+\varepsilon_G}
 >\frac{2576}{10001}>\frac{32}{125}.
\]
The overlap sign trace is therefore $-2$, so its native index is $1$.

Finally $\|H(\phi)-H(\psi)\|\le|\phi-\psi|$ on the circular phase
parameter, and every phase lies within $\pi/16<11/56$ of a mesh point.
Consequently the gap is greater than
$32/125-11/56=417/7000$ on the entire phase circle.
No eigenvalue can cross zero, so the index is fixed. The exact finite
clock Fourier decomposition, including both stated temporal spin choices,
then gives $k$ unit-index blocks on the $3k$-layer cover. This is a
finite direct-sum theorem and introduces no continuum limit.
\end{proof}

\emph{If the stated exact dyadic witnesses are supplied and replayed}, the
resulting certificate removes the inherited arithmetic-model premise from
the corresponding all-link neighborhoods, positive finite topological-sector
bounds, overlap locality and local chiral charts. Until that replay, these
claims retain their original conditional certificate boundary.
Their existing link, action, mass, measure and domain hypotheses remain.
It does not identify an occurring physical topological sector, assign its
complete action weight, or prove a physical Yang--Mills mass gap.


\subsection{Historical floating-product residual evidence}
\label{p10:app:v199-inertia-certificate}
The following evidence is retained as a historical record. When the exact
certificate conditions of Appendix~\ref{p10:app:v256-exact-inertia-certificate}
are actually replayed, they supersede the arithmetic-model premise.
Until their cited witness files are available, the historical
floating-product acceptance model remains a distinct, conditional
computational evidence route; it cannot be replaced by an unreplayed claim.

This historical appendix gave the computational evidence used in
Theorems~\ref{p10:thm:v199-native-unit-component} and
\ref{p10:thm:v199-native-clock-covers}.  The accompanying scripts
\texttt{Paper00-NativeUnitClutchCertificate-v1\_0.py} and
\texttt{Paper00-NativeClockCoverCertificate-v1\_0.py} specify every matrix
entry, radical enclosure and product-error bound.  They use NumPy and
SciPy; their results are retained in
\texttt{Paper00-GeometricCertificateResults-v1\_0.json}.

The unit matrix is exactly \(H_*=(H_a+\sqrt3H_b)/16\).
Both coefficient matrices are Hermitian Gaussian-integer matrices, and
their componentwise absolute row sums are bounded by 192 and 64.  The
midpoint links need only rational quaternion entries and \(\sqrt3\).
Enclose that radical between dyadic endpoints whose squares are checked
using exact integers.  Treat the computed eigenvector and eigenvalue arrays
\(Q,D\) as exact dyadic witnesses, not as trusted exact eigenpairs.
Explicit standard IEEE binary64 matrix-product error bounds give
\[
 \|Q^\dagger Q-I\|_\infty<2.316\times10^{-6},\quad
 \|Q^\dagger H_*Q-D\|_\infty<5.879\times10^{-5},\quad
 \min_i|D_{ii}|>0.25777829.
\]
Both exact error matrices are Hermitian, so these infinity bounds dominate
their spectral norms.  The Gram bound makes \(Q\) invertible.  The
congruence bound keeps every sign of \(D\), whose negative/positive counts
are \((433,431)\); Sylvester inertia proves the same counts for \(H_*\).
Moreover, writing the two bounds as \(\varepsilon_G,\varepsilon_C\),
\[
 \min|\operatorname{Spec}H_*|
 \ge\frac{\min_i|D_{ii}|-\varepsilon_C}{1+\varepsilon_G}
 >0.2577189>0.256.
\]
The script uses \(u=2^{-53}\), conservative complex-product
\(\gamma_{8n}=8nu/(1-8nu)\), \(n=864\), explicit component-norm
bounds, and a further one-percent inflation of the scalar error estimates.
No underflow/overflow or nonstandard arithmetic is admitted by this model.
This is an a posteriori certificate under that model, not a proof-assistant
verification of a BLAS implementation.

For \(\phi_j=j\pi/8\), the second script encloses
\(\sqrt2\) and \(\sqrt{2\pm\sqrt2}\) by exact rational square checks.
The phase-component enclosure widths are below \(2^{-49}\).  The exact
sample matrices have component-infinity norm below 40; a conservative
assembly error bound is \(10^{-10}\).  Each of the 16 samples satisfies
the same inertia and a residual-certified gap greater than \(0.256\).
The script checks Gram error below \(10^{-4}\) and congruence error below
\(10^{-3}\) before accepting a sample; the actual recorded margins are
stronger.  The phase Lipschitz estimate and exact finite Fourier
decomposition in the main proof then establish every phase and every finite
clock cover without extrapolating a list of tested cover degrees.


% Source Paper VII-A, original line 7105; inherited proof bodies unchanged.
\subsection{Topology-resolving comparison and the sampling boundary}
\label{p10:app:v199-topology-comparison}
For Hermitian representatives on the same identified finite carrier,
\(\operatorname{gap}H_{\rm ref}\ge\gamma>0\) and
\(\|H_{\rm native}-H_{\rm ref}\|<\gamma\) imply equal sign traces:
the straight-line homotopy cannot cross zero.  To turn this into an index
match, the reference must already carry the correct topology and chirality
convention.  Different carrier dimensions require an explicit stabilization
by zero-index gapped blocks, not an untyped comparison of Dirac operators.

On an embedded \(SU(2)\) clutch branch, a specified simplicial readout
subdivision whose vertex images \(q_i\) lie in a common open hemisphere
has the continuous quaternion interpolant
\[
 g_{\rm int}(x)=\frac{\sum_i\ell_i(x)q_i}{|\sum_i\ell_i(x)q_i|}.
\]
The hemisphere normal pairs positively with the numerator, which therefore
never vanishes; the formula agrees on common faces.  If the original clutch
and this interpolant stay in the same hemisphere on each simplex, normalized
straight-line interpolation is a homotopy.  A sufficient check is a
Lipschitz bound \(L_g\) and mesh diameter \(h\) with \(L_gh<\pi/2\).
This is a readout-interpolation criterion, not spatial refinement of the
fixed physical momentum carrier or a general \(SU(3)\) admissibility theorem.

Clutch samples alone do not determine the integer.  Every Hurwitz vertex in
\(2T\) obeys \(q^{12}=1\), so the continuous maps \(q^n\) and
\(q^{n+12}\) have identical vertex values.  Their degrees differ by 12:
the primitive generator of \(H^3(SU(2);\mathbb Z)\) is additive under
group multiplication, so the power map multiplies it by the exponent.
This alias concerns clutch-vertex samples only.  The complete clutched
time-dependent links in the clock-cover theorem retain the continuous
bundle owner and do not identify these histories.


% Source Paper VII-A, original line 7138; inherited proof bodies unchanged.
\subsection{Optional Euclidean Cartan Yang--Mills branch}
\label{p10:app:v199-cartan-yang-mills}
The mathematical transgression, cylinder functional, minimizing kink and
homogeneous fluctuation spectrum for the Cartan family are owned by
Paper~I-C, G11.  Paper~VII-A uses that theorem bank only on the additional
Euclidean pure-Yang--Mills branch specified below; it does not rederive the
underlying Cartan variational problem.

Use the fundamental \(2\oplus1\) embedding and the left representative
\(A_f=f\Theta\), \(\Theta=g^{-1}dg\).  Transport the Paper-I-B orientation
and trace convention
\(-\int_{S^3}\operatorname{tr}\Theta^3/(24\pi^2)=1\), and define
\[
 N_{\rm CS}(A)=\frac1{8\pi^2}\int_{S^3}
 \operatorname{tr}(A\wedge dA+\tfrac23A^3).
\]
This orientation is to be transported consistently if the right
representative or the direction of its parameter is used.

\begin{proposition}[Imported transgression and smooth clutched unit configuration]
\label{p10:prop:v199-cartan-transgression}
The Cartan family satisfies
\[
 N_{\rm CS}(A_f)=3f^2-2f^3,\qquad
 Q[A_f]=\frac1{8\pi^2}\int\operatorname{tr}(F_f\wedge F_f)
       =[3f^2-2f^3]_{\tau_0}^{\tau_1}.
\]
A smooth \(0\)-to-\(1\) profile constant near the two ends glues to a
smooth charge-one connection on a clutched \(S^3\times S^1\) bundle.
\end{proposition}
\begin{proof}
This is Paper~I-C, Proposition~\ref{p1b:prop:g11-transgression}
\cite{Kugelmann2026Paper1B}, after the displayed fundamental embedding,
trace and orientation are identified.  No additional Paper-VII-A derivation is
used.
\end{proof}
A periodic profile in one fixed trivialization has net charge zero.
The static Cartan anchor has \(N_{\rm CS}=1/2\) and \(Q=0\); that relative
Chern--Simons value is not a selected strong-CP angle.  The clutched smooth
configuration is not automatically a Yang--Mills critical point.

On the additional Euclidean cylinder metric
\(ds^2=R^2(d\tau^2+d\Omega_3^2)\), take the classical pure gauge
constituent \(S_{\rm YM}=-g^{-2}\int\operatorname{tr}(F\wedge *F)\),
with \(T_a=-i\sigma_a/2\).  The radius here belongs to this transported
position-metric branch; it is not a hierarchy scale or momentum cutoff.
The Euclidean lift of succession is used, while \(S_E^1\) supplies no
extra spacetime coordinate.

\begin{proposition}[Imported cylinder minimum and homogeneous fluctuations]
\label{p10:prop:v199-yang-mills-kink}
The homogeneous family obeys
\[
 S_{\rm YM}[f]=\frac{12\pi^2}{g^2}\int
 [\dot f^2+4f^2(1-f)^2]\,d\tau,\qquad
 \ddot f=4f(1-f)(1-2f).
\]
On the infinite cylinder with unit boundary charge,
\[
 f_*(\tau)=\frac{1+\tanh(\tau-\tau_0)}2,\quad
 S_{\rm YM}[f_*]=\frac{8\pi^2}{g^2}.
\]
The homogeneous quadratic fluctuation operator is
\[
 \mathcal L=-\partial_\tau^2+4-6\operatorname{sech}^2z,
 \qquad z=\tau-\tau_0,
\]
with simple bound eigenvalues \(0,3\), normalized eigenfunctions
\(\sqrt3\operatorname{sech}^2z/2\) and
\(\sqrt{3/2}\operatorname{sech}z\tanh z\), and absolutely continuous
spectrum \([4,\infty)\).
\end{proposition}
\begin{proof}
Paper~I-C, Theorem~\ref{p1b:thm:g11-functional} and
Proposition~\ref{p1b:prop:g11-fluctuations}
\cite{Kugelmann2026Paper1B}, give the complete variational and spectral
statements under exactly the cylinder, trace and homogeneous-profile data
fixed above.  This imports the mathematical owner result rather than
repeating its proof in Paper~VII-A.
\end{proof}

In proper-length units these operator eigenvalues scale as \(R^{-2}\).
The midpoint \(f=1/2\) instead has homogeneous operator
\(-\partial_\tau^2-2\) and an unstable constant direction.
In the comparison coordinate \(r=e^\tau\), the kink is
\(f=r^2/(r^2+r_0^2)\), displaying the unselected scale/translation modulus.
The common-heat value \(g_3^2=2\pi^2\) gives a classical dimensionless
action four on precisely that matching branch.  It is not a tunnelling
probability, a controlled quantum weight or a populated-sector prediction.
The homogeneous fluctuation operator is not the full gauge-fixed determinant.
The infinite-cylinder kink is not an exactly periodic primitive-circle
solution or a stationary solution of an unspecified full nonlocal action.

% BEGIN M8 PREPARED EUCLIDEAN CHIRAL CALCULATION

% Source Paper VII-A, original line 7232; inherited proof bodies unchanged.
\subsubsection{The complete chiral zero-mode count on the optional kink}
\label{p10:app:m8td-kink-dirac}
This extension uses the same additional infinite Euclidean cylinder and
pure Yang--Mills kink as Proposition~\ref{p10:prop:v199-yang-mills-kink}.
It does not assert a primitive-circle periodic solution or select a
physical transition history. Import the spatial doublet operator from
Theorem~\ref{p3f:thm:m8td-hodge-dirac} with the displayed left gauge and
orientation.

\begin{proposition}[Unique normalizable chiral kink mode]
\label{p10:prop:m8td-kink-zero}
For $z=\tau-\tau_0$, $f_*=(1+\tanh z)/2$, define dimensionless blocks
\[
 W=\partial_\tau+R\mathcal D_{f_*},\qquad
 W^\dagger=-\partial_\tau+R\mathcal D_{f_*}.
\]
On the natural $L^2$ cylinder realization, per doublet,
\begin{equation}
 \ker W=0,\quad\dim\ker W^\dagger=1,\quad\operatorname{ind}W=-1,
 \qquad
 \Psi_0=\frac{\operatorname{sech}^{3/2}z}{\pi^{3/2}R^2}\xi_0.
 \label{p10:eq:m8td-kink-zero}
\end{equation}
Here $\xi_0$ is pointwise unit and the volume measure is
$R^4d\tau\,d\mu_{S_1^3}$.
\end{proposition}
\begin{proof}
The preceding Yang--Mills reduction supplies
$f_*=(1+\tanh(\tau-\tau_0))/2$. Its doublet Dirac family is the
one in Paper~1B, Theorem~\ref{p1b:thm:dirac-spectrum}; hence
Proposition~\ref{p1b:prop:kink} \cite{Kugelmann2026Paper1B}
applies on the same natural $L^2$ cylinder domain. Its adjoint convention
is $W^\dagger=-\partial_\tau+R\mathcal D_{f_*}$ and its volume measure
is exactly $R^4d\tau\,d\mu_{S_1^3}$. It therefore gives the two
kernels, index and normalization in \eqref{p10:eq:m8td-kink-zero}.
\end{proof}

On the singlet put $A=\partial_z+\tfrac32\tanh z$. The partners are
\[
 A^\dagger A=-\partial_z^2+\tfrac94-\tfrac{15}4\operatorname{sech}^2z,
 \qquad AA^\dagger=-\partial_z^2+\tfrac94-\tfrac34\operatorname{sech}^2z.
\]
The first has bound values $0,2$ and the second the single bound value
$2$, with continuum beginning at $9/4$ for each. Profiles are
$\operatorname{sech}^{3/2}z$, $\operatorname{sech}^{1/2}z\tanh z$ and,
for the partner, $\operatorname{sech}^{1/2}z$. Factoring once more at
parameter $1/2$ gives a partner potential everywhere at least its continuum
threshold, excluding extra bound states. Physical squared values include
$R^{-2}$. Reversing the profile exchanges the chiral roles; a gauge
transformation of the same connection does not. The color singlet
spectator has no cylinder zero. These are fermion operators, distinct
from the homogeneous gauge fluctuation operator and from a bosonic RAR
critical boundary Hessian.
% END M8 PREPARED EUCLIDEAN CHIRAL CALCULATION



% Source Paper VII-A, original line 7288; inherited proof bodies unchanged.
\subsection[Supplementary free-adjoint spectral and interaction-selection corollaries]{Supplementary free-adjoint spectral\\and interaction-selection corollaries}
\label{p10:app:v245-free-adjoint-corollaries}
The results in this section are diagnostics of the quadratic massless adjoint
branch.  They are not substitutions for the interacting singlet QCD
Hamiltonian used in the main proof chain.

\begin{proposition}[Exact singlet-projected free-adjoint thermal trace]
\label{p10:prop:v245-singlet-projected-free-gluon-trace}
Let \(x=e^{-\beta\Eper}\).  On the regular bosonic free-adjoint branch, level
\(n\ge2\) contains \(2(n^2-1)\) spatial/helicity copies of the adjoint
representation.  The exact color-singlet thermal trace is
\[
 \boxed{
 Z_{g,0}^{\rm sing}(\beta)
 =\int_{SU(3)}dg\,
 \prod_{n=2}^{\infty}
 \det_{\mathbf8}\!\left(I-x^n\operatorname{Ad}_g\right)^{-2(n^2-1)}.}
\]
At \(g=e\) this reduces to the unprojected colored trace
\[
 Z_{g,0}^{\rm unproj}
 =\prod_{n=2}^{\infty}(1-x^n)^{-16(n^2-1)}.
\]
On an actual finite occupation carrier the exact formula is instead the
finite Haar-projected trace
\[
 \boxed{
 Z_{g,L,N}^{\rm sing}
 =\int_{SU(3)}dg\,
 \operatorname{Tr}_{\mathcal F_{L,N}}
 \left[\Gamma(\operatorname{Ad}_g)e^{-\beta d\Gamma(H_g)}\right].}
\]
Hence the factor sixteen counts unprojected local adjoint helicity modes; it
must not replace the physical singlet projection at finite QCD resolution.
\end{proposition}

\begin{proposition}[Free adjoint spectral zeta]
\label{p10:prop:v245-free-gluon-spectral-zeta}
The unprojected free transverse adjoint spectral zeta is eight copies of the
photon coexact zeta,
\[
 \boxed{\zeta_g^{\rm free}(s)
 =16R_x^{2s}\bigl[\zeta_R(2s-2)-\zeta_R(2s)\bigr].}
\]
Under zeta regularization its corresponding free transverse reference value is
\[
 \boxed{E_{0,g}^{\zeta}
 =\frac{11}{15}\frac{\hbar c_{\rm op}}{R_x}.}
\]
This is an unprojected free spectral invariant, not the interacting singlet
QCD vacuum energy.
\end{proposition}

\begin{proposition}[Triple representation filter for the cubic gluon vertex]
\label{p10:prop:v245-cubic-gluon-selection-filters}
Expand a free interaction-picture connection in temporal characters, coexact
\(S^3\) harmonics and a Cartan--Weyl color basis.  A nonzero cubic coefficient
must pass the following independent necessary filters:
\begin{enumerate}[label=\textnormal{(\roman*)},leftmargin=2.1em]
\item the color weights appearing in
      \(Q(X,[Y,Z])\) sum to zero, including the zero Cartan weights;
\item the spatial representations obey
\[
 \operatorname{Hom}_{SU(2)_L\times SU(2)_R}
 (R_1\otimes R_2\otimes R_3,\mathbf1)\ne0;
\]
\item the primitive-cell/Fourier average of the exact Floquet interaction
      of Theorem~\ref{p10:thm:v247-singlet-floquet-qcd} obeys
\[
 \boxed{\epsilon_1n_1+\epsilon_2n_2+\epsilon_3n_3=0},
 \qquad \epsilon_i=\pm1.
\]
\end{enumerate}
The third condition is precisely the zero Fourier sector of the exact
one-turn Floquet interaction.  It does not say that an off-resonant
instantaneous vertex vanishes; nonzero Fourier sectors remain in the full
time-ordered unitary and can contribute through higher-order/virtual
processes.
\end{proposition}



% Source Paper VII-A, original line 7369; inherited proof bodies unchanged.
\subsection{Finite representation tails and controlled block truncation}
\label{p10:app:v199-heat-tails}
For \(SU(3)\) labels \((p,q)\in\mathbb N^2\),
\[
 d_{pq}=\frac{(p+1)(q+1)(p+q+2)}2,\quad
 C_{pq}=\frac{p^2+q^2+pq+3p+3q}{3}.
\]
Truncate at total degree \(p+q\le L\).  Uniformly on the group,
\[
 |K_t-K_t^{(L)}|\le\sum_{p+q>L}d_{pq}^2e^{-tC_{pq}/2}
 \le\sum_{k=L+1}^\infty a_k(t),\quad
 a_k(t)=\frac{(k+1)(k+2)^6}{64}e^{-tk(k+4)/8}.
\]
Indeed at \(p+q=k\), \(d_{pq}\le(k+2)^3/8\) and
\(C_{pq}\ge k(k+4)/4\), with \(k+1\) representations.  The ratio
\[
 \frac{a_{k+1}}{a_k}=\frac{k+2}{k+1}
 \left(\frac{k+3}{k+2}\right)^6e^{-t(2k+5)/8}
\]
decreases in \(k\).  If its value \(r_L(t)\) at \(k=L+1\) is below one,
the explicit bound is
\begin{equation}
 \epsilon_L(t)\le\frac{a_{L+1}(t)}{1-r_L(t)}.
 \label{p10:eq:v199-representation-tail}
\end{equation}
For the \(j\)-th Casimir moment multiply the bounding sequence by
\([k(k+3)/3]^j\), using the upper Casimir bound and the corresponding
modified ratio.  This controls the stiffness and quartic coefficient too.
At the displayed isotropic time the total-degree-17 heat-kernel bound is
below \(3.96\times10^{-117}\); the root solver and moment evaluation use
the same convergent sums.

For a finite block, \(|K_t^{(L)}|\le K_t(e)\) as well.  Telescoping the
face product and integrating normalized Haar measure therefore yields
\begin{equation}
 \|F_{\mathcal B}-F_{\mathcal B}^{(L)}\|_\infty
 \le\sum_f\epsilon_f\prod_{g\ne f}K_{t_g}(e)=:\delta.
 \label{p10:eq:v199-block-error}
\end{equation}
If \(m_{\mathcal B}>0\) is a lower bound for the full boundary weight and
\(\delta<m_{\mathcal B}\), the approximant is positive and
\[
 \|\log F_{\mathcal B}-\log F_{\mathcal B}^{(L)}\|_\infty
 \le-\log(1-\delta/m_{\mathcal B}).
\]
A lower bound follows from the minima of the positive face kernels; where
\(K_t(e)<2\), the elementary \(K_t(g)\ge2-K_t(e)>0\) suffices for a
face.  Without the error condition partial character sums are not assumed
pointwise positive.  These bounds may grow with the block and deteriorate
as face times tend to zero; they do not establish a uniform RG norm transfer.

% BEGIN V202 HEAT CERTIFICATE AND OPTIONAL VECTORLIKE BRANCH

% Source Paper VII-A, original line 7421; inherited proof bodies unchanged.
\subsubsection{Rational native heat-root and tail certificate}
\label{p10:app:v202-native-heat-certificate}
Put $t_0=14383/2500$ and $\beta=4/\pi^2$. The inequality
$h(t_0)>1/(2\pi^2)$ is equivalent to
\[
 \sum_{(p,q)\ne(0,0)}d_{pq}^2(C_{pq}-\beta)e^{-t_0C_{pq}/2}>\beta.
\]
Every summand is positive. Terms with $1\le p+q\le3$ alone leave a
rationally certified margin greater than $28/10^6$, using an upper bound
for $\beta$ and lower bounds for the exponentials. Strict decrease of
$h$ therefore gives $t_{\rm iso}>t_0$. The packaged standard-library
check uses Machin's arctangent identity and alternating-series bounds for
$\pi$. For $x\ge0$, with $S_n(x)=\sum_{j=0}^nx^j/j!$, it uses
\[
 S_n(x)\le e^x\le S_n(x)+
 \frac{x^{n+1}/(n+1)!}{1-x/(n+2)},\qquad n+2>x.
\]
Apply these exact rational bounds to the $k=18$ term and ratio of
\eqref{p10:eq:v199-representation-tail}. The tail at $t_0$ is below
$4\,10^{-117}$; its decimal upper-bound diagnostic is
$3.970950225\,10^{-117}$. Thus no rounded heat-root output is used as
an exact premise. Increasing the degree controls every fixed finite
block; the bounds are not size-independent RG interaction norms.


% Source Paper VII-A, original line 7445; inherited proof bodies unchanged.
\subsubsection{An optional real massive vectorlike determinant}
\label{p10:app:v202-vectorlike}
This supplementary branch assumes the declared real vectorlike insertion
$D_m=mI+(1-a_Rm/2)D_{\rm ov}$ with $m>0$, $a_R>0$. It does not
select the physical mass or a chiral measure.
\begin{proposition}[Positive determinant on the declared vectorlike branch]
\label{p10:prop:v202-vectorlike-positive}
For that finite operator,
\[
 \sigma_{\min}(D_m)\ge\min(m,2/a_R)>0,\qquad \det D_m>0.
\]
\end{proposition}
\begin{proof}
For this overlap matrix $V=\Gamma_5\epsilon$ is unitary and
$\Gamma_5V\Gamma_5=V^\dagger$. The specified massive insertion is
$D_m=(a_R^{-1}+m/2)I+(a_R^{-1}-m/2)V$. Apply Paper~1B,
Proposition~\ref{p1b:prop:positive-massive-branch}
\cite{Kugelmann2026Paper1B}, with $a=a_R$ and the given $m>0$.
\end{proof}
Constant positive flavor masses give a product of positive determinants.
On this declared branch the determinant may multiply a positive finite
gauge weight in Theorem~\ref{p12:thm:v202-bounded-source-domain}.
Theta/CP phase selection, physical masses, a chiral measure and reflection
positivity are not supplied by this algebraic determinant result.
% END V202 HEAT CERTIFICATE AND OPTIONAL VECTORLIKE BRANCH

% Overlap/Cayley scalar-source block moved into the main chiral chain in v2.164.


% BEGIN G16-MASSIVE-OVERLAP-RESOLVENT v1.67

% Source Paper VII-A, original line 7475; inherited proof bodies unchanged.
\subsubsection{massive-overlap resolvent and interacting-propagator invertibility}
\label{p10:sec:v167-massive-overlap-resolvent}

\begin{theorem}[Uniform accretive bound for the massive overlap resolvent]
\label{p10:thm:v167-overlap-resolvent-accretive}
On the declared vectorlike massive-overlap branch, write
\[
 D_m=aI+bV,\qquad
 a=a_R^{-1}+\frac m2,\qquad
 b=a_R^{-1}-\frac m2,
\]
with $V=\Gamma_5\epsilon_H$ unitary and $m,a_R>0$.  Then
\[
 \boxed{
 \frac{D_m^{-1}+D_m^{-1\dagger}}2
 \succeq
 \frac{1}{\max(m,2/a_R)}I .}
\]
Consequently, for any positive normalized interacting gauge measure on this
finite branch,
\[
 S_L(m):=\langle D_m^{-1}\rangle_L
\]
satisfies
\[
 \boxed{
 \frac{S_L+S_L^\dagger}{2}
 \succeq
 \frac{1}{\max(m,2/a_R)}I,}
\qquad
 \boxed{\|S_L^{-1}\|\le \max(m,2/a_R).}
\]
In particular the full finite interacting propagator is automatically
invertible on the declared positive massive branch.
\end{theorem}
\begin{proof}
For an eigenvalue $e^{i\theta}$ of $V$,
\[
 \operatorname{Re}\frac1{a+b e^{i\theta}}
 =\frac{a+b\cos\theta}{a^2+b^2+2ab\cos\theta}.
\]
Since $a>|b|$, this is minimized at the endpoint selected by the sign of $b$
and has minimum $(a+|b|)^{-1}$.  Because
$a+|b|=\max(m,2/a_R)$, functional calculus gives the first operator
inequality.  Positive normalized averaging preserves the Hermitian-part
inequality.  If $c=1/\max(m,2/a_R)$, then
$\operatorname{Re}\langle x,S_Lx\rangle\ge c\|x\|^2$, hence
$\|S_Lx\|\ge c\|x\|$; finite dimensionality gives invertibility and the
stated inverse bound.
\end{proof}

\begin{corollary}[Positive uniform-mass source interval]
\label{p10:cor:v167-positive-mass-source-interval}
For a real uniform scalar-source displacement $m\mapsto m+J$ with
$|J|<m$, the same branch remains massive and
\[
 \sigma_{\min}(D_{m+J})\ge \min(m+J,2/a_R)>0,
 \qquad \det D_{m+J}>0.
\]
Thus the native uniform mass source has a connected positive finite source
interval around $J=0$.  More general local or momentum-carrying source charts
retain their separately stated positivity/noncriticality hypotheses.
\end{corollary}

\begin{remark}[G16 scope]
\label{p10:rem:v167-g16-scope}
The theorem removes propagator invertibility as an independent finite G16
premise.  It does not evaluate a named continuum-scheme coefficient, prove a
global phase-crossing source chart, or select the physical fermion-family
member.
\end{remark}
% END G16-MASSIVE-OVERLAP-RESOLVENT v1.67

% BEGIN VERIFIED QCD LITERATURE TRANSFERS v2.210

% Source Paper VII-A, original line 7993; inherited proof bodies unchanged.
\subsection{covariant spectral bands on a supplied global Dirac domain}
\label{p10:app:v2186-covariant-spectral-band}

This appendix tests a finite representation alternative that is not consumed
by the native positive proof chain.  It assumes a global compact
self-adjoint elliptic Dirac realization and its gauge-domain covariance;
the local quark/Dirac germ alone does not establish these premises.
The conclusion is finite spectral covariance and an exterior-Fock/CAR lift.
It supplies no physical cutoff, primitive position lattice, extensive
same-action family, compatible pressure or condensate.

\begin{proposition}[Gauge-related finite Dirac bands and their CAR lift]
\label{p10:prop:v2186-covariant-dirac-band-car}
Let $D_A$ be a self-adjoint Dirac operator with compact resolvent on a
supplied compact spinor/color bundle domain.  Suppose each admitted gauge
map is unitary, maps the operator domains, and satisfies
\[
 D_{A^g}=M_gD_AM_g^{-1},\qquad
 M_{g_2g_1}=M_{g_2}M_{g_1}.
\]
For a finite comparison index $r>0$ define
\[
 E_r(A)=\mathbf1_{[0,r^2]}(D_A^2),\quad
 \mathcal K_r(A)=\operatorname{Ran}E_r(A),\quad
 D_r(A)=D_A|_{\mathcal K_r(A)}.
\]
Then $\mathcal K_r(A)$ is finite dimensional and
\[
 E_r(A^g)=M_gE_r(A)M_g^{-1},\qquad
 G_g:=M_g|_{\mathcal K_r(A)}:
 \mathcal K_r(A)\xrightarrow{\sim}\mathcal K_r(A^g)
\]
is unitary with
$D_r(A^g)=G_gD_r(A)G_g^{-1}$.  These maps compose exactly between the
gauge-related fibers.  On
\[
 \mathcal F_-(\mathcal K_r(A))
 =\bigoplus_{n=0}^{\dim\mathcal K_r(A)}\Lambda^n\mathcal K_r(A)
\]
the canonical creation/annihilation maps satisfy exact CAR and
$\bigoplus_n\Lambda^nG_g$ intertwines their gauge-related realizations.
\end{proposition}
\begin{proof}
Compact resolvent gives only finitely many eigenvalues of $D_A$ in
$[-r,r]$, counted with multiplicity.  Unitary functional calculus gives the
projector covariance and hence the restricted unitary intertwiner.
Restriction preserves multiplication of the $M_g$ maps because the target
of each restriction is the next admitted fiber.  Exterior multiplication
and its Hilbert adjoint give
$\{c(f),c^\dagger(h)\}=\langle f,h\rangle I$ and the vanishing equal-type
anticommutators.  Unitarity preserves these inner products, and naturality
of exterior powers gives the claimed intertwiner.
\end{proof}

This construction uses a map between $\mathcal K_r(A)$ and
$\mathcal K_r(A^g)$, rather than the compressed multiplier
$E_r(A)M_gE_r(A)$ on one fixed band.  It therefore is outside the
fixed-band premise of
Proposition~\ref{p10:prop:v2180-compressed-local-gauge-no-go}.
It does not identify either construction with the actual QCD regulator.

\begin{remark}[Source charts, physical license and limiting scope]
\label{p10:rem:v2186-band-source-scope}
If a supplied source family $L(J)=D_{A,J}^2$ has a differentiable
resolvent and an isolated retained spectral cluster on a source chart, its
Riesz projector has constant finite rank there.  For a counterclockwise
contour $\mathcal C$ enclosing that cluster,
\[
 E(J)=\frac1{2\pi i}\oint_{\mathcal C}(z-L(J))^{-1}\,dz,
\]
and differentiation retains the variation of the projector as well as that
of the restricted operator.  It is not legitimate to freeze a
source-dependent $E(J)$ when differentiating the determinant/source word.
An appropriate unitary trivialization and the complete same-owner
source/sea prescription are required for comparison with a fixed finite
source chart.

The index $r$ here is a mathematical comparison index.  Assigning it a
physical sharp-frontier role or identifying it with $\Lambda_F$ requires an
independent licensed theorem.  Likewise the inherited quark field/mode and
statistics assignment must be verified before the mathematical CAR lift is
used as a physical descendant of
Theorem~\ref{p4:thm:finite-second-quantization}.  A large-extent family
requires its global domains, density/symbol and source correspondence;
pressure convergence and nonzero ordered response are further certificates.
None of those physical or analytic obligations follows from the proposition.
\end{remark}

% END INLINE SOURCE: Paper7A-CovariantSpectralBand-Appendix-v1_0.tex

% END XFE-EXTENSIVE-FERMION-DIAGNOSTICS v2.180

% BEGIN NATIVE CONSTRUCTOR AND TOPOLOGY DIAGNOSTICS v2.189

% BEGIN INLINE SOURCE: Paper7A-NativeClockSourceComponent-Appendix-v1_0.tex
% Framework v2.189: diagnostics outside the positive clock-source chain.


% BEGIN FRAMEWORK v2.323: EXACT PETERSEN HADRON SOURCE-COMPRESSION CERTIFICATE
\section{Petersen mirror compression, coherence, and color-center certificates}
\label{p13:sec:v323-petersen-proof}
The following is an exact finite check on the \emph{specified}
Petersen--Hopf geometry.  It does not compute a gauge-invariant
hadronic pole.  Take the ten unoriented axis vectors $u_{ij}$ of
Paper~I-A, the fifteen oriented disjoint-pair edges $e=(v,w)$,
and their signed normalized cross-product columns
\[
 k_e=\frac32\operatorname{sgn}(u_v\cdot u_w)(u_v\times u_w).
\]
Form $K=[k_e]$, its Galois conjugate $K^\sigma$ under
$\sqrt5\mapsto-\sqrt5$, and the two cycle-space projectors
\[
 \Pi_3=\frac15K^{\mathsf T}K,\quad
 \Pi_{3'}=\frac15(K^\sigma)^{\mathsf T}K^\sigma,\quad
 D=\Pi_3-\Pi_{3'}.
\]
The five core projectors $P_m$ select the three Petersen edges whose
four constituent labels omit $m$; they are mutually orthogonal and
sum to the identity on $\mathbb R^{15}$.

\begin{proposition}[Exact support-compression characteristic polynomials]
\label{p13:prop:v323-petersen-charpolys}
For $S\subset\{1,2,3,4,5\}$, put
$D_S=P_SDP_S|_{P_SW}$ and $k=|S|$.  Its characteristic polynomial
is independent of the particular $k$-subset and equals
\[
\begin{array}{c|l}
 k&\det(tI-D_S)\\\hline
 0&1\\
 1&t^3\\
 2&t^2(t^2-3/20)^2\\
 3&t^3(t^2-1/5)(t^2-7/20)^2\\
 4&t^6(t^2-3/5)^3\\
 5&t^9(t^2-1)^3.
\end{array}
\]
For every $m$, $P_mDP_m=0$; for $m\ne n$ the inter-core block
$P_mDP_n$ has singular values
$\sqrt{3/20},\sqrt{3/20},0$.  In particular
$D^3=D$ and the global difference has ranks $3+3$ and zero trace.
\end{proposition}
\begin{proof}
Insert the coordinates of the ten Hopf axes into $K$ and
$K^\sigma$, then form the displayed rank-three projectors.  The
single-core diagonal blocks vanish because each of $\Pi_3$ and
$\Pi_{3'}$ compresses to $P_m/5$.  The inter-core $3\times3$ Gram
matrices have eigenvalues $3/20,3/20,0$.  Direct exact determinants
of one representative $k$-subspace give the table.  Every $k$-subset
lies in one orbit of the signed $A_5$ action, and $D$ commutes with
that action; this completes the proof for all 32 subsets.  The
accompanying executable $\mathbb Q(\sqrt5)$ certificate separately
checks all 60 even permutations and all ten inter-core blocks.
\end{proof}

\begin{corollary}[A sharp coherence bound, not a hadron spectrum]
\label{p13:cor:v323-petersen-coherence-bound}
For any normalized state $\rho$ with support in $P_SW$,
\[
 |\operatorname{Tr}(\rho D)|\leq\|D_S\|_{\rm op}.
\]
The sharp bounds are $0$, $\sqrt{15}/10$,
$\sqrt{35}/10$, $\sqrt{15}/5$, and $1$ for $k=1,2,3,4,5$.
Equality is attainable by a normalized eigenvector of the appropriate
compressed operator.  If $\rho=\sum_mP_m\rho P_m$, the left hand
side is exactly zero, regardless of the number of populated cores.
\end{corollary}
\begin{proof}
The variational spectral bound is sharp for an extremal eigenvector.
Block-incoherent states give zero by $P_mDP_m=0$.
\end{proof}

\begin{remark}[Same-action interpretation firewall]
This finite certificate does not identify a core as a quark, the
$k=2,3$ supports as mesons or baryons, the edge operators as physical
$SU(3)$ links, or the numerical eigenvalues as energies.  A source-owned
color-singlet map, a nonzero admissible inter-core matrix element, and
the full compiled Hamiltonian/Schur pole are required before making a
hadronic claim.  The central $\mathbb Z_2$ sign is not the
$\mathbb Z_3$ center of $SU(3)_c$.
\end{remark}
% END FRAMEWORK v2.323: EXACT PETERSEN HADRON SOURCE-COMPRESSION CERTIFICATE


\section{Chiral, topology, source diagnostics and baryon comparisons}
\label{p13:sec:v237-d3}
% Source Paper VII-A, original line 6620; inherited proof bodies unchanged.
% Former appendix marker: technical material is now placed in VII-D.


% BEGIN INLINE SOURCE: Paper7A-ChiralDiagnostics-Appendix-v1_0.tex
% Framework v2.164: supplementary calculations, with native and comparison domains explicit.

% Source Paper VII-A, original line 6625; inherited proof bodies unchanged.
\subsection{Chiral geometry, trial coherence and finite execution diagnostics}
\label{p10:app:chiral-diagnostics}
These calculations support or delimit the main source certificates.
The free and trial branches are comparisons on the declared momentum
carrier, and the anomaly branch requires additional continuum/infrared
hypotheses. None evaluates the actual interacting ordered condensate.


% Source Paper VII-A, original line 6632; inherited proof bodies unchanged.
\subsubsection{Protected finite rank, zero sets and the clock}
\label{p10:app:protected-chiral-rank}
If a geometric operator has a protected rank $r$ and a refinement-uniform
complementary gap $g>0$, then for $0<m<g$ its near-zero count is at most
$r$. Dividing by $\mathcal V_L$ and taking extent first makes this
contribution vanish when $r$ is fixed and $\mathcal V_L\to\infty$.
This does not bound the rank or density of the full interacting QCD
operator, which is a different consumed object.

For a multiplication-symbol comparison on the unit $S^3_{\rm mom}$,
Haar probability gives a coordinate $q_3$ density
$2\sqrt{1-q_3^2}/\pi$. The symbol $K_1=\kappa q_3\tau_1$ therefore
has a nonzero paired pushforward density at zero and, with trace over
both eigenvalues, scalar response tending to $4/\kappa$.
The chirally permitted second coefficient
$K_2=\kappa(q_3\tau_1+q_2\tau_2)$ instead has near-zero
density proportional to $|\lambda|$ and vanishing response.
A separately stipulated continuous-clock completion can turn the
one-coefficient comparison into a response of order
$m\log(\kappa/m)$, still vanishing while its susceptibility diverges.
These comparisons change a declared symbol/clock problem; they do not
identify it with the complete QCD generator or its trace weight.

The indefinite Dirac complement admits both signs of its Schur shift:
on $\operatorname{diag}(0,\sigma g)$ with off-diagonal perturbation $v$,
\[
 \lambda_{\rm low}
 =\frac{\sigma}{2}\bigl(g-\sqrt{g^2+4v^2}\bigr)
 =-\sigma v^2/g+O(v^4/g^3),\quad \sigma=\pm1.
\]
This differs from the positive Hamiltonian pivot below its complement
in Theorem~\ref{p10:thm:complete-schur-scalar-source}.
For a heat trace with $Z'(0)=0$,
$(-\log Z)''(0)=-Z''(0)/Z(0)$: the curvature sign reverses.
Neither isolated functional is the complete physical effective action.
The elementary odd scalar and the composite chiral bilinear retain their
distinct owners. A local/Wilsonian negative Hessian is not an assertion
about the exact convex Legendre potential.

At a symmetric finite source with $\mathcal C_L(0)=0$, a uniform bound
$|\partial_h\mathcal C_L(h)|\le M$ on one common neighborhood gives
$|\mathcal C_L(h)|\le Mh$, hence zero ordered response.
Failure of that regularity is necessary for a positive limit, but a
susceptibility diverging only logarithmically can still have vanishing
area in the shrinking source interval.


% Source Paper VII-A, original line 6678; inherited proof bodies unchanged.
\subsubsection{Conditional CAR energy-moment simplification}
\label{p10:app:clifford-chiral-moment}
For $H_f=d\Gamma(h_f)$ and a fixed numerical, or gauge-central,
one-particle source $O=d\Gamma(S)$, the CAR algebra gives
\[
 [O,[H_f,O]]=d\Gamma([S,[h_f,S]]).
\]
If the actual canonical source has $S=\beta=\beta^\dagger$,
$\beta^2=I$, and $\{\beta,h_f\}=0$, this becomes $-4H_f$.
Pure gauge terms commute with $O$. A retained color-current term
$\sum_{a,b}J_aG_{ab}J_b$, $J_a=d\Gamma(T_a)$, commutes if
$[S,T_a]=0$ and $G_{ab}$ commutes with $O$.
On that complete branch $M_1=-2\langle H_f^{\rm raw}\rangle$.
Restore any normal-order constant in the raw energy; the commutator
itself is unchanged by a scalar shift. Noncentral source coefficients
or a moving frame require the full commutator and contacts.
If additionally $\|h_f\|\le\Lambda$ uniformly and
$d_f/\mathcal N$ is bounded, then $M_1\le2d_f\Lambda=O(\mathcal N)$.
Compact momentum alone does not bound all gauge amplitudes or prove
that mode/trace estimate. The quadratic variance lower bound remains
a calculation of the actual state.


% Source Paper VII-A, original line 6700; inherited proof bodies unchanged.
\subsubsection{Fixed-frontier free Clifford comparison}
\label{p10:app:free-chiral-comparison}
On the explicitly supplied free Clifford branch, extending the owned
polynomial momentum trace comparison to the continuous spectral
functions below requires uniformly bounded full operator norm and
normalized identity trace. Uniform polynomial approximation then
licenses the extension; it is not an identification with the interacting
QCD Hamiltonian.
With positive free energy $E=cp$, occupied negative projector $P_-$
and source $\beta$, finite CAR calculus gives
$M_0=\operatorname{Tr}P_-$ and $M_1=2\sum_{E>0}E$.
The fixed-frontier per-flavor scalar response is
\[
 C_{\rm free}(h)=\frac{N_ch}{2\pi^2\hbar^3c^3}
 \left[cp_F\sqrt{c^2p_F^2+h^2}
       -h^2\operatorname{arsinh}(cp_F/h)\right]\longrightarrow0.
\]
Radial integration also gives
\[
 \chi_{\rm free}(0)=\frac{N_cp_F^2}{2\pi^2\hbar^3c},\quad
 \frac{M_0}{N_fV}=\frac{N_cp_F^3}{3\pi^2\hbar^3},\quad
 \frac{M_1}{N_fV}=\frac{N_ccp_F^4}{2\pi^2\hbar^3},\quad
 \Delta=\frac32cp_F.
\]
Here $h$ is the energy-valued source probe, not an empirical mass
assignment. The Hamiltonian spatial normalization and a Euclidean
clock-volume normalization agree only through their licensed transfer.


% Source Paper VII-A, original line 6728; inherited proof bodies unchanged.
\subsubsection{Globally shared flavor-orbit trial}
\label{p10:app:shared-flavor-chiral-trial}
On each comparison pair channel $k$, let
\[
 \gamma_k(U)=\frac12
 \begin{pmatrix}(1-\xi_k)I_3&\eta_kU\\
                 \eta_kU^\dagger&(1+\xi_k)I_3\end{pmatrix},
 \quad U\in SU(3),\quad \xi_k^2+\eta_k^2=1.
\]
Tensor with $I_{N_c}$ and use the off-diagonal unit scalar source.
This is a positive projector covariance. Complete color multiplets give
a globally color-singlet Slater trial. The mean scalar source is
$N_c\sum_k\eta_k\operatorname{ReTr}U$.
Haar restoration uses
$\int\operatorname{ReTr}U=0$ and
$\int(\operatorname{ReTr}U)^2=1/2$, since
$\int|\operatorname{Tr}U|^2=1$ and $\int(\operatorname{Tr}U)^2=0$.
Its exact variance is
\[
 M_0=N_c\sum_k(3-\tfrac32\eta_k^2)
             +\frac{N_c^2}{2}\left(\sum_k\eta_k\right)^2.
\]
The first term equals the Wick term of the averaged covariance.
Thus $\Xi_{\rm shared}=N_c^2(\sum_k\eta_k)^2/2$.
Independent orientation restoration gives instead
$\Xi_{\rm independent}=N_c^2\sum_k\eta_k^2/2$.
The common orientation illustrates a quadratic mixed-channel return.
It supplies neither the actual invariant ground nor its amplitude.

For the comparison pair energy
$h_{0,k}=\operatorname{diag}(E_kI_3,-E_kI_3)$, $E_k>0$,
the cost is $3N_c\sum_kE_k(1-\xi_k)$.
Minimizing cost minus the aligned source gives
$\eta_k=h/\sqrt{E_k^2+h^2}$ and
$\xi_k=E_k/\sqrt{E_k^2+h^2}$.
Without a zero-energy atom its coherence density vanishes as $h\downarrow0$.
An assigned fixed extensive coherence generally costs extensive free
energy. Therefore the trial coefficient alone is no energy certificate.

With $\operatorname{Tr}T^aT^b=\delta^{ab}/2$, the complete
color exchange $X_c=\sum_aT^a\otimes(-T^{a*})$ acts as
$-4/3$ on $\mathbf1$ and $1/6$ on $\mathbf8$ in
$\mathbf3\otimes\bar{\mathbf3}$. This follows from
$C_{\rm total}=8I/3+2X_c$, with singlet Casimir zero and octet Casimir
three. Separate $\operatorname{Tr}T^a=0$ contractions erase this
return. The color identity supplies no new four-fermion coefficient,
propagator normalization or critical coupling.


% Source Paper VII-A, original line 6776; inherited proof bodies unchanged.
\subsubsection{Conditional three-flavor anomaly comparison}
\label{p10:app:conditional-chiral-anomaly}
On an additional confined four-dimensional anomaly-matching branch
with ordinary massless color-singlet composite fermions and well-defined
integer anomaly sums, integer baryon charge and
$t_L+t_R=0\pmod3$ imply
\[
 1=\sum_r\ell_rB_r\,\mathcal I(r_L)d(r_R),\qquad
 \mathcal I(\mathbf3)=1.
\]
The ultraviolet coefficient is $N_cB_{\rm quark}=3(1/3)=1$.
In a zero-triality irrep, $Y=\operatorname{diag}(1,1,-2)$ has
weights $3n$ with $\sum n=0$, so
$\mathcal I(r)=\operatorname{Tr}_rY^2/6
 =3\sum n^2/2\in3\mathbb Z$.
For nonzero triality,
$d(a,b)=(a+1)(b+1)(a+b+2)/2$ is divisible by three:
the distinct residues of $a+1,b+1$ force a zero factor.
Together with integral $SU(3)$ indices, every infrared term is divisible
by three, contradicting the unit ultraviolet coefficient.
This is the prime-factor comparison of
\cite{CiambrielloEtAl2024ChiralStrategy}.

The finite singlet projector does not establish that continuum,
confinement or infrared spectrum. Chiral symmetry breaking under those
antecedents does not itself compute the native bilinear response.
The discrete anomaly in \cite{Tanizaki2018AnomalyConstraint}
excludes the naive Stern phase on its stated content; additional
massless/topological content remains a separate possibility.
No universal numerical lower bound on $\mathcal C_L$ follows here.


% Source Paper VII-A, original line 6807; inherited proof bodies unchanged.
\subsubsection{Finite numerical energy/source certificate}
\label{p10:app:numerical-chiral-certificate}
\begin{proposition}[Error-carrying finite response lower bound]
\label{p10:prop:numerical-chiral-certificate}
Let $\widetilde H_L,\widetilde O_L$ approximate the complete canonical
operators with certified norm errors $\epsilon_H,\epsilon_O$.
For $\ell_L\le E_{0,L}$ and a normalized physical vector $v$, put
\[
 q_L^-=\frac{\langle v,\widetilde O_Lv\rangle-\epsilon_O}{\mathcal N_L},
 \qquad
 \delta_L^+=
 \frac{\langle v,\widetilde H_Lv\rangle+\epsilon_H-\ell_L}{\mathcal N_L}.
\]
Then $\mathcal C_L(h)\ge q_L^- -\delta_L^+/h$.
\end{proposition}
\begin{proof}
The operator errors bound the exact source below by $q_L^-$
and exact excess energy above by $\delta_L^+$. Apply
Theorem~\ref{p4:thm:finite-energy-source-certificate}.
\end{proof}
The companion interacting program supplies outward-rounded Gershgorin
ground-energy bounds and exact rational trial quadratic forms for the
represented binary floats. A finite eigenvector residual alone is not
a ground-energy lower-bound certificate. Identification with a physical
packet still consumes its provenance, positive kinetic normalization and
operator error hypotheses; uniform compatible readout limits remain open.

Four historically named comparison programs are intended to check
geometric Schur/source identities, source moments, full CAR/Haar/color
contractions and interacting Schur/energy formulas. Their executable
bytes and specific input/output provenance are not present under the named
paths in the four currently supplied framework archives. The programs are:
\begin{itemize}
\item \path{Paper00-MomentumCartanChiralDensityChecks-v1_0.py}
\item \path{Paper00-ChiralSourceMomentChecks-v1_0.py}
\item \path{Paper00-ChiralFourPointChecks-v1_0.py}
\item \path{Paper00-InteractingChiralCertificateChecks-v1_0.py}
\end{itemize}
Their finite numerical outputs are not replay-certified by this release;
their physical condensate status also remains open. The accompanying v1.5
exploration and v1.3 reproduction reports retain historical source
provenance and are not substitutes for the absent program/input/output bytes.
The finite low-energy-state route has related mathematical precedents
in \cite{KomaTasaki1994SymmetryBreaking,Tasaki2018LongRangeOrder};
their lattice/locality hypotheses do not supply the missing native
moment bound or physical QCD sequence.

% END INLINE SOURCE: Paper7A-ChiralDiagnostics-Appendix-v1_0.tex



% BEGIN INLINE SOURCE: Paper7A-ReadoutDomainComparisons-Appendix-v1_0.tex
% Framework v2.187: supplementary domain tests, not physical constructors.

% Source Paper VII-A, original line 8208; inherited proof bodies unchanged.
\subsection{Spatial quark-constructor comparison and source diagnostics}
\label{p10:app:v2189-native-spatial-quark-comparison}

The main construction is on the already licensed clocked quark-action/CAR
branch. Its proof does not invent a fermion/statistics assignment from a
spinor shell. The geometric bundle, reciprocal spin transport, global
compact Dirac domain and orbital restriction are positive native inputs.
The physical fermionic branch retains the named domain of V-A; the one-body
mode assignment and exterior algebra are inherited through IV. A particular
cofinal physical history is not selected. Each admitted finite frontier has
the displayed constructor, and the cofinal reciprocal geometric comparison
is a family of those words where that comparison is claimed.


% Source Paper VII-A, original line 8221; inherited proof bodies unchanged.
\subsubsection{Orbital completion and exact operator type}
The component-orbital band is not a new $|D|$ spectral cutoff. For one spatial
rank-two spinor it has
\[
 \mathcal K_m^{\rm orb}
 =\bigoplus_{k=0}^m\mathcal H_{k,+}
 \oplus\bigoplus_{k=0}^{m-1}\mathcal H_{k,-},\qquad
 \dim\mathcal K_m^{\rm orb}=2D_m.
\]
The symmetric complete Dirac-shell comparison through $k=m$ has dimension
$2(m+1)(m+2)(m+3)/3$. Its difference is the negative boundary shell,
of rank $(m+1)(m+2)$. Both have the same leading density, but this equality
does not identify their determinants, source contacts, sea allocations or
anomalies. A physical use of the symmetric completion needs its actual
restriction word. The positive construction instead uses the component
band already printed in III-F's free metric Dirac specialization.

The rank-$72$ quark fibre is the physical $P_+$ factorization, with rank-$18$
left and right internal quark factors and rank-$2$ external Weyl factors.
The analytic quark heat carrier has rank $144$ because Euclidean first-order
Clifford action exchanges $P_\pm$. Equal-time clocked Clifford action is
even and preserves $P_+$; adding the analytic partner as independent physical
CAR modes would double the physical count without an occurrence theorem.
The charged source remains
$R_s^{\rm ch}=\sum_rm_{s,r}^{\rm fin}\iota_{s,r}$ with its actual incidences.
No diagonal source, source singular values, physical pole or rest mass is
substituted for this packet.


% Source Paper VII-A, original line 8249; inherited proof bodies unchanged.
\subsubsection{Gauge covariance retains the field embedding}
The construction gives exact covariance of the word with its actual
embedding: $(a,\Phi,J_m)\mapsto(a^g,\Phi^g,M_gJ_m)$. It does not prove that
$P_mM_gP_m$ is a representation on one fixed component band, nor furnish
an $a$-alone canonical cutoff selection. The exact native quaternion
$SU(3)$ map $g(q)=\operatorname{diag}(q_0I-iq_a\sigma_a,1)$ provides an
independent test: on the primitive complete component band, exact Haar
integration gives
\[
 \operatorname{Tr}\{I-(P_1M_gP_1)^\dagger(P_1M_gP_1)\}=6.
\]
Transporting the actual embedding instead preserves its Gram matrix and
its $q_0$ source matrix exactly. Independently reassembling
$a_a^g=-(L_ag)g^\dagger$ and differentiating $gJ_1$ gives the same
retained covariant derivatives with exact rational zero residual. Thus the
positive covariance proof consumes a real transported fibre, not a false
fixed-band multiplier identity.


% Source Paper VII-A, original line 8267; inherited proof bodies unchanged.
\subsubsection{Actual QCD matching is a consumer of this constructor}
The native global field-restriction word now has explicit modes, dual ports,
CAR maps, action matrices, source jets and density. It need not be a primitive
position lattice. Its match to a particular VII-A finite or ordered readout
must preserve the actual field, gauge/constraint reduction, central/sea
allocation, normalization and clock. Finite Hurwitz point sampling is not
exact $L^2$ Galerkin compression, as
Theorem~\ref{p10:thm:hurwitz-shell-dirac-compression} already states.
The finite retained edge holonomies do not determine a global smooth
connection compression: a smooth variation supported away from all retained
edges leaves those holonomies fixed but can change $P c(a)P$.
This is a comparison test, not a new global geometry or generic CAR debt.

The normalized bilinear dictionary must likewise be transported before
equating its derivative to the ordered common scalar source. On the overlap
branch the scalar insertion is
$R_{\rm ov}=(I-\Gamma_5\operatorname{sign}H)/2$; it is not automatically
the local Hamiltonian mass insertion or a raw identity. The latter includes
its clock-hermitizing, field/dual, kinetic and sea conventions. The stationary
VII-A condensate normalizes by $N_fV$ on its stated $N_f=3$ branch and
consumes a matched clock/ground-state passage. The full generated quark
carrier includes all charged quark roles; no unlicensed light-flavor spectral
projection is introduced merely to turn the rank-$72$ count into that
low-energy branch.


% Source Paper VII-A, original line 8292; inherited proof bodies unchanged.
\subsubsection{Source bounds and finite geometric execution}
For a fixed bounded numerical Hermitian source symbol,
$\|S_m\|\le\|S\|_\infty$ and
$\|d\Gamma(S_m)\|\le\operatorname{Tr}|S_m|
 \le72D_m\|S\|_\infty$.
Together with the proved density this supplies the one-body source/density
antecedents of the existing fixed-source CAR variance bounds. Noncentral
quantum source coefficients, singlet reduction, actual central derivatives
and full source contacts retain their own complete word. The displayed bound
does not silently bound an uncomputed central source derivative or prove a
quadratic-volume connected return.

The first frontier has $D_1=5$ and $360$ physical quark component modes.
In the unit Haar pairing its orbital basis is
$1,2q_0,2q_1,2q_2,2q_3$. The native $q_0$ multiplication source has only
the two entries $Q_{01}=Q_{10}=1/2$, hence
$\operatorname{Tr}Q^2=1/2$. For the actual external axial grading on
$\mathcal Q_+$ this gives the finite geometric source Gram
$\operatorname{Tr}S^2=36$. Its canonical CAR tracial comparison variance
is $9$. This is an exact closed trace/readout of the newly exhibited source
word, not the actual QCD state, a physical susceptibility value or a
condensate. The same exact execution checks the round component Dirac and
the nonflat midpoint Born doublet block. Its protected geometric zero is
retained as a geometric zero; the actual charged source and other coupled
sectors are not erased when interpreting the physical Hamiltonian.

A faithful finite Gibbs word has strictly positive Kubo source curvature
when its actual source is noncentral on that consumed support. This criterion
does not by itself prove noncentrality after the full local Gauss/singlet
reduction or any thermodynamic lower bound. In particular, extensive mode
count and bounded source norm do not establish the action-weighted shrinking
chiral window, the variational low-energy source states, a nonzero finite
ordered condensate, Goldstone residue or GMOR theorem hypotheses.

% END INLINE SOURCE: Paper7A-NativeSpatialQuarkConstructor-Appendix-v1_0.tex


% BEGIN INLINE SOURCE: Paper7A-NativeWeightedTopology-Appendix-v1_0.tex
% Framework v2.189: supplementary topology-weight diagnostics.

% Source Paper VII-A, original line 8331; inherited proof bodies unchanged.
\subsection{Finite topological weight normalization and source diagnostics}
\label{p10:app:v2189-topological-weight-diagnostics}

The main positive chain proves the Haar floor, native link and face counts,
complete vectorlike determinant bounds, zero-index witness and normalized
source conclusions. This appendix preserves their normalization checks,
the detailed sign estimate and the excluded weight/domain cases. It adds
no mass member, action selector, physical transfer state or limit hypothesis.


% Source Paper VII-A, original line 8340; inherited proof bodies unchanged.
\subsubsection{Normalized torus and link bookkeeping}
For $SU(3)$, with $\theta_3=-\theta_1-\theta_2$, the normalized Weyl density
of a class function is
\[
 \frac1{6(2\pi)^2}
 \prod_{i<j}|e^{i\theta_i}-e^{i\theta_j}|^2\,d\theta_1d\theta_2.
\]
Differentiating conjugation on a complex root plane gives multiplication by
$e^{i(\theta_i-\theta_j)}-1$; its real determinant is the squared chord.
Expanding the Vandermonde and its conjugate yields $36$ permutation pairs.
Only the six identical pairs have zero two-angle Fourier character: the
three exponent differences sum to zero, and torus restriction makes the
character trivial only when those differences coincide, hence all vanish.
The normalized torus integral is six. This verifies the prefactor without
an extra $\sqrt3$ coordinate factor.

Lemma~\ref{p10:lem:next187-su3-haar-ball} gives, for
$0<\epsilon\le1/256$,
\[
 q(\epsilon)=\frac{1225}{3}\left(\frac7{22}\right)^8
                 \left(\frac\epsilon4\right)^8.
\]
Its rectangle has angle-difference minima $7\delta/2$, $5\delta$, and
$\delta$, and area $\delta^2/2$, where $\delta=\epsilon/4$.
The squared Vandermonde floor is $19600\delta^6/\pi^6$.
All eigenvalue chords lie strictly inside the operator ball even on the
rectangle's boundary. The strict rational comparison uses
\[
 \frac{22}{7}-\pi
 =\int_0^1\frac{x^4(1-x)^4}{1+x^2}\,dx>0.
\]
Translation preserves operator distance and Haar mass. Product Haar on
$360$ independent links therefore gives full-ball mass greater than
$q(\epsilon)^{360}$; no independence of face holonomies is asserted.
At the two radii used in the main chain,
\[
 q_*:=q(1/256)
  =\frac{7061881225}{199022581174732616540136492402475008},
 \qquad q(1/4096)=2^{-32}q_*.
\]
With $q_*=n/d$, the decimal floor in
Theorem~\ref{p10:thm:next187-native-heat-interval} is precisely the integer
inequality
\[
 4\,7^{1152}n^{720}10^{18796}>18^{1152}d^{720}.
\]
The integer certificate also verifies
$4(7/18)^{1152}q_*^{720}<10^{-18795}$; neither comparison uses a rounded
logarithm as a premise.

The Hurwitz graph's $96$ unoriented edges are independent link variables;
the $192$ directed edges include their adjoints. The cyclic three-layer
product has $288$ spatial and $72$ temporal independent links. It has
$288$ spatial triangles and $288$ temporal ribbons, including the $96$
closing ribbons. Closing holonomies are identity only on the restricted
Cartan center of Proposition~\ref{p10:prop:v199-restricted-heat-dynamics}.
They remain heat factors under arbitrary all-link perturbations. The
temporal one-complex supplies no additional filled temporal triangle.


% Source Paper VII-A, original line 8399; inherited proof bodies unchanged.
\subsubsection{Exact heat bounds and characteristic modulus}
The degree-$17$ relative heat-tail certificate for $N\le96$ is not a
certificate for the full $576$-face three-layer product. The main topology
proof uses the exact pointwise $a,b_H$ heat bounds, with every closing
ribbon retained. It does not extrapolate a truncation error beyond its
printed face count.

Positive unit-sector probabilities imply a real return to one precisely
at $2\pi\mathbb Z$ and least positive period $2\pi$. They do not imply
$|\phi(\theta)|<1$ everywhere else. For example, the diagnostic law with
probability $1/2$ at each of $+1$ and $-1$ has $\phi(\pi)=-1$.
This is a counterexample to the stronger modulus inference, not a proposed
native ensemble. The proved real-return inequality needs no zero-sector mass.
The exact massless law in
Theorem~\ref{p10:thm:v2189-massless-topological-endpoint} is instead
supported on zero, so its characteristic function is identically one.


% Source Paper VII-A, original line 8416; inherited proof bodies unchanged.
\subsubsection{The actual zero-index witness and sign-continuity constants}
The center in Theorem~\ref{p10:thm:v2189-zero-index-overlap-witness}
uses actual $SU(3)$ links: spatial and interior temporal links are identity,
and the closing gauge transporter is $\zeta I_3$, with
$\zeta=e^{2\pi i/3}$ and $\det(\zeta I_3)=1$.
The retained fermion spin sign is a separate factor; it is not selected
anew and is not a gauge link or physical Fourier source.

In each temporal Fourier block, the exact stencil of
Theorem~\ref{p1b:thm:overlap} has temporal derivative $i\sin\omega$ and
temporal Wilson term $1-\cos\omega$. At $\rho=1$ the Hermitian part is
$W_s-\cos\omega I$. The center is time independent and its temporal
transporter central, so temporal diagonalization is exact on the full
Hurwitz spatial stencil. The proof does not assume $[D_s,W_s]=0$.
With $A=\Gamma_4(D_s+W_s-\cos\omega I)$ and $s=\sin\omega$,
\[
 K=\Gamma_4(A+i sI),\qquad K^\dagger K=A^2+s^2I,
 \qquad A^\dagger=A.
\]
The imaginary part of $(A+i sI)(A^2+s^2I)^{-1/2}$ is definite for
$s\ne0$. It cannot equal the real expectation of $-\Gamma_4$ on a
putative $V=-1$ vector. This excludes exact overlap zeros without a
proxy lattice or a commuting spatial-Wilson assumption.

The source's unit-stencil diagonal norm is $|4-\rho|=3$, each of eight
spatial blocks has norm $3/4$, and each of two temporal blocks has norm
one. Thus $\|H\|\le11$. For both retained spin signs,
$\min|\sin\omega|=\sin(\pi/9)>2/9$, and the main argument gives
$\sigma_{\min}(D_{\rm ov}/(2/a_R))>1/99$ at the center.
The analytic ball constants are
\[
 \epsilon_0=1/4096,\quad b=8\epsilon_0=1/512,\quad
 g=2/9-b=1015/4608,
\]
\[
 \frac{b}{2g}=\frac9{2030},\qquad
 \eta=\frac1{99}-\frac9{2030}=\frac{1139}{200970}>0.
\]
The block row and column perturbation estimate is independent of the
center's index, so it applies here without importing the clock-center's
index-$k$ premise.

For completeness, the source sign estimate follows directly from its
convergent difference-resolvent identity. If Hermitian $H,K$ have gaps
at least $g>0$,
\[
 \operatorname{sign}H-\operatorname{sign}K
 =\frac1\pi\int_{\mathbb R}
      [(H-itI)^{-1}-(K-itI)^{-1}]\,dt.
\]
The resolvent identity bounds the integrand by
$\|H-K\|/(g^2+t^2)$, whose norm integral is $\pi\|H-K\|/g$.
Therefore the normalized $D_{\rm ov}/(2/a_R)$ difference is bounded by
$\|H-K\|/(2g)$. This reproduces
Corollary~\ref{p10:cor:m8td-overlap-continuity} and requires no
straight-path sign derivative assumption. The zero-index center and
neighborhood proof has no IEEE residual premise. The separate
nonzero-index centers are now also certified exactly by
Theorem~\ref{p10:thm:v256-exact-native-inertia}.


% Source Paper VII-A, original line 8475; inherited proof bodies unchanged.
\subsubsection{Mass normalization, finite source jets and domain guards}
For $x=a_Rm/2$ and $c=2/a_R$,
\[
 D_m/c=(I+V)/2+x(I-V)/2.
\]
At $x=1$ this is exactly the identity, so its determinant is constant and
its normalized gauge law is pure heat. Any positive gauge-independent
normalization factor per flavor likewise cancels. The residual $x$
dependence does not cancel by changing the overlap coordinate while holding
physical $m$ fixed. A positive constant flavor matrix acting in an
independent flavor factor and commuting with $V$ reduces by diagonalization
to the scalar theorem; a noncommuting, spatially varying, chiral or complex
mass insertion needs its own statement.

At a fixed configuration the $V=-1$ zero modes contribute $x^{n_0}$;
nonreal conjugate pairs contribute positive quadratic factors. The index
gives $n_0\ge|\nu|$ and $n_0\equiv\nu\pmod2$. It does not guarantee
$n_0=|\nu|$: extra opposite-chirality pairs raise the source order.
Integrating the finite polynomial against a fixed mass-independent positive
remaining weight gives the one-sided mass-source derivatives at zero from
the retained positive branch. The polynomial has an algebraic extension;
no separately licensed negative-mass probability chart is inferred.
An arbitrary external Grassmann insertion needs its own source word,
zero-mode saturation and normalization. Its coefficient is not automatically
a positive probability weight or a numerical anomaly coefficient.

The complete domain $\Omega_{\rm gap}$ is an explicitly declared finite
representative of the generated gapped family for these bounded integrals;
no occurring physical transfer is asserted to select it. The following
scope distinctions preserve the complete parent used by each result:
\begin{itemize}
\item A domain restricted to nonzero-index sectors can have zero unsaturated
massless partition function. The exact normalized massless endpoint uses
the positive partition contribution from its retained zero-index ball.
\item An additional factor that vanishes on the unit-index balls can remove
their activity. Strict positivity alone supplies no quantitative global
floor and ceiling. Known bounds $0<g_-\le G\le g_+$ multiply the massive
probability floor by $g_-/g_+$ and the suppression ceiling by $g_+/g_-$.
\item A source-dependent state or action factor changes mass derivatives and
their hierarchy. A determinant phase or configuration-dependent chiral
measure can give a signed law; positive charge variance is not automatic
for that law.
\item The common finite scalar derivative tends to zero for at least two
simultaneously vanishing flavors on the stated fixed domain, whereas one
vanishing flavor allows a linear anomaly term. Neither statement proves
an ordered condensate or a thermodynamic spectral density. The inverse
ball mass and determinant floor are not uniform thermodynamic estimates.
\end{itemize}
The finite interval concerns $\chi V_L$. Physical masses and their transport,
the volume map, a chiral measure, reflection positivity, a continuum
susceptibility or pole interpretation, and an ordered SSB limit remain
their separately typed consumers.

% END INLINE SOURCE: Paper7A-NativeWeightedTopology-Appendix-v1_0.tex

% END NATIVE CONSTRUCTOR AND TOPOLOGY DIAGNOSTICS v2.189



% BEGIN GEOMETRY-FIRST CONSEQUENCES v2.201
\setlength{\emergencystretch}{2em}


\section{Hadron blocks, LFRep, thermal and nuclear auxiliary calculations}
\label{p13:sec:v237-d4}
% Source Paper VII-B, original line 2220; inherited proof bodies unchanged.
\subsubsection{Explicit constructed finite blocks}
\label{p11:subsec:explicit-constructed-finite-blocks}

The closure ledger above reduces the remaining QCD work to finite spectral
execution.  The first executable blocks can be written once the finite
Hamiltonian/source member and the indicated projectors are supplied.  Their
algebraic form uses the derived color factors, but their numerical entries,
topological weights, and chiral data are outputs of that supplied family and
of the regulated charge-covariance and double-limit correlator calculations;
none is fixed by a native contraction constant or multiplicity count.

\begin{theorem}[Compiled primitive-support transfer block]
\label{p11:thm:explicit-primitive-support-transfer-block}

Let \(P_{\chi,L}\) be the explicitly constructed orthogonal primitive
singlet channel projectors, with no pairwise orthogonality assumed.
Their canonical combined-support \emph{orthogonal projection} is
\[
P_{\rm prim,L}:=\operatorname{supp}\left(\sum_\chi P_{\chi,L}\right)
=\mathbf1_{(0,\infty)}\left(\sum_\chi P_{\chi,L}\right),\qquad
\operatorname{ran}P_{\rm prim,L}=\sum_\chi\operatorname{ran}P_{\chi,L}.
\]
It equals \(\sum_\chi P_{\chi,L}\) when the channel ranges are pairwise
orthogonal, but that is not needed. For \(t>0\), the exact finite block is
\[
K_{\rm prim,L}(t)=P_{\rm prim,L}e^{-t(H_L-E_{0,L})}P_{\rm prim,L}.
\]
If each channel projector reduces \(H_L\), and \(E_{0,L}\) is its least
energy, the restriction of this block to \(\operatorname{ran}P_{\rm prim,L}\)
is strictly positive with spectrum in \((0,1]\), and the retained energy
eigenvalues are \(-t^{-1}\log\lambda\). The full-space compressed
operator may additionally have zero eigenvalues on the orthogonal
complement; one never takes their logarithms. The block need not be
diagonal in support labels and is not fixed by \(C^2\delta\).
\phantomsection\label{p11:eq:explicit-primitive-transfer-block}
\phantomsection\label{p11:eq:explicit-primitive-hamiltonian-block}
\phantomsection\label{p11:eq:explicit-pure-glue-scalar-transfer-block}
\phantomsection\label{p11:eq:explicit-pure-glue-scalar-hamiltonian-block}
\end{theorem}
\begin{proof}

Set \(S_L=\sum_\chi P_{\chi,L}\ge0\). Finite spectral calculus gives
the orthogonal projection \(P_{\rm prim,L}=\mathbf1_{(0,\infty)}(S_L)\),
with range exactly the span of the primitive channel ranges.
If \([P_{\chi,L},H_L]=0\) for every channel, then \([S_L,H_L]=0\)
and hence \([P_{\rm prim,L},H_L]=0\), even when the channel
projectors do not commute with each other. On the combined range,
\(0<e^{-t(H_L-E_0)}\le I\) by finite positive spectral calculus,
so all retained transfer eigenvalues lie in \((0,1]\). The operator
vanishes on the orthogonal complement. Support-label diagonality does
not follow from reduction of the \emph{combined} range: the individual
channels may overlap or mix inside it and require diagonalization of
the resulting finite restriction.
\end{proof}

\paragraph{Parent of finite sector and measure execution.}
Paper~VII-A, Corollary~\ref{p10:cor:finite-qcd-parent-readouts}, locates the
covered finite hadron/nuclear extractions and winding weights in the existing
complete QCD source-transfer parent.  The following finite traces,
susceptibilities and residues execute those retained maps.  They are not new
sector/measure constructors.  Keep any action/state family and its compatible
readout image when a unique member has not been established.  Physical
isolation, ordered limits and open-channel continuation are distinct
claim-domain certificates; topology counts do not fix the dynamical weights.

\begin{theorem}[Finite topological generating block]
\label{p11:thm:explicit-positive-topological-winding-block}

For the regulated integer charge $\nu_L$ on its admitted native index domain,
retain the same-action sector coefficients $Z_{L,\nu}\ge0$, their complete
normalization $Z_L(0)>0$, and the state/thermal datum where consumed.  Define
the auxiliary characteristic source
\[
 p_{L,\nu}=\frac{Z_{L,\nu}}{Z_L(0)},\qquad
 Z_L(\theta)=\sum_\nu Z_{L,\nu}e^{i\nu\theta},\qquad
 \sum_\nu p_{L,\nu}=1,
\]
and
\[
 \chi_{t,L}=-V_L^{-1}\partial_\theta^2\log Z_L(\theta)|_0
 =V_L^{-1}\operatorname{Var}_0(\nu_L).
\]
The winding weights and all mixed covariances are readouts of the complete
packet; no universal count multiplier or independently constructed physical
$H_L(\theta)$ is required.  A physical strong-CP action coupling remains a
separate licensed sector-character claim.
\phantomsection\label{p11:eq:explicit-mixed-shell-topological-enhancement}
\phantomsection\label{p11:eq:explicit-topological-block-moments}
\phantomsection\label{p11:eq:explicit-positive-topological-transfer}
\end{theorem}
\begin{proof}

Expand the normalized generating function as \(Z_L(\theta)/Z_L(0)=\sum_\nu p_{L,\nu}e^{i\nu\theta}\).  Differentiating twice at zero gives the variance.  Positivity of the probabilities gives \(\chi_{t,L}\ge0\), while their magnitudes remain dynamical.
\end{proof}

\begin{corollary}[Paper-V-C second-moment evolution applied to Paper-VII-B data]
\label{p11:cor:explicit-second-moment-dglap-blocks}
\leavevmode\par\noindent
Applying the imported Paper-V-C DGLAP operator to the Paper-VII-B singlet boundary moments gives the matrix
\begin{equation}
\label{p11:eq:explicit-dglap-second-moment-block}
A_{n_f}^{\rm BH}=\begin{pmatrix}
-16/9&n_f/3\\[0.2em]
16/9&-n_f/3
\end{pmatrix}.
\end{equation}
For \(n_f=3\),
\begin{equation}
A_3^{\rm BH}=\begin{pmatrix}
-1.777777778&1\\
1.777777778&-1
\end{pmatrix},
\qquad
\operatorname{Spec}A_3^{\rm BH}=\left\{0,-\frac{25}{9}\right\}.
\end{equation}
The zero eigenvector gives \((Q_\infty,G_\infty)=(9/25,16/25)\).  Across thresholds,
\begin{equation}
\label{p11:eq:dglap-threshold-fixed-points}
\begin{array}{c|c|c|c}
n_f&\lambda_{\rm rel}&Q_\infty&G_\infty\\ \hline
3&-25/9&9/25&16/25\\
4&-28/9&3/7&4/7\\
5&-31/9&15/31&16/31\\
6&-34/9&9/17&8/17
\end{array}
\end{equation}
\end{corollary}
\begin{proof}
Equation~\eqref{p11:eq:explicit-dglap-second-moment-block} is the second Mellin moment of the fixed-color leading DGLAP kernels imported from Paper~V-C.  Its columns sum to zero,
expressing momentum conservation.  Solving
\[
  A_{n_f}(Q,G)^T=0,
  \qquad
  Q+G=1,
\]
gives
\[
  Q_\infty=\frac{3n_f}{16+3n_f},
  \qquad
  G_\infty=\frac{16}{16+3n_f}.
\]
The nonzero eigenvalue is
\[
  -\left(\frac{16}{9}+\frac{n_f}{3}\right).
\]
\end{proof}

\begin{corollary}[Leading pseudoscalar mass matrix from typed chiral data]
\label{p11:cor:explicit-leading-pseudoscalar-blocks}

Once a physical quark-mass operator, the chiral condensate/decay constant, and a topological anomaly term are supplied, the leading pseudoscalar mass operator is the usual chiral Hessian.  For off-diagonal flavor \(i\bar j\),
\[
M_{ij}^2=B_0(m_i+m_j),
\]
while the neutral sector is the corresponding flavor Hessian plus the computed singlet-anomaly entry.  All pole values are eigenvalues of that matrix; no numerical pseudoscalar spectrum is exported before these inputs are computed.
\phantomsection\label{p11:eq:explicit-pseudoscalar-octet-mass-map}
\phantomsection\label{p11:eq:native-pseudoscalar-spectrum-numerical}
\phantomsection\label{p11:eq:explicit-native-neutral-spectrum-summary}
\phantomsection\label{p11:eq:neutral-reclosure-coefficient}
\phantomsection\label{p11:eq:numerical-neutral-pseudoscalar-pole-readout}
\phantomsection\label{p11:eq:native-fks-y-factor}
\phantomsection\label{p11:eq:native-singlet-axial-renormalization}
\phantomsection\label{p11:eq:native-fks-eta-eta-prime-matrix}
\phantomsection\label{p11:eq:native-eta-eta-prime-closed-poles}
\phantomsection\label{p11:eq:explicit-soft-pion-scattering-block}
\end{corollary}
\begin{proof}

Expand the chiral potential to quadratic order about the chosen vacuum.  The mass spurion gives \(B_0(m_i+m_j)\) for off-diagonal generators, and the anomalous singlet source adds only to the singlet-direction Hessian at this order.  Diagonalization is finite linear algebra.  Current Paper~VII-A leaves the required physical inputs as correlator/source targets.
\end{proof}

\begin{theorem}[Geometric heavy-source form factor and physical-kernel boundary]
\label{p11:thm:finite-hurwitz-heavy-source-kernel-package}

The finite Hurwitz coordinate set determines exact finite source densities and their Fourier form factors.  These dimensionless trigonometric identities may be used as source shapes.  A physical heavy-quark potential additionally requires the typed quark masses, gauge-action/propagator member, and the position-space pairing/area map; the geometry alone does not determine a Coulomb coefficient or string tension.
\phantomsection\label{p11:eq:hurwitz-heavy-source-density}
\phantomsection\label{p11:eq:hurwitz-heavy-source-form-factor}
\phantomsection\label{p11:eq:hurwitz-heavy-source-swave-density-kernel}
\phantomsection\label{p11:eq:hurwitz-heavy-kinetic-symbol}
\phantomsection\label{p11:eq:finite-hurwitz-coulomb-kernel}
\phantomsection\label{p11:eq:finite-hurwitz-flux-profile}
\phantomsection\label{p11:eq:finite-hurwitz-raw-flux-length}
\phantomsection\label{p11:eq:finite-hurwitz-flux-average-A}
\phantomsection\label{p11:eq:finite-hurwitz-flux-subtraction}
\phantomsection\label{p11:eq:finite-hurwitz-heavy-kernel-substitutions}
\end{theorem}
\begin{proof}

Fourier transformation of a finite point measure is an exact finite sum, so the geometric form factors follow without dynamics.  Multiplying a candidate propagator or area functional by an arbitrary positive response coefficient leaves those source form factors unchanged while changing the physical potential, proving that the latter requires additional operator data.
\end{proof}

\begin{corollary}[Conditional heavy 1S variational bound]
\label{p11:cor:leading-finite-heavy-1s-variational-readout}

For any supplied self-adjoint heavy-source Hamiltonian \(H_{Q\bar Q,L}\) and normalized trial family \(\psi_\beta\),
\[
E_{0,L}\le \inf_\beta\langle\psi_\beta,H_{Q\bar Q,L}\psi_\beta\rangle.
\]
This Rayleigh--Ritz bound is parameter-free once the Hamiltonian is supplied; no absolute heavy pole follows from the Hurwitz source form factor alone.
\phantomsection\label{p11:eq:leading-finite-heavy-1s-functional}
\phantomsection\label{p11:eq:finite-heavy-1s-spin-average}
\end{corollary}
\begin{proof}

The Rayleigh--Ritz principle gives the inequality.  The trial parameter is varied rather than fitted.  The numerical Hamiltonian still depends on the physical heavy-mass and interaction operators, which are not fixed by source geometry alone.
\end{proof}

% Source Paper VII-B, original line 2400; inherited proof bodies unchanged.
\subsubsection{Conditional downstream blocks from the physical quark-mass operator}
\label{p11:subsec:native-downstream-e40-blocks}

Current Paper~VII-A does not export a quark mass from a visible-shell count.  It instead inherits the actual microscopic quark-source member from G0 and transports that already owned operator to \(\mu\) through the G16 RG/threshold map.  The following chiral and scalar-response statements are therefore conditioned on that inherited renormalized mass/source operator \(M_q^{\rm phys}(\mu)\), on occurrence of the controlled chiral regime, and on the required correlator/low-energy limits.  No \({\Ksh}=40\) shell coordinate or finite count is identified with a quark mass in this section, and no physical-pole certificate is required merely to define the running chiral spurion.

\begin{theorem}[Criterion for removing Kaplan--Manohar ambiguity]
\label{p11:thm:kaplan-manohar-ambiguity-absent}

The Kaplan--Manohar reparametrization freedom is removed as a physical selector if the physical quark-mass/source operator is microscopically fixed by an injective upstream operator map and transported in a fixed renormalization convention.  Without such microscopic ownership, low-order chiral data alone do not select the representative.
\phantomsection\label{p11:eq:km-lambda-zero}
\end{theorem}
\begin{proof}

The ambiguity is precisely an invariance of the truncated chiral description under a redefinition of the mass spurion accompanied by low-energy-constant shifts.  A microscopically fixed mass/source operator breaks this equivalence by fixing which spurion represents the inherited source member.  Paper~V-A supplies an injective source-owner map; after an actual G0 member occurs, Paper~V-B/VII-A transport that member through the certified G16 RG map.  Hence the Kaplan--Manohar parameter is not an additional physical selector on this native branch.  This statement does not determine the low-energy constants, the controlled chiral regime, correlator limits, or higher-order truncation errors.
\end{proof}

\begin{corollary}[Dashen relation without an absolute electromagnetic coefficient]
\label{p11:cor:universal-dashen-lift-native}

In the leading chiral \(SU(3)\) limit, the Dashen theorem equates the charged pion and kaon electromagnetic mass-squared lifts while neutral leading lifts vanish up to the stated approximation.  The common coefficient is a dynamical electromagnetic current-correlator low-energy constant and is not fixed to unity by dimensional analysis.
\phantomsection\label{p11:eq:universal-dashen-lift-native}
\phantomsection\label{p11:eq:dashen-native-relations}
\end{corollary}
\begin{proof}

The leading chiral electromagnetic operator has one common coefficient multiplying the charge commutator structure.  Symmetry therefore relates splittings, but multiplying that operator by any dimensionless constant preserves all symmetries.  Hence the relation is derived while the absolute coefficient must be computed/matched.
\end{proof}

\begin{corollary}[PCAC pseudoscalar-density relation]
\label{p11:cor:pcac-pseudoscalar-density-native}

For a pseudoscalar state \(P_{ij}\) in the domain where PCAC holds,
\[
(m_i+m_j)\langle0|\bar q_i i\gamma_5q_j|P_{ij}\rangle=F_{P_{ij}}M_{P_{ij}}^2
\]
(up to the stated renormalization/sign convention).  The relation is exact at the operator-identity level; its numerical matrix element requires the computed \(F_P,M_P,m_i,m_j\).
\phantomsection\label{p11:eq:pcac-density-native}
\phantomsection\label{p11:eq:pseudoscalar-density-universal-native}
\end{corollary}
\begin{proof}

Take the vacuum-to-one-particle matrix element of the axial Ward identity \(\partial^\mu A_\mu=(m_i+m_j)P\) and use \(\langle0|A_\mu|P(p)\rangle=iF_Pp_\mu\).  No count-derived chiral constant is needed.
\end{proof}

\begin{corollary}[Feynman--Hellmann Goldstone sigma terms]
\label{p11:cor:goldstone-sigma-native}

For a differentiable pseudoscalar eigenbranch,
\[
\sigma_{q}^{(P)}=m_q\frac{\partial M_P}{\partial m_q}
\]
with derivatives taken with respect to the physical mass-source operator.  Leading chiral formulas follow by differentiating the leading mass matrix; no fixed numerical sigma term is implied before that operator is supplied.
\phantomsection\label{p11:eq:goldstone-sigma-native}
\phantomsection\label{p11:eq:mixed-neutral-sigma-native}
\end{corollary}
\begin{proof}

This is the Feynman--Hellmann theorem applied to the mass-source parameter.  The derivative equals the corresponding scalar matrix element.  Substitution of any chiral approximation is conditional on its domain and parameters.
\end{proof}

\begin{corollary}[Soft Goldstone kernel and scattering boundary]
\label{p11:cor:native-soft-pipi-pik-scattering}

The leading amputated soft \(\pi\pi\) and \(\pi K\) kernels are the standard
current-algebra/leading-chiral expressions in terms of the computed
pseudoscalar masses and decay constants.  They become asymptotic scattering
amplitudes or scattering lengths only on the applicable
\textnormal{(SM1)--(SM4)} domain of
Definition~\ref{p12:def:claim-indexed-smat-package}.  Paper~VII-B retains the
symbolic kernel relations but withdraws numerical values based on obsolete
\(F_\pi\) and mass inputs.
\phantomsection\label{p11:eq:native-pipi-scattering-lengths-numeric}
\phantomsection\label{p11:eq:native-pik-scattering-lengths}
\end{corollary}
\begin{proof}

The \(O(p^2)\) chiral action and crossing/isospin projection give the amputated
kernel.  LSZ identifies it with a scattering amplitude only under the named
asymptotic clauses.  Since the chiral coefficients are the actual correlator
outputs, neither statement determines those outputs.
\end{proof}

\begin{corollary}[Strong eta-to-three-pion source]
\label{p11:cor:native-eta-to-three-pi-source}

At leading chiral order the amputated strong \(\eta3\pi\) kernel is
proportional to the isospin-breaking scalar source
\(B_0(m_d-m_u)/F_\pi^2\), with the standard kinematic polynomial and neutral
mixing factors.  It becomes a physical \(\eta\to3\pi\) decay amplitude only
after an asymptotic interpretation is chosen: either a stable-parent
approximation on the applicable \textnormal{(SM1)--(SM4)} domain, or an
unstable-pole definition using \textnormal{(SM1)--(SM3)} and
\textnormal{(SM5)}, with \textnormal{(SM4)} for any LSZ-normalized stable
external legs.  A decay rate additionally uses the phase-space and
channel-normalization parts of \textnormal{(SM7)} in
Definition~\ref{p12:def:claim-indexed-smat-package}.
The kernel is a computed chiral-source readout, not a geometric constant.
\phantomsection\label{p11:eq:native-eta3pi-source}
\phantomsection\label{p11:eq:native-eta3pi-amplitude}
\end{corollary}
\begin{proof}

The leading chiral mass spurion is the only CP-even strong source linear in
\(m_d-m_u\) with the required flavor transformation.  Expanding the chiral
action gives the amputated kernel.  The claim-indexed asymptotic clauses alone
promote it to a decay amplitude or rate.  Current Paper~VII-A does not fix the
physical quark-mass operator or \(B_0,F_\pi\) numerically.
\end{proof}

\begin{corollary}[Neutral transition form factors from charge traces and computed mixing]
\label{p11:cor:native-neutral-transition-ff}

Anomaly/factorization charge traces fix the flavor coefficients of neutral pseudoscalar transition form factors.  The physical normalizations require the computed decay constants and the neutral-state mixing matrix.  Therefore the charge-trace relations survive, while predecessor numerical transition constants are withdrawn.
\phantomsection\label{p11:eq:native-neutral-charge-coeff}
\phantomsection\label{p11:eq:native-transition-ff-asymptotics}
\end{corollary}
\begin{proof}

The electromagnetic anomaly is linear in the flavor generator weighted by \(Q_f^2\), giving fixed trace coefficients.  Changing the normalized physical eigenvectors or decay constants changes the amplitude without changing those traces, so the latter do not close the numerical form factor.
\end{proof}

\begin{corollary}[Charged-pseudoscalar weak kernels from current matrix elements]
\label{p11:cor:native-charged-pseudoscalar-weak-kernels}

For each charged pseudoscalar, \(\langle0|A_\mu|P(p)\rangle=iF_Pp_\mu\) defines the QCD decay constant.  Semileptonic vector form factors satisfy their current-normalization and Ademollo--Gatto constraints in the appropriate symmetry limit.  Absolute kernels require the computed \(F_P\), masses and form-factor slopes.
\phantomsection\label{p11:eq:native-leptonic-qcd-kernel}
\phantomsection\label{p11:eq:native-kell3-normalization}
\end{corollary}
\begin{proof}

These statements are Ward/current-algebra consequences.  The Ademollo--Gatto theorem removes first-order symmetry breaking from the protected vector normalization, but it does not determine higher-order corrections or decay constants.
\end{proof}

\begin{corollary}[Exact two-pseudoscalar kinematic thresholds]
\label{p11:cor:native-two-pseudoscalar-thresholds}

For stable pseudoscalar poles with masses \(M_a,M_b\), the bottom of the
noninteracting two-particle spectral sum in the rest frame is exactly
\(E_{ab}^{\rm thr}=M_a+M_b\).  Numerical thresholds are therefore downstream
spectral readouts, not independent constants.  Calling the associated
continued singularity an \(S\)-matrix branch point additionally uses
clause \textnormal{(SM5)} of
Definition~\ref{p12:def:claim-indexed-smat-package}.

\end{corollary}
\begin{proof}

At zero total three-momentum the minimum of
\(\sqrt{M_a^2+\mathbf p^2}+\sqrt{M_b^2+\mathbf p^2}\) occurs at
\(\mathbf p=0\), giving the mass sum.  This kinematic spectral statement is
independent of the separate \(S\)-matrix interpretation.
\end{proof}

\begin{corollary}[Feynman--Hellmann flavor shifts replace additive flavor floors]
\label{p11:cor:primitive-flavor-floors-native}

For a differentiable hadron eigenbranch along a physical mass-source path \(M_q(s)\),
\[
M_h(1)-M_h(0)=\int_0^1\!ds\;\langle h,s|\partial_sH[M_q(s)]|h,s\rangle.
\]
Thus flavor splittings are source-response integrals.  No universal rule \(M_{\rm floor}+\sum m_q\) follows from color support.
\phantomsection\label{p11:eq:primitive-meson-flavor-floor}
\phantomsection\label{p11:eq:primitive-decuplet-equal-spacing-native}
\end{corollary}
\begin{proof}

Integrate the Feynman--Hellmann derivative of the isolated eigenvalue along the path.  Sea, spin, glue and orbital contributions can change with the source and are automatically included in the matrix element; color representation alone does not make the integrand equal to a valence count.
\end{proof}

\begin{corollary}[Conditional pseudoscalar HRG formulas]
\label{p11:cor:pseudoscalar-hrg-native}

Given a specified set of physical pseudoscalar pole masses and conserved
charges, the dilute noninteracting HRG pressure and charge cumulants are the
corresponding Bose-gas sums.  Interactions may be added by a virial/scattering
expansion only when its phase shifts satisfy the applicable channel package in
Definition~\ref{p12:def:claim-indexed-smat-package}.  No fixed \(T_c\),
condensate ratio, or mass list is derived from the finite carrier alone.
\phantomsection\label{p11:eq:pseudoscalar-hrg-pressure-native}
\phantomsection\label{p11:eq:native-ps-hrg-tc-numerics}
\phantomsection\label{p11:eq:native-condensate-ps-hrg}
\phantomsection\label{p11:eq:native-lowT-charge-susceptibilities}
\end{corollary}
\begin{proof}

The ideal-gas partition function is a product over species and momentum modes; logarithmic differentiation gives pressure and cumulants.  The approximation is conditional on dilution and on the supplied spectrum.  The finite-carrier geometry does not select the thermodynamic temperature at which the approximation ceases to apply.
\end{proof}


% Source Paper VII-B, original line 2917; inherited proof bodies unchanged.
\subsubsection{Geometric lift, pure-glue representatives, and leading nuclear blocks}
\label{p11:subsec:geometric-lift-glue-nuclear-blocks}

The primitive support blocks above are one-support readouts.  The next finite step is the geometric lift from one support to relative multi-support geometry.  This subsection records the first such lift.  These entries are not comparison fits: they use only the Hurwitz group law, the singlet color projectors, the supplied finite Hamiltonian/source member, and the baryon/chiral constants fixed earlier.  Whenever a statement is only the leading finite-cell block, this is stated explicitly.

\begin{theorem}[Hurwitz relative-coordinate and convolution lift]
\label{p11:thm:hurwitz-relative-support-lift}

For \(G=2T\), the quotient identities \(G^2/G_{\rm diag}\cong G\) and \(G^3/G_{\rm diag}\cong G^2\) are exact.  Given any supplied class kernels \(k_1,k_2,k_3\), the corresponding relative transfers are the finite left/right convolutions determined by the group law.  Their spectra are products of the actual Fourier/representation eigenvalues of those kernels; no universal \(q_C=1-C^2\delta\) is implied.
\phantomsection\label{p11:eq:two-support-relative-coordinate}
\phantomsection\label{p11:eq:relative-two-support-transfer}
\phantomsection\label{p11:eq:equal-support-relative-eigenvalues}
\phantomsection\label{p11:eq:three-support-relative-coordinate}
\phantomsection\label{p11:eq:three-support-relative-transfer}
\end{theorem}
\begin{proof}

The relative coordinates \(x^{-1}y\) and \((x_1^{-1}x_2,x_1^{-1}x_3)\) are invariants of the diagonal left action.  The transition law follows by substituting independent group increments.  Peter--Weyl diagonalizes class convolution and shows that spectral products depend on the supplied kernel eigenvalues, completing the exact geometric lift.
\end{proof}


\begin{theorem}[Pure-glue relative geometry and fibre multiplicity]
\label{p11:thm:geometric-pure-glue-relative-block-corrected}

The two-support Hurwitz quotient and the adjoint--curvature fibre multiplicity are exact finite geometry.  If a relative transfer block \(K_J^{\rm rel}\) is supplied, the trivial fibre lift \(K_J^{\rm rel}\otimes I_{32}\) repeats eigenvalues and multiplies source traces, not pole energies.  Absolute glue energies must come from the compiled Yang--Mills Hamiltonian/stiffness and are not fixed by a shell coefficient.
\phantomsection\label{p11:eq:ss-relative-primitive-energies-corrected}
\phantomsection\label{p11:eq:primitive-glue-relative-representatives-corrected}
\phantomsection\label{p11:eq:curvature-fibre-transfer-tensor-product}
\end{theorem}
\begin{proof}

Tensoring an operator with an identity repeats every eigenvalue with multiplicity equal to the fibre dimension.  Hence multiplicity affects degeneracy and unnormalized traces but not eigenvalue values.  Any physical energy scale must therefore enter through the Hamiltonian quadratic form, as in current Paper~VII-A.
\end{proof}

\begin{theorem}[Exact Hodge kernel and generalized Yang--Mills stiffness]
\label{p11:thm:finite-plaquette-hodge-curvature-hessian}
\label{p11:thm:bilinear-curvature-heat-time}

The dimensionless face-Hodge operators
\[
\Delta_2=d_1d_1^\dagger+d_2^\dagger d_2,
\qquad
\mathcal B_2=\Delta_2\otimes I+I\otimes\Delta_2
\]
are exact geometric operators.  A physical curvature energy is determined by the supplied Yang--Mills quadratic form \(Q_{\YM,L}\).  Its stiffness relative to the Hodge metric is the generalized Rayleigh quotient
\[
\kappa_{\YM,L}=\inf_{F\ne0}\frac{Q_{\YM,L}[F]}{\langle F,\Delta_{2,H}F\rangle},
\]
not a universal \(\Lambda_{\rm BH}\).
\phantomsection\label{p11:eq:face-hodge-curvature-operator}
\phantomsection\label{p11:eq:bilinear-curvature-hessian}
\phantomsection\label{p11:eq:curvature-stiffness-kappa-lambdabh}
\end{theorem}
\begin{proof}

The Hodge operators are constructed solely from incidence matrices.  Current Paper~VII-A proves their coercive geometric role and defines physical stiffness through the actual Hamiltonian quadratic form.  Rescaling the Yang--Mills action coefficient changes \(Q_{\YM,L}\) and \(\kappa_{\YM,L}\) without changing \(\Delta_2\), proving the separation.
\end{proof}

\begin{corollary}[Cell-local scalar Hodge moment]
\label{p11:cor:cell-local-scalar-curvature-stiffness}

For the normalized cell-local scalar source \(v_c\) obtained from the eight triangular faces of an octahedral cell, the exact dimensionless Hodge bilinear has the computed local Rayleigh moment
\[
\frac{\langle v_c,\mathcal B_2v_c\rangle}{\langle v_c,v_c\rangle}=10
\]
for the canonical source convention used here.  The corresponding physical energy is obtained only after solving the generalized Yang--Mills quadratic-form problem; \(10\) is not an energy in MeV.
\phantomsection\label{p11:eq:cell-local-scalar-source}
\phantomsection\label{p11:eq:cell-local-scalar-projected-hessian}
\phantomsection\label{p11:eq:cell-local-scalar-stiffness-energy}
\end{corollary}
\begin{proof}

The diagonal face-Hodge contribution is five per factor, so the diagonal bilinear contribution is ten for the canonical diagonal source.  Normalizing the source gives the stated dimensionless moment.  A physical stiffness requires the independent Hamiltonian quadratic form, so multiplying by an obsolete universal scale is not licensed.
\end{proof}


\begin{corollary}[Conditional finite face-incidence form-factor ratios]
\label{p11:cor:finite-face-incidence-glue-residues}
Let the \emph{actually defined} normalized triangular face-incidence
operator $F$ act on a finite Hilbert space with pairwise orthogonal
adjacency eigenprojections $P_1,P_2,P_3$. If their exact source
incidence calculation verifies
\[
 F P_j=\lambda_jP_j,\qquad
 (\lambda_1,\lambda_2,\lambda_3)
 =\left(\frac23,\frac13,\frac16\right),
\]
then for equally normalized vectors $v_j\in\operatorname{ran}P_j$
with nonzero common norm, the squared geometric form factors satisfy
$\|Fv_1\|^2:\|Fv_2\|^2:\|Fv_3\|^2=16:4:1$.
The current source does not specify the normalized incidence matrix,
its adjacency eigenprojections and a check of the three asserted
\(\lambda_j\) for that actual operator. Hence the statement is
conditional; the factor values are targets of the exact finite-source
certificate, \emph{not yet an independently derived spectrum} in this release.
Interpreting them as physical $0^{++},2^{++},0^{-+}$ poles additionally
requires source/operator projection and genuine spectral overlap in the
compiled Yang--Mills Hamiltonian.
\phantomsection\label{p11:eq:triangle-incidence-source-operator}
\phantomsection\label{p11:eq:glue-face-form-factors}
\phantomsection\label{p11:eq:glue-reduced-residue-ratio}
\end{corollary}
\begin{proof}
From $FP_j=\lambda_jP_j$ one obtains
$\|Fv_j\|^2=|\lambda_j|^2\|v_j\|^2$ on each specified
projection. The eigenvalue ratios are $4:2:1$, so their squares are
$16:4:1$. This proof checks only the spectral-algebra implication and
cannot replace an exact computation of the source-owned incidence matrix,
normalization and eigenspaces. A form factor is not automatically a
Hamiltonian pole residue.
\end{proof}


\begin{corollary}[Conditional string tension and glue ratios]
\label{p11:cor:native-string-tension-and-relative-glue-ratios}
For a specified pure-glue or certified unscreened fundamental Wilson sector,
the reciprocal-position metric area and Wilson holonomy are already
generated.  If the actual complete action/measure admits the same-owner APBI
heat presentation of Corollary~\ref{p10:cor:v2158-apbi-string-tension-reduction},
its uniform strict majorant gives
\[
 |\langle W(C)\rangle|
 \le C_0\exp\!\left(c_\partial P(C)-\sigma_*A_{g_x}(C)\right),
 \qquad \sigma_*:=\frac{\log z}{A_+}>0.
\]
APBC is the square-face sufficient specialization.  A direct same-owner
multiface/flux surface theorem can instead supply a positive lower bound on
the stated sector; no new area-pairing coefficient is needed.  The bound is
not a numerical evaluation of the tension, and it is not an automatic
extension to a determinant-weighted full-QCD screened ground sector.
Glueball-to-positive-string ratios require both the relevant pole spectrum
and the actual positive tension on this domain.
\phantomsection\label{p11:eq:native-string-tension-derived-corrected}
\phantomsection\label{p11:eq:compact-glue-string-normalized-ratios}
\phantomsection\label{p11:eq:compact-glue-string-normalized-ratios-numeric}
\end{corollary}
\begin{proof}
The position area and Wilson source are inherited.  The APBI theorem supplies
the displayed bound on its same-action unscreened heat branch, while
Theorem~\ref{p10:thm:hurwitz-static-potential-floor} records the separate direct
flux route with its full admissible-sector hypotheses.  Neither route
supplies its unproved extensive owner/bound or evaluates the complete
ground-state tension merely from color counts.
\end{proof}

\begin{theorem}[Light meson support block as a spectral compression]
\label{p11:thm:light-qqbar-relative-meson-block}

Color algebra fixes the singlet \(q\bar q\) support and its \(J^{PC}\)/flavor decomposition.  For each such sector, the relative meson spectrum is the spectrum of the corresponding compression of the supplied \(H_L-E_{0,L}\).  The Hurwitz relative-coordinate construction may organize basis states, but it does not fix the transfer eigenvalues by \(1-C_F^2\delta\).
\phantomsection\label{p11:eq:qqbar-relative-energy-tower}
\phantomsection\label{p11:eq:qqbar-native-channel-list}
\end{theorem}
\begin{proof}

The singlet projector and quantum-number projectors define a finite invariant block.  The spectral theorem yields its energy levels.  Independent changes of the radial/spin interaction preserve the color factor but change the energies, proving that the latter are dynamical.
\end{proof}

\begin{theorem}[Determinant baryon geometry and spectral multiplets]
\label{p11:thm:three-quark-determinant-geometry-baryon-multiplets}

The three-quark determinant fixes the color singlet and the associated exchange symmetry.  Spin/flavor/orbital multiplets are obtained by decomposing the remaining state space.  Their masses and hyperfine splittings are eigenvalues of the corresponding baryon Hamiltonian blocks, not powers of a universal determinant transfer number.
\phantomsection\label{p11:eq:determinant-color-weight-cb}
\phantomsection\label{p11:eq:baryon-relative-exponent-spectrum}
\phantomsection\label{p11:eq:baryon-su6-ground-decomp}
\end{theorem}
\begin{proof}

Exterior antisymmetry of color is exact.  The remaining representation decomposition is group theory.  The Hamiltonian restricted to each irrep block is an arbitrary symmetry-compatible Hermitian matrix supplied by the compiler, so its eigenvalues require actual matrix elements rather than a count rule.
\end{proof}

\begin{corollary}[Baryon current/scalar package as computed matrix elements]
\label{p11:cor:leading-baryon-current-scalar-response-package}

The leading baryon readout package consists of conserved charges, current form factors, axial/tensor reduced matrix elements, magnetic slopes, and scalar Feynman--Hellmann derivatives evaluated on the computed baryon eigenstates.  Symmetry relates members within a multiplet but does not supply universal numerical dressing factors.
\phantomsection\label{p11:eq:baryon-current-package-numbers}
\phantomsection\label{p11:eq:fd-leading-baryon-package}
\end{corollary}
\begin{proof}

Ward identities, Wigner--Eckart, and Feynman--Hellmann give the stated decomposition.  Each remaining reduced element is a matrix element of a specified operator, so it is determined by the compiled state/operator pair rather than by a separate phenomenological coefficient.
\end{proof}

\begin{theorem}[OPE plus computed short-range nuclear operators]
\label{p11:thm:native-ope-and-nn-contact-block}

In the low-energy two-nucleon domain, the long-range one-pion-exchange operator has the standard form determined by the computed \(g_A,F_\pi,m_\pi\).  The short-range interaction is the most general Hermitian contact operator allowed by the stated symmetries at the chosen EFT order, with coefficients obtained by matching QCD multi-baryon matrix elements.  No fixed matrix \(\operatorname{diag}(6,14/3)\) or \(1-R\delta\) transfer law follows from color counting.
\phantomsection\label{p11:eq:native-ope-momentum-kernel}
\phantomsection\label{p11:eq:ope-coefficient-number}
\phantomsection\label{p11:eq:nn-contact-r-matrix}
\phantomsection\label{p11:eq:nn-contact-transfer-matrix}
\phantomsection\label{p11:eq:nn-contact-hamiltonian-matrix}
\phantomsection\label{p11:eq:contact-only-threshold-lengths}
\end{theorem}
\begin{proof}

Pion exchange follows from the two pion--nucleon vertices and the pion propagator.  Locality and symmetries classify the contact operator basis but do not determine its reduced coefficients.  Distinct coefficients preserve all color/spin labels while changing scattering lengths, which proves the need for QCD matching.
\end{proof}

\begin{proposition}[Deuteron as a conditional two-body spectral problem]
\label{p11:prop:leading-deuteron-contact-ope-block}

Given a fully specified finite/continuum two-nucleon Hamiltonian \(H_D\) containing the computed contact and pion-exchange operators, a deuteron exists exactly when \(H_D\) has an isolated eigenvalue below the two-nucleon threshold.  Its binding energy and D-state probability are spectral/eigenvector readouts; no fixed MeV or percentage is asserted without the execution matrix.
\phantomsection\label{p11:eq:leading-deuteron-coupled-hamiltonian}
\phantomsection\label{p11:eq:deuteron-eigenvalue-readout}
\phantomsection\label{p11:eq:leading-deuteron-estimate}
\end{proposition}
\begin{proof}

This is the definition of a bound state for a self-adjoint two-body Hamiltonian.  The min--max principle determines the bound energy, and the squared norm of the D-wave projection determines the D-state probability.  The current corpus does not supply the complete contact kernel needed for a numerical execution.
\end{proof}



% Source Paper VII-B, original line 3304; inherited proof bodies unchanged.
\subsubsection{Additional native validation signatures from Blocks D--F}
\label{p11:subsec:additional-native-validation-signatures-blocks-d-f}

The preceding Blocks D--F do not add phenomenological inputs.  They give native
comparison signatures: quantities which must be computed from the same finite
Hamiltonian, source dictionary, charge operators, and transfer cumulants before
any external comparison is made.

\begin{corollary}[Distribution moment relations and the kaon odd-moment target]
\label{p11:cor:distribution-moment-fingerprint-kaon-odd}

Pion exchange symmetry forces its odd Gegenbauer moments to vanish in the isospin limit.  A kaon odd moment is allowed by flavor asymmetry but is a matrix element of the typed light-front operator and therefore a computation target, not a finite-geometry constant.
\phantomsection\label{p11:eq:distribution-moment-fingerprint-general}
\phantomsection\label{p11:eq:a1-xi-moment-fingerprint}
\phantomsection\label{p11:eq:pion-odd-moment-zero-fingerprint}
\phantomsection\label{p11:eq:kaon-odd-moment-native-readout}
\end{corollary}
\begin{proof}

The pion statement follows from \(x\leftrightarrow1-x\) symmetry and orthogonality of odd Gegenbauer polynomials.  No corresponding symmetry exchanges strange and light quarks in a kaon, so an odd coefficient is permitted; its value requires the physical distribution operator.
\end{proof}

\begin{corollary}[Single-trace covariance identities without a fixed thermal lock]
\label{p11:cor:single-trace-thermal-lock-signature}

At finite level, derivatives of \(\log Z_L\) with respect to temperature and source parameters are cumulants/covariances in the same Gibbs state.  This creates exact correlation identities among thermal responses.  It does not force their extrema to occur at one common temperature.
\phantomsection\label{p11:eq:thermal-lock-vector}
\phantomsection\label{p11:eq:thermal-lock-response-definitions}
\end{corollary}
\begin{proof}

Differentiate the normalized finite trace.  First derivatives give expectation values and second mixed derivatives give connected covariances.  Distinct analytic functions can have extrema at distinct temperatures, so a common trace is insufficient to prove a universal thermal lock.
\end{proof}

\begin{corollary}[Logarithmic transfer distinguishes computed contact channels]
\label{p11:cor:nn-contact-ratio-log-hamiltonian-distinction}

For any positive computed finite contact-channel transfer eigenvalues \(0<\lambda_1,\lambda_2\le1\), their energy difference is
\[
\Delta E=-t^{-1}\log(\lambda_2/\lambda_1).
\]
No universal raw contact ratio follows before the contact transfer itself is computed.
\phantomsection\label{p11:eq:nn-contact-raw-ratio-nine-seven}
\phantomsection\label{p11:eq:nn-contact-transfer-exact-forms}
\phantomsection\label{p11:eq:nn-contact-hamiltonian-ratio-exact}
\end{corollary}
\begin{proof}

Functional calculus gives \(E_i=-t^{-1}\log\lambda_i\); subtraction yields the formula.  Since the predecessor contact entries were not symmetry-fixed, their ratio cannot be retained as a theorem.
\end{proof}

\begin{corollary}[Scalar--glue correlated-residue validation target]
\label{p11:cor:scalar-glue-correlated-residue-validation-target}
In the scalar singlet channel the pole data are a correlated Schur readout,
not a list of independent glueball and quarkonium masses or fitted mixing angles.
Let \(P\) project onto the compact scalar curvature/quarkonium support and \(Q=1-P\) onto open channels.  Then
\begin{equation}
\label{p11:eq:scalar-glue-effective-hamiltonian-signature}
H_{{\rm eff},L}^{(0^{++})}(z)=PHP+PHQ(z-QHQ)^{-1}QHP,
\qquad z\notin\sigma(QHQ),
\end{equation}
and its finite poles solve
\begin{equation}
\label{p11:eq:scalar-glue-signature-pole-equation}
\det\!\bigl[z\mathbf1-H_{{\rm eff},L}^{(0^{++})}(z)\bigr]=0.
\end{equation}
They are real and have residues obtained by compressing orthogonal spectral
projectors.  After the separately
proved open-channel limit and continuation, a simple pole \(s_p\) of
\(D^{\rm II}(s)\) has correlated matrix residue
\begin{equation}
\label{p11:eq:scalar-glue-correlated-residues-signature}
\operatorname*{Res}_{s=s_p}D^{\rm II}(s)^{-1}
=\frac{v_pw_p^\dagger}{w_p^\dagger(D^{\rm II})'(s_p)v_p}.
\end{equation}
Consequently the observable scalar pattern is a residue vector across the scalar
singlet sources and open channels.  It may validate the framework only after the
finite Schur matrix is executed; it may not be used to choose \(M_G,\xi,\rho_G\), or any scalar mixing angle.
\end{corollary}
\begin{proof}
This is the scalar specialization of
Theorem~\ref{p11:thm:block-b-finite-pole-dressing},
Theorem~\ref{p11:thm:schur-real-part-sheets-residues}, and
Theorem~\ref{p11:thm:isolated-resonance-pole-residue-limit}.  The finite Schur
complement gives the real spectral poles; the separately continued inverse
operator gives the left/right residue formula.  Since every finite entry is
produced by the compiler and source dictionary, no independent scalar residue
or mixing parameter remains.
\end{proof}

\begin{proposition}[Mixed \texorpdfstring{\(c\bar b\)}{c bbar} hyperfine validation target]
\label{p11:prop:mixed-cb-hyperfine-validation-target}
Let Paper~V-B fix the charm and bottom quark slots before QCD transport.  In Paper~VII-B the mixed heavy-flavour meson blocks are then
\begin{equation}
\label{p11:eq:bc-blocks-native-spectral-target}
H_{c\bar b}^{0^-}=P_{c\bar b,0^-}H_{\rm QCD}^{\rm fin}P_{c\bar b,0^-},
\qquad
H_{c\bar b}^{1^-}=P_{c\bar b,1^-}H_{\rm QCD}^{\rm fin}P_{c\bar b,1^-}.
\end{equation}
The native hyperfine target is
\begin{equation}
\label{p11:eq:bc-hyperfine-native-target}
\Delta_{\rm hf}^{c\bar b,\rm BH}
=
M_{B_c^\ast}^{\rm BH}-M_{B_c}^{\rm BH},
\qquad
M_{B_c^{(\ast)}}^{\rm BH}\in\operatorname{Spec}(H_{c\bar b}^{J^P}).
\end{equation}
This is a sharp validation target because it combines the upstream Paper~V-B charm--bottom slot geometry with the downstream Paper~VII-B QCD transport.  It is not a new input and it is not a number until the two finite heavy-flavour blocks are diagonalized.
\end{proposition}
\begin{proof}
The flavor projectors \(P_{c\bar b,J^P}\) are fixed by the quark slots, spin/parity labels, and finite singlet projector.  Restricting the already compiled finite Hamiltonian to these two sectors gives finite matrices.  The hyperfine splitting is the spectral difference of the lowest vector and pseudoscalar eigenvalues.  Since neither the charm/bottom slots nor the QCD transport entries are selected from the measured \(B_c\) spectrum, comparison cannot backflow into the construction.
\end{proof}



% Source Paper VII-B, original line 4775; inherited proof bodies unchanged.
% Former appendix marker: technical material is now placed in VII-D.


% Source Paper VII-B, original line 4777; inherited proof bodies unchanged.
\subsection{Optional exact finite LFRep positive-measure realization}
\label{p11:app:v264-exact-lfrep-positive-measures}
This appendix retains the stronger exact finite occupation-space realization.
It is not a premise of the main physical light-ray/twist-two chain.  When the
full certificate \(\LFRep_L\) is supplied, the following theorems provide a
commuting fraction simplex, positive number/momentum POVMs, an exact finite
support proof on \([0,1]\), and positive-measure convergence criteria.  These
finite statements do not by themselves imply positivity or support of a
scheme-renormalized continuum PDF.  On this appendix branch, clause~\textnormal{(R3)} is assumed and we write
\[
 P_L^+:=P_{F,L}^+=P_{\rm phys,L}^+,
 \qquad E_{p_h,L}=E^{\rm phys}_{p_h,L}.
\]

\begin{theorem}[Constructed total momentum-fraction operators]
\label{p11:thm:constructed-light-front-fraction-operators}
On the fixed sector in Definition~\ref{p11:def:hadronic-distribution-package},
define
\begin{equation}
\label{p11:eq:constructed-total-fraction-operators}
X_{i|h,L}:=\frac{1}{p_h^+}
E_{p_h,L}P_{i,L}^+E_{p_h,L} .
\end{equation}
Then the \(X_{i|h,L}\) are commuting positive self-adjoint operators and
\begin{equation}
\label{p11:eq:fraction-operator-simplex}
0\le X_{i|h,L}\le E_{p_h,L},\qquad
[X_{i|h,L},X_{j|h,L}]=0,qquad
\sum_iX_{i|h,L}=E_{p_h,L}.
\end{equation}
Consequently every spectrum lies in \([0,1]\).  These operators are constructed
from longitudinal translation and species data; no independent
``momentum-fraction observable'' and no phenomenological normalizer is added.
\end{theorem}
\begin{proof}
The commuting positive one-body operators \(p_L^+\pi_{i,L}\) have commuting
positive second quantizations.  Their sum is \(P_L^+\).  Compression by the
commuting fixed-momentum projector preserves positivity and commutativity.
Equation~\eqref{p11:eq:lf-fixed-total-plus-sector} therefore gives
\(\sum_iX_{i|h,L}=E_{p_h,L}\); positivity of every summand implies
\(X_{i|h,L}\le E_{p_h,L}\).  The spectral inclusion follows from the last
operator inequality.
\end{proof}

\begin{theorem}[Parton number and momentum-fraction operator measures]
\label{p11:thm:light-front-fraction-povms}
For a Borel set \(B\subset[0,1]\), define on the same sector
\begin{align}
\label{p11:eq:parton-number-povm}
\mathsf N_{i|h,L}(B)
&:=E_{p_h,L}\,\mathrm d\Gamma\!\left(
\pi_{i,L}\mathbf 1_B(p_L^+/p_h^+)\right)E_{p_h,L},\\
\label{p11:eq:parton-momentum-povm}
\mathsf M_{i|h,L}(B)
&:=\frac{1}{p_h^+}E_{p_h,L}\,\mathrm d\Gamma\!\left(
p_L^+\pi_{i,L}\mathbf 1_B(p_L^+/p_h^+)\right)E_{p_h,L}.
\end{align}
They are commuting positive operator-valued measures supported on \([0,1]\).
The counting measure \(\mathsf N_i\) need not be normalized, whereas the family
\(\{\mathsf M_i\}_{i\in\mathcal I}\) is a normalized POVM on
\(\mathcal I\times[0,1]\).  Moreover,
\begin{equation}
\label{p11:eq:fraction-povm-identities}
\mathsf M_{i|h,L}(B)=\int_Bx\,\mathsf N_{i|h,L}(dx),\qquad
\mathsf M_{i|h,L}([0,1])=X_{i|h,L},\qquad
\sum_i\mathsf M_{i|h,L}([0,1])=E_{p_h,L}.
\end{equation}
Thus
\begin{equation}
\label{p11:eq:finite-cut-pdf-measures}
\nu_{i,h,L}(B):=\langle h,L;p_h^+|\mathsf N_{i|h,L}(B)|h,L;p_h^+\rangle,
\quad
\mu_{i,h,L}(B):=\langle h,L;p_h^+|\mathsf M_{i|h,L}(B)|h,L;p_h^+\rangle
\end{equation}
are respectively the finite-cut number and momentum measures, with
\(\mu_{i,h,L}(dx)=x\nu_{i,h,L}(dx)\).  The spectral measure of the single
total-share operator \(X_{i|h,L}\) instead describes fluctuations of the
\emph{total} species share and is not the one-parton PDF.
\end{theorem}
\begin{proof}
Finite spectral calculus makes \(B\mapsto\mathbf1_B(p_L^+/p_h^+)\) a
projection-valued measure.  Second quantization, multiplication by the
commuting positive operator \(p_L^+\pi_{i,L}\), and compression preserve
countable additivity and positivity.  In a fixed positive-total-momentum
occupation configuration, no occupied one-parton momentum exceeds \(p_h^+\),
which gives support in \([0,1]\).  Functional calculus gives the first identity
in \eqref{p11:eq:fraction-povm-identities}; the other two follow from
Theorem~\ref{p11:thm:constructed-light-front-fraction-operators}.  A zero mode
appears at \(x=0\) in \(\mathsf N\) and contributes zero to \(\mathsf M\), so
the formulas retain rather than discard it.
\end{proof}

\begin{theorem}[Optional exact-LFRep Wilson/POVM identification]
\label{p11:thm:v264-optional-finite-wilson-povm-identification}
Assume clause~\textnormal{(R5)} of \(\LFRep_L\) and the alignment in
\(\LFState_{h,L}\) intertwine the finite mode expansion with the
gauge-invariant good-field bilocals along the derived path
\(z_\lambda:=\gamma_{+,z_0}(\lambda)\):
\begin{equation}
\label{p11:eq:finite-light-ray-bilocal}
\mathcal O_{q,L}(\lambda)
=\bar\psi_{q,+,L}(z_\lambda)\gamma^+
W_{+,L}(\lambda,0)\psi_{q,+,L}(z_0)
\end{equation}
and the analogous \(F^{+\perp}W_+F^+{}_{\perp}\) gluon operator.  Then their
finite Fourier coefficients are the number measures
\(\nu_{i,h,L}\), and their local Mellin moments agree with the moments of
\(\mathsf N_{i|h,L}\).  This proves the finite-cut PDF identification.

A continuum collinear PDF is obtained only under
clauses~\textnormal{(L2)--(L4)} of \(\LFLim\): the null helix has the stated
local-flat/decompactifying light-ray limit, the color holonomies converge, and
the renormalized operator limit commutes with the moment/Fourier operation.
TMDs additionally require clause~\textnormal{(L5)}.  After
such subtraction, a scheme-dependent renormalized PDF is an operator-valued
distribution and need not itself be the spectral measure of a positive
commuting operator, even though the finite-cut operator measures above are positive.
\end{theorem}
\begin{proof}
Expand the good fields in the finite longitudinal momentum eigenbasis.  Finite
Fourier orthogonality along the represented flow parameter \(\lambda\) turns
the separation on \(\gamma_{+,z_0}\) into the spectral
projectors \(\mathbf1_B(p_L^+/p_h^+)\); normal ordering separates quark and
antiquark occupation terms.  The transporter makes the bilocal gauge invariant,
and the assumed intertwiner identifies the resulting coefficients with
\eqref{p11:eq:parton-number-povm}.  Differentiation at \(\lambda=0\) gives the
Mellin-moment identities.  Ultraviolet light-ray and rapidity divergences arise
only upon removal of the cut and therefore require the separately stated
renormalization limits; finite spectral positivity does not pass automatically
through arbitrary scheme subtraction~\cite{Collins2011Foundations}.
\end{proof}

\begin{theorem}[Optional exact-LFRep finite PDF sum-rule realization]
\label{p11:thm:v264-optional-finite-pdf-sum-rules}
Under Theorem~\ref{p11:thm:v264-optional-finite-wilson-povm-identification}, let \(h\) have net
flavor charges \(N_f^h\), baryon number \(B_h\), and electric charge \(Q_h\).
The finite-cut measures obey
\begin{align}
\label{p11:eq:pdf-net-flavor-sum-rule}
\nu_{q_f,h,L}([0,1])-\nu_{\bar q_f,h,L}([0,1])&=N_f^h,\\
\label{p11:eq:pdf-baryon-sum-rule}
\frac13\sum_f\bigl(\nu_{q_f,h,L}-\nu_{\bar q_f,h,L}\bigr)([0,1])&=B_h,\\
\label{p11:eq:pdf-charge-sum-rule}
\sum_f Q_f\bigl(\nu_{q_f,h,L}-\nu_{\bar q_f,h,L}\bigr)([0,1])&=Q_h,\\
\label{p11:eq:pdf-momentum-sum-rule}
\sum_{i=q,\bar q,g}\mu_{i,h,L}([0,1])&=1.
\end{align}
At finite $L$, the momentum identity is algebraic operator closure from the
constructed species decomposition and the fixed-total-plus sector; it consumes
no additional Noether or total-momentum premise.  If the renormalized continuum
limit preserves the corresponding finite quantum Ward identities and total
energy--momentum identity under the state/regulator and QFT4/QFT10 controls of
Paper~VII-C, Proposition~\ref{p12:prop:v248-finite-quantum-ward-antecedent} and
Theorem~\ref{p12:thm:finite-identity-preservation}, these become the standard
PDF integrals \(\int(q_f-\bar q_f)=N_f^h\) and
\(\sum_i\int_0^1x f_{i/h}(x,\mu)\,dx=1\).
\end{theorem}
\begin{proof}
The first identity is the occupation realization of the conserved net-flavor
operator; Theorem~\ref{p10:thm:vector-flavor-charge-conservation} fixes its
eigenvalue.  Baryon number and charge are the displayed linear combinations.
The last identity is the expectation of the final equality in
\eqref{p11:eq:fraction-povm-identities}; it is therefore already closed at
finite level once the species decomposition and fixed-$P_L^+$ state are
supplied.  Theorem~\ref{p11:thm:v134-fixed-plus-sector-preservation} shows
that the total-momentum sector itself is preserved by the finite Hamiltonian
flow.  Continuum transport is a separate quantum-Ward/renormalization
certification problem and must retain every anomaly, contact, and regulator
term controlled by Paper~VII-C.  No scattering hypothesis is used.
\end{proof}

\begin{theorem}[Hausdorff moment limit for the momentum measures]
\label{p11:thm:hausdorff-pdf-moment-limit}
Suppose that for every species \(i\) and every integer \(r\ge0\),
\begin{equation}
\label{p11:eq:hausdorff-moment-convergence}
m_{i,r}:=\lim_{L\to\infty}\int_{[0,1]}x^r\,\mu_{i,h,L}(dx)
\end{equation}
exists.  Then there is a unique positive measure \(\mu_{i,h}\) on \([0,1]\)
with these moments, and the full sequence converges weakly:
\(\mu_{i,h,L}\Rightarrow\mu_{i,h}\).  Moreover
\(\sum_i\mu_{i,h}([0,1])=1\).  Identification
\(\mu_{i,h}(dx)=x f_{i/h}(x,\mu)\,dx\) additionally requires the renormalized
Wilson-line convergence in Theorem~\ref{p11:thm:v264-optional-finite-wilson-povm-identification};
absolute continuity is not a consequence of moment convergence.  The analogous
statement for number measures \(\nu_{i,h,L}\) requires a uniform total-number
bound, which is an extra small-\(x\)/zero-mode condition.
\end{theorem}
\begin{proof}
The momentum sum rule bounds every positive \(\mu_{i,h,L}\) by one in total
mass.  Compactness of \([0,1]\) therefore gives weakly convergent subsequences.
Every subsequential limit has moments \(m_{i,r}\).  If two finite measures on
\([0,1]\) have all the same moments, they agree on polynomials; the Weierstrass
theorem and the Riesz representation theorem then make them equal on every
continuous test function.  The Hausdorff moment problem is therefore
determinate, so all subsequences have the same limit and the full sequence
converges.  Taking the zeroth moments in the finite sum rule gives the limiting
normalization.  Neither density with respect to Lebesgue measure nor control of
\(\nu=\mu/x\) at \(x=0\) follows from these arguments.
\end{proof}

\begin{theorem}[Bounded-Lipschitz refinement limit for momentum measures]
\label{p11:thm:bounded-lipschitz-pdf-limit}
Assume the finite species set \(\mathcal I\), the finite momentum sum rule, and
for every \(i\in\mathcal I\) the Cauchy condition
\begin{equation}
\label{p11:eq:bounded-lipschitz-cauchy}
 \lim_{L,M\to\infty}
 \sup_{\substack{\|f\|_\infty\le1\\
                  \operatorname{Lip}(f)\le1}}
 \left|
 \int f\,d\mu_{i,h,L}-\int f\,d\mu_{i,h,M}
 \right|=0 .
\end{equation}
Then there is a unique positive measure \(\mu_{i,h}\) on \([0,1]\) for each
species such that
\[
 \mu_{i,h,L}\Rightarrow\mu_{i,h},
 \qquad
 \sum_{i\in\mathcal I}\mu_{i,h}([0,1])=1.
\]
All momentum moments converge as a consequence. Identification
\[
\mu_{i,h}(dx)=x f_{i/h}(x,\mu)\,dx
\]
still requires clauses~\textnormal{(L2)--(L4)} of \(\LFLim\), and absolute
continuity is not implied. Number-measure convergence separately requires
uniform small-\(x\)/zero-mode control.

Equivalently, clause~\textnormal{(L1)} may establish
\eqref{p11:eq:bounded-lipschitz-cauchy} through the positive maps
\[
 \Phi_{i,L}(f):=\int f(x)\,\mathsf M_{i|h,L}(dx).
\]
Each component is completely positive and subunital, while the combined map
on \(C(\mathcal I\times[0,1])\) is unital completely positive.  Comparing
different \(L\) still requires the declared refinement channels and compatible
\(\LFState_{h,L}\) states; the common codomain is not inferred.
\end{theorem}

\begin{proof}
The sum rule bounds each positive measure by one in total mass.  On compact
\([0,1]\), bounded sets of positive measures are weakly sequentially compact,
and the bounded-Lipschitz metric metrizes weak convergence.  Every subsequence
therefore has a weakly convergent subsubsequence, while the Cauchy condition
forces any two such limits to have zero bounded-Lipschitz distance.  Hence the
full sequence has one weak limit.  Testing with the constant function and
summing over the finite species set transports the finite normalization to the
limit.  Every monomial is continuous on \([0,1]\), so its integral converges.
Positive maps from a commutative \(C^*\)-algebra are completely positive; the
finite operator-measure normalization makes the combined map unital.  None of
these facts supplies the separately typed renormalized light-ray or
small-\(x\) limits.
\end{proof}




% Source Paper VII-B, original line 5032; inherited proof bodies unchanged.
\subsection{Compression-commutator defect for the G35 compatibility test}
\label{p11:app:v199-compression-defect}
\begin{proposition}[Exact finite symmetry leakage]
\label{p11:prop:v199-compression-defect}
For finite self-adjoint \(H,G\), an orthogonal projection \(P\), and
\(Q=I-P\),
\[
 [PHP,PGP]=P[H,G]P-PHQGP+PGQHP.
\]
Thus \([H,G]=0\) and \([P,G]=0\) suffice for exact inherited symmetry.
In the parent-symmetric case the defect norm is at most
\[
 \|PHQ\|\,\|QGP\|+\|PGQ\|\,\|QHP\|.
\]
\end{proposition}
\begin{proof}
Paper~1B, Proposition~\ref{p1b:result:compression-defect}
\cite{Kugelmann2026Paper1B}, with $A=H$ and $B=G$ gives the displayed
identity for this compression. Submultiplicativity gives the norm bound;
$[P,G]=0$ removes both leakage terms. For unbounded parents the identity
requires a common domain on which all displayed products are defined.
\end{proof}
This is a diagnostic after the native parent generator is constructed.
Defining an unrelated finite exponential \(e^{itG}\) proves none of the
required QCD/null-flow or species relations.

% BEGIN VERIFIED QCD READOUT CERTIFICATES v2.210


\section{Analytic refinement, continuum gauge and renormalization estimates}
\label{p13:sec:v237-d5}
% Source Paper VII-C, original line 2864; inherited proof bodies unchanged.
\subsubsection{Action-first effective-action realization of QFT4}
\label{p12:subsec:action-first-effective-rg}

On the optional action-first gauge branch, the finite coarse map is not an
independent lattice input.  It is the evaluation-dual restriction map of the
inherited reciprocal-position path refinement.

\begin{theorem}[Inherited reciprocal-path refinement determines the finite holonomy coarse map]
\label{p12:thm:v251-inherited-path-coarse-map}
Let $L\le M$ be two finite reciprocal-position presentation levels of the
already constituted position role.  Let $\mathsf{Path}_L$ and
$\mathsf{Path}_M$ be compatible oriented finite path groupoids used only as
auxiliary reciprocal-position presentations, and let
\[
 \jmath_{LM}:\mathsf{Path}_L\hookrightarrow\mathsf{Path}_M
\]
be their genuine subdivision/refinement inclusion.  Thus an oriented coarse
edge $e$ has a retained fine representative
\[
 \jmath_{LM}(e)=e_1^{\epsilon_1}\cdots e_k^{\epsilon_k},
 \qquad \epsilon_i\in\{\pm1\}.
\]
For a compact gauge group $G$ and a fine holonomy configuration
$U_M:\mathsf{Path}_M\to G$ respecting concatenation and inversion, define
\begin{equation}
\label{p12:eq:v251-holonomy-restriction}
 \boxed{\pi_{M\to L}(U_M):=U_M\circ\jmath_{LM}.}
\end{equation}
Then this is the unique coarse configuration map satisfying
\begin{equation}
\label{p12:eq:v251-evaluation-square}
 \operatorname{Hol}_{\pi_{M\to L}U_M}(\gamma)
 =\operatorname{Hol}_{U_M}(\jmath_{LM}\gamma)
 \qquad(\gamma\in\mathsf{Path}_L).
\end{equation}
It obeys:
\begin{enumerate}[label=\textnormal{(\roman*)},leftmargin=3em]
\item exact composition,
$\pi_{N\to L}=\pi_{M\to L}\circ\pi_{N\to M}$ whenever
$\jmath_{LN}=\jmath_{MN}\circ\jmath_{LM}$;
\item the edge-product formula
\begin{equation}
\label{p12:eq:v251-edge-product}
 (\pi_{M\to L}U)_e
 =U_{e_1}^{\epsilon_1}\cdots U_{e_k}^{\epsilon_k};
\end{equation}
\item gauge equivariance, with all intermediate vertex transformations
cancelling along a subdivided coarse path;
\item on the admitted cellular subdivision branch in which distinct
independent coarse edges use interior-disjoint fine edge chains, projective
product Haar measure,
\begin{equation}
\label{p12:eq:v251-haar-projective}
 (\pi_{M\to L})_*\mu_{0,M}=\mu_{0,L};
\end{equation}
\item if the inherited path refinement intertwines Euclidean reflection,
$\jmath_{LM}\Theta_L=\Theta_M\jmath_{LM}$, then
$\pi_{M\to L}\Theta_M=\Theta_L\pi_{M\to L}$;
\item for every bounded coarse cylinder observable $F$,
$\alpha_{LM}F:=F\circ\pi_{M\to L}$ is the reciprocal coarse-to-fine test
inclusion, and its state dual is
\begin{equation}
\label{p12:eq:v251-state-test-dual}
 \omega_M\longmapsto\omega_M\circ\alpha_{LM}
 \equiv (\pi_{M\to L})_*\omega_M .
\end{equation}
\end{enumerate}
The objects $\mathsf{Path}_K$ and $\jmath_{LM}$ are inherited presentation
structure of the already licensed reciprocal-position role; they are not a
new primitive position lattice, physical scale, or independent carrier.
\end{theorem}
\begin{proof}
Equation~\eqref{p12:eq:v251-holonomy-restriction} is genuine restriction of
the fine holonomy evaluation to an inherited subpresentation.  Hence
\eqref{p12:eq:v251-evaluation-square} determines the map uniquely because the
retained path evaluations separate finite holonomy configurations.  The
composition law is contravariance of precomposition.  Concatenation and
inversion give \eqref{p12:eq:v251-edge-product}; the endpoint gauge factors
survive while every internal factor cancels, proving gauge equivariance.
Normalized Haar measure is idempotent under convolution, so the product of
the independent fine edge variables along each subdivided coarse edge is
again Haar distributed; interior-disjoint subdivision chains give
\eqref{p12:eq:v251-haar-projective} jointly.  Reflection compatibility follows
by precomposition with the stated path intertwiner.  Finally,
$\int F\,d(\pi_*\omega)=\int(F\circ\pi)\,d\omega$, giving
\eqref{p12:eq:v251-state-test-dual}.

The typing statement uses the Paper~II-A same-role inheritance rule for a
genuine structured restriction together with its evaluation-dual
co-realization theorem: no new physical sort is introduced by retaining the
finite path presentation, and no compression such as $P_KAP_K$ is being
inserted.  Thus the refinement incidence, its holonomy restriction, the
coarse-observable pullback and the state pushforward are co-realized
presentations of one inherited relation rather than independent physical
constructors.
\end{proof}

\begin{corollary}[Finite same-owner geometric--quantum block square]
\label{p12:cor:v251-finite-same-owner-block-square}
Let $\mathcal G^{\rm blk}_{M\to L}$ denote the complete finite action-first
block constructor using the same inherited $\jmath_{LM}$ and
$\pi_{M\to L}$ of Theorem~\ref{p12:thm:v251-inherited-path-coarse-map}:
retain the coarse presentation, transport the gauge constraints and coarse
sources by the same restriction, and integrate the fine fibre against the
projective reference measure.  On the positive Haar branch,
\begin{equation}
\label{p12:eq:v251-finite-same-owner-square}
 \boxed{
 \mathcal Q_L\,\mathcal G^{\rm blk}_{M\to L}
 =\mathcal R_{M\to L}\,\mathcal Q_M,
 \qquad
 \mathcal R_{M\to L}=[(\pi_{M\to L})_*].}
\end{equation}
The identity holds in the declared essential quotient.  Source derivatives
are the conditional insertions of
Theorem~\ref{p12:thm:v199-rgas-disintegration}; gauge constraints descend by
equivariance; and reflection positivity is preserved when the inherited path incidence
intertwines reflection and its observable pullback carries the coarse
positive-time algebra into the fine positive-time algebra.
\end{corollary}
\begin{proof}
By Theorem~\ref{p12:thm:v251-inherited-path-coarse-map}, the geometric block
and the state/test comparison use one and the same inherited incidence and
projective Haar system.  Fibre integration therefore gives exactly the
pushforward measure.  Writing that pushforward relative to the coarse
reference measure produces the effective action of
Theorem~\ref{p12:thm:v199-rgas-disintegration}, so quantizing the complete
blocked datum and pushing forward the fine quantum ray are the same measure
class.  The source, constraint, reflection and essential-quotient statements
follow from the same restriction/pushforward operation and the cited finite
blocking transport results.
\end{proof}

For a positive fine action measure
$d\mu_M=Z_M^{-1}e^{-S_M}d\mu_{0,M}$, define the exact coarse pushforward
$\mu_{M\downarrow L}=(\pi_{M\to L})_*\mu_M$ using this inherited map.
Because the reference Haar measure is projective, the coarse measure has a
positive Radon--Nikodym density and hence an exact effective action
$S^{\rm eff}_{M\downarrow L}$, unique up to a vacuum constant.


\begin{theorem}[Exact finite blocking, score and connected response]
\label{p12:thm:v199-rgas-disintegration}
Let the finite configuration spaces be standard Borel, with projective
probability reference measures
\((\pi_{M\to L})_*\nu_M=\nu_L\), as on the admitted finite Haar branch.
For \(w_M=e^{-S_M}>0\) integrable, use a regular conditional distribution
\(\nu_M(dy\mid x)\) to define
\begin{equation}
\label{p12:eq:v199-rgas-disintegration}
 w_L(x)=\int e^{-S_M(y)}\,\nu_M(dy\mid x),\qquad
 S_L^{\rm eff}(x)=-\log w_L(x).
\end{equation}
The weight is finite and positive almost everywhere.  The partition
function is preserved; blocking through an intermediate level equals
direct blocking almost everywhere.  If maps and reference measures are
parameter independent and differentiation under the conditional integral
is justified, then
\begin{equation}
\label{p12:eq:v199-rgas-score}
 \partial_a S_L^{\rm eff}=\langle \partial_aS_M\rangle_x,\qquad
 \partial_a\partial_b S_L^{\rm eff}
 =\langle\partial_a\partial_b S_M\rangle_x
   -\operatorname{Cov}_x(\partial_aS_M,\partial_bS_M).
\end{equation}
Conditional averages use the weight \(e^{-S_M}/w_L(x)\).
\end{theorem}
\begin{proof}
Use the conditional blocking construction of Paper~1B,
Theorems~\ref{p1b:thm:blocking-transport} and~\ref{p1b:thm:blocking}
\cite{Kugelmann2026Paper1B}, on the given projective reference measures.
Conditional expectation preserves the integral of \(w_M\); positivity and
integrability give the almost-everywhere assertions.  The tower property
for the sigma-algebras of the composed coarse maps gives
\(\mathbb E[\mathbb E[w_M\mid L]\mid K]=\mathbb E[w_M\mid K]\).
Differentiate \(-\log w_L\) once for the first response formula.
Differentiating that normalized average gives the expected second
derivative minus its connected covariance.
\end{proof}

\begin{corollary}[Source and reflection transport under finite blocking]
\label{p12:cor:v199-rgas-sources-reflection}
For \(S_M-J^aO_a\) with the same domination conditions, the first source
derivative of \(S_L^{\rm eff}[J]\) is \(-\langle O_a\rangle_x\) and its
second is \(-\operatorname{Cov}_x(O_a,O_b)\).  Coarse pulled-back sources
factor through the block integral; fine sources generate conditional
composite insertions.  If the fine state is reflection positive, the
coarse map commutes with reflection, and it pulls positive-time coarse
observables into the fine positive-time algebra, the blocked state is
reflection positive.
\end{corollary}
\begin{proof}
Apply \eqref{p12:eq:v199-rgas-score} to the source deformation.
The coarse reflection quadratic form at \(A\) equals the fine form at
\(A\circ\pi\), which is nonnegative by the two compatibility conditions.
\end{proof}
The removable vacuum constant is independent of every surviving field.
If the metric survives, \(c\int\sqrt g\) is an action coordinate and is
retained.  Exact finite blocking does not supply the uniform source and
complete-action bounds needed by the following continuum theorem.



% BEGIN V203 POSITIVE BLOCKING COMPARISON
\begin{theorem}[Action oscillation and sharp positive-state comparison]
\label{p12:thm:v203-blocking-oscillation}
On the real positive finite-measure branch of
Theorem~\ref{p12:thm:v199-rgas-disintegration}, use the same coarse map and
reference conditional measures for
$\mathcal B(S)(x)=-\log\int e^{-S(y)}\nu(dy\mid x)$.
Assume both weights are integrable with nonzero fibre integrals almost
everywhere. If $m\le h=T-S\le M$ essentially, then
\[
 m\le\mathcal B(T)-\mathcal B(S)\le M,\qquad
 \operatorname{osc}(\mathcal B(T)-\mathcal B(S))
 \le\operatorname{osc}(T-S),
\]
where $\operatorname{osc}h=\operatorname{ess\,sup}h-
\operatorname{ess\,inf}h$. With $\operatorname{TV}=\tfrac12\|\cdot\|_1$,
the normalized states obey the sharp bound
\[
 \operatorname{TV}(\mu_S,\mu_T)
 \le\tanh\bigl(\operatorname{osc}(T-S)/4\bigr).
\]
The same bound holds after exact blocking.
\end{theorem}
\begin{proof}
The positive weights and common conditional measures in
Theorem~\ref{p12:thm:v199-rgas-disintegration} are precisely the data of
Paper~1B, Theorem~\ref{p1b:thm:blocking}
\cite{Kugelmann2026Paper1B}. Its range contraction and sharp
total-variation estimate give the assertions, with the same essential
oscillation and $\operatorname{TV}=\tfrac12\|\cdot\|_1$ conventions.
\end{proof}

\begin{corollary}[Summable action defects give fixed-coarse state convergence]
\label{p12:cor:v203-summable-action-defects}
For a countable tower with exact compatible blocking, suppose the positive
finite actions satisfy
\[
 \operatorname{osc}(\mathcal B_{n+1\to n}S_{n+1}-S_n)\le\epsilon_n.
\]
For $\ell\le M<N$,
\[
 \operatorname{osc}(\mathcal B_{N\to\ell}S_N
                   -\mathcal B_{M\to\ell}S_M)
 \le\sum_{n=M}^{N-1}\epsilon_n.
\]
Consequently $\sum_n\epsilon_n<\infty$ gives total-variation convergence
at each fixed coarse level. The limiting marginal states are consistent.
For bounded classical observables $|O_i|\le M_i$, comparison of any two
such states gives
\[
 |\kappa_k^\mu(O_1,\ldots,O_k)-\kappa_k^\nu(O_1,\ldots,O_k)|
 \le2\operatorname{TV}(\mu,\nu)
       \left(\sum_{r=1}^k r!\,{k\brace r}\right)\prod_i M_i,
\]
where ${k\brace r}$ are Stirling numbers of the second kind.
\end{corollary}
\begin{proof}
Use the actual compatible tower and its stated adjacent defect bounds
in Paper~1B, Theorem~\ref{p1b:thm:blocking}
\cite{Kugelmann2026Paper1B}. Iterating its oscillation contraction
gives the finite sum from $M$ to $N-1$; its summable-tail and cumulant
conclusions give the remaining claims. For a particular tower the
required sum is the sum of that tower's own $\epsilon_n$.
\end{proof}
Adding the same real bounded source tilt to both compared actions preserves
$h$, so the comparison bound is uniform on any common integrable tilted
domain. All surviving fields enter the configuration variable: a term
$c\int\sqrt g$ is not a removable constant when $g$ survives.
This supplies a quantitative sufficient criterion for bounded QFT4/state
comparison and for propagation of a proposed quantum RG defect. Native
$\epsilon_n$ summability, volume uniformity, complete-action and mixed-polymer
bounds are still to be proved; an extensive local error can have growing
oscillation. Complex chiral weights require a different argument. The given
geometric UV fixed point, interacting quantum intertwiner, essential
critical surface and physical trajectory retain their distinct statuses.

\begin{theorem}[Exact presentation change commutes with finite blocking]
\label{p12:thm:v204-blocking-presentation}
At each existing finite level let $F_n:X_n\to\widetilde X_n$ be an
invertible measurable presentation change, with measurable inverse, and
transport the reference measures, actions, sources and coarse maps:
\[
 \widetilde\nu_n=F_{n*}\nu_n,\qquad
 \widetilde S_n=S_n\circ F_n^{-1},\qquad
 \widetilde\pi_{n,m}=F_m\pi_{n,m}F_n^{-1}.
\]
Then the exact conditional blocking of
Theorem~\ref{p12:thm:v199-rgas-disintegration} obeys
\[
 \widetilde{\mathcal B}_{n\to m}(\widetilde S_n)
   =(\mathcal B_{n\to m}S_n)\circ F_m^{-1}
\]
almost everywhere in compatible versions. Normalized states, bounded
correlators, their total-variation comparison distances and the
action-oscillation defects of
Corollary~\ref{p12:cor:v203-summable-action-defects} are unchanged.
Thus exact presentation transport contributes zero blocking discrepancy.
\end{theorem}
\begin{proof}
Apply Paper~1B, Theorem~\ref{p1b:thm:blocking-transport}
\cite{Kugelmann2026Paper1B}, to the present reference measures, coarse
maps and bimeasurable changes $F_n$. Their transported data are exactly
those in the theorem. The almost-everywhere conditional identity and
preservation of integrals, total variation and essential oscillations
therefore give the stated zero presentation discrepancy.
\end{proof}
The isometric image presentations established in Papers~II-C--IV are such
changes on their existing finite configuration variables, with all other
variables retained. They do not themselves change a refinement level or
replace genuine integration over removed degrees of freedom. Consequently
native inter-level action errors and their required uniformity still require
the actual action calculation; neither sampling nor a new coordinate measure
may be charged as an independent approximation. If a different reference
measure is used, its Jacobian belongs to the transported action.

% END V203 POSITIVE BLOCKING COMPARISON

\begin{theorem}[Effective-action stability implies bounded QFT4]
\label{p12:thm:effective-action-stability-qft4}
Suppose the exact coarse action admits the decomposition
\[
 S^{\rm eff}_{M\downarrow L}
 =S_L+\Delta S_{ML}+\text{const},
\]
where the unique CP-even marginal coordinate and the CP-odd coordinate have
already been matched to their owned coarse values and the remaining
$\Delta S_{ML}$ is local/quasilocal.  If for every fixed bounded local
cylinder observable $A$
\[
 \sup_{M\succeq L}\sup_{0\le s\le1}
 \left|\operatorname{Cov}^{(M)}_{L,s}(A,\Delta S_{ML})\right|
 \le C_A\varepsilon_L,
 \qquad \varepsilon_L\to0,
\]
then
\[
 \boxed{
 \sup_{M\succeq L}
 |\omega_M(A\circ\pi_{M\to L})-\omega_L(A)|
 \le C_A\varepsilon_L\to0.}
\]
Exact composition of the pushforwards gives refinement-path independence.
If the coarse maps commute with Euclidean reflection, reflection positivity
is preserved under pushforward.
\end{theorem}
\begin{proof}
By definition of the pushforward effective state,
$\omega_M(A\circ\pi)=\omega^{\rm eff}_{M\downarrow L}(A)$.
Interpolate between $S_L$ and $S^{\rm eff}_{M\downarrow L}$; the derivative
is the covariance with $\Delta S_{ML}$.  Integrating the uniform bound gives
the estimate.  Exact composition of the $\pi$ maps gives the common-refinement
argument, while reflection compatibility lets every coarse positive-time
cylinder be pulled back to a positive-time fine cylinder.
\end{proof}


\begin{theorem}[Uniform-source effective-action stability]
\label{p12:thm:v185-uniform-source-effective-action-stability}
Under the hypotheses of Theorem~\ref{p12:thm:effective-action-stability-qft4},
let
\[
 \Psi_L(J)=\sum_{r=1}^N J_r\mathcal O_{r,L}
\]
be a finite family of bounded local source insertions compatible with the
coarse maps, so that
$\Psi_M(J)=\Psi_L(J)\circ\pi_{M\to L}$ on their common cylinder domain.  Put
\[
 \widehat Z_L(J):=\omega_L(e^{\Psi_L(J)}).
\]
Assume there is a fixed complex polydisc $\mathbb D_r^N$ on which the
interpolation covariance bound of
Theorem~\ref{p12:thm:effective-action-stability-qft4} holds uniformly for the
bounded observable $A_J=e^{\Psi_L(J)}$:
\[
 \sup_{M\succeq L}\sup_{0\le s\le1}\sup_{J\in\mathbb D_r^N}
 \left|\operatorname{Cov}^{(M)}_{L,s}(A_J,\Delta S_{ML})\right|
 \le C_r\varepsilon_L,
 \qquad \varepsilon_L\to0.
\]
Then
\[
 \boxed{
 \sup_{M\succeq L}\sup_{J\in\mathbb D_r^N}
 |\widehat Z_M(J)-\widehat Z_L(J)|\to0.}
\]
Consequently $\widehat Z_L$ converges uniformly on every smaller polydisc.
If the family $\widehat Z_L$ is uniformly bounded away from zero on one
such smaller polydisc, then $W_L=\log\widehat Z_L$ converges there after
choosing the branch fixed by
$W_L(0)=0$, and Cauchy differentiation gives convergence of every fixed mixed
source moment and connected cumulant generated by this source family.
\end{theorem}
\begin{proof}
Source compatibility makes
$A_J\circ\pi_{M\to L}=e^{\Psi_M(J)}$.  Theorem~\ref{p12:thm:effective-action-stability-qft4}
applied uniformly to $A_J$ gives the displayed Cauchy estimate for
$\widehat Z_L$.  Uniform convergence on compact sub-polydiscs and the common
analytic source domain allow Cauchy's integral formula to be applied to each
mixed derivative.  Where the limiting functional is nonzero, the normalized
branch $W(0)=0$ is analytic and derivatives of $W_L$ are the connected
cumulants.  No conclusion is drawn for an unbounded or $L$-rescaled composite
field unless its regulator normalization and a common source domain have
already been proved.
\end{proof}

\begin{remark}[Semigroup and covariance consequences]
\label{p12:rem:action-first-semigroup-consequence}
If the finite transfer is constructed from the same reflection-positive
action measure, finite-time transfer matrix elements are Euclidean cylinder
expectations.  Hence QFT4 control of the reflected/time-separated bounded
cylinder family implies the corresponding finite-time semigroup Cauchy
matrix elements without a second physical intertwining premise.  The same
interpolation theorem also controls the interacting part of local-flat
covariance whenever the chart-to-flat action difference is a uniformly
irrelevant/quasilocal defect.
\end{remark}


% BEGIN v2.36 SOURCE-UNIFORM CONTINUUM COMPRESSION
\begin{definition}[Source-uniform refinement certificate]
\label{p12:def:v236-surg}
The symbol \(\mathsf{SURG}\) denotes the hypotheses of
Theorem~\ref{p12:thm:v185-uniform-source-effective-action-stability} together
with the already declared refinement coherence and, for claims of local-flat
interacting covariance, the required uniform chart-to-flat defect bound.  It
is an analytic refinement certificate, not new sector physics.
\end{definition}

\begin{corollary}[Source-owned continuum package]
\label{p12:cor:v236-source-owned-continuum}
On a finite sector package whose bounded local source insertions belong to one
common source/action lineage, \(\mathsf{SURG}\) gives the canonical QFT4 local
state and all fixed mixed source moments/cumulants.  Finite reflection
positivity then gives QFT5, uniform bounded local representatives give the
bounded part of QFT7, and the chart-defect clause gives the interacting local
part of QFT6.  Any local field realized as a regulator-normalized derivative
of the same uniformly controlled source family belongs to the
\(\mathrm{QFT10}_{\rm source}\) class and its fixed mixed moments descend by
Cauchy differentiation.

This corollary does not cover an independently normalized unbounded composite,
light-ray/rapidity operator or other residue not proved to be a derivative of
the same source owner.  Those remain in
\(\mathrm{QFT10}_{\rm ext}\) and retain their explicit UCS/mixed-moment
certificate.  QFT8 nontriviality and QFT9 gap/controlled-gapless claims also
remain claim-specific.
\end{corollary}

\begin{corollary}[Quantum Ward/anomaly transport is dependent]
\label{p12:cor:v236-ward-dependent}
Let a finite regulated fermionic/action owner satisfy its exact Ward or
anomaly identity, including all Jacobian, contact and regulator terms.  If
every term lies in the \(\mathsf{SURG}\)-controlled source class or the
required \(\mathrm{QFT10}_{\rm ext}\) class, then
Theorem~\ref{p12:thm:finite-identity-preservation} transports that identity to
the continuum.  Thus continuum quantum Ward/anomaly control is a dependent
consequence of the finite regulated owner plus the continuum certificate; it
is not an independent physical sector or selector.
\end{corollary}
% END v2.36 SOURCE-UNIFORM CONTINUUM COMPRESSION



% Source Paper VII-C, original line 3329; inherited proof bodies unchanged.
\subsubsection{Scale-split action-first certificate}
\label{p12:subsec:g78-msrg-certificate}

The action-first route requires two different analytic coordinate systems.
The ultraviolet heat-kernel density concentrates near flat connections and is
not perturbative about product Haar.  After a quantitative crossover to a
strong-heat domain, a Haar-centered polymer expansion becomes effective.  The
following theorem fixes that infrared domain before the scale-split parent is
stated.

\begin{definition}[Heat-polymer activity norm]
\label{p12:def:heat-polymer-activity-norm}
Let $X_L$ be a finite oriented cell complex with edge set $E_L$ and two-face
set $F_L$.  Faces are adjacent when their boundary edge supports intersect.
For normalized Haar measure and the central heat kernel satisfying
$\partial_tK_t=\frac12\Delta K_t$, put
\[
 u_t(g):=K_t(g)-1,
 \qquad
 \varepsilon_*:=\sup_{L,f}\|u_{t_{f,L}}\|_\infty.
\]
Let $N_n(f)$ be the number of connected $n$-face polymers whose edge support
intersects that of $f$.  For $a>0$ define
\begin{equation}
\label{p12:eq:heat-polymer-activity}
 \mathcal Q_a
 :=\sup_{L,f}\sum_{n\ge1}N_n(f)e^{an}\varepsilon_*^n.
\end{equation}
\end{definition}

\begin{theorem}[Uniform source-stable strong-heat domain]
\label{p12:thm:source-stable-heat-polymer-domain}
Consider the finite action measure
\[
 d\mu_L(U)=Z_L^{-1}
 \prod_{f\in F_L}K_{t_{f,L}}(U_f)
 \prod_{e\in E_L}dU_e.
\]
If $\mathcal Q_a\le a$ with a strict margin, uniformly in level and volume,
then the partition function and every bounded local source functional admit
an absolutely convergent polymer expansion, uniformly for fixed boundary
holonomies.  The logarithm is analytic in all sources that consume less than
the strict activity margin, and its derivatives obey connected-tree bounds.
If the same activity norm carries an exponential spacetime-distance weight,
the connected-tree bound carries the corresponding decay weight.

If, more concretely,
\begin{equation}
\label{p12:eq:heat-polymer-geometric-count}
 N_n(f)\le B_{\rm geo}C_{\rm geo}^{n-1},
\end{equation}
then the explicit sufficient inequalities are
\begin{equation}
\label{p12:eq:heat-polymer-kp-criterion}
 \boxed{
 \frac{B_{\rm geo}\varepsilon_*e^a}
 {1-C_{\rm geo}\varepsilon_*e^a}\le a,
 \qquad
 C_{\rm geo}\varepsilon_*e^a<1.}
\end{equation}
For $SU(3)$,
\begin{equation}
\label{p12:eq:su3-heat-character-majorant}
 \varepsilon(t)
 \le R_0(t)
 :=\sum_{(p,q)\ne(0,0)}
 \left[\frac{(p+1)(q+1)(p+q+2)}2\right]^2
 \exp\!\left[-\frac t6
 (p^2+q^2+pq+3p+3q)\right].
\end{equation}
Thus \eqref{p12:eq:heat-polymer-kp-criterion} is a directly executable
representation-theoretic certificate once the heat-time normalization and
cell incidences are fixed.
\end{theorem}
\begin{proof}
Write every face factor as $1+u_{t_f}(U_f)$ and expand the product.  A selected
face set splits into connected components in the face-overlap graph, and
product Haar integration factorizes between components with disjoint edge
supports.  Hence the components form a hard-core polymer gas with activity
\[
 z_L(\gamma)=
 \int\prod_{f\in\gamma}u_{t_{f,L}}(U_f)
 \prod_{e\in E(\gamma)}dU_e,
 \qquad
 |z_L(\gamma)|\le\varepsilon_*^{|\gamma|}.
\]
Equation~\eqref{p12:eq:heat-polymer-activity} majorizes the incompatible-
polymer sum in the Koteck\'y--Preiss criterion
\cite{KoteckyPreiss1986}; the strict form gives absolute convergence and
analytic source differentiation.  Fixed boundary holonomies change only the
arguments of boundary face factors and not their sup norms.  The standard
tree-graph reorganization gives the connected-tree estimate, and inserting a
distance weight gives the weighted version.  Summing
\eqref{p12:eq:heat-polymer-geometric-count} as a geometric series gives
\eqref{p12:eq:heat-polymer-kp-criterion}.

For the final claim, Peter--Weyl theory gives
\[
 K_t(g)=1+\sum_{\lambda\ne0}d_\lambda
 e^{-tC_2(\lambda)/2}\chi_\lambda(g),
\]
and $|\chi_\lambda(g)|\le d_\lambda$.  Substitution of
$d_{p,q}=(p+1)(q+1)(p+q+2)/2$ and
$C_2(p,q)=(p^2+q^2+pq+3p+3q)/3$ proves
\eqref{p12:eq:su3-heat-character-majorant}.
\end{proof}

\begin{corollary}[Quasilocal defect stability inside the polymer domain]
\label{p12:cor:polymer-defect-stability}
Let $D_L=\sum_Xd_{X,L}$ and define
\[
 \|D_L\|_\mu
 :=\sup_x\sum_{X\ni x}
 e^{\mu\operatorname{diam}(X)}\|d_{X,L}\|_\infty.
\]
If every interpolation $S_L+sD_L$, $0\le s\le1$, remains in one strict
source-stable domain of Theorem~\ref{p12:thm:source-stable-heat-polymer-domain},
then every fixed bounded local $A$ satisfies
\begin{equation}
\label{p12:eq:polymer-defect-stability}
 |\langle A\rangle_{S_L+D_L}-\langle A\rangle_{S_L}|
 \le C_A\|D_L\|_\mu.
\end{equation}
Consequently an exact coarse-action remainder of order $O(a_L^\delta)$ in
this norm gives the same QFT4/RGSTAB rate.
\end{corollary}
\begin{proof}
Differentiate along the interpolation:
$\frac d{ds}\langle A\rangle_s=-\operatorname{Cov}_s(A,D_L)$.
The weighted two-insertion tree bound from
Theorem~\ref{p12:thm:source-stable-heat-polymer-domain}, followed by the sum
over $X$, bounds the covariance by $C_A\|D_L\|_\mu$ uniformly in $s$.
Integration proves \eqref{p12:eq:polymer-defect-stability}; the final statement
is Theorem~\ref{p12:thm:effective-action-stability-qft4}.
\end{proof}

\begin{proposition}[No one-regime Haar-polymer route to a finite-mass continuum]
\label{p12:prop:no-one-regime-haar-polymer-continuum}
For $SU(3)$, a level-uniform Haar-centered criterion of the form
\eqref{p12:eq:heat-polymer-activity} cannot persist when the face heat times
tend to zero.  Conversely, if the heat times stay in a fixed strict
strong-heat domain with graph-distance decay rate $m_{\rm graph}>0$ while the
physical cell size $a_L$ tends to zero, connected correlations at distinct
fixed physical points vanish as
\[
 O\!\left(e^{-m_{\rm graph}r/a_L}\right),
\]
so the physical inverse correlation length diverges.  A nontrivial finite-mass
continuum therefore requires a scale-split analysis with a correlation length
diverging in lattice units.
\end{proposition}
\begin{proof}
The compact-group heat asymptotic gives
$K_t(e)\sim c_Gt^{-\dim(G)/2}$; hence for $SU(3)$,
$\|K_t-1\|_\infty\ge K_t(e)-1\asymp t^{-4}\to\infty$ as $t\downarrow0$.
This violates every fixed positive bound in
\eqref{p12:eq:heat-polymer-activity}.  In the fixed strong-heat case, supports
at physical separation $r$ have graph distance of order $r/a_L$, so the
stated decay and infinite physical inverse correlation length follow.
\end{proof}

\begin{definition}[Scale-split G7/G8 multiscale certificate]
\label{p12:def:g78-msrg}
The certificate $\mathrm{G78\!-\!MSRG}$ consists of the following
five statements for one physical $SU(3)$ action/transfer family.
\begin{description}[leftmargin=4.2em,style=nextline]
\item[\textnormal{M0}] The spectral quadratic coefficient, compact-group face
normalization, exact coarse pushforwards, reflection-positive transfer, and
source insertions belong to one explicitly bridged action owner.
\item[\textnormal{M1}] Above a fixed physical matching scale, a gauge-fixed
small-field/large-field block RG has the correctly normalized asymptotically
free marginal flow and a weighted quasilocal remainder satisfying
\[
 \|R_{j+1}\|_{\rm UV}
 \le\rho\|R_j\|_{\rm UV}+Cg_j^p,
 \qquad 0<\rho<1,
\]
with summable large-field tails and the same estimates for the required local
sources.
\item[\textnormal{M2}] At one finite physical scale $j_*$, the exact blocked
multi-face action enters a strict source-stable infrared domain:
\[
 \mathcal Q_{\rm IR}(S_{j_*}^{\rm eff})<q_*.
\]
\item[\textnormal{M3}] Below $j_*$, exact blocking, boundary conditions,
sources, and irrelevant defects preserve a strict distance-weighted polymer
margin in physical Euclidean units.
\item[\textnormal{M4}] The same action is reflection positive, centered local
vectors are cyclic in the selected vacuum sector, and the transfer generator
satisfies a Paper-IV weighted-conjugation/DLOC estimate or an independently
proved uniformly quasilocal-logarithm criterion.
\end{description}
\end{definition}

The ultraviolet form of M1 is the action-specific completion demanded by the
multiscale lattice-gauge program; the rigorous finite-volume constructions of
Ba\l aban provide relevant theorem architecture but do not discharge M1--M2
for the present action \cite{Balaban1987RGI,Balaban1988RGII}.

\begin{theorem}[Scale-split common-parent implication]
\label{p12:thm:g78-msrg-consequence}
For the actual action-first sequence,
\[
 \boxed{
 \mathrm{G78\!-\!MSRG}
 \Longrightarrow
 \mathrm{G7b}+\mathrm{G8A}+\mathrm{G8B}+\mathrm{G8C}.}
\]
Equivalently, M0--M4 give bounded QFT4/RGSTAB, a positive physical transfer
gap on the selected cyclic sector, dynamical locality, and all-order
source-connected tree/Ursell estimates.  Non-Gaussianity and unbounded residue
field reconstruction remain separate QFT8/QFT10 obligations.
\end{theorem}
\begin{proof}
Iterating M1 controls all ultraviolet telescoping differences and source
derivatives.  M2 places the exact crossover action, including its generated
multi-face terms, inside the strict M3 domain.  Theorem~
\ref{p12:thm:source-stable-heat-polymer-domain} and Corollary~
\ref{p12:cor:polymer-defect-stability} then control the infrared blocks,
arbitrary connected source derivatives, and the exact effective-action
defects.  Exact composition gives refinement-path independence.  Reflection
positivity identifies physical-time cylinder correlations with matrix
elements of the transfer semigroup.  M3's fixed-physical-time decay excludes
spectral support below a positive threshold on centered local vectors, and
M4 cyclicity extends the exclusion to the selected vacuum-orthogonal sector.
The M4 weighted-conjugation estimate gives G8B.  None of these implications
forces a surviving higher connected continuum moment or the normalization and
Cauchy control required for an unbounded composite field.
\end{proof}



% Source Paper VII-C, original line 3560; inherited proof bodies unchanged.
\subsubsection{Quantitative norm interface for the ultraviolet--infrared crossover}
\label{p12:subsec:v198-rg-norm-interface}

The statements below refine the certificate of
Definition~\ref{p12:def:g78-msrg} without asserting that the actual blocked
action has entered its infrared domain.  The index $j$ denotes successive
blocking steps on a specified finite interval.  The physical evaluation
scale, the refinement cutoff and this index must be related by declared
maps; they are not identified by notation.  A small-field expansion alone
does not control the complete interaction norm.  Appendix~
\ref{p12:app:v198-norm-domain-diagnostics} gives the exact domain obstruction
and records the numerical comparison checkpoint.

\begin{definition}[Complete-action crossover norm interface]
\label{p12:def:v198-complete-action-norm-interface}
Work on one M0 action/transfer family.  On a branch whose retained CP-even
isotropic $F^2$ stiffness is a positive scalar $\kappa_j$, match that
coordinate by Paper~VII-A's heat-Hessian theorem:
\[
 t_j=h_{SU(3)}^{-1}(\kappa_j),\qquad
 S_j^{\rm heat}=\sum_f s_{t_j}(U_f),\qquad
 D_j=S_j^{\rm eff}-S_j^{\rm heat}.
\]
Fix the vacuum-constant convention and a local interaction decomposition
$D_j=\sum_Xd_{X,j}$ before evaluating its norm.  Thus
$\Pi_{F^2}D_j=0$ on the matched coordinate, but no bound on $D_j$ follows
from that identity.  An anisotropic branch must retain its stiffness
operator or face coordinates explicitly; it is not covered by replacing
them with an arbitrary scalar average.

Let $X$ be the actual finite face support of an interaction, let $X\sim f$
mean edge-support overlap with $f$, and set
\begin{equation}
\label{p12:eq:v198-crossover-defect-norm}
 r_j=\sup_f\sum_{X\sim f}e^{|X|/4}\|d_{X,j}\|_\infty.
\end{equation}
The supremum is over the full compact configuration domain, with uniformity
in the volumes, boundary holonomies and finer regulators required by the
claim at fixed physical matching scale.  The decomposition, its support
weights and its normalization are part of the norm interface.  Hessian
matching makes the heat representative unique on the stated branch; it
does not make an arbitrary decomposition of $D_j$ canonical.

An ultraviolet norm-transfer certificate consists of a bound
\begin{equation}
\label{p12:eq:v198-uv-ir-norm-transfer}
 r_j\le C_j^{\rm tr}\|R_j\|_{\rm UV}+\eta_j,
 \qquad C_j^{\rm tr},\eta_j\ge0,
\end{equation}
proved for the exact blocked action.  It includes the discarded large-field
contribution, all generated multi-face terms and the stated uniformity
domains.  Required bounded local sources must satisfy compatible bounds on
a fixed source polydisc.  Claims of spatial or temporal decay additionally
require the physical-distance weights used by their target theorem.
\end{definition}

\begin{proposition}[Finite recurrence bound after norm transfer]
\label{p12:prop:v198-finite-rg-recurrence}
Suppose, on $0\le j<N$, that the complete ultraviolet remainder satisfies
\[
 \|R_{j+1}\|_{\rm UV}\le\rho_j\|R_j\|_{\rm UV}+A_jg_j^4,
 \qquad \rho_j,A_j\ge0,
\]
and that \eqref{p12:eq:v198-uv-ir-norm-transfer} holds at $N$.  Then
\begin{equation}
\label{p12:eq:v198-finite-crossover-budget}
 r_N\le C_N^{\rm tr}\left[
 \left(\prod_{i=0}^{N-1}\rho_i\right)\|R_0\|_{\rm UV}
 +\sum_{n=0}^{N-1}A_ng_n^4
       \prod_{i=n+1}^{N-1}\rho_i\right]+\eta_N.
\end{equation}
Empty products are one.  The fourth-power forcing is an explicit hypothesis
about this norm on the complete blocked action, not a consequence of a
local Taylor coefficient.  No ultraviolet asymptotic estimate is extended
through a strongly coupled interval without a bound on that interval.
\end{proposition}
\begin{proof}
Induction on $N$ gives the bracketed bound on $\|R_N\|_{\rm UV}$.
Substitution into the norm-transfer inequality proves the claim.  In
particular the estimate is finite-interval and does not assume that an
asymptotically free indexing convention agrees with the blocking direction.
\end{proof}

\begin{proposition}[Heat activity bounded by the matched stiffness]
\label{p12:prop:v198-stiffness-activity}
In the normalization of
Theorem~\ref{p10:thm:v185-heat-hessian-inversion}, let
$\kappa=h_{SU(3)}(t)$.  For $0<\kappa<1/6$,
\begin{equation}
\label{p12:eq:v198-stiffness-activity}
 R_0(t)\le\overline\varepsilon(\kappa)
 :=\frac{6\kappa}{1-6\kappa}.
\end{equation}
Consequently, if $N_n(f)\le BC^{n-1}$ uniformly on the actual cell
presentation and $C e^{1/4}\overline\varepsilon(\kappa)<1$, the heat part
of the activity criterion at $a=1/4$ is bounded by
\begin{equation}
\label{p12:eq:v198-heat-budget}
 \overline Q(\kappa)
 =\frac{B e^{1/4}\overline\varepsilon(\kappa)}
 {1-C e^{1/4}\overline\varepsilon(\kappa)}.
\end{equation}
The constants $B,C$ belong to that cell presentation, not merely to a
comparison geometry or the compact momentum support.
\end{proposition}
\begin{proof}
Every nontrivial $SU(3)$ representation has $C_2\ge4/3$ by the displayed
Casimir formula in \eqref{p12:eq:su3-heat-character-majorant}.  Since
$\dim SU(3)=8$, the heat-Hessian formula gives
\[
 \kappa\ge\frac18\frac{(4/3)R_0(t)}{1+R_0(t)}
 =\frac{R_0(t)}{6(1+R_0(t))}.
\]
Rearrangement yields the first claim; substitution into the monotone
geometric-series majorant proves the second.
\end{proof}

\begin{proposition}[Conditional complete-action stopping budget]
\label{p12:prop:v198-conditional-stopping-budget}
In addition to the preceding hypotheses, suppose a polymer construction for
the actual multi-face action proves the complete activity domination
\begin{equation}
\label{p12:eq:v198-complete-activity-domination}
 \mathcal Q_{\rm IR}(S_j^{\rm heat}+sD_j;J)
 \le\overline Q(\kappa_j)+r_je^{r_j}+q_{\rm src},
 \quad 0\le s\le1,\quad J\in\mathbb D_{\rm src},
\end{equation}
uniformly on a common source polydisc, where $q_{\rm src}\ge0$ is the
source budget.  This premise must include mixed heat/defect polymers and
their incompatibility weights; it is not inferred from
\eqref{p12:eq:v198-crossover-defect-norm} alone.  If
\begin{equation}
\label{p12:eq:v198-stopping-budget}
 \overline Q(\kappa_j)+r_je^{r_j}+q_{\rm src}<\frac14,
\end{equation}
then the action enters the corresponding strict source-stable M2 domain.
For a fixed admissible source budget and $r\ge0$, put
\[
 m=\frac14-re^r-q_{\rm src}>0,\qquad
 \varepsilon_{\max}=\frac{me^{-1/4}}{B+Cm}.
\]
A sufficient stiffness bound is
\begin{equation}
\label{p12:eq:v198-stiffness-threshold}
 0<\kappa<\frac{\varepsilon_{\max}}{6(1+\varepsilon_{\max})}.
\end{equation}
\end{proposition}
\begin{proof}
Complete activity domination and the strict inequality give the stipulated
infrared criterion, with the same source domain.  Solving
$B e^{1/4}\varepsilon/(1-Ce^{1/4}\varepsilon)<m$ gives
$\varepsilon<me^{-1/4}/(B+Cm)$; the denominator is then positive.
Solving $6\kappa/(1-6\kappa)<\varepsilon_{\max}$ proves the final formula.
\end{proof}

\begin{remark}[Claim-specific premise reduction]
\label{p12:rem:v198-premise-reduction-scope}
An enhanced M1 which proves the full norm transfer, complete activity
domination and finite stopping inequality supplies M2 as a consequence.
The existence of those bounds is still an action-specific obligation.
For bounded local states, retain the covariance/comparison conditions of
Theorem~\ref{p12:thm:effective-action-stability-qft4}. For a fixed real
bounded-cylinder source tuple in positive states, Theorems
\ref{p12:thm:v202-bounded-source-domain} and
\ref{p12:thm:v202-bounded-source-inheritance} derive a common zero-free
domain and cumulant convergence from bounded state convergence, without
the full M3 package. The broader source/covariance domains in
Theorem~\ref{p12:thm:v185-uniform-source-effective-action-stability}
and actual matching-scale comparison remain separate obligations.  Positive
physical gap, deep-infrared distance decay, dynamical locality, unbounded
fields and non-Gaussianity keep their own hypotheses.  Likewise, centering
a cyclic algebra does not prove that a smaller selected local algebra is
cyclic.  Reflection-compatible pushforward preserves reflection positivity
of the same action but does not identify distinct transfer owners.
\end{remark}

\begin{definition}[Born--Hopf refinement-compatibility map]
\label{p12:def:bh-renormalization-map}
Let \(L\preceq L'\).  A Born--Hopf refinement-compatibility map is a completely positive unital comparison map
\[
\mathcal R_{L'\to L}:\A_{L'}(\mathcal O_{L'})\longrightarrow\A_L(\mathcal O_L)
\]
satisfying
\begin{equation}
\label{p12:eq:renormalization-expectation-compatibility}
\left|
\omega_{L'}(A_{L'})-\omega_L(\mathcal R_{L'\to L}A_{L'})
\right|
\le C_A a_L^\eta
\end{equation}
for local representatives of fixed continuum cylinder observables.
\end{definition}


\begin{definition}[Two-scale gapped refinement comparison]
\label{p12:def:two-scale-gapped-comparison}
For \(L\preceq M\), let \(J_{ML}:\mathcal H_L\to\mathcal H_M\) be the native refinement isometry on a gapped block.  The block satisfies two-scale gapped comparison if
\[
\|J_{ML}K_L-K_MJ_{ML}\|\le\eta_L,
\qquad
\|J_{ML}A_L-A_MJ_{ML}\|\le\xi_L(A),
\]
for every bounded local cylinder observable \(A\), uniformly for sufficiently fine \(M\succeq L\), with \(\eta_L\to0\) and \(\xi_L(A)\to0\), and if the non-vacuum transfer spectrum is separated from the vacuum by a uniform gap.  In this gapped-block criterion, \(K_L\) denotes the fixed-physical-time sector transfer used to isolate the vacuum projection; it is not the infinitesimal refinement transfer \(K_L^{\mathrm{ref}}=e^{-a_LH_L}\).
\end{definition}

\begin{theorem}[Two-scale gapped comparison with controlled ground block]
\label{p12:thm:two-scale-comparison-qft4}
In addition to Definition~\ref{p12:def:two-scale-gapped-comparison},
assume the fixed-physical-time transfers are positive self-adjoint
contractions with eigenspace projections $P_L,P_M$ at eigenvalue one,
ranks $r_L,r_M\ge1$, and non-ground spectrum bounded by $q<1$.
Write $\rho_L=P_L/r_L$, $\rho_M=P_M/r_M$ for the \emph{canonical
normalized ground-block states}. If the transfer defect is $\eta_L$,
observable defect $\xi_L(A)$, and
$\varepsilon_L=\eta_L/(1-q)<1$, then $r_M\ge r_L$ and
\[
 \boxed{\|\rho_M-J_{ML}\rho_LJ_{ML}^\dagger\|_1
 \le2\varepsilon_L+2(1-r_L/r_M)}.
\]
Consequently
\[
 |\omega_M(A_M)-\omega_L(A_L)|
 \le\|A_M\|[2\varepsilon_L+2(1-r_L/r_M)]+\xi_L(A).
\]
The QFT4 Cauchy property follows when these errors vanish uniformly
on directed tails, including the rank-ratio defect
$\sup_{M\succeq L}(1-r_L/r_M)\to0$.
In particular a cofinally constant ground multiplicity, not
necessarily rank one, suffices \emph{for the chosen normalized
block states}. The raw gap/intertwining conditions alone do not select
a physical state on a degenerate vacuum space.
\end{theorem}
\begin{proof}
Since $K_LP_L=P_L$, the intertwining bound implies
$\|(I-P_M)J_{ML}P_L\|\le\eta_L/(1-q)=\varepsilon_L$.
The inequality $\varepsilon_L<1$ makes $P_M$ injective on the
$r_L$-dimensional transported ground space, so $r_M\ge r_L$.
Let $Q=JP_LJ^\dagger$ and $Q'$ be the projection onto
$P_M(\operatorname{ran}Q)$. Principal-angle estimates give
$\|Q-Q'\|_1\le2r_L\varepsilon_L$.
Since $Q'\le P_M$ and has rank $r_L$,
$\|P_M/r_M-Q'/r_L\|_1=2(1-r_L/r_M)$.
The triangle inequality gives the first display, and
trace/operator-norm duality plus
$\|J^\dagger A_MJ-A_L\|\le\xi_L(A)$ gives the second.
Uniform pairwise comparison and common refinements yield the Cauchy
net precisely under the stated rank/state guards.
\end{proof}





% Source Paper VII-C, original line 3800; inherited proof bodies unchanged.
\subsubsection{Local-state convergence from same-side vacuum control}
\label{p12:sec:qft4-subladder}
The comparison subladder uses the following retained labels:
\begin{enumerate}[label=\textnormal{(QFT4\alph*)},leftmargin=3.9em]
\item \label{p12:item:qft4a} the directed comparison tower;
\item \label{p12:item:qft4b} the common bounded cylinder algebra;
\item \label{p12:item:qft4c} uniform representative bounds;
\item \label{p12:item:qft4d} on the positive-gap comparison branch, an
actually proved rank-one ground projection and uniform fixed-physical-time
contraction. This is not a premise on gap-free cluster-state existence.
\end{enumerate}

\begin{definition}[Cross-vacuum diagnostic and same-side comparison residuals]
\label{p12:def:local-vacuum-comparison-seminorm}
Let $J=J_{ML}:\mathcal H_L\to\mathcal H_M$ be the inherited isometry,
$\Omega_L,\Omega_M$ normalized vacuum vectors, and $A_L,A_M$ the actual
representatives of one bounded question. For a cross-level map $D$ retain
the diagnostic
\begin{equation}
\label{p12:eq:local-vacuum-seminorm}
 \|D\|_{\omega,\mathcal O}
 =\sup_{\|B\|,\|C\|\le1}
 |\langle B_M\Omega_M,D C_L\Omega_L\rangle|.
\end{equation}
This diagnostic alone is not a state-comparison norm: it can vanish when
the two vacuum lines are orthogonal. The sufficient residuals below instead
use the same transported vacuum on both sides:
\begin{equation}
\label{p12:eq:qft4-defects}
 \eta^K_{ML}:=\|(JK_L-K_MJ)\Omega_L\|,\qquad
 \eta^A_{ML}:=\|(JA_L-A_MJ)\Omega_L\|.
\end{equation}
Here $K_L=e^{-t_0(H_L-E_{0,L})}$ uses one fixed physical comparison time
$t_0>0$. The shifted Hamiltonian residual is
$\eta^H_{ML}:=\|[J(H_L-E_{0,L})-(H_M-E_{0,M})J]\Omega_L\|$.
All residuals are evaluated on the actual finite carriers and source marks.
\end{definition}

\begin{definition}[Sufficient same-side refinement compatibility]
\label{p12:def:admissible-local-vacuum-compatibility}
On a uniformly isolated rank-one vacuum branch require, cofinally in $L$
and uniformly for sufficiently fine $M\succeq L$,
\begin{equation}
\label{p12:eq:qft4e-transfer-lv}
 \eta^K_{ML}\longrightarrow0,
\end{equation}
and, for each fixed bounded cylinder product consumed by the claim,
\begin{equation}
\label{p12:eq:qft4f-observable-lv}
 \eta^A_{ML}\longrightarrow0.
\end{equation}
These are a sufficient QFT4e--QFT4f certificate, not consequences of the
cross-vacuum diagnostic or of smooth symbol approximation. A Hamiltonian
residual tending to zero may replace the transfer residual by the theorem
below. Direct state naturality remains an alternative without a gap.
\end{definition}

\begin{lemma}[QFT4g: quantitative vacuum-line alignment]
\label{p12:lem:qft4g-projection-stability}
Let $P_L=|\Omega_L\rangle\langle\Omega_L|$,
$P_M=|\Omega_M\rangle\langle\Omega_M|$, $K_M P_M=P_M$, and
$0\le K_M|_{P_M^\perp}\le q_*I$ with $q_*<1$, while $K_L\Omega_L=\Omega_L$.
Put $v=J\Omega_L$ and $r=\|(I-P_M)v\|$. Then
\begin{equation}
\label{p12:eq:qft4g-projection-stability}
 \boxed{r=\|(P_MJ-JP_L)\Omega_L\|
       \le {\eta^K_{ML}\over1-q_*}.}
\end{equation}
If the shifted fine Hamiltonian satisfies
$H_M-E_{0,M}\ge\gamma_*(I-P_M)$, then also
$r\le\eta^H_{ML}/\gamma_*$. These statements control the compared vacuum
line; they do not assert norm-intertwining on every excited vector.
\end{lemma}
\begin{proof}
Since $(I-K_M)v=(JK_L-K_MJ)\Omega_L$, spectral calculus on $P_M^\perp$
gives $(1-q_*)r\le\|(I-K_M)v\|$. For the Hamiltonian version use
$(H_M-E_{0,M})v=-[J(H_L-E_{0,L})-(H_M-E_{0,M})J]\Omega_L$ and the
spectral lower bound on $P_M^\perp$. The projection identity follows from
$P_L\Omega_L=\Omega_L$. Neither estimate divides by an unproved vacuum
overlap; instead it forces $|\langle\Omega_M,v\rangle|^2=1-r^2\to1$.
\end{proof}

\begin{theorem}[QFT4h--QFT4j: an error-carrying local-state bound]
\label{p12:thm:qft4hij-local-state-convergence}
Under the preceding vacuum-line hypotheses,
\begin{equation}
\label{p12:eq:qft4h-cauchy-bound}
 \boxed{|\omega_M(A_M)-\omega_L(A_L)|
 \le 2\|A_M\|r+\eta^A_{ML}
 \le {2\|A_M\|\eta^K_{ML}\over1-q_*}+\eta^A_{ML}.}
\end{equation}
One may replace the first fraction by
$2\|A_M\|\eta^H_{ML}/\gamma_*$ on the Hamiltonian-gap branch.
Uniform boundedness and the stated cofinal residual limits therefore make
each controlled cylinder expectation Cauchy and refinement-path independent.
The positive normalized limit descends through the actual compatible product
and null quotient to the common cylinder algebra.
\end{theorem}
\begin{proof}
The trace distance of the pure states of $v=J\Omega_L$ and $\Omega_M$ is
$2\sqrt{1-|\langle\Omega_M,v\rangle|^2}=2r$. Their expectations of $A_M$
differ by at most $2\|A_M\|r$. The remaining difference is
$\langle v,(A_MJ-JA_L)\Omega_L\rangle$, bounded by $\eta^A_{ML}$.
The previous lemma gives the residual bounds. For unrelated fine levels use
a common refinement twice. Apply this to every fixed product required for
positivity and representative independence; the compatible quotient then
inherits a positive normalized state.
\end{proof}

\begin{proposition}[Mixed-state comparison without a gap]
\label{p12:prop:v2211-mixed-state-refinement-bound}
For any finite densities $\rho_L,\rho_M$ and isometry $J$,
\[
 |\operatorname{Tr}(\rho_M A_M)-\operatorname{Tr}(\rho_L A_L)|
 \le\|A_M\|\|\rho_M-J\rho_LJ^\dagger\|_1
 +|\operatorname{Tr}\rho_L(J^\dagger A_MJ-A_L)|.
\]
Thus uniform vanishing of these two actual defects on the controlled cylinder
family is another sufficient QFT4 certificate, including degenerate symmetric
ground densities. It requires no positive-gap assertion.
\end{proposition}
\begin{proof}
Add and subtract $\operatorname{Tr}(J\rho_LJ^\dagger A_M)$ and use
trace/operator norm duality. The common-refinement argument is the same as
above. Direct asymptotic state naturality is a further valid route when this
stronger trace-norm comparison is unavailable.
\end{proof}

The older cross-vacuum diagnostic cannot replace these hypotheses. For
$J=I$, the diagonal algebra on $\mathbb C^2$ and alternating vacua $e_1,e_2$,
take $K_L=\operatorname{diag}(1,q)$ or $\operatorname{diag}(q,1)$ and
$A=\operatorname{diag}(1,-1)$. Every cross-vacuum transfer/observable
diagnostic is zero, although the state values alternate. The same-side
transfer residual is instead $1-q$ between opposite vacua and correctly
fails the sufficient certificate. This is a test of the comparison lemma,
not a claim that the native action realizes that alternating family.


% Source Paper VII-C, original line 3937; inherited proof bodies unchanged.
\subsubsection{Tier-2 estimates for bounded spectral-window data}
\label{p12:subsec:tier2-spectral-window-estimates}

\begin{definition}[Bounded spectral-window presentation class]
\label{p12:def:tier2-spectral-window-class}
Smooth symbols may obey spectral approximation estimates
\begin{equation}
\label{p12:eq:spectral-tail-bound}
 \|(1-P_L)a\|_\infty\le C_{a,s}L^{-s}\|a\|_{C^s}.
\end{equation}
For state convergence the separate Hamiltonian-on-vacuum estimate is
\begin{equation}
\label{p12:eq:hamiltonian-tail-bound}
 \eta^H_{ML}\le C_{H,s}L^{-s}.
\end{equation}
The latter is an actual interacting comparison estimate. It does not follow
from the former symbol tail. The compact momentum ceiling remains fixed.
\end{definition}

\begin{proposition}[Same-side Duhamel bound and Toeplitz no-go]
\label{p12:prop:tier2-qft4-estimates}
For shifted nonnegative Hamiltonians and a normalized coarse vacuum,
\begin{equation}
\label{p12:eq:tier2-transfer-estimate}
 \boxed{\eta^K_{ML}\le t_0\eta^H_{ML}.}
\end{equation}
The required observable comparison is the actual residual
\begin{equation}
\label{p12:eq:tier2-observable-estimate}
 \eta^A_{ML}=\|(JA_L-A_MJ)\Omega_L\|.
\end{equation}
Smooth symbol approximation alone does not prove its vanishing.
\end{proposition}
\begin{proof}
Writing the shifted Hamiltonians as $A_L,A_M$, Duhamel gives
\[
 Je^{-t_0 A_L}-e^{-t_0 A_M}J
 =-\int_0^{t_0}e^{-(t_0-u)A_M}(JA_L-A_MJ)e^{-uA_L}\,du.
\]
On $\Omega_L$, $e^{-uA_L}\Omega_L=\Omega_L$ and the other semigroup is
a contraction, proving the estimate without a stability assumption on a
cross-vacuum seminorm. For the symbol warning, the highest cutoff mode
$e_L=\sqrt{L+1}D^{L/2}_{L/2,L/2}$ and fixed symbol
$f=D^{1/2}_{1/2,1/2}$ satisfy
$fe_L=\sqrt{(L+1)/(L+2)}e_{L+1}$. The symbol tail is zero while the
operator intertwining defect on this mode tends to one. Whether the actual
vacuum occupies that leaking direction is a separate calculation; no
vacuum-residual lower bound is inferred from this example.
\end{proof}

\begin{theorem}[Actual same-side defects close bounded QFT4]
\label{p12:thm:tier2-closes-qft4-bounded}
Suppose the actual bounded cylinder representatives have
$\eta^A_{ML}\to0$ and $\eta^H_{ML}\to0$ cofinally, uniformly for finer
$M\succeq L$, and the actual fine rank-one vacuum satisfies a uniform
$\gamma_*>0$. Then
\[
 |\omega_M(A_M)-\omega_L(A_L)|
 \le {2\|A_M\|\over\gamma_*}\eta^H_{ML}+\eta^A_{ML}\longrightarrow0.
\]
This proves QFT4 on the controlled cylinder products. Symbol density and
quadrature do not supply these residuals. When a uniform gap is not claimed,
the mixed-state or direct state-naturality certificate can instead be used.
\end{theorem}
\begin{proof}
Apply the vacuum-line Hamiltonian estimate and the local-state bound, then
pass to a common refinement for path independence.
\end{proof}

\begin{proposition}[Positive finite quadrature exists and inherits spectral approximation error]
\label{p12:prop:tier2-quadrature}
Let \((X,\mu)\) be a compact probability space and let
\(V_L\subset C(X)\) be any finite-dimensional real function space containing
the constants.  There exist nodes \(x_{\alpha,L}\in X\) and positive
normalized weights \(w_{\alpha,L}\) such that
\[
\sum_\alpha w_{\alpha,L}v(x_{\alpha,L})=\int_Xv\,d\mu
\qquad\forall v\in V_L.
\]
In particular, for \(X=S^3\) take the real and imaginary parts of retained
Peter--Weyl products needed at level \(L\).  For a smooth \(f\) on \(S^3\),
if \(P_Lf\in V_L\) and
\(\|(1-P_L)f\|_\infty\le C_sL^{-s}\|f\|_{C^s}\), then
\begin{equation}
\label{p12:eq:tier2-quadrature-estimate}
\left|\sum_\alpha w_{\alpha,L}f(x_{\alpha,L})-\int_{S^3}f\,d\mu\right|
\le 2C_sL^{-s}\|f\|_{C^s}.
\end{equation}
Thus the existence of a positive exact quadrature is derived from finite-dimensional convexity rather than inserted as a position lattice.
\end{proposition}
\begin{proof}
Choose a basis \(1,v_2,\ldots,v_d\) of \(V_L\) and map \(x\mapsto(v_2(x),\ldots,v_d(x))\in\mathbb R^{d-1}\).  The moment vector \((\int v_2d\mu,\ldots,\int v_dd\mu)\) is the barycenter of the image of the probability measure and therefore lies in its convex hull.  By finite-dimensional convexity (Carath\'eodory's theorem), it is a convex combination of finitely many evaluation vectors.  The corresponding positive coefficients sum to one because the constant function is included, proving exact quadrature.  Apply exactness to \(P_Lf\); positivity and normalization bound both the quadrature and integral of \((1-P_L)f\) by its sup norm, giving the factor two in \eqref{p12:eq:tier2-quadrature-estimate}.
\end{proof}


\begin{corollary}[Quadrature-presentation invariance of the retained Yang--Mills Hessian]
\label{p12:cor:v185-ym-quadrature-presentation-invariance}
Let $\mathscr R_L$ be any finite real retained curvature module on a compact
reciprocal-position chart/quadrature domain $(X_L,\mu_L)$, let $q_C$ be the
Casimir-normalized invariant metric of Paper~VII-A,
Theorem~\ref{p10:thm:v185-heat-hessian-inversion}, and let
$V_L^{\rm YM}\subset C(X_L)$ contain the constants and every scalar density
\[
 x\longmapsto
 \sum_{\mu<\nu}q_C(F_{\mu\nu}(x),G_{\mu\nu}(x)),
 \qquad F,G\in\mathscr R_L.
\]
Any two positive normalized quadratures that are exact on
$V_L^{\rm YM}$ induce the same bilinear Yang--Mills quadratic form on
$\mathscr R_L$.  Therefore the finite Hessian on the retained physical
curvature module is unique even when its auxiliary node/weight presentation
is not.  On a face chart satisfying the hypotheses of Paper~VII-A,
Corollary~\ref{p10:cor:v185-common-heat-face-stiffness}, every such exact
presentation yields the same quadratic action and hence the same target
stiffness operator on the retained image; a unique physical position lattice
is not required.
\end{corollary}
\begin{proof}
For $F,G\in\mathscr R_L$ the displayed scalar density lies in
$V_L^{\rm YM}$.  Exactness of either quadrature therefore makes its weighted
sum equal to the same continuum/chart integral.  Polarization gives equality of the two
induced bilinear forms for every pair $F,G$.  The final statement is only a
change of coordinates on this already fixed finite bilinear form; it does not
promote the auxiliary quadrature nodes to fundamental position data.
\end{proof}

\begin{lemma}[Derived local-flat geometry and the actual covariance comparison]
\label{p12:lem:tier2-flat-chart-covariance}
Assume the reciprocal-position radius \(R_{x,L}\to\infty\).  Fix one positive reference length \(\ell_0\) already available from the finite geometric normalization and choose
\[
\rho_L:=\sqrt{R_{x,L}\ell_0}.
\]
Then \(\rho_L\to\infty\) and \(\rho_L/R_{x,L}=\sqrt{\ell_0/R_{x,L}}\to0\).  On these charts the metric and volume-form defects from the flat Euclidean model are \(O(\rho_L^2/R_{x,L}^2)\).  For a bounded local observable define the total expectation-covariance defect
\[
\varepsilon^{\rm cov}_L(A,g)
:=|\omega_L((gA)_L)-\omega_L(A_L)|.
\]
Then local flat \(E(4)\) covariance of the continuum state follows if and only if, in addition to the derived geometric exhaustion and representative convergence, \(\varepsilon^{\rm cov}_L(A,g)\to0\) for every fixed compactly supported local question and local Euclidean motion.
\begin{equation}
\label{p12:eq:tier2-flat-chart-estimate}
\text{geometry contributes }O(\rho_L^2/R_{x,L}^2),\quad
\begin{gathered}\text{canonical state term: controlled by action defect;}\\\text{actual representative comparison remains.}\end{gathered}
\end{equation}
\end{lemma}
\begin{proof}
Normal coordinates on a round sphere of radius \(R_{x,L}\) satisfy \(g_{ij}(y)=\delta_{ij}+O(|y|^2/R_{x,L}^2)\) and \(\sqrt{\det g}=1+O(|y|^2/R_{x,L}^2)\).  Choosing \(\rho_L=\sqrt{R_{x,L}\ell_0}\) gives a dimensionally consistent exhausting sequence with vanishing relative curvature error.  This proves the geometric part without an extra axiom.  A metric approximation alone does not compare the actual transported observable words. On an equivariant normalized full-ground/Gibbs branch, Theorem~\ref{p12:thm:native-bh-covariance} and Corollary~\ref{p12:cor:v227-covariance-provider} bound the total defect by the representative residual, plus the controlled state term on approximate branches. This supplies the missing sufficient implementation without an independent state selector. Arbitrary states or Euclidean motions lacking the actual comparison unitary retain their correlation comparison. Convergence of the displayed total defect is exactly the expectation-covariance assertion once the geometric errors vanish.
\end{proof}

\begin{lemma}[Conditional asymptotic locality from operator-window overlap]
\label{p12:lem:tier2-asymptotic-locality}
Let \(A_L(f)=W_{f,L}X_LW_{f,L}\) and
\(B_L(g)=W_{g,L}Y_LW_{g,L}\) be actual \emph{even}, gauge-admitted
bounded local representatives in a common finite operator representation,
where the window operators are self-adjoint contractions,
\(\|X_L\|\le M_A\) and \(\|Y_L\|\le M_B\) uniformly in \(L\).
Let \(d=\operatorname{dist}(\operatorname{supp}f,
\operatorname{supp}g)>0\). If their \emph{operator} overlap is
verified to obey
\[
 \|W_{f,L}W_{g,L}\|\le C_N(1+Ld)^{-N}
\]
with constants independent of \(L\), then
\begin{equation}
\label{p12:eq:tier2-asymptotic-locality}
\boxed{\|[A_L(f),B_L(g)]\|_{\omega_L}
 \le2M_AM_BC_N(1+Ld)^{-N}.}
\end{equation}
If the actual window operators have orthogonal supports, the commutator
vanishes exactly. The scalar smooth-window kernel estimate
\(|K_L(x,y)|\le C'_NL^3(1+Ld(x,y))^{-N}\) can help establish the
operator-overlap premise only after the source-defined operator realization,
normalization, support, constraints and uniform kernel-to-operator norm
estimates are supplied. It is not by itself an interacting locality theorem.
For odd fields a graded commutator and the actual parity/constraint
implementation are required.
\end{lemma}
\begin{proof}
The products factor as
\[
\begin{aligned}
A_LB_L&=W_fX_L(W_fW_g)Y_LW_g,\\
B_LA_L&=W_gY_L(W_gW_f)X_LW_f.
\end{aligned}
\]
Submultiplicativity and self-adjointness of the windows give
\[
\|[A_L,B_L]\|\le2M_AM_B\|W_fW_g\|.
\]
The selected-state/GNS seminorm is bounded by operator norm. Insert the
stated overlap estimate. Orthogonality gives both products zero.
No step infers an operator overlap from a scalar integral kernel in the
absence of a compatible actual operator construction.
\end{proof}


\textbf{Tier-2 consequence.}
Peter--Weyl approximation gives the symbol-level bounded sector, Proposition~\ref{p12:prop:tier2-quadrature} derives positive exact finite quadratures, and Lemma~\ref{p12:lem:tier2-flat-chart-covariance} derives the geometric local-flat exhaustion.  None of these geometric/symbol results proves QFT4 for the interacting tower.  QFT4 is closed only after the actual observable and Hamiltonian/semigroup cross-level defects tend to zero as in Theorem~\ref{p12:thm:tier2-closes-qft4-bounded}.  Likewise, spectral UCS requires the separate normalization and direct tree-domination hypotheses of Theorem~\ref{p12:thm:native-spectral-ucs}.


\begin{proposition}[RG flow as finite refinement bookkeeping, with identifiability guard]
\label{p12:prop:rg-not-new-input}
If the renormalization maps of Definition~\ref{p12:def:bh-renormalization-map}
exist and satisfy QFT4, then the running coefficients occurring in a
specified finite representative action can be \emph{bookkept} as coordinates
of that representative system and its refinement maps; they are not an
independent continuum postulate. QFT4 compatibility of normalized local
expectations does \emph{not}, on its own, uniquely determine every
numerical finite coupling. Such uniqueness additionally requires a fixed
operator/counterterm scheme and an injective coupling-to-observable map
(after quotienting invisible/redundant coordinates) on the declared query
class, or actual finite-action coefficients supplied by the existing owner.
\end{proposition}
\begin{proof}
Refinement maps relate the given finite representatives and therefore
specify how \emph{already known} action coordinates are transported.
They cannot in general invert a noninjective observable map. For example,
\(H\mapsto H+cI\) leaves every normalized finite Gibbs state
\(e^{-\beta H}/\operatorname{Tr}(e^{-\beta H})\) and all Heisenberg
commutators unchanged; the coefficient of the identity cannot be inferred
from these normalized local expectations, however exactly they satisfy
QFT4. On the declared injective quotient and fixed scheme, equality of
all determining expectation coordinates fixes the corresponding coupling
coordinates. This does not add a new continuum action input.
\end{proof}

\begin{theorem}[RG flow as resolution-boundary response]
\label{p12:thm:v2140-rg-boundary-response}
Assume the finite refinement maps of
Definition~\ref{p12:def:bh-renormalization-map} and the finite positive shell
calculus of Paper~I-C.  Let $k$ denote the location of a spectral/resolution
boundary, and let $\Gamma_k^{\partial}$ be the effective action on the
retained side after the complementary finite modes are integrated out.  On a
differentiable finite bosonic representative its exact modified-Legendre
response is
\begin{equation}
\label{p12:eq:v2140-rg-boundary-response}
 \boxed{
 \partial_{\ell}\Gamma_k^{\partial}
 =\frac12\operatorname{Tr}
 \left[(\Gamma_k^{\partial(2)}+R_k)^{-1}\partial_{\ell}R_k\right],}
 \qquad \ell=\log k,
\end{equation}
with the trace restricted to the actual finite support.  Whenever the
occurring finite theory has the graded boson/fermion/gauge-quotient owners of
Papers~III-F--IV, the Gaussian-order version is the corresponding finite
supertrace.  If
\[
 \Gamma_k^{\partial}=\sum_a g_a(k)\mathcal O_a+\cdots
\]
in a fixed retained operator basis, then
\begin{equation}
\label{p12:eq:v2140-beta-projection}
 \boxed{
 \beta_a(k):=\partial_{\log k}g_a(k)
 =\Pi_a\!\left(\partial_{\log k}\Gamma_k^{\partial}\right).}
\end{equation}
Thus running couplings and wave-function factors are coordinates of a moving
resolution-boundary response of the same finite representative system.
\end{theorem}

\begin{proof}
Paper~I-C proves the finite shell covariance and the exact modified Legendre
flow on the minimal affine support.  The blocking maps here identify the
retained support at scale $k$, so changing $k$ changes the eliminated shell
and therefore the effective boundary functional, with
\eqref{p12:eq:v2140-rg-boundary-response} as its differential response.
Paper~IV fixes the relative determinant signs/powers of finite bosonic and
paired-Grassmann blocks, while Paper~III-F supplies the finite gauge-orbit
quotient response when that owner occurs; replacing the ordinary trace by the
resulting finite supertrace is therefore bookkeeping of the already typed
blocks, not a new interaction.  Projection onto a fixed operator basis gives
\eqref{p12:eq:v2140-beta-projection} by definition of its coordinates.
\end{proof}

\begin{corollary}[Standard-Model gauge realization]
\label{p12:cor:v2140-sm-one-loop-boundary-response}
On the Paper--V-A \(\mathsf R_5^{\rm SSQFT}\) gauge branch, the
$F_i^2$ projection of the Gaussian resolution-boundary response is
\[
 \boxed{(b_1,b_2,b_3)=\left(\frac{41}{10},-\frac{19}{6},-7\right).}
\]
The universal factors $-11/3$, $2/3$, and $1/3$ are the graded logarithmic
heat response of the one-form/quotient, Weyl, and complex-scalar blocks,
respectively; they are not independent finite-representative inputs.
\end{corollary}

\begin{remark}[Scale boundary is not history time]
The parameter $k$ labels resolution/refinement.  The theorem does not identify
$\partial_{\log k}$ with physical time evolution, does not choose a
cosmological scale history $k(t)$, and does not determine whether a later
physical scale variable is Born-even or Born-reflected.  Those are separate
owner/history questions.
\end{remark}

\begin{remark}[Analytic closure domain]
The remaining interacting continuum content is concentrated in actual cross-level local-state/semigroup comparison (QFT4), interacting local-flat covariance (QFT6), the optional uniform-gap inequality \(\gamma_*>0\) (QFT9), and regulator-normalized mixed-moment convergence for unbounded residues (QFT10).  Peter--Weyl density, positive finite quadrature, finite reflection-positivity passage, bounded-cylinder regularity, and the conditional gap-transfer theorem are already derived.  Direct spectral UCS is a useful sufficient bound for field moments but requires a uniform residue envelope and an all-order tree/cumulant estimate; neither follows from finite dimensionality or pair mixing.
\end{remark}


% BEGIN EXACT-GRADED-BOUNDARY-RESPONSE v2.98

% Source Paper VII-C, original line 6413; inherited proof bodies unchanged.
% Former appendix marker: technical material is now placed in VII-D.


% Source Paper VII-C, original line 6415; inherited proof bodies unchanged.
% Former logical part: Appendix: G7/G8 Analytic Obstructions and Optional Certificates; integrated into the new thematic argument.




% Source Paper VII-C, original line 6419; inherited proof bodies unchanged.
\subsection{RG comparison diagnostics and a sufficient spectral criterion}
\label{p12:app:v199-rgas-diagnostics}

\begin{proposition}[Finite negative essential spectrum under compactness]
\label{p12:prop:v199-rgas-spectral-criterion}
On a Hilbert space of essential perturbations, suppose \(B_0\) is
self-adjoint, bounded below, and has compact resolvent, and \(K\) is
bounded and self-adjoint.  Then \(B=B_0+K\), on the domain of \(B_0\),
has only finitely many eigenvalues at or below zero, counting multiplicity.
\end{proposition}
\begin{proof}
A bounded perturbation preserves self-adjointness, a lower bound and
compact resolvent.  The increasingly ordered \(\lambda_j(B_0)\) tend to
\(+\infty\).  The min--max inequality
\(\lambda_j(B)\ge\lambda_j(B_0)-\|K\|\) proves the assertion.
\end{proof}
Applying the lemma to essential RG requires proving these properties
for its actual stability operator.  It bounds linear relevant and
marginal candidates with the main-text sign convention; marginal
nonlinear dynamics must still be treated.  Compactness of a spacetime
Laplacian alone does not establish these hypotheses.

\begin{proposition}[Classical response plateau with a fixed-point-shaped beta]
\label{p12:prop:v199-rgas-response-plateau}
For the comparison action \(F(R)=(16\pi G_0)^{-1}(R+\beta R^2)\),
with \(G_0,\beta,\xi>0\), define a tensor-response coupling along
\(R=\xi k^2\) by
\[
 g_{\rm resp}(k)=\frac{G_0k^2}{1+2\beta\xi k^2},\qquad
 g_\infty=\frac{G_0}{2\beta\xi}.
\]
Then \(\partial_{\log k}g_{\rm resp}
=2g_{\rm resp}(1-g_{\rm resp}/g_\infty)\), without any quantum block
integration in its definition.
\end{proposition}
\begin{proof}
Differentiate the rational function; both sides equal
\(2G_0k^2/(1+2\beta\xi k^2)^2\).
\end{proof}
This comparison explains why a response plateau cannot replace the
quantum intertwining identity.  It changes none of the supplied UV
geometry or full Laurent endpoint data.



% Source Paper VII-C, original line 6463; inherited proof bodies unchanged.
\subsection{Crossover norm diagnostics and conditional comparison checkpoint}
\label{p12:app:v198-norm-domain-diagnostics}

This appendix supports the norm interface in the main text.  Its examples
test invalid inference patterns; they do not refute the actual heat action.
Its numerical checkpoint uses the explicitly identified affine-$F_4$
comparison counts and does not prove that a native RG trajectory reaches it.

For $0<g\le1$, the family $F_g(Y)=g^4Y^4$ is $O(g^4)$ on each fixed
bounded $Y$ domain, while
\[
 \sup_{|Y|\le g^{-1}}|F_g(Y)|=1.
\]
Thus a small-field expansion in $U=\exp(gY)$ does not establish a uniform
bound on the full compact $U$ domain.  Nor does the absence of a quadratic
mismatch bound a quartic coefficient: $aY^4$ has zero Hessian at the origin
for every $a>0$.  These examples locate the missing norm-transfer estimate;
they do not determine the native regulator's forcing order.

For the affine-$F_4$ comparison of
Appendix~\ref{p12:app:optional-affine-f4-heat-certificate}, take
$B=76$, $C=300$ and $a=1/4$.  At zero defect and zero source budget,
\eqref{p12:eq:v198-stiffness-threshold} gives
\[
 \kappa_{\max}=2.1462414083\ldots\times10^{-4}.
\]
The reported handoff checkpoint uses
\[
\begin{gathered}
t_{\rm Hess}=h_{SU(3)}^{-1}(1/(2\pi^2)),\\
R_0(3t_{\rm Hess})=1.81144085249\ldots\times10^{-4}.
\end{gathered}
\]
The geometric-series heat budget is approximately $0.0190031153$, leaving
approximately $0.2309968847$ for defects and sources.  Under a proved
complete-activity bound with that heat budget, the zero-source defect
ceiling is approximately $0.190860649$, the positive solution of
$r e^r=0.2309968847\ldots$.  The rounding is illustrative: use the exact
strict inequality, with certified numerical error and source reserve, in
any stopping certificate.  A nonzero source budget consumes part of this
margin.  The universal zero-heat, zero-source ceiling for this particular
budget is $r<W(1/4)$; it is not a necessary physical threshold.

No assumption that every branch has a three-face collar is made.  Two-face
heat convolution is exact on its incidence domain; multi-face branch
interactions must remain in the complete activity construction.  The
momentum 24-cell's incidence is not used as the reciprocal-position
blocking incidence.


% Source Paper VII-C, original line 6508; inherited proof bodies unchanged.
\subsection{Valence-three obstruction to product heat-kernel projectivity}
\label{p12:app:valence-three-heat-projectivity-obstruction}

The shortest exact witness already occurs for $U(1)$.  With normalized Haar
measure and
\[
 K_t(\theta)=\sum_{n\in\ZZ}e^{-tn^2/2}e^{in\theta},
\]
let one eliminated edge angle $x$ occur in three face angles.  Then
\begin{equation}
\label{p12:eq:u1-three-face-heat-integral}
\begin{aligned}
 I_t(\alpha,\beta,\gamma)
 &:={1\over2\pi}\int_0^{2\pi}
 K_t(\alpha+x)K_t(\beta+x)K_t(\gamma+x)\,dx\\
 &=\sum_{n+m+\ell=0}
 e^{-t(n^2+m^2+\ell^2)/2}
 e^{i(n\alpha+m\beta+\ell\gamma)}.
\end{aligned}
\end{equation}
The constraint $n+m+\ell=0$ is a joint interaction.  In particular the
Fourier coefficient $(1,0,0)$ vanishes in
\eqref{p12:eq:u1-three-face-heat-integral}, whereas it is nonzero in every
product $K_{T_1}(\alpha)K_{T_2}(\beta)K_{T_3}(\gamma)$ with finite positive
$T_i$.  Thus no renormalization of three one-face heat times reproduces the
pushforward.

For $SU(3)$ the same obstruction is representation-theoretic.  Three-factor
Haar integration retains invariant tensors rather than one diagonal
convolution channel.  Since
$\mathbf3\otimes\mathbf3\otimes\mathbf3\supset\mathbf1$, the invariant
epsilon tensor produces a joint three-face term.  Restricting the external
holonomies to center elements makes the separation especially transparent:
Haar integration removes every single-variable nonzero-triality mode while
the joint triality mode $(1,1,1)$ survives.  This proves the scope statement
in Proposition~\ref{p10:prop:multiface-heat-blocking}.  Degree-two
convolution identities remain exact; what fails is their unrestricted
four-dimensional extrapolation.


% Source Paper VII-C, original line 6547; inherited proof bodies unchanged.
\subsection{Optional affine-\texorpdfstring{$F_4$}{F4} comparison member and explicit infrared margin}
\label{p12:app:optional-affine-f4-heat-certificate}

This section is an optional action-first comparison branch.  It is not a
claim that spectral refinement itself creates a physical spacetime lattice,
and it is not imported into the native theorem without M0.

Let $a_L=2^{-L}a_0$ and form the affine Coxeter arrangement
\[
 \mathscr H_L=
 \{x:\langle\alpha^\vee,x\rangle=ka_L\}_{
 \alpha^\vee\in\Phi^\vee(F_4),\ k\in\ZZ}.
\]
The level-$L$ hyperplanes are contained in those at level $L+1$, so the
alcove complexes form exact subdivisions.  They have finitely many local
types modulo affine-Weyl transport, and each two-cell is triangular.  The link
of an edge is a rank-three spherical Coxeter complex; the largest possibility
is $B_3$.  Its group has order $48$, and the three vertex orbits have sizes
\[
 48/8=6,
 \qquad 48/4=12,
 \qquad 48/6=8.
\]
Hence at most $26$ triangles meet an edge.  At most
$1+3(26-1)=76$ faces overlap one fixed triangle, so the face-overlap graph has
degree at most $75$.  Rooted plane-tree encoding then gives the safe bounds
\begin{equation}
\label{p12:eq:affine-f4-polymer-counts}
 B_{\rm geo}=76,
 \qquad C_{\rm geo}=4\cdot75=300.
\end{equation}

\begin{proposition}[Optional $D_4$ Voronoi constant-curvature comparison]
\label{p12:prop:v2158-d4-voronoi-constant-curvature}
On the standard flat $D_4$ Voronoi honeycomb, let $a$ be the edge length of a
regular $24$-cell.  Its four-volume and triangular two-face area are
\[
 V_{24}=2a^4,
 \qquad
 A_\triangle=\frac{\sqrt3}{4}a^2.
\]
For the $96$ triangular faces of one cell, let $b_f$ be the unit simple
bivector in the face plane.  The exact $24$-cell orbit identity is
\[
 \boxed{\sum_{f=1}^{96}b_f\otimes b_f=16I_{\Lambda^2\mathbb R^4}.}
\]
Since three $24$-cells meet at every triangular face of the
$\{3,4,3,3\}$ honeycomb, the global face weight
\[
 w_\triangle=\frac{3V_{24}}{16}=\frac38a^4
\]
reproduces the quadratic norm of every constant Lie-algebra-valued two-form.
Hence
\[
 \frac{w_\triangle}{A_\triangle^2}=2,
 \qquad
 \kappa_{D_4}=\frac1{\pi^2}
\]
on the native color normalization $g_3^2=2\pi^2$, and the matching group heat
time is
\[
 \boxed{t_{D_4}=h_{SU(3)}^{-1}(1/\pi^2)
 =4.4209988618076\ldots .}
\]
The honeycomb vertex figure is the tesseract $\{4,3,3\}$, so four triangular
faces meet at each edge.  On the equal-face-time rooted-tree comparison at
$z=1.1$,
\[
 n=vz(1+N)^2,\qquad
 F_N=(1+n)^3-1,
\]
\[
 x=uz(1+N+B)^2,\qquad
 F_B=\frac{(1+N+2x)^3-(1+N-x)^3}{3},
\]
with $u=e^{-2t_{D_4}/3}$ and $v=e^{-3t_{D_4}/2}$.  The pair
$(N_*,B_*)=(0.01,0.8)$ is strict:
\[
 F_N=0.004444017911\ldots<0.01,
 \qquad
 F_B=0.707432606033\ldots<0.8.
\]
This proposition is a flat constant-curvature comparison diagnostic only.  It
does not prove a $D_4$-supported exact quadrature for the complete retained
$S_x^3\times\mathbb R$ curvature module, does not identify the optional
honeycomb with the native transfer-first finite QCD owner, and does not
discharge M0 or APBI.
\end{proposition}
\begin{proof}
In the standard $24$-cell realization the $96$ triangular face bivectors form
one full reflection-group orbit.  Direct orbit summation gives
$\sum_f b_f\otimes b_f=16I$; tracing checks $96=6\cdot16$.  Sharing each
triangle among three cells gives the stated global weight.  The stiffness
then follows from Paper~VII-A's exact-cell formula, and monotonic
heat-Hessian inversion gives the displayed root.  The tesseract vertex figure
gives four faces per edge, and direct substitution gives the strict recursion
values.
\end{proof}

Before adding any optional equality between differently typed heat
parameters, the same-owner Hessian bridge of Paper~VII-A gives a distinct
source-matched group-heat value on an isotropic exact cell presentation:
\[
 \boxed{
 t_{\rm Hess}=h_{SU(3)}^{-1}\!\left(\frac1{2\pi^2}\right)
 =5.7532958466\ldots .}
\]
Using the representation-theoretic majorant of
Theorem~\ref{p12:thm:source-stable-heat-polymer-domain} gives
\[
 R_0(t_{\rm Hess})=0.405089\ldots .
\]
For the safe affine-$F_4$ bounds below and $a=1/4$, the explicit
Koteck\'y--Preiss sufficient threshold is
\[
 \varepsilon_{\rm crit}
 =\frac{ae^{-a}}{B_{\rm geo}+aC_{\rm geo}}
 =0.00128940527\ldots,
\]
which is reached by the $SU(3)$ majorant at
\[
 t_{\rm crit}=14.3159577\ldots .
\]
Hence the source-matched native point is not already in this particular
Haar-centered strong-heat domain.  This is a crossover diagnostic only: it
does not prove an RG trajectory and does not replace M1--M2.

Now add the explicitly optional standard heat-time member
\[
 \mathsf{HK}_{\rm std}:\qquad
 t_f=t_*=g_3^2=2\pi^2.
\]
The equality is a regulator convention on this branch; Paper~VII-A's
Remark~\ref{p10:rem:three-distinct-heat-parameters} proves that it is not a
source-native consequence.  For $r=p+q$, the two fundamental representations
contribute
\[
 18e^{-4\pi^2/3}
 \simeq3.4688554154\times10^{-5}.
\]
For $r\ge2$,
\[
 C_2(p,q)\ge {r^2\over4}+r,
 \qquad
 d_{p,q}^2\le{(r+2)^6\over4},
\]
and there are $r+1$ pairs with sum $r$.  Therefore
\begin{equation}
\label{p12:eq:optional-standard-heat-r0-bound}
\begin{aligned}
 R_0(2\pi^2)
 &\le18e^{-4\pi^2/3}
 +\sum_{r=2}^\infty
 { (r+1)(r+2)^6\over4}
 e^{-\pi^2(r^2/4+r)}\\
 &<3.4689\times10^{-5}.
\end{aligned}
\end{equation}
For an exact arithmetic enclosure, use
$\pi>103993/33102$ and finite Taylor lower sums for each positive
exponential.  Sixty terms bound the fundamental contribution and the tail
ratio, and eighty terms bound the first tail term.  They give respectively
$3.468855433\times10^{-5}$, $4.250863\times10^{-10}$, and a subsequent-term
ratio below $1.154\times10^{-9}$, proving
\eqref{p12:eq:optional-standard-heat-r0-bound} without a truncated
representation sum.

Choose $a=1/4$.  Using $e^{1/4}<1.285$, put
$x=R_0(2\pi^2)e^{1/4}<4.4576\times10^{-5}$.  Equations
\eqref{p12:eq:heat-polymer-kp-criterion} and
\eqref{p12:eq:affine-f4-polymer-counts} then give
\[
 300x<0.013373,
 \qquad
 {76x\over1-300x}<0.003434<{1\over4}.
\]
Thus the bare one-face action on this optional member has unallocated
activity margin greater than $0.2465$.  This closes the bare infrared
inequality only.  M2 must still put the generated multi-face crossover action
and all required sources inside that margin, and M1/M4 remain independent.


% Source Paper VII-C, original line 6728; inherited proof bodies unchanged.
\subsection{Comparison diagnostic: fixed strong heat under a shrinking position-lattice regulator}
\label{p12:app:fixed-strong-heat-scaling-boundary}


oindent	extbf{Scope.}  This subsection is a conventional comparison diagnostic for an auxiliary position-lattice branch with physical cell size \(a_L	o0\).  It is not a native obstruction to the Born continuum, whose ultraviolet momentum ceiling remains fixed while reciprocal-position extent grows.

Suppose a strict level-independent strong-heat polymer estimate gives
$|\langle A_L;B_L\rangle_c|\le C_{A,B}e^{-m_{\rm graph}d_L}$ with
$m_{\rm graph}>0$.  If the physical cell size is $a_L$ and the supports stay
a fixed physical distance $r>0$ apart, then $d_L\asymp r/a_L$ and
\[
 |\langle A_L;B_L\rangle_c|
 \le C_{A,B}e^{-m_{\rm graph}r/a_L}\longrightarrow0.
\]
The physical inverse correlation length is of order
$m_{\rm graph}/a_L\to\infty$.  A finite positive continuum mass instead
requires $m_{{\rm graph},L}\sim a_Lm_{\rm phys}$, so the correlation length
diverges in lattice units.  This is the critical scaling delegated to M1--M2,
not a consequence of the fixed strong-heat theorem.


% Source Paper VII-C, original line 9285; inherited proof bodies unchanged.
\subsection{Uniform momentum-factor tests across a supplied refinement}
\label{p12:sec:v226-refinement-coefficients}
\begin{proposition}[A fixed momentum inverse lifts uniform operator bounds]
\label{p12:prop:v226-refinement-coefficients}
For any admitted refinement family with frozen intrinsic momentum factor,
let actual residuals satisfy
$A_n\in\operatorname{End}(\mathcal K_n)\otimes\mathcal A_D$.
Then
\[
 A_n=\sum_{j=1}^{137}q_j(A_n)\otimes D_j,\qquad
 \|A_n\|\le\sum_j\|q_j(A_n)\|\,\|D_j\|.
\]
If the right-hand side tends to zero, the complete residual tends to zero
in operator norm, independently of the growing outer dimensions. The same
statement holds for each actually consumed derivative residual.
\end{proposition}
\begin{proof}
Use Corollary~\ref{p1b:cor:v226-source-jets} on every actual finite carrier.
The 137 duals belong to one fixed finite momentum algebra, so their norms
are independent of the refinement index. The asserted convergence follows
from the printed operator-valued coefficient bound.
\end{proof}
Scalar expectations on a nonseparating state family do not imply these
operator-valued norm hypotheses. The new degree-six coordinate presentation
therefore does not prove uniform position-band estimates, strict locality,
reflection positivity, a continuum mass gap, limiting scattering or a
renormalized physical current. Each claim retains its existing source,
state, action, measure and ordered-limit certificates. The compact principal
bundle descent problem also remains the actual one-bit/pullback frontier
of Theorem~\ref{p1b:thm:v226-principal-two-types}, not a result inferred from
stable bundle counts or the completed mixed-$U(2)$ classification.

% END FRAMEWORK v2.226: KINEMATIC CLOSURES

% BEGIN FRAMEWORK v2.227: DOWNWARD CLOSURES


\section{Selector subsystem proofs and finite theory criteria}
\label{p13:sec:v237-d6}
% Source Paper VII-C, original line 6748; inherited proof bodies unchanged.
% Former logical part: Appendix: Selector Subsystems and Ownership Boundary; integrated into the new thematic argument.

\noindent\textbf{Orientation.}
The appendix classifies ordinary theories as selector presentations of the finite quantum calculus already owned by Paper~IV; it does not establish a second finite-quantum foundation.  A terminal ownership note records that finite nuclear closed-shell and split-grammar readouts belong to Paper~VII-B and are not part of the Paper-VII-C continuum certificate.

% BEGIN FINITE-INTERNAL-MOMENTUM-TRACE v2.102

% Source Paper VII-C, original line 6754; inherited proof bodies unchanged.
\subsection{Finite-internal extension of the momentum-ball trace}
\label{p12:app:v2102-finite-internal-momentum-trace}

\begin{corollary}[Factorized finite-internal momentum-ball trace]
\label{p12:cor:v2102-finite-internal-momentum-trace}
Assume a downstream sector has the explicitly factorized presentation
\[
 \mathcal P_x(m)\otimes V_{\rm int},
\]
with fixed finite inherited internal carrier $V_{\rm int}$ as in
Corollary~\ref{p12:cor:v259-resolution-growth-no-internal-species}.  Let
$B\in\End(V_{\rm int})$ be fixed and let $F$ be any fixed noncommutative
polynomial in the scaled spatial-transport triple.  Then the exact finite trace
factorization and Theorem~\ref{p12:thm:spatial-transport-momentum-ball-trace}
give
\[
 \boxed{
 \lim_{a\to\infty}\frac1{V_a}
 \operatorname{Tr}_{\mathcal P_x(m(a))\otimes V_{\rm int}}
 \!\left(F(\Pi^{(a)})\otimes B\right)
 =\operatorname{Tr}_{V_{\rm int}}(B)
 \int_{|p|\le\hbar\Lambda_F}
 \frac{d^3p}{(2\pi\hbar)^3}F_{\rm comm}(p).}
\]
In particular, for an internal projector $P$ of rank $r$ the external mode
density in that selected finite internal block is multiplied by $r$.
\end{corollary}

\begin{proof}
At every finite level,
\[
 \operatorname{Tr}_{\mathcal P_x\otimes V_{\rm int}}
 \left(F(\Pi^{(a)})\otimes B\right)
 =\operatorname{Tr}_{\mathcal P_x}F(\Pi^{(a)})\,
  \operatorname{Tr}_{V_{\rm int}}B.
\]
The internal factor is independent of reciprocal-position refinement, while
the external normalized trace converges by
Theorem~\ref{p12:thm:spatial-transport-momentum-ball-trace}.  Taking the limit
proves the result.
\end{proof}

\begin{remark}[Scope]
\label{p12:rem:v2102-finite-internal-momentum-trace-scope}
This is a trace-level product-sector corollary.  It does not assert
factorization of the physical state or measure, an interacting thermal
quasiparticle spectrum, a Debye/screening coefficient, a new particle species,
or a new RG flow.  Mixed $x$--internal operators and nonfactorized
states/measures retain their existing analytic obligations.
\end{remark}
% END FINITE-INTERNAL-MOMENTUM-TRACE v2.102


% Source Paper VII-C, original line 6863; inherited proof bodies unchanged.
\subsection{Finite selector subsystems and the Born--Hopf theory criterion}
\label{p12:app:finite-selector-subsystems}


\begin{table}[htbp]
\centering
\small
\begin{tabularx}{\textwidth}{@{}p{0.55\textwidth}X@{}}
\toprule
\textbf{Selector-subsystem step} & \textbf{Discharged by} \\
\midrule
Paper-IV finite carrier import & Definition~\ref{p12:def:finite-bh-carrier-algebra}. \\
Selector datum & Definition~\ref{p12:def:selector-datum}. \\
Accessible commutant algebra & Definition~\ref{p12:def:accessible-commutant-algebra}. \\
Selector presentation over the Paper-IV envelope & Proposition~\ref{p12:thm:finite-selector-commutant-subsystem}. \\
Imported Paper-IV selector decomposition & Proposition~\ref{p12:thm:wedderburn-selector-form}. \\
Hidden degrees are multiplicity spaces & Corollary~\ref{p12:cor:hidden-dof-multiplicity}. \\
Order-reversing selector lattice & Theorem~\ref{p12:thm:order-reversing-selector-lattice}. \\
Subsystem logic & Remark~\ref{p12:rem:subsystem-logic}. \\
\bottomrule
\end{tabularx}
\caption[Finite selector subsystem ledger]{Finite selector subsystem ledger.}
\label{p12:tab:finite-selector-subsystem-ledger}
\TableNotes{The appendix records how ordinary theories appear as finite subsystem presentations rather than competitors to the Born--Hopf framework.  A selector narrows the accessible algebra; it does not enlarge the finite carrier by fiat.}
\end{table}

This appendix records the finite subsystem principle used implicitly throughout the series.  Its finite algebra, state calculus, Born rule, channel theory, conservation law, gauge projection, thermal structure, and bosonic boundary defect are imported from Paper~IV.  The Paper-VII-C contribution is the classification of supplied selector data and the determination of which selected finite presentations admit the controlled continuum reconstruction of the main text.  The principle is not that every mathematically writable continuum theory is fundamental.  It is that every admissible finite physical theory with specified energy, symmetry, gauge, environment, or observable restrictions is represented inside the Paper-IV carrier by a selector algebra, a commutant, an invariant compression, or a controlled continuum refinement.

Thus ordinary theories are not competitors to the Born--Hopf framework.  They are sector presentations:
\[
\begin{gathered}
\text{finite carrier}+\text{selector data}
\\
\Longrightarrow
\text{subsystem / quotient / fixed-point sector / continuum reconstruction}.
\end{gathered}
\]
The GP2 filter is essential: a sector is derived only when the selector data are native to the framework or already produced by earlier papers.  If the selector is supplied externally, the sector is selected or interface-level rather than fundamental.


% Source Paper VII-C, original line 6902; inherited proof bodies unchanged.
\subsubsection{Selector data and accessible algebras}

\begin{definition}[Imported Paper-IV finite carrier convention]
\label{p12:def:finite-bh-carrier-algebra}
Paper~IV, Theorem~\ref{p4:thm:full-finite-observable-algebra}, fixes, before any sector, apparatus, environment, or superselection
selector is supplied, the maximal resolved finite operator envelope
\begin{equation}
\label{p12:eq:maximal-carrier-algebra}
\mathcal A_C:=\End(\mathcal H_C).
\end{equation}
Paper~VII-C adopts this as inherited notation.  A finite physical presentation in
this appendix is a selected subpresentation of this Paper-IV pair; the full
algebra and its state/Born calculus are not rederived here.
\end{definition}

\begin{definition}[Selector datum]
\label{p12:def:selector-datum}
A selector datum on \(\mathcal H_C\) is a finite collection
\[
\mathcal S=\{S_1,\ldots,S_k\}\subseteq\mathcal A_C
\]
or a finite compact symmetry representation
\[
U:G\to U(\mathcal H_C)
\]
together with any selected Hamiltonian, charge, projection, gauge condition, environment, pointer algebra, transfer operator, boundary condition, or local comparison datum that is intended to define a physical subsystem.

The selector algebra generated by \(\mathcal S\) is
\begin{equation}
\label{p12:eq:selector-algebra}
\mathcal B_{\mathcal S}:=C^\ast(\mathcal S)\subseteq\mathcal A_C.
\end{equation}
If a compact group representation \(U:G\to U(\mathcal H_C)\) is included, then \(\mathcal B_{\mathcal S}\) also contains the finite-dimensional algebra generated by \(U(G)\).
\end{definition}

\begin{definition}[Accessible commutant algebra]
\label{p12:def:accessible-commutant-algebra}
The observable algebra compatible with the selector datum \(\mathcal S\) is the commutant
\begin{equation}
\label{p12:eq:accessible-commutant}
\mathcal A_{\mathcal S}
:=
\mathcal B_{\mathcal S}'
=
\{a\in\mathcal A_C:[a,b]=0\ \forall b\in\mathcal B_{\mathcal S}\}.
\end{equation}
It is the algebra of finite observables that do not disturb, violate, or resolve beyond the selected data.
\end{definition}

\begin{proposition}[Selector presentation over the Paper-IV envelope]
\label{p12:thm:finite-selector-commutant-subsystem}
Let \(\mathcal H_C\) carry the Paper-IV finite quantum envelope and let
\(\mathcal S\) be a selector datum.  Then
\[
(\mathcal H_C,\mathcal A_{\mathcal S})
\]
is a selected finite presentation of that envelope.  If
\(P\in Z(\mathcal B_{\mathcal S})\) is a central selector-sector projection,
then
\begin{equation}
\label{p12:eq:compressed-selector-subsystem}
(P\mathcal H_C,\;P\mathcal A_{\mathcal S}P)
\end{equation}
is the corresponding compressed sector presentation.
\end{proposition}

\begin{proof}
Paper~IV, Theorem~\ref{p4:thm:full-finite-observable-algebra} and Lemma~\ref{p4:lem:proper-subalgebras-require-selectors}, proves that the unreduced finite algebra is the full endomorphism
envelope and that every proper non-scalar subalgebra carries selector data.
The commutant and central compression displayed here therefore classify a
supplied selector inside that already closed finite quantum calculus; they do
not define a new algebraic foundation.
\end{proof}

\begin{proposition}[Imported Paper-IV selector decomposition]
\label{p12:thm:wedderburn-selector-form}
For the supplied finite selector algebra \(\mathcal B_{\mathcal S}\), the
Paper-IV selector-subalgebra result, Lemma~\ref{p4:lem:proper-subalgebras-require-selectors}, gives a finite decomposition
\begin{equation}
\label{p12:eq:wedderburn-carrier-decomposition}
\mathcal H_C\cong\bigoplus_{\alpha}V_\alpha\otimes M_\alpha
\end{equation}
with
\begin{equation}
\label{p12:eq:selector-wedderburn}
\mathcal B_{\mathcal S}\cong
\bigoplus_{\alpha}\End(V_\alpha)\otimes\mathbf1_{M_\alpha},
\end{equation}
and
\begin{equation}
\label{p12:eq:commutant-wedderburn}
\mathcal A_{\mathcal S}=\mathcal B_{\mathcal S}'\cong
\bigoplus_{\alpha}\mathbf1_{V_\alpha}\otimes\End(M_\alpha).
\end{equation}
Paper~VII-C uses the multiplicity spaces \(M_\alpha\) only as selector-classification data.
\end{proposition}

\begin{proof}
This is the finite Wedderburn decomposition already used in Paper~IV, Lemma~\ref{p4:lem:proper-subalgebras-require-selectors}, to prove
that proper finite subalgebras require selectors.  The present proposition
records the imported normal form for the continuum-witness classification.
\end{proof}

\begin{corollary}[Hidden degrees of freedom are multiplicity spaces]
\label{p12:cor:hidden-dof-multiplicity}
A selected finite subsystem contains unresolved hidden degrees of freedom precisely when at least one multiplicity space \(M_\alpha\) satisfies
\[
\dim M_\alpha>1.
\]
If every \(M_\alpha\) is one-dimensional, the commutant is abelian.  If the commutant is only \(\mathbb C\mathbf1\), the selector datum leaves no nontrivial finite subsystem.
\end{corollary}

\begin{proof}
By Proposition~\ref{p12:thm:wedderburn-selector-form}, all accessible noncommutative freedom after selection is carried by the matrix factors \(\End(M_\alpha)\).
\end{proof}

\begin{theorem}[Order-reversing selector lattice]
\label{p12:thm:order-reversing-selector-lattice}
If two selector data satisfy
\[
\mathcal B_{\mathcal S_1}\subseteq\mathcal B_{\mathcal S_2},
\]
then their accessible algebras satisfy
\begin{equation}
\label{p12:eq:order-reversing-selector}
\mathcal A_{\mathcal S_2}
=
\mathcal B_{\mathcal S_2}'
\subseteq
\mathcal B_{\mathcal S_1}'
=
\mathcal A_{\mathcal S_1}.
\end{equation}
Thus adding selector assumptions reduces the accessible subsystem.
\end{theorem}

\begin{proof}
If \(a\) commutes with every element of the larger algebra \(\mathcal B_{\mathcal S_2}\), then it commutes with every element of the smaller algebra \(\mathcal B_{\mathcal S_1}\).  Hence
\[
\mathcal B_{\mathcal S_2}'\subseteq\mathcal B_{\mathcal S_1}'.
\]
\end{proof}

\begin{remark}[Subsystem logic]
\label{p12:rem:subsystem-logic}
The framework therefore organizes finite theories by a Galois-type relation:
\[
\text{more assumptions}
\quad\Longleftrightarrow\quad
\text{smaller accessible algebra}.
\]
This is the precise sense in which finite physical theories appear as Born--Hopf subsystems.  A theory is specified by what it selects, what it ignores, and what commutant remains.
\end{remark}


% Source Paper VII-C, original line 7056; inherited proof bodies unchanged.
\subsubsection{Energy, symmetry, and finite superselection}

\begin{definition}[Energy--symmetry selector]
\label{p12:def:energy-symmetry-selector}
An energy--symmetry selector consists of a self-adjoint finite Hamiltonian
\[
H=H^\ast\in\End(\mathcal H_C)
\]
and a compact unitary symmetry representation
\[
U:G\to U(\mathcal H_C).
\]
It is exact when
\begin{equation}
\label{p12:eq:exact-energy-symmetry}
[H,U(g)]=0
\qquad
\forall g\in G.
\end{equation}
\end{definition}

\begin{theorem}[Finite energy--symmetry sector theorem]
\label{p12:thm:finite-energy-symmetry-sector}
Let \((H,G)\) be an exact energy--symmetry selector.  Then
\[
\mathcal H_C
\cong
\bigoplus_{\lambda,\rho}
V_\rho\otimes M_{\lambda,\rho},
\]
where \(\lambda\) runs over energy eigenvalues and \(\rho\) over irreducible \(G\)-types.  The energy-preserving and \(G\)-invariant observable algebra is
\begin{equation}
\label{p12:eq:energy-symmetry-observable-algebra}
\mathcal A_{H,G}
=
\{a\in\End(\mathcal H_C):[a,H]=0,\ [a,U(g)]=0\ \forall g\in G\}
\end{equation}
and has the finite form
\begin{equation}
\label{p12:eq:energy-symmetry-wedderburn}
\mathcal A_{H,G}
\cong
\bigoplus_{\lambda,\rho}
\mathbf1_{V_\rho}\otimes\End(M_{\lambda,\rho}).
\end{equation}
\end{theorem}

\begin{proof}
Since \(H\) commutes with \(U(G)\), the finite spectral decomposition of \(H\) and the finite isotypic decomposition of \(U\) may be simultaneously refined.  Applying Proposition~\ref{p12:thm:wedderburn-selector-form} to the algebra generated by \(H\) and \(U(G)\) gives the result.
\end{proof}

\begin{remark}[Finite generator and conservation status]
\label{p12:cor:finite-noether}
Paper~IV now obtains finite one-parameter unitary generators from the
finite-dimensional generator theorem rather than using finite Stone as an
independent physical input.  Its conservation statement remains the typed
conditional theorem: if a finite self-adjoint generator \(Q\) commutes with
the Hamiltonian, \([H,Q]=0\), then
\[
\alpha_t(Q)=e^{itH}Qe^{-itH}=Q.
\]
On a native same-owner branch the symmetry-of-dynamics antecedent may already
be derived upstream; for an externally supplied \(H,Q\) it remains an explicit
hypothesis.  Paper~VII-C uses the resulting finite conservation law only to
identify conserved selector variables and their controlled continuum limits.
\end{remark}

\begin{definition}[Stabilizer after symmetry breaking]
\label{p12:def:stabilizer-after-breaking}
If a selected Hamiltonian or order datum \(X\) does not preserve a full compact symmetry \(G\), the surviving stabilizer is
\begin{equation}
\label{p12:eq:surviving-stabilizer}
G_X
:=
\{g\in G:U(g)XU(g)^{-1}=X\}.
\end{equation}
\end{definition}

\begin{proposition}[Symmetry breaking as selector reduction]
\label{p12:prop:symmetry-breaking-selector-reduction}
Symmetry breaking in a finite Born--Hopf subsystem is the replacement
\[
G\longrightarrow G_X.
\]
The exact selected subsystem is determined by the selector algebra generated by the order datum \(X\) together with the surviving symmetry representation \(U(G_X)\), not by the broken part of \(G\).
\end{proposition}

\begin{proof}
By definition, precisely the elements of \(G_X\) preserve the selected datum \(X\).  Only these transformations commute with the selected structure and therefore survive as exact symmetries of the subsystem.
\end{proof}


% Source Paper VII-C, original line 7147; inherited proof bodies unchanged.
\subsubsection{Gauge selectors and invariant physical sectors}

\begin{definition}[Imported Paper-IV finite gauge selector]
\label{p12:def:finite-gauge-selector}
Let \(G\) be compact and let \(U:G\to U(\mathcal H_C)\) be a finite unitary
representation.  Paper~IV, Theorem~\ref{p4:thm:finite-gauge-constraint-closure}, supplies the Haar projector and invariant algebra
\begin{equation}
\label{p12:eq:gauge-haar-projection}
P_G:=\int_G U(g)\,dg,
\end{equation}
\begin{equation}
\label{p12:eq:gauge-physical-carrier}
\mathcal H_C^G:=P_G\mathcal H_C,
\end{equation}
\begin{equation}
\label{p12:eq:gauge-invariant-observables}
\mathcal A_C^G:=\{a\in\mathcal A_C:U(g)aU(g)^{-1}=a\ \forall g\in G\}.
\end{equation}
These formulas are inherited finite-quantum data; Paper~VII-C asks only whether
compatible towers of such invariant sectors satisfy the continuum certificate.
\end{definition}

\begin{proposition}[Gauge-sector classification consequence]
\label{p12:thm:finite-gauge-projection}
If \(H=H^\ast\in\mathcal A_C^G\), then Paper~IV, Theorem~\ref{p4:thm:finite-gauge-constraint-closure},
implies
\[
e^{-itH}\mathcal H_C^G\subseteq\mathcal H_C^G.
\]
Hence
\[
(\mathcal H_C^G,\;P_G\mathcal A_C^GP_G,\;H|_{\mathcal H_C^G})
\]
is the selected finite gauge-sector presentation used by the Paper-VII-C
continuum ledger.
\end{proposition}

\begin{proof}
This is a direct application of the finite gauge-constraint closure proved in
Paper~IV.  No independent gauge-state or evolution theorem is introduced here.
\end{proof}

\begin{remark}[QCD and singlet sectors]
\label{p12:rem:qcd-singlet-sector}
The QCD singlet sector used in Paper~VII-B is an instance of this imported gauge-sector proposition.  The Paper~V-A color carrier is imported as upstream theorem-level data, the Paper~V-C singlet support is the finite gauge selector used by the compiled strong sector, and the nonperturbative Wilson/transfer sector is a gauge-invariant subsystem.  Perturbative gluon Fock language is a later low-boundary or continuum interface, not the primitive definition of the finite color sector.
\end{remark}


% Source Paper VII-C, original line 7194; inherited proof bodies unchanged.
\subsubsection{Subsystem status classes: native, selected, interface, no-witness}

\begin{table}[htbp]
\centering
\small
\begin{tabularx}{\textwidth}{@{}
>{\raggedright\arraybackslash}p{0.55\textwidth}
>{\raggedright\arraybackslash}X
@{}}
\toprule
\textbf{Criterion or subsystem claim} & \textbf{Discharged by} \\
\midrule
Subsystem status classes
& Definition~\ref{p12:def:subsystem-status-classes}. \\

GP2 filter for new sectors
& Proposition~\ref{p12:prop:gp2-filter-new-sectors}. \\

Theory-subsystem classification
& Theorem~\ref{p12:thm:theory-subsystem-classification}. \\

Ordinary theories as subsystem presentations
& Remark~\ref{p12:prop:qm-subsystem} and Propositions~\ref{p12:prop:classical-subsystem}--\ref{p12:prop:condensed-matter-subsystems}. \\

Controlled continuum subsystem
& Definition~\ref{p12:def:controlled-continuum-subsystem}. \\

Continuum theories as subsystem limits
& Theorem~\ref{p12:thm:continuum-theories-subsystem-limits}. \\

Paper-IV finite CCR boundary import and Paper-VII-C limit criterion
& Remarks~\ref{p12:lem:no-exact-finite-global-bosonic-ccr} and
\ref{p12:prop:truncated-one-mode-ccr-defect}, followed by
Corollary~\ref{p12:cor:ccr-regular-subsystem-condition}. \\

Finite-witness criterion for new sectors
& Theorem~\ref{p12:thm:finite-witness-criterion-new-sectors}. \\

Finite-witness requirement for BSM sectors
& Corollary~\ref{p12:cor:finite-witness-bsm-sectors}. \\

Finite discovery algorithm
& Definition~\ref{p12:def:finite-discovery-algorithm} and
Theorem~\ref{p12:thm:finite-resolution-discovery-coverage}. \\

Finite witness target ledger
& Definition~\ref{p12:def:finite-witness-target-transcript},
Table~\ref{p12:tab:finite-discovery-target-ledger}, and
Proposition~\ref{p12:prop:finite-witness-target-priority}. \\
\bottomrule
\end{tabularx}
\caption[Theory criterion and new-physics ledger]{Theory criterion and new-physics ledger.}
\label{p12:tab:theory-criterion-new-physics-ledger}
\TableNotes{This table turns the final subsystem discussion into a result roadmap.  It separates native finite presentations, selected finite subsystems, interface presentations, and unsupported sector insertions.}
\end{table}

\begin{definition}[Subsystem status classes]
\label{p12:def:subsystem-status-classes}
A Born--Hopf subsystem has one of the following statuses.

\begin{enumerate}[label=(\roman*), leftmargin=2em]
\item \textbf{Native subsystem.}
The selector data are generated by GP1--GP2 and the previous papers.

\item \textbf{Selected subsystem.}
The carrier and admissibility rules are native, while a finite Hamiltonian, state, preparation, environment, boundary condition, or sector selector is supplied on a stated theorem domain.

\item \textbf{Interface subsystem.}
An experimentally supplied apparatus, thermodynamic bath, material Hamiltonian, coarse-graining, or continuum presentation is represented inside the finite carrier.

\item \textbf{No-witness insertion.}
A proposed subsystem is not admitted as new fundamental structure when it requires a new continuum input, arbitrary extra field, extra gauge representation, preferred lattice, preferred Lorentzian branch, or unmotivated selector not generated by the Born--Hopf chain.
\end{enumerate}
\end{definition}

\begin{proposition}[GP2 finite-witness admissibility rule]
\label{p12:prop:gp2-filter-new-sectors}
Within the constitutive GP2 definition of the Born--Hopf framework, a proposed new \emph{native} sector is admissible only if it is one of the following:
\begin{enumerate}[label=(\alph*), leftmargin=2em]
\item a nontrivial commutant or multiplicity space left by native selector data;
\item a native invariant or gauge-projected sector;
\item a native topological, winding, or transfer sector;
\item a finite-boundary or finite-resolution correction to an already present carrier;
\item a controlled continuum or large-\(m_x\) refinement of finite data;
\item a declared nonminimal parent-completion slot with explicit activation data.
\end{enumerate}
Otherwise it has no native witness and is not admitted as a new fundamental sector.
\end{proposition}
\begin{proof}
This is a consequence of the framework's definition of ``native'': GP2 admits only structure generated from the finite carrier, already derived selectors, controlled refinement, or an explicitly declared nonminimal completion.  The listed mechanisms are the corresponding finite witness classes used by the framework.  Therefore a proposed sector outside them is, by definition, an external insertion rather than a native sector.  This proposition is a constitutive classification rule; it is not an empirical theorem excluding theories formulated outside the Born--Hopf framework.
\end{proof}


% Source Paper VII-C, original line 7286; inherited proof bodies unchanged.
\subsubsection{Ordinary theories as subsystem presentations}

\begin{theorem}[Finite-resolution subsystem classification principle]
\label{p12:thm:theory-subsystem-classification}
Let a finite-resolution subsystem presentation \(T\) inside the Paper-IV envelope be specified by finite state spaces, finite observables, energy or transfer data, compact symmetry or gauge data, and finite experimental resolution.  Then \(T\) is represented in the Born--Hopf framework by one of the following finite selector presentations:
\begin{enumerate}[label=(\roman*), leftmargin=2em]
\item commutant subsystem;
\item central sector compression;
\item gauge fixed-point sector;
\item stabilizer subsystem;
\item coarse-grained subsystem;
\item controlled continuum refinement.
\end{enumerate}
If the selector data of \(T\) are native, the theory is a derived Born--Hopf sector.  If the selector data are external, the theory is a selected or interface sector.
\end{theorem}

\begin{proof}
Finite physical data generate selector data inside the Paper-IV envelope whenever they are represented at finite resolution.  Proposition~\ref{p12:thm:finite-selector-commutant-subsystem}, Theorem~\ref{p12:thm:finite-energy-symmetry-sector}, and Proposition~\ref{p12:thm:finite-gauge-projection} classify the corresponding imported finite presentations.  Continuum theories arise only after the independent refinement obligations of the main text are satisfied.  The status distinction follows from Definition~\ref{p12:def:subsystem-status-classes}.
\end{proof}

\begin{remark}[Quantum mechanics is a Paper-IV import]
\label{p12:prop:qm-subsystem}
The finite quantum presentation used throughout Paper~VII-C is
\[
\mathcal A_C=\End(\mathcal H_C),
\qquad
\omega(a)=\Tr(\rho a),
\]
with states and Born probabilities as in Paper~IV, Theorems~\ref{p4:thm:density-matrix} and~\ref{p4:thm:born}, and with instruments, channels, composition,
and reduced dynamics exactly as proved there.  Compression,
tensor-factor selection, commutant selection, and environment selection create
subsystem presentations of that calculus.  Paper~VII-C makes no independent
finite-quantum postulate or proof.
\end{remark}

\begin{proposition}[Classical probability and conditional classical mechanics]
\label{p12:prop:classical-subsystem}
Classical probability is a commutative selected subalgebra
\(\mathcal C\subseteq\mathcal A_C\) with
\(\mathcal C\simeq\mathbb C^n\). A selected classical Hamiltonian
mechanics limit requires more than asymptotic commutativity: it additionally
requires a specified correspondence family \(Q_L\) from a Poisson algebra
\((\mathcal P,\{\cdot,\cdot\})\) into the actual finite observables,
a convergence topology \(\tau\), a scaling \(\hbar_L\), and
compatible Hamiltonians \(H_L\) such that, on the stated common domain,
\[
 Q_L(ab)-Q_L(a)Q_L(b)\xrightarrow{\tau}0,
 \qquad
 \frac{1}{i\hbar_L}[Q_L(a),Q_L(b)]
   -Q_L(\{a,b\})\xrightarrow{\tau}0,
\]
\[
 \frac{i}{\hbar_L}[H_L,Q_L(a)]
   -Q_L(\{a,h\})\xrightarrow{\tau}0.
\]
On any source-verified branch satisfying these comparisons, classical
Hamiltonian dynamics is a controlled selected subsystem of finite quantum
dynamics, not a primitive replacement for the carrier. Commutativity by
itself supplies neither a Poisson bracket nor this dynamical limit.
\end{proposition}
\begin{proof}
A finite commutative \(C^*\)-algebra is the algebra of functions on a
finite set and gives classical probability. For mechanics the displayed
commutator and Hamiltonian comparisons are precisely the missing
algebraic and dynamical correspondence checks. Their limits transport the
Heisenberg derivation to the stated Poisson Hamiltonian derivation on its
common domain. The commutative algebra alone does not provide these maps or
limits.
\end{proof}

\begin{remark}[Thermodynamics is a Paper-IV finite-quantum import]
\label{p12:prop:thermodynamics-subsystem}
For a supplied finite Hamiltonian \(H\), Paper~IV, Theorem~\ref{p4:thm:finite-thermal-kms-kinetic-interface}, owns the finite Gibbs/KMS and
entropy formulas
\[
\rho_\beta=\frac{e^{-\beta H}}{\Tr(e^{-\beta H})},
\qquad
S(\rho)=-\Tr(\rho\log\rho),
\]
as well as the channel/open-system origin of coarse-grained irreversibility.
Paper~VII-C uses only refinement families of these finite thermal states when a
continuum thermal sector is under consideration.
\end{remark}

\begin{corollary}[Finite KMS target versus dynamical thermalization]
\label{p12:cor:k16-finite-kms-target-versus-thermalization}
Let \(\mathcal H_{\rm out}^{\rm ord}\) be a supplied finite occupied carrier
from the Paper--VI-B \(K=16\) interface, let
\(\mathcal A_{\rm ord}=\End(\mathcal H_{\rm out}^{\rm ord})\), and supply a
self-adjoint \(H_{\rm ord}\) and \(\beta>0\).  Then the normalized
\(\beta\)-KMS state on \(\mathcal A_{\rm ord}\) is uniquely
\begin{equation}
\label{p12:eq:k16-unique-finite-kms-target}
 \rho_{\beta,\rm ord}
 =\frac{e^{-\beta H_{\rm ord}}}
        {\Tr(e^{-\beta H_{\rm ord}})}.
\end{equation}
If the physical observable algebra is instead a direct sum of superselection
blocks, the KMS condition fixes this density within each block but does not
fix the central mixture weights.  In either case, the static KMS statement
does not imply that the produced state \(\rho_{\rm prod}\) thermalizes.  Such a
claim still requires a specified state evolution \(\Phi_t\), sector-weight
data where applicable, an energy/entropy condition selecting \(\beta\), and a
proved convergence statement such as
\begin{equation}
\label{p12:eq:k16-thermalization-obligation}
 \lim_{t\to\infty}
 \left\|\Phi_t(\rho_{\rm prod})-\rho_{\beta,\rm ord}\right\|_1=0.
\end{equation}
Consequently no independent scalar \(\mathrm{THERM16}\) temperature selector
is present.  On the displayed simple-algebra domain, a supplied \(\beta\) has
the unique KMS target above.  Selection of a physical \(\beta\) from the
cosmological/energy history, superselection central weights, collision or
kinetic dynamics, entropy production, and convergence are separately typed
claim obligations rather than components of one missing THERM16 selector.
\end{corollary}

\begin{proof}
Paper~IV, Theorem~\ref{p4:thm:finite-simple-kms-uniqueness}, gives uniqueness
on the simple matrix algebra and the central-weight caveat on direct sums.
The KMS identity is an equilibrium characterization and contains no evolution
from \(\rho_{\rm prod}\).  Equation~\eqref{p12:eq:k16-thermalization-obligation}
therefore requires the additional dynamical data listed above.
\end{proof}

\begin{corollary}[Imported special-relativity subsystem compatibility]
\label{p12:prop:sr-subsystem}
Import the Paper~IV minimal finite free-sector dispersion theorem,
Theorem~\ref{p4:thm:finite-one-particle-dispersion}, together with the
Papers~III--IV Lorentzian chain closure,
Theorem~\ref{p3:thm:p3-p4-lorentzian-chain-closure}.  For every stable isolated
intrinsic block in that theorem domain, the Paper~VII-C large-\(m_x\) local
continuum presentation contains the corresponding local special-relativistic
free subsystem.  In local flat-chart notation its mass shell is
\[
E^2=c^2|p|^2+m^2c^4,
\]
and its preserving spacetime transformations are the Lorentz/Poincar\'e
transformations of the imported Lorentzian quadratic form.  Paper~VII-C does not
rederive the finite dispersion relation, select the Lorentzian sign, or claim
that special relativity is the unique global theory selected by GP1--GP2.
\end{corollary}

\begin{proof}
Paper~IV proves the positive-frequency affine second-order dispersion relation
on the stated stable isolated intrinsic finite block.  Paper~III composes that
finite-field result with the scalar-cover and positive-spatial-operator chain
to constitute the one-time Lorentzian sign.  Theorem~\ref{p12:thm:local-chart-os-reconstruction}
then supplies the local continuum chart and reconstructed field presentation.
Restricting that presentation to the imported free block yields the displayed
special-relativistic subsystem.  No Paper-VII-C selector, dispersion theorem, or
signature theorem is added.
\end{proof}

\begin{proposition}[Imported GR-subsystem compatibility]
\label{p12:prop:gr-subsystem}
Papers~III and VI-A already fix the gravitational action, metric field equation, scalar-response constitution, and their closed-cosmology realization.  Paper~VII-C does not derive those objects again.  It proves only that the reconstructed local continuum presentation is compatible with their low-curvature Einstein--Hilbert subsystem,
\[
S_{\rm eff}
\sim
\int (R-2\Lambda)\sqrt{-g}\,d^4x
+
S_{\rm matter}.
\]
Here \(S_{\rm matter}\) is subsystem notation for the already-fixed reciprocal-position leakage/matter readout.  No gravitational coefficient, field equation, curvature branch, or cosmological response is selected in Paper~VII-C.
\end{proposition}

\begin{proof}
The Einstein--Hilbert truncation and metric dynamics are Paper--III imports, and their physical closed-cosmology realization is a Paper--VI-A import.  The Paper--VII-C OS/Wightman reconstruction supplies local fields, covariance, and local-chart compatibility for the inherited finite tower.  Restricting that reconstructed presentation to the imported low-curvature metric subsystem yields the displayed form without constituting a new gravity theory.
\end{proof}

\begin{proposition}[Gauge theory and Yang--Mills]
\label{p12:prop:gauge-ym-subsystem}
Gauge theory is the finite local comparison calculus of compact internal selector data.  The invariant physical carrier is obtained by the Haar projection \(P_G\), and the native Born--Hopf continuum Yang--Mills sector is recovered only after the Paper~VII-C refinement, locality, positivity, and susceptibility obligations are satisfied.
\end{proposition}

\begin{proof}
The finite gauge statement is the Paper-IV import recorded in Proposition~\ref{p12:thm:finite-gauge-projection}.  A continuum gauge field is not inserted as a primitive; it is reconstructed from finite local comparison and Wilson/curvature-residue data under the main Paper~VII-C obligations.
\end{proof}

\begin{proposition}[Standard Model and QCD sectors]
\label{p12:prop:sm-qcd-subsystems}
The Standard Model and QCD sectors are not arbitrary continuum insertions.  In the series they arise from finite internal carrier data, chirality, anomaly cancellation, scalar stabilizer structure, finite color, singlet projection, transfer positivity, and continuum reconstruction.  Additional elementary particles or gauge groups are not admissible unless a native commutant, parent-completion slot, or controlled finite witness is exhibited.
\end{proposition}

\begin{proof}
The finite internal data and their closure are supplied before Paper~VII-C.  Paper~VII-C reconstructs the continuum QFT layer from those finite data.  Corollary~\ref{p12:cor:v259-resolution-growth-no-internal-species} additionally rules out using reciprocal-position refinement itself as a hidden species generator: larger $m_x$ changes external mode resolution, not the inherited internal carrier.  Therefore new particle content cannot be added at the continuum level without adding an otherwise licensed finite witness or selector datum, which is excluded by GP2 when absent.
\end{proof}

\begin{proposition}[Effective field theories]
\label{p12:prop:eft-subsystem}
An effective field theory is an energy-window subsystem.  If \(P_{\le E}\) is the projection onto a finite energy window, then
\[
\mathcal A_{\le E}
=
P_{\le E}\mathcal A_CP_{\le E}
\]
with the compressed dynamics gives the finite EFT presentation.  RG flow is comparison between such presentations over nested or overlapping windows, not a new primitive.
\end{proposition}

\begin{proof}
Energy selection is a central case of the selector-presentation proposition.  Compressing to a spectral window gives a finite subsystem.  Comparing such subsystems under refinement gives the renormalization maps used in the main text.
\end{proof}

\begin{proposition}[Finite continuity law and conditional hydrodynamic interface]
\label{p12:prop:hydrodynamics-subsystem}
Suppose the actual finite action supplies a local (possibly graded-even)
operator presentation $\mathcal H=\bigotimes_{x\in V}\mathcal H_x$,
$H=\sum_Xh_X$, and $Q=\sum_xq_x$ with local self-adjoint charges.
For $x\in X$ define
$e_x^X=(i/\hbar)[h_X,q_x]$ and the antisymmetric current
\[
 J_{xy}^X=(e_y^X-e_x^X)/|X|\quad(x,y\in X),\qquad
 s_x=\sum_{X\ni x}\frac{i}{\hbar|X|}[h_X,Q_X],
 \quad Q_X=\sum_{y\in X}q_y.
\]
Then, with $J_{xy}=\sum_{X\supset\{x,y\}}J_{xy}^X$,
\[
 \boxed{\frac{dq_x}{dt}+\sum_yJ_{xy}=s_x.}
\]
If every interaction is locally charge preserving,
$[h_X,Q_X]=0$, the law is source-free. For every finite region
$A\subset V$,
\[
 \dot Q_A=-\sum_{x\in A,\,y\notin A}J_{xy}+\sum_{x\in A}s_x,
 \quad
 \|\dot Q_A\|\le\frac2\hbar
 \sum_{X\cap A\ne\emptyset,\,X\not\subset A}
 \|h_X\|\,\|Q_{X\cap A}\|
\]
under local interaction charge preservation.
Global $[H,Q]=0$ alone does not supply this presentation, nor does
continuity alone prove a diffusion, viscosity or constitutive law.
Hydrodynamic observables require the occurring local charges,
state/scale selection and long-wavelength response certificates.
\end{proposition}
\begin{proof}
For each interaction $X$, $\dot q_x=e_x^X$ and
$\sum_{y\in X}J_{xy}^X=|X|^{-1}\sum_{y\in X}e_y^X-e_x^X$.
Summing over $X$ gives the displayed continuity identity. Under local
conservation, $s_x=0$; antisymmetry cancels internal currents on summing
over a region and leaves only crossing interactions. The bound follows
from $\|[h_X,Q_{X\cap A}]\|\le2\|h_X\|\|Q_{X\cap A}\|$.
No constitutive law is used.
\end{proof}


\begin{proposition}[Condensed-matter and apparatus theories]
\label{p12:prop:condensed-matter-subsystems}
Condensed-matter models, apparatus pointer models, spin chains, band Hamiltonians, and material EFTs are external interface subsystems unless their Hamiltonians and symmetry data are derived natively.  They are represented by finite selector data inside the carrier, but they are not new fundamental sectors.
\end{proposition}

\begin{proof}
Such models specify external Hamiltonians, environments, lattices, boundary conditions, or material symmetries.  These are valid finite selector data and hence define subsystems, but their selectors are supplied by preparation or material context rather than by GP1--GP2 alone.
\end{proof}


% Source Paper VII-C, original line 7499; inherited proof bodies unchanged.
\subsubsection{Continuum reconstruction as a subsystem limit}

\begin{definition}[Controlled continuum subsystem]
\label{p12:def:controlled-continuum-subsystem}
A controlled continuum subsystem is a refinement family of finite subsystems
\[
(\mathcal H_L,\mathcal A_L,\omega_L,H_L)_{L\in\mathcal R}
\]
with comparison maps to a common local cylinder algebra such that the Paper~VII-C obligations hold: local state compactness, positivity, covariance, asymptotic locality, spectral control, and the required sector-specific Ward or Schwinger--Dyson identities.
\end{definition}

\begin{theorem}[Continuum theories are recovered subsystem limits]
\label{p12:thm:continuum-theories-subsystem-limits}
A continuum QFT or continuum geometric field theory is admissible in the Born--Hopf framework only as a controlled continuum subsystem.  Its continuum Hilbert space, local fields, distributions, path-integral notation, propagators, and Ward identities are reconstructed outputs, not primitive inputs.
\end{theorem}

\begin{proof}
This is the main Paper~VII-C reconstruction principle applied to the finite selector subsystems of this appendix.  The finite carriers and local comparison maps are primary.  The continuum theory is the limit object obtained after the relevant obligations are satisfied.
\end{proof}

\begin{remark}[Continuum reconstruction boundary]
\label{p12:rem:continuum-reconstruction-boundary}
A proposed continuum field cannot bypass the finite-witness criterion by being written directly as a Lagrangian.  It must have a finite carrier witness, selector data, admissible local comparison maps, and the required positivity/locality/covariance limits.
\end{remark}


% Source Paper VII-C, original line 7524; inherited proof bodies unchanged.
\subsubsection{Finite bosonic CCR boundary defect as one subsystem criterion}

The finite CCR boundary defect is one important example of the general subsystem principle.  It shows how an apparent continuum theory can fail to be an ordinary subsystem if its finite boundary terms do not vanish.

\begin{remark}[Imported Paper-IV obstruction to exact finite bosonic CCR]
\label{p12:lem:no-exact-finite-global-bosonic-ccr}
Paper~IV, Proposition~\ref{p4:prop:no-exact-ccr}, proves by the finite trace identity that no operators on a
finite-dimensional Hilbert space satisfy \([a,a^\dagger]=\mathbf1\).  This
obstruction is imported here and is not a Paper-VII-C result.
\end{remark}

\begin{remark}[Imported Paper-IV truncated one-mode CCR defect]
\label{p12:prop:truncated-one-mode-ccr-defect}
Paper~IV, Proposition~\ref{p4:prop:explicit-finite-ccr-boundary-defect},
proves on \(\operatorname{span}\{|0\rangle,\ldots,|N\rangle\}\) the
exact finite boundary formula
\begin{equation}
\label{p12:eq:finite-ccr-boundary-defect-short}
[a,a^\dagger]=\mathbf1-(N+1)|N\rangle\langle N|.
\end{equation}
Paper~VII-C uses this imported identity only to formulate and test a continuum
regularity condition for selected refinement families.
\end{remark}

\begin{corollary}[CCR-regular subsystem condition]
\label{p12:cor:ccr-regular-subsystem-condition}
A finite bosonic refinement family recovers ordinary CCR only if the occupation-boundary defect vanishes in the selected states.  In the one-mode case this requires
\[
(N+1)\omega_N(|N\rangle\langle N|)\to0.
\]
If this fails, the subsystem is a finite-saturation bosonic sector rather than ordinary bosonic QFT.
\end{corollary}
\begin{proof}
\ProofPlan{We read the truncated one-mode CCR defect in selected states and use its boundary term to distinguish ordinary CCR recovery from finite-saturation bosonic sectors.}
\begin{proofsteps}
\item \textbf{Truncated one-mode CCR defect.}
By the imported Paper-IV boundary identity recorded in Remark~\ref{p12:prop:truncated-one-mode-ccr-defect},
\[
[a,a^\dagger]
=
\mathbf1-(N+1)|N\rangle\langle N|.
\]

\item \textbf{Expectation in selected subsystem states.}
Taking expectation in a selected state \(\omega_N\) gives
\[
\omega_N([a,a^\dagger])
=
1-(N+1)\omega_N(|N\rangle\langle N|).
\]

\item \textbf{Vanishing boundary-defect condition.}
Ordinary CCR recovery in the selected subsystem requires the commutator expectation to converge to the continuum value \(1\).  Therefore the defect term must vanish:
\[
(N+1)\omega_N(|N\rangle\langle N|)\to0.
\]

\item \textbf{Finite-saturation alternative.}
If this condition fails, then the limiting selected states retain a nonzero contribution from the top occupation projector.  The selected subsystem therefore does not realize the ordinary bosonic CCR; it retains a finite occupation-boundary saturation defect.  Hence it is a finite-saturation bosonic sector rather than ordinary bosonic QFT.
\end{proofsteps}
\end{proof}

\begin{remark}[New physics from CCR failure]
\label{p12:rem:new-physics-ccr-failure}
The CCR defect does not introduce a new boson.  It modifies the high-occupation behavior of already present bosonic modes.  Thus high-occupation anomalies should first be tested against finite-boundary saturation before postulating extra elementary fields.
\end{remark}

\begin{proposition}[Executable CCR-boundary witness numerics]
\label{p12:prop:executable-ccr-boundary-witness-numerics}
Let
\[
 p_N:=\omega_N(|N\rangle\langle N|),
 \qquad
 \epsilon_{\rm CCR}^{(N)}=(N+1)p_N.
\]
Then the selected commutator readout is
\[
 \omega_N([a,a^\dagger])=1-\epsilon_{\rm CCR}^{(N)},
\]
and the following canonical cutoff states give the closed finite witness values:
\begin{center}
\scriptsize
\renewcommand{\arraystretch}{1.12}
\begin{tabularx}{0.96\textwidth}{@{}
>{\raggedright\arraybackslash}p{0.27\textwidth}
>{\raggedright\arraybackslash}p{0.30\textwidth}
>{\raggedright\arraybackslash}p{0.18\textwidth}
>{\raggedright\arraybackslash}X
@{}}
\toprule
Selected finite state & Top weight \(p_N\) & \(\epsilon_{\rm CCR}^{(N)}\) & Boundary-commutator verdict \\
\midrule
Boundary-top state & \(1\) & \(N+1\) & Divergent finite-boundary saturation. \\
Uniform window on \(0,\ldots,N\) & \(1/(N+1)\) & \(1\) & Persistent defect; commutator expectation equals zero. \\
Polynomial boundary leakage & \((N+1)^{-\gamma}\) & \((N+1)^{1-\gamma}\) & Expectation defect vanishes only for \(\gamma>1\); finite nonzero limit at \(\gamma=1\). \\
Truncated geometric tail \((0<r<1)\) & \((1-r)r^N/(1-r^{N+1})\) & \((N+1)(1-r)r^N/(1-r^{N+1})\) & Expectation defect vanishes exponentially; full Weyl regularity is not inferred. \\
\bottomrule
\end{tabularx}
\end{center}
For example, for \(N=8,32,128\) the polynomial case \(p_N=(N+1)^{-2}\) gives
\[
\epsilon_{\rm CCR}^{(N)}=1/9,\quad 1/33,\quad 1/129,
\]
while the geometric case \(r=1/2\) gives approximately
\[
1.761\times10^{-2},\qquad3.842\times10^{-9},\qquad1.895\times10^{-37}.
\]
These are evaluations of the closed boundary formula, not new fitted parameters.
\end{proposition}

\begin{proof}
The first identity is the imported Paper-IV boundary formula in Remark~\ref{p12:prop:truncated-one-mode-ccr-defect} evaluated in the selected state \(\omega_N\).  The table substitutes the displayed top weights into \(\epsilon_{\rm CCR}^{(N)}=(N+1)p_N\).  The displayed boundary-commutator verdicts follow from the exact identity: vanishing of \(\epsilon_{\rm CCR}^{(N)}\) is equivalent to convergence of this \emph{one commutator expectation} to one. Corollary~\ref{p12:cor:ccr-regular-subsystem-condition} establishes this as a necessary, not sufficient, condition for full CCR recovery.  The three cutoff examples are direct substitutions at \(N=8,32,128\).\qedhere
\end{proof}
\begin{remark}[Boundary expectation does not imply Weyl regularity]
The finite numbers in Proposition~\ref{p12:prop:executable-ccr-boundary-witness-numerics}
are only boundary-commutator diagnostics. Convergence of the regular Weyl
representation, multiword moments or local field net additionally requires
the distinct state, tightness, domain and Weyl-continuity hypotheses in the
later source theorems; none is supplied by \(\epsilon_{\rm CCR}^{(N)}\to0\)
alone.
\end{remark}


% Source Paper VII-C, original line 7637; inherited proof bodies unchanged.
\subsubsection{Finite-witness test for new sectors}

\begin{definition}[Native new-physics witness]
\label{p12:def:native-new-physics-witness}
A native new-physics witness is a finite Born--Hopf structure that survives the GP2 filter and is not already identified with a known subsystem.  It must be one of:
\begin{enumerate}[label=(\roman*), leftmargin=2em]
\item a nontrivial native commutant factor \(\End(M_\alpha)\);
\item a native topological or winding sector;
\item a finite boundary-defect sector;
\item a native scalar, metric, or transfer-response correction;
\item a controlled gapless or continuum-refinement sector;
\item a declared nonminimal parent-completion slot with explicit activation data.
\end{enumerate}
\end{definition}



\begin{corollary}[Finite witness transcript closure]
\label{p12:cor:finite-witness-transcript-closure}
A proposed new sector passes the Paper~VII-C criterion only after its native witness has been written as an executable finite witness transcript in the sense of Definition~\ref{p12:def:finite-witness-target-transcript}.  In particular, continuum notation for a new particle, gauge group, scalar, axionlike field, sterile sector, or extra family is not an admissible substitute for the finite tuple \(\mathfrak W_{\rm fin}\).
\end{corollary}
\begin{proof}
Definition~\ref{p12:def:native-new-physics-witness} lists the allowed witness types, but the witness becomes testable only after the finite carrier, selector algebra, commutant, Wedderburn multiplicities, readout residue, and null condition are specified as \(\mathfrak W_{\rm fin}\) in Definition~\ref{p12:def:finite-witness-target-transcript}.  Table~\ref{p12:tab:finite-discovery-target-ledger} gives the corresponding executable target classes.  Therefore a continuum name without this finite transcript has no GP2 witness and is rejected by the finite-witness criterion.
\end{proof}

\begin{theorem}[Finite-witness criterion for new sectors]
\label{p12:thm:finite-witness-criterion-new-sectors}
Let \(X\) be a proposed new physical sector.  Then \(X\) is admissible only if it has a native new-physics witness in the sense of Definition~\ref{p12:def:native-new-physics-witness}.  If no such witness exists, \(X\) is a no-witness insertion.
\end{theorem}

\begin{proof}
By Proposition~\ref{p12:thm:finite-selector-commutant-subsystem}, any finite sector must be generated by selector data, commutants, invariant projections, or compressions.  By Theorem~\ref{p12:thm:continuum-theories-subsystem-limits}, any continuum theory must be a controlled continuum subsystem of finite data.  Corollary~\ref{p12:cor:finite-witness-transcript-closure} upgrades the witness requirement to the executable transcript \(\mathfrak W_{\rm fin}\).  Hence a proposed new sector without one of the native witnesses of Definition~\ref{p12:def:native-new-physics-witness} and its finite transcript is an insertion.
\end{proof}


\begin{corollary}[Finite-witness requirement for BSM sectors]
\label{p12:cor:finite-witness-bsm-sectors}
Within the Born--Hopf framework, BSM sectors such as new light particles, gauge groups, scalar fields, sterile sectors, axionlike fields, or extra generations require a native finite witness or an explicitly declared nonminimal parent-completion slot.  They are not added solely by continuum naming.
\end{corollary}

\begin{proof}
Each listed object changes the finite carrier content, gauge representation content, or selector algebra.  Without a native witness, this is an ungenerated insertion.
\end{proof}

\begin{corollary}[Native finite witness classes]
\label{p12:cor:where-new-physics-search}
Within the minimal framework, the following are admissible classes of finite residual witnesses before any new continuum field is named:
\begin{enumerate}[label=(\alph*), leftmargin=2em]
\item finite scalar/metric response residues;
\item determinant or ruler-transport residues;
\item low-mode or large-angle finite harmonic residues;
\item finite topological or winding sectors;
\item finite bosonic boundary defects;
\item nontrivial residual commutants or multiplicity spaces;
\item explicitly declared parent-completion slots.
\end{enumerate}
This is a classification of admissible witness types, not a theorem imposing a priority order for experimental or theoretical searches.
\end{corollary}
\begin{proof}
Each listed item is one of the witness mechanisms admitted by Definition~\ref{p12:def:native-new-physics-witness} and Proposition~\ref{p12:prop:gp2-filter-new-sectors}.  Those definitions classify what counts as a native finite witness; they do not define a total ordering among the classes.  Hence the mathematically valid conclusion is the classification itself, while the former phrase ``searched first'' is removed as an unsupported methodological preference.
\end{proof}


% Source Paper VII-C, original line 7699; inherited proof bodies unchanged.
\subsubsection{Discovery algorithm}

\begin{definition}[Finite discovery algorithm]
\label{p12:def:finite-discovery-algorithm}
Given a proposed physical regime, the Born--Hopf discovery algorithm is:

\begin{enumerate}[label=\textbf{Step \arabic*.}, leftmargin=3em]
\item Identify the finite carrier \(\mathcal H_C\) and maximal algebra \(\mathcal A_C=\End(\mathcal H_C)\).
\item Identify all native selector data \(\mathcal S\): energy, transfer, symmetry, gauge, scalar, topology, environment, local chart, or boundary data.
\item Form the selector algebra
\[
\mathcal B_{\mathcal S}=C^\ast(\mathcal S).
\]
\item Compute the commutant
\[
\mathcal A_{\mathcal S}=\mathcal B_{\mathcal S}'.
\]
\item Decompose
\[
\mathcal H_C\simeq\bigoplus_\alpha V_\alpha\otimes M_\alpha.
\]
\item Classify every \(M_\alpha\): known spin, charge, color, weak isospin, family, mass/mixing, gauge singlet, topology, boundary, response, or unidentified.
\item If an unidentified \(M_\alpha\), boundary defect, topological block, or response residue remains, test it as a native new-physics witness.
\item If a continuum theory is requested, verify the Paper~VII-C refinement obligations before naming a continuum field theory.
\end{enumerate}
\end{definition}

\begin{theorem}[Finite-resolution discovery coverage]
\label{p12:thm:finite-resolution-discovery-coverage}
At fixed finite resolution, the discovery algorithm finds all selector-compatible finite subsystem degrees of freedom.
\end{theorem}

\begin{proof}
At fixed finite resolution, all operators lie in the finite matrix algebra \(\mathcal A_C\).  The selector-compatible observables are exactly the commutant \(\mathcal B_{\mathcal S}'\).  By the Wedderburn theorem, all degrees of freedom in the commutant are exactly the multiplicity factors \(M_\alpha\).  Therefore no additional selector-compatible finite degree of freedom exists outside the algorithm.
\end{proof}

\begin{remark}[What the algorithm cannot do by itself]
\label{p12:rem:algorithm-limits}
The finite discovery algorithm classifies finite residual structure.  It does not by itself prove that a residual sector is realized in nature, dynamically populated, light, stable, or experimentally visible.  Those require additional mass, transfer, coupling, state-population, or boundary-weight calculations.
\end{remark}


% Source Paper VII-C, original line 7740; inherited proof bodies unchanged.
\subsubsection{Finite discovery targets and residual witness ledger}
\label{p12:subsec:finite-discovery-targets}

The discovery algorithm becomes useful only after the admissible witness types
are turned into executable finite readout targets.  This subsection records the
minimal target ledger used by the Born--Hopf framework.  It is not a catalogue
of new particles.  It is a readout-chronology rule for deciding where an unidentified
finite residue may first be tested before any continuum BSM language is allowed.

\begin{definition}[Executable finite witness transcript]
\label{p12:def:finite-witness-target-transcript}
An executable finite witness transcript for a proposed residual phenomenon is a
tuple
\begin{equation}
\label{p12:eq:finite-witness-target-transcript}
\mathfrak W_{\rm fin}
=
\left(
\mathcal H_C,
\mathcal S,
\mathcal B_{\mathcal S},
\mathcal A_{\mathcal S},
\{M_\alpha\},
\mathcal R_{\rm wit},
\mathcal N_{\rm null}
\right),
\end{equation}
where \(\mathcal H_C\) is the finite carrier, \(\mathcal S\) is the native
selector data, \(\mathcal B_{\mathcal S}=C^\ast(\mathcal S)\),
\(\mathcal A_{\mathcal S}=\mathcal B_{\mathcal S}'\), and
\(\{M_\alpha\}\) are the Wedderburn multiplicity factors.  The object
\(\mathcal R_{\rm wit}\) is one of the native witness readouts of
Definition~\ref{p12:def:native-new-physics-witness}, and
\(\mathcal N_{\rm null}\) is the corresponding null condition under which the
residue reduces to an already-known subsystem or disappears.
\end{definition}

\begin{table}[htbp]
\centering
\scriptsize
\begin{tabularx}{\textwidth}{@{}
>{\raggedright\arraybackslash}p{0.19\textwidth}
>{\raggedright\arraybackslash}p{0.27\textwidth}
>{\raggedright\arraybackslash}p{0.27\textwidth}
>{\raggedright\arraybackslash}X
@{}}
\toprule
\textbf{Target class} & \textbf{Native finite witness} & \textbf{Null condition} & \textbf{Interpretation if nonzero} \\
\midrule
Finite bosonic boundary defect
& \(\epsilon_{\rm CCR}^{(N)}=(N+1)\omega_N(|N\rangle\langle N|)\)
& \(\epsilon_{\rm CCR}^{(N)}\to0\)
& ordinary bosonic CCR is recovered only in the selected states; a persistent value is a finite-saturation residue, not a new elementary boson. \\

Residual commutant or multiplicity space
& an unidentified factor \(M_\alpha\) in
\(\mathcal H_C\simeq\bigoplus_\alpha V_\alpha\otimes M_\alpha\)
& every \(M_\alpha\) is already classified by known spin, charge, color, weak isospin, family, mass/mixing, gauge, topology, or response data
& a genuine hidden finite degree of freedom, requiring separate mass, transfer, stability, and population calculations before any particle interpretation. \\

Topological or winding block
& finite index/winding readouts such as \(\nu\), \(\langle\nu^2\rangle/V\), \(b_{2k}\), or a finite anomaly-response block
& the finite winding/index response is already exhausted by the known gauge or anomaly sector
& a native topological residual; it must be tested before introducing axionlike, instanton-phenomenology, or other continuum topological fields. \\

Scalar, metric, or transfer-response residue
& a finite scalar response, determinant-ruler response, transfer-kernel residue, or low-mode harmonic remainder fixed by the carrier
& the response is absorbed by the already-derived scalar/metric/transport subsystem on its theorem domain
& a native correction to propagation, ruler transport, or response laws rather than an inserted dark-sector particle. \\

Controlled gapless or continuum-refinement sector
& a finite local subsystem whose comparison limit satisfies the Paper--VII-C refinement obligations
& the sufficient same-side vacuum/state-naturality, positivity, covariance, and locality defects vanish in the admissible refinement
& a legitimate continuum presentation of finite data; failure of the obligations is not evidence for a new continuum input. \\

Nonminimal parent-completion slot
& an explicitly declared parent-completion datum with activation data and no comparison feedback into the minimal carrier
& the activation datum is absent or projected out by the minimal split phase
& selected nonminimal physics; it is outside the minimal core unless the parent-completion slot is separately constituted. \\
\bottomrule
\end{tabularx}
\caption[Finite discovery target ledger]{Finite discovery target ledger.}
\label{p12:tab:finite-discovery-target-ledger}
\TableNotes{The table refines the witness classes of
Definition~\ref{p12:def:native-new-physics-witness} into executable targets.
The null condition must be stated before comparison.  A nonzero witness does
not by itself prove a new realized particle; it identifies the finite residue
that must next receive mass, transfer, coupling, stability, population, or
boundary-weight analysis.}
\end{table}

\begin{proposition}[Executable CCR boundary witness]
\label{p12:prop:finite-witness-target-priority}
The finite bosonic CCR boundary defect is an exactly executable native witness:
\begin{equation}
\label{p12:eq:epsilon-ccr-target}
\epsilon_{\rm CCR}^{(N)}:=(N+1)\omega_N(|N\rangle\langle N|),
\qquad
\omega_N([a,a^\dagger])=1-\epsilon_{\rm CCR}^{(N)}.
\end{equation}
A persistent value
\begin{equation}
\label{p12:eq:epsilon-ccr-nonzero}
\limsup_{N\to\infty}\epsilon_{\rm CCR}^{(N)}>0
\end{equation}
is a finite-saturation witness and prevents ordinary CCR recovery in that selected state family.  This algebraic fact does not impose a research-priority ordering relative to other finite witnesses and does not imply a new elementary boson.
\end{proposition}
\begin{proof}
The exact truncated commutator is \([a,a^\dagger]=I-(N+1)|N\rangle\langle N|\).  Taking \(\omega_N\) gives \eqref{p12:eq:epsilon-ccr-target}.  Corollary~\ref{p12:cor:ccr-regular-subsystem-condition} proves that ordinary CCR recovery requires \(\epsilon_{\rm CCR}^{(N)}\to0\); hence a positive limsup is precisely a finite-saturation obstruction.  No comparison among unrelated witness classes follows from this identity, so the former word ``first'' was a methodological preference rather than a theorem and is removed.
\end{proof}

\begin{corollary}[No particle claim from a witness alone]
\label{p12:cor:no-particle-claim-from-witness-alone}
A nonzero entry in Table~\ref{p12:tab:finite-discovery-target-ledger} licenses
only the next finite calculation: mass or gap, transfer amplitude, coupling,
stability, population, and boundary-weight analysis.  It does not license naming
a new continuum particle or adding a new Lagrangian sector unless the required
finite witness transcript has first been completed.
\end{corollary}

\begin{proof}
This is the operational form of Remark~\ref{p12:rem:algorithm-limits} and
Corollary~\ref{p12:cor:finite-witness-bsm-sectors}.  The algorithm detects
residual finite structure; realization, visibility, and continuum naming require
the additional finite data listed in the statement.
\end{proof}


% Source Paper VII-C, original line 7867; inherited proof bodies unchanged.
\subsubsection{SR and GR revisited as examples}

\begin{remark}[SR is contained but not uniquely selected]
\label{p12:rem:sr-contained-not-unique}
The framework contains SR as the local stable free Lorentzian subsystem, but GP1--GP2 do not select SR as the whole theory.  The selector data needed to isolate SR include stability, isolation, free second-order dynamics, positive scalar-cover time, and a local large-\(m_x\) flat chart.  Without those selectors, the full Born--Hopf carrier contains more structure than SR.
\end{remark}

\begin{remark}[GR is contained as a low-curvature metric subsystem]
\label{p12:rem:gr-contained-low-curvature}
The framework contains GR as the low-curvature metric response subsystem already derived in Paper~III and realized gravitationally in Paper~VI-A; Paper~VII-C only certifies compatibility of its reconstructed local continuum fields with that imported Einstein--Hilbert truncation.  It does not identify the full Born--Hopf theory with GR.  Finite scalar, determinant-ruler, compact-momentum, and spectral-response corrections belong to the larger theory.
\end{remark}


% Source Paper VII-C, original line 7879; inherited proof bodies unchanged.
\subsubsection{Summary of the subsystem criterion}

The subsystem principle can be summarized as
\[
\boxed{
\text{finite carrier}
+
\text{selector algebra}
\quad\Longrightarrow\quad
\text{commutant subsystem}.
}
\]
Gauge sectors are invariant compressions:
\[
\boxed{
\mathcal H_{\rm phys}=P_G\mathcal H_C.
}
\]
Energy and symmetry sectors are finite Wedderburn blocks:
\[
\boxed{
\mathcal H_C
=
\bigoplus_{\lambda,\rho}
V_\rho\otimes M_{\lambda,\rho}.
}
\]
Continuum theories are controlled subsystem limits:
\[
\boxed{
\text{continuum QFT}
=
\text{finite local subsystem}
+
\text{Paper VII-C refinement obligations}.
}
\]
New physics is admissible only when a native finite witness remains:
\[
\boxed{
\begin{gathered}
\text{new sector}
\Rightarrow
\text{native finite witness},\\
\text{i.e. commutant / topology / boundary / response / refinement / parent completion}.
\end{gathered}
}
\]

Thus the same selector machinery also provides a finite classification criterion for admissible physical subtheories.  Ordinary theories appear as sectors, compressions, fixed-point algebras, low-boundary regimes, or continuum reconstructions of the finite carrier.  Extra continuum structures require the corresponding finite witness criterion.



\section{Comparison branches, conditional numerical targets and scope checks}
\label{p13:sec:v237-d7}
The previously recorded $E_{36}$/$E_{40}$ split-grammar formulas and nuclear transfer ansatz are preserved explicitly as non-derivational phenomenological targets. Any future promotion requires a computed common-action QCD operator and a specified physical observable map.
% Source Paper VII-A, original line 6855; inherited proof bodies unchanged.
\subsection{Conditional screening and source-regularity comparisons}
\label{p10:app:v2187-readout-domain-comparisons}
These comparisons protect the domains of the main readout statements.
The screening argument assumes actual same-owner static-source trial states
and an energy estimate that this paper does not construct.  The regularity
example tests a mathematical implication and is not a QCD action member.

\begin{proposition}[Conditional static-light variational screening test]
\label{p10:prop:v2187-conditional-static-light-screening}
For a declared full-QCD static-source Hamiltonian $H_R^{\rm stat}$ at separation
$R$, use one common static self-energy subtraction and let $V(R)$ be its
ground-energy readout in that convention.  Suppose its complete admissible
domain contains normalized two-static-light-singlet trial vectors $\psi_R$
with
\[
 \langle\psi_R,H_R^{\rm stat}\psi_R\rangle
 \le B_{\rm sl}+\varepsilon_R,
 \qquad B_{\rm sl}<\infty,\qquad \varepsilon_R=o(R),
\]
and suppose $V(R)\ge-B_{\rm low}$ with $B_{\rm low}<\infty$ independent of $R$.
Then any asymptotic slope defined by these ground-state readouts is zero:
\[
 \lim_{R\to\infty}\frac{V(R)}{R}=0.
\]
\end{proposition}
\begin{proof}
Rayleigh--Ritz gives $V(R)\le B_{\rm sl}+\varepsilon_R$.  Dividing the upper and
lower bounds by $R$ and using $\varepsilon_R/R\to0$ proves the statement.
The color pairing $C\otimes C^*=\mathbf1\oplus\mathbf8$ permits a singlet
channel; it does not supply the trial vectors, their energy bound, locality,
mixing, or overlap with a selected Wilson source.  A Wilson-source extraction
equals this complete ground-state slope only after its spectral-support and
ordered-limit identification is established.
\end{proof}

The proposition neither contradicts the pure-glue obstruction of
Theorem~\ref{p10:thm:pure-glue-center-obstruction-screening} nor proves native
full-QCD screening.  It shows why a positive unscreened area theorem cannot
be imposed on every full-QCD ground-state or finite hadron observable without
its sector hypothesis.

\begin{remark}[Twice differentiable even source shapes]
\label{p10:rem:v2187-c2-source-remainder-comparison}
The even function
\[
 E(\theta)=E(0)+\tfrac12\chi\theta^2+|\theta|^3
\]
is $C^2$ near zero with $E''(0)=\chi$, but its remainder is not
$O(\theta^4)$.  Thus the $o(\theta^2)$ remainder in
Corollary~\ref{p10:cor:native-full-qcd-theta-shape} uses precisely its stated
second-order regularity; the stronger fourth-order remainder needs the
additional regularity printed there.
\end{remark}

% END INLINE SOURCE: Paper7A-ReadoutDomainComparisons-Appendix-v1_0.tex



% Source Paper VII-A, original line 8089; inherited proof bodies unchanged.
\subsection{Native clock-sector response and source-domain comparisons}
\label{p10:app:v2189-native-clock-comparison}

The positive clock construction is in
Section~\ref{p10:sec:v2189-native-clock-source-component}.
This appendix retains its fixed-background positive vectorlike branch and
actual source/volume conventions. It distinguishes its temporal components
from a finite ordered condensate, a uniform sharp-band source chart and an
identification from retained link data alone. These comparisons are not new
prerequisites for already covered finite massive QCD readouts.

\begin{proposition}[Clock index density is not a finite condensate]
\label{p10:prop:v2189-clock-zero-mode-scope}
The actual native clock sectors have $\nu_k=k$, hence at least $k$ exact
overlap zero modes. On their fixed-background positive scalar branch,
\[
 \mathcal V_{4,k}^{-1}\operatorname{Tr}(D_{k,m}^{-1}R_k)
 \ge\frac1{\mathcal V_{4,1}m}.
\]
Thus these fixed sectors have no finite ordinary massless scalar response
limit. For $N_f$ equal positive masses $m$, the determinant retains the
zero-mode factor $m^{N_fn_{0,k}}$; for unequal masses the factor is
$\prod_{f=1}^{N_f}m_f^{n_{0,k}}$. Actual full action/measure weights must
be retained before interpreting a normalized ensemble response.
\end{proposition}
\begin{proof}
Use Theorem~\ref{p10:thm:v190-generated-4d-gw-regulator}, the actual
index of Theorem~\ref{p10:thm:v199-native-clock-covers}, and the exact
Cayley response of
Theorem~\ref{p10:thm:v2158-overlap-cayley-scalar-response}.
Every zero mode has unit scalar insertion and contributes $1/m$; every
paired contribution is nonnegative. Divide by
$\mathcal V_{4,k}=k\mathcal V_{4,1}$.
The determinant zero-mode factor and normalized-sector cancellation are
exactly Corollary~\ref{p10:cor:v2158-finite-volume-zero-mode-firewall}.
The uniform clock gap is a Wilson sign-kernel gap, not a physical fermion
mass gap or many-body transfer gap.
\end{proof}


% Source Paper VII-A, original line 8128; inherited proof bodies unchanged.
\subsubsection{Temporal extent and the distinct spatial constructor}
The clock propositions supply geometric gluing, compatible temporal
Euclidean CAR and positive-mass periodic scalar-source components. They do
not assign independent physical equal-time quark modes to time layers.
A fixed finite parent and its restrictions have rank bounded by that parent,
so that construction alone cannot give positive spatial mode density when
the spatial volume diverges. The distinct
Theorem~\ref{p10:thm:v2189-native-orbital-quark-constructor} instead uses
the already constituted global quark field and complete native orbital
bands; Proposition~\ref{p10:prop:v2189-spatial-car-density} proves their
spatial density and exact exterior-CAR maps.
The actual QCD field/dual, source, gauge/constraint, central/sea and clock
comparison remains a consumer of that constructor. No full physical
$\Sigma>0$, residual-pseudoscalar susceptibility, Goldstone or GMOR limit
is inferred from temporal degree or spatial rank alone.


% Source Paper VII-A, original line 8144; inherited proof bodies unchanged.
\subsubsection{Sharp spectral bands need their actual source chart}
Even a supplied global compact self-adjoint elliptic Dirac owner with exact
gauge-domain covariance does not by itself give a common differentiable
sharp-band source chart. On a two-dimensional chiral block, set
\[
 D_L(h)=(r-\varepsilon_L+h)\sigma_1,
 \qquad r>0,\quad 0<\varepsilon_L<r,\quad\varepsilon_L\downarrow0.
\]
Then $E_L(h)=\mathbf1_{[0,r^2]}(D_L(h)^2)$ has rank two for
$0\le h<\varepsilon_L$ and rank zero for $h>\varepsilon_L$ sufficiently
close to that crossing. Thus no level-independent nonzero source interval
is automatic. This finite example is self-adjoint, has compact resolvent
and trivial exact gauge covariance. It only refutes the inference of a
uniform source chart from those properties.
It does not prohibit a licensed component restriction or a sharp
restriction with a proved cluster gap and retained projector jets.

The precise source-chart scope is
Remark~\ref{p10:rem:v2186-band-source-scope}.
Differentiating the actual field/dual word retains all three terms
\[
 \partial_h(J^\dagger D J)
 =J_h^\dagger DJ+J^\dagger D_hJ+J^\dagger DJ_h,
\]
and the kinetic, interaction, projector and central/sea contacts of
Theorem~\ref{p3f:thm:v1148-finite-metric-constructor}.


% Source Paper VII-A, original line 8171; inherited proof bodies unchanged.
\subsubsection{Retained holonomy does not determine global compression}
Theorem~\ref{p10:thm:hurwitz-shell-dirac-compression} licenses a
point stencil with first-symbol consistency and explicitly distinguishes
it from exact $L^2$ Galerkin compression. A finite retained edge list leaves
an open region disjoint from those edges. A smooth connection variation
supported there preserves every retained edge holonomy while it can change
$P c(\delta a)P$ in a global local-Dirac compression. This is the scope of
Proposition~\ref{p10:prop:v2180-finite-holonomy-dirac-no-go}.
It excludes an identification from that finite link list alone, rather
than native global Dirac geometry or an actual same-owner source comparison.


% Source Paper VII-A, original line 8182; inherited proof bodies unchanged.
\subsubsection{Global geometry and the licensed quantum field word}
The canonical global pullback supplies the existing bundle gluing, and
III-F already specializes its generated joint operator to $S_R^3$ with
the unique spin structure and self-adjoint elliptic domains.
The joint-operator theorem alone does not select a physical action, field
member or normalization. The positive spatial construction instead consumes
the already licensed quark-action/statistics branch and its actual
finite-metric field/dual word. It therefore does not infer a regulated
physical measure from global geometry alone or require a new arbitrary
sharp physical cutoff. Its remaining identification with a particular QCD
source/sea/clock word has the scope stated in
Appendix~\ref{p10:app:v2189-native-spatial-quark-comparison}.


% Source Paper VII-A, original line 8195; inherited proof bodies unchanged.
\subsubsection{Divisor cofinality of compatible clock covers}
The sequence $3^j$ is unbounded but not cofinal in the odd-integer divisor
order, since $5$ never divides $3^j$. The positive clock construction uses
$k_j=(2j+1)!!$. Every odd positive integer divides a sufficiently large
odd double factorial, and each consecutive quotient is odd, preserving
the fixed antiperiodic sign. This is clock-cover cofinality, not spatial
cofinality.

% END INLINE SOURCE: Paper7A-NativeClockSourceComponent-Appendix-v1_0.tex


% BEGIN INLINE SOURCE: Paper7A-NativeSpatialQuarkConstructor-Appendix-v1_0.tex
% Framework v2.189: native spatial quark constructor appendix module.

% Source Paper VII-B, original line 4326; inherited proof bodies unchanged.
\subsubsection{Hurwitz split-grammar phenomenological targets (not yet derived)}
\label{p11:subsec:split-grammar-numerical-diagnostics}

\textbf{Diagnostic-only boundary.}
This subsection preserves numerical coincidences from the predecessor for reproducibility, but it is not part of the native derivation chain and supplies no evidence for a unique QCD operator map.  None of the rows below may be used to choose a source member, transfer coefficient, mass operator, or nuclear interaction.

The preceding finite count is independent of numerical spectroscopy.  Nevertheless, the auxiliary inverse split-grammar ledger gives a useful check on whether the same finite supports that produce the shell-count ranks also give diagnostic proximity to established QCD and nuclear scales.  These rows are recorded here only as diagnostics.  They are not used as inputs in the proof of Theorem~\ref{p11:thm:native-boundary-reduced-shell-count}, and the particle rows already fixed upstream by Paper~V-B are not counted as independent confirmations.

The historical ledger used the representation-specific Paper-V-B label-carrier ratio
\[
        \eta_{\rm lab}=\frac{24}{25}
\]
and the shell scales \(E_{32},E_{36},E_{40},E_{56}\).  The predecessor interpreted these diagnostics through the following role assignment (not derived in current Paper~VII-B):
\[
        E_{40}
        \quad\text{carries the visible first-family }u,d\text{ split grammar,}
\]
whereas
\[
        E_{36}
        \quad\text{carries the composite QCD/hadronic interface.}
\]
The displayed association is retained solely to reproduce the predecessor's arithmetic.  Current Paper~VII-B neither promotes it to an operator ownership statement nor uses it in any theorem.

The strongest nontrivial rows in the ledger are:
\[
\begin{array}{c|c|c|c|c}
\text{target} & \text{grammar expression} & \text{value} & \text{target value} & \text{relative error}\\
\hline
\sigma_{\pi N}
&
E_{36}\eta_{\rm lab}/\sqrt{24}
&
43.0285207
&
43.02
&
1.98\times10^{-4}
\\
m_0
&
E_{36}\bigl(\sqrt4+\sqrt4\bigr)
&
878.316
&
878.0384434
&
3.16\times10^{-4}
\\
M_{\rho/\omega}
&
E_{36}\eta_{\rm lab}^{1/2}\sqrt{13}
&
775.707688
&
773.655
&
2.65\times10^{-3}
\\
M_\phi
&
E_{36}\eta_{\rm lab}\sqrt{24}
&
1032.68450
&
1026.65
&
5.88\times10^{-3}
\\
M_N^{\rm iso}
&
E_{36}\eta_{\rm lab}^{-3/2}\bigl(\sqrt4+\sqrt4\bigr)
&
933.778660
&
939.603
&
6.20\times10^{-3}
\\
\sigma_{sN}
&
E_{36}\eta_{\rm lab}^{-1}/\sqrt{24}
&
46.6889330
&
46.335
&
7.64\times10^{-3}
\\
B_d
&
E_{40}\sqrt2
&
2.20064404
&
2.224566
&
1.08\times10^{-2}
\\
B_\alpha
&
\displaystyle
\frac{E_{36}\eta_{\rm lab}^{-3/2}}{\sqrt6+\sqrt{12}+\sqrt6}
&
27.9137153
&
28.295674
&
1.35\times10^{-2}
\end{array}
\]
All masses and binding energies in the table are in MeV.  The displayed expressions have been parenthesized as finite path sums; for example \(E_{36}(\sqrt4+\sqrt4)\) is a two-vector path, not \(E_{36}\sqrt4+\sqrt4\).

These numerical coincidences were historically interpreted through the following separation of roles; current Paper~VII-B does not derive that interpretation.  The \(E_{40}\) split
\[
        6=1+5,\qquad 5=2+3
\]
is the visible first-family grammar.  It gives the light \(u,d\) branch content and, after passing through
\[
        p=uud,\qquad n=udd,
\]
the projected baryonic support invariants
\[
        B_2,\quad T_3,\quad B_2\oplus T_3,\quad 3B_2,\quad \Phi_4(B_2\oplus T_3),
\]
with ranks
\[
        2,\quad3,\quad5,\quad6,\quad20.
\]
These are the ranks that appear in the low and medium closed-shell stiffness benchmarks.  By contrast, the \(E_{36}\) rows cluster around scalar, vector, topological, and nucleon-scale QCD quantities.  This clustering is only a retrospective comparison pattern.  It does not establish that \(E_{36}\) owns a QCD mass operator or that \(E_{40}\) owns a light-flavor mass source.

The predecessor also recorded the following conjectural nuclear-transfer ansatz:
\[
        E(2_1^+;Z,N)^2
        \approx
        E_{40}^{\,2}\,
        \frac{k_{\rm sel}(Z,N)}
        {D_{32}(Z,N)} ,
\]
where \(k_{\rm sel}(Z,N)\) is selected from the projected \(E_{40}\)-baryonic branch ranks, while \(D_{32}(Z,N)\) denotes the finite \(E_{32}\)-level proton/neutron shell-transfer softening away from exact closure.  At exact closed shell \(D_{32}=1\).  Away from closure, valence proton and neutron pairs activate finite transfer channels and lower the first \(2^+\) energy; the dual amplitude-side readout is expected to increase accordingly.

This remains a historical conjecture, not a theorem or completed nuclear-structure derivation.  To promote it beyond a diagnostic, the finite \(A\)-body transfer matrix defined in Paper~VII-B must derive the map
\[
        (Z,N)\longmapsto k_{\rm sel}(Z,N)
\]
and the softening denominator \(D_{32}(Z,N)\) without fitting them from the observed \(2^+\) energies.




% Source Paper VII-C, original line 8365; inherited proof bodies unchanged.
\subsection{continuum selector synchronization}
\label{p12:sec:v286-safeguard-continuum-selector-sync}
Paper VII-C is a reconstruction and certification layer.  It neither reselects
the actual finite fermion source/mass member nor introduces an independent
thermal scalar.  On the rooted charged branch the finite source assembly is
already generated by Paper~V-A,
Corollary~\ref{p5:cor:v397-charged-preterminal-assembly-closure};
it needs no second G0 member or noncentral assembly selector.  Other declared
source branches retain their own upstream certificates.  G16 supplies
RG/pole transport when required, and a supplied inverse
temperature has the finite KMS target theorem on its stated domain.  Physical
history and thermalization remain separately scoped dynamical claims.
% END SAFEGUARD-CONTINUUM-SELECTOR-SYNC v2.86


% BEGIN GAP-TAXONOMY-CONTINUUM v2.87
\begin{remark}[Finite internal source gap versus continuum sector gap]
\label{p12:rem:v287-finite-source-versus-sector-gap}
A positive finite/internal fermion source gap is not the same claim as a
positive gap of the entire interacting charged Hamiltonian sector.  Fixed
internal full rank survives ordinary external position refinement as
$I_{\mathcal P_x(L)}\otimes T$, but a sector may still contain massless gauge
quanta.  QFT9-type sector-gap inheritance and a sharp charged-particle pole
therefore remain separately typed analytic claims.  This distinction prevents
a finite fermion gap from being over-promoted into a photon-containing sector
gap.
\end{remark}
% END GAP-TAXONOMY-CONTINUUM v2.87


% BEGIN SOURCE-DUAL-G16 v2.88

% Source Paper VII-C, original line 8396; inherited proof bodies unchanged.
\subsection{source-dual G16 comparison and finite normalization}
\label{p12:sec:v288-source-dual-g16}

\begin{definition}[Native scalar-source dual pairing]
\label{p12:def:v288-native-scalar-source-dual}
Let $T^{\rm BH}$ be an already occurring canonically normalized finite
fermion mass/source coordinate and let $\mathcal O_m^{\rm BH}$ be its native
source derivative in the minimal source dual of
Theorem~\ref{p12:thm:minimal-source-quotient}.  Their source insertion is the
pairing
\[
 \mathfrak I_m=\langle T^{\rm BH},\mathcal O_m^{\rm BH}\rangle.
\]
This definition transports rather than selects $T^{\rm BH}$.  Its upstream
source certificate is reused: the rooted charged assembly is already closed,
while a distinct generic or neutral target must declare its own source domain.
No claim-specific interacting comparison or physical pole is evaluated here.
\end{definition}

\begin{proposition}[Contragredient finite-to-continuum source transport]
\label{p12:prop:v288-contragredient-mass-source-transport}
Suppose a named continuum scheme $S$ represents the same scalar insertion by
\[
 \mathcal O_m^S(\mu)=C_S(\mu)\mathcal O_m^{\rm BH}
\]
on the claim-relevant finite source space.  Invariance of the pairing forces
\[
 \boxed{M^S(\mu)=C_S(\mu)^{-\dagger}T^{\rm BH}}.
\]
For a one-dimensional or flavor-blind scalar block this reduces to
$m_q^S(\mu)=C_S(\mu)^{-1}t_q^{\rm BH}$.
\end{proposition}

\begin{corollary}[Mass RGE as dual operator transport]
\label{p12:cor:v288-mass-rge-dual-transport}
Let
\[
 \Gamma_S(\mu):=\mu\,\frac{dC_S}{d\mu}C_S^{-1}.
\]
Then
\[
 \boxed{\mu\frac{dM^S}{d\mu}=-\Gamma_S^\dagger M^S}.
\]
Thus ordinary running changes the continuum coordinate of an already owned
finite source; it does not generate a family hierarchy or repair missing
source rank.
\end{corollary}

\begin{proposition}[Finite-normalization degeneracy]
\label{p12:prop:v288-finite-normalization-degeneracy}
For every constant invertible finite comparison map $K$, replacing
$C_S(\mu)$ by $KC_S(\mu)$ leaves its logarithmic anomalous-dimension equation
unchanged up to the corresponding similarity action.  In the scalar case,
$C_S(\mu)\mapsto\kappa C_S(\mu)$ leaves
$d\log C_S/d\log\mu$ unchanged.  Hence beta functions and anomalous dimensions
fix relative running but cannot by themselves determine the finite
BH-to-$S$ normalization or the scale at which $C_S(\mu)=I$.
\end{proposition}

\begin{remark}[Split G16 frontier]
\label{p12:rem:v288-split-g16-frontier}
For running-mass claims, distinguish the RG/threshold-flow certificate from
the finite scalar-source comparison certificate.  The former transports a
fixed scheme coordinate once its boundary value is supplied; the latter
fixes the native-BH-to-named-scheme finite normalization.  A pole certificate
remains a third, claim-scoped continuation/isolation/residue problem.  No
numerical $2\,{\rm GeV}$ or $E_{32}/\sqrt{240}$ identification is asserted by
this release.
\end{remark}
% END SOURCE-DUAL-G16 v2.88

% BEGIN G16-NATIVE-SCALAR-NORMALIZATION v2.89

% Source Paper VII-C, original line 8468; inherited proof bodies unchanged.
\subsection{native scalar normalization and finite scheme-vertex frontier}
\label{p12:sec:v289-native-scalar-normalization}

\begin{proposition}[Native overlap scalar-source normalization]
\label{p12:prop:v289-native-overlap-scalar-normalization}
On the declared massive overlap branch of Paper~VII-A v1.65, the native source
derivative is
\[
 \boxed{\mathcal O_m^{\rm BH}=\frac12(I-\Gamma_5\epsilon_H),}
\]
and is independent of the overall overlap normalization $a_R$.  Hence $a_R$
is not an unresolved finite mass/source normalization parameter on this
branch.
\end{proposition}

\begin{corollary}[Tree-level source coordinate is fixed internally]
\label{p12:cor:v289-tree-source-coordinate}
Because the scalar insertion is the derivative of the same finite massive
Dirac operator with respect to its source coordinate, there is no independent
tree-level rescaling of that native mass coordinate.  Any named-scheme
conversion therefore has the form
\[
 C_{S\leftarrow {\rm BH}}(\mu)=I+\text{finite interacting corrections}
\]
in a common external-leg normalization.  This statement does not evaluate the
interacting correction.
\end{corollary}

\begin{remark}[Historical v2.89 frontier, superseded below]
\label{p12:rem:v289-refined-g16-frontier}
At release v2.89 the native scalar normalization was already closed while the
finite interacting scheme vertex remained the next structural target.  Release
v2.93 below supplies that constructor on its certified noncritical comparison
domain.  The finite-normalization-degeneracy proposition remains relevant: RG
running transports a finite comparison value but does not manufacture it, and
no $E_{32}/\sqrt{240}$ unity-crossing theorem is asserted.
\end{remark}
% END G16-NATIVE-SCALAR-NORMALIZATION v2.89


% BEGIN G16-FINITE-SCHEME-CONSTRUCTOR-CLOSURE v2.93

% Source Paper VII-C, original line 8509; inherited proof bodies unchanged.
\subsection{finite G16 constructor closure and native/source-scheme separation}
\label{p12:sec:v293-g16-finite-constructor-closure}

Paper~VII-A, Theorem~\ref{p10:thm:v167-overlap-resolvent-accretive}, proves that
on the positive massive-overlap branch the full finite interacting propagator
is invertible with a uniform accretive lower bound.  The rooted source-transport
module then derives the complete normalized-measure scalar response,
\[
 \Lambda_{m,L}=\partial_mS_L^{-1},
\]
including the determinant/sea covariance term, and defines the finite
BH-to-named-scheme comparison on every common noncritical chart by
\[
 \boxed{
 C_{S\leftarrow{\rm BH}}
 =V_SV_{\rm BH}^{-1}.}
\]
The comparison is an owned coordinate transformation of an already generated
source insertion; it is not a new mass-family selector.

\begin{theorem}[G16 structural status]
\label{p12:thm:v293-g16-structural-status}
On the declared positive massive-overlap branch and on every certified common
noncritical comparison chart,
\[
 \boxed{\texttt{G16\_NATIVE\_SCALAR\_NORMALIZATION=CLOSED\_ANALYTIC}},
\]
\[
 \boxed{\texttt{G16\_INTERACTING\_VERTEX\_CONSTRUCTOR=CLOSED\_FINITE}},
\]
\[
 \boxed{\texttt{G16\_FINITE\_SCHEME\_CONSTRUCTOR=CLOSED\_ON\_DOMAIN}}.
\]
For a claim that requests a specific external coordinate,
\[
 \boxed{\texttt{G16\_NAMED\_SCHEME\_VALUE}[S,\mu,\kappa]=\texttt{EVALUATION}},
\]
and a pole/mass-shell statement separately consumes
\[
 \boxed{\texttt{G16\_POLE=CLAIM\_SCOPED}}.
\]
The historical aggregate token
\texttt{G16\_FINITE\_SCHEME\_VERTEX} is retired as too coarse because it
conflated a closed constructor with its claim-specific numerical evaluation.
\end{theorem}

\begin{corollary}[Native physical source evolution bypasses named-scheme coordinates]
\label{p12:cor:v293-native-source-evolution-bypasses-scheme}
The exact finite coarse-pushforward relation
\[
 \partial_aS_L^{\rm eff}=\langle\partial_aS_M\rangle_x
\]
transports an owned native source through the effective action under its stated
hypotheses.  Therefore an internal physical calculation does not require an
intermediate $\overline{\rm MS}$ or other named-scheme coordinate unless the
claim explicitly asks for that coordinate.  This does not close the separate
physical running-scale, threshold, kinetic-history or phase-selection problems
of downstream cosmology.
\end{corollary}
% END G16-FINITE-SCHEME-CONSTRUCTOR-CLOSURE v2.93

% BEGIN CHARGED-SOURCE-CONTINUUM-SYNC v2.94

% Source Paper VII-C, original line 8646; inherited proof bodies unchanged.
\subsection{electromagnetic no-free-matching synchronization}
\label{p12:sec:v2100-g17-no-free-matching-sync}

\begin{corollary}[Retained-source naturality forbids a second electromagnetic normalization]
\label{p12:cor:v2100-g17-no-free-matching}
On the same-source branch of Paper~V-C Release v2.75, exact blocking
naturality,
\[
 D_A^2\mathscr R_{NL}S_N[A_Q]=D_A^2S^{\rm eff}_{N\downarrow L}[A_Q],
\]
transports the electromagnetic Hessian without changing the retained source
coordinate.  Therefore a finite threshold or interacting correction belongs
inside the resulting $K_{Q,L}$ and $Z_Q(L,q^2)$; it cannot be represented as
an additional independently fitted rescaling of $A_Q$ while claiming the same
physical descendant.  In particular, named-scheme finite matching constants
are coordinate data for an effective-action comparison and do not constitute
a second native electromagnetic source normalization.
\end{corollary}

\begin{corollary}[Narrowed Thomson frontier]
\label{p12:cor:v2100-thomson-frontier}
After the no-free-matching theorem, the unresolved physical endpoint is the
controlled evaluation of
\[
 \boxed{
 Z_Q^{\rm phys}(0)
 =\lim_{q^2\to0}\lim_{k\to0}Z_Q(k,q^2),}
\]
with the order replaced by the common joint limit when uniformity proves their
interchange.  The current-weighted theorem and its uniform-transport
corollary supply a sufficient temporal-chart criterion, but the actual QCD
current measure, contact contribution, joint external-momentum regularity and
continuum/refinement bounds remain claim-consumed obligations.  No further
absolute source-normalization constant is part of this frontier.
\end{corollary}
% END G17-ENDPOINT-RESPONSE-SYNC v2.100

% framework v2.178: exact source coordinate retained as descendant of finite K16 production owner.

% BEGIN INLINE SOURCE: Paper7C-K16ExactSourceTransport-Main-v1_1.tex
% Framework v2.177: exact native K16 gauge-source transport.



% BEGIN SOURCE-INTEGRATED S3 FINITE RESPONSE v1_59 (18 November 2026)
\section[S$^3$ finite-response projectors and conditional action-owner tests]{S$^3$ finite-response projectors\\and conditional action-owner tests}
\label{p13:sec:s3-finite-response-v159}

This section is an exact \emph{finite mathematical comparison} on the
marked massless Hurwitz $24$-cell spin stencil, not a construction of the
unselected physical position-sector spectral owner.  The original
matrix, oriented incidence, integer witnesses, exact certificates and
independent replay code are retained in
\texttt{Paper00-S3DAGInvariantBridge-v1\_16.zip}.  In the propositions
below, the words ``fixed stencil'', ``fixed color Cartan background'',
and ``fixed carrier'' are essential hypotheses.  Upstream geometric
results, existing physical constructors and their scopes are unchanged.

\begin{definition}[Comparison curl, spin stencil and blind sector]
\label{p13:def:s3-comparison-curvature-v159}
Let $B_1:\mathbb R^{96}\to\mathbb R^{24}$ and
$B_2:\mathbb R^{96}\to\mathbb R^{96}$ be the oriented
one-boundary and two-boundary operators of the actual Hurwitz cell
complex, and put $L=B_2B_2^\mathsf T$.
Let $D_H(U)$ denote the \emph{massless} $48\times48$ round-spin
stencil on the fixed $24$-root comparison chart, with the standard
fundamental SU(3) Cartan path on its $96$ links; its unperturbed
spin-only squared operator $X_0=D_H(I)^2$ has eigenvalues $4,16$.
Define $P_{\mathrm b}$ as the exact orthogonal rank-eight
\emph{link-space} projector specified by the companion integer
construction (equivalently
$P_{\mathrm b}=P_5(L)(I-Z_4/4)$ with the archived commuting
integer matrix $Z_4$).  No physical gauge-field or QCD Hamiltonian
occurrence is asserted by this definition.
\end{definition}

\begin{proposition}[Exact transverse curvature and comparison rank]
\label{p13:prop:s3-comparison-rank-v159}
On the fixed chart of Definition~\ref{p13:def:s3-comparison-curvature-v159},
\[
B_1B_2=0,\quad
\operatorname{rank}B_1=23,\quad
\operatorname{rank}B_2=73,\quad
B_1P_{\mathrm b}=0,\quad LP_{\mathrm b}=5P_{\mathrm b}.
\]
The stipulated uniform positive triangular-face action has quadratic
Hessian $H_{\rm face}=\kappa L\otimes I_8$ for $\kappa>0$,
with 23 gauge-zero and 73 positive link directions per color.
The specified spin-stencil scalar heat action has 31 zero and
65 positive link directions per color, for every positive heat time;
eight of its zeros are precisely $\operatorname{ran}P_{\mathrm b}$.
\end{proposition}
\begin{proof}
The chain-complex identity, ranks and projectors follow by exact
integer elimination and matrix multiplication on the oriented
24-cell incidence arrays.  Diagonalization in their commuting
integer projector algebra yields the stated Hessian ranks and
strict positivity of the 65 nonzero heat modes; the exact matrices
and replay checks are part of the cited finite companion.
No numerical rank threshold or full physical-action identification
is used.
\end{proof}

\begin{theorem}[Exact Duhamel--contact cancellation on eight modes]
\label{p13:thm:s3-eight-mode-contact-v159}
In the massless spin-only comparison write $X(U)=D_H(U)^2$,
$S_\tau(U)=\tfrac12\operatorname{Tr}e^{-\tau X(U)}$,
and $P_4,P_{16}$ for the unperturbed spin spectral projectors.
For all real link tangents $u,v\in\operatorname{ran}P_{\mathrm b}$,
\[
\operatorname{Tr}(X'_uX'_v)=\frac{208}{5}\langle u,v\rangle,\quad
\operatorname{Tr}(P_4X''_{uv})=\frac{52}{15}\langle u,v\rangle,\quad
\operatorname{Tr}(P_{16}X''_{uv})=-\frac{52}{15}\langle u,v\rangle.
\]
The first-jet blocks $P_4X'_uP_4$ and
$P_{16}X'_uP_{16}$ vanish but $X'_u\ne0$ if $u\ne0$.
For every $\tau>0$ the complete second heat response is
\[
\begin{aligned}
G_\tau(u,v)&=\frac{26}{15}\tau
  (e^{-4\tau}-e^{-16\tau})\langle u,v\rangle,\\
C_\tau(u,v)&=-G_\tau(u,v),\\
\boxed{D^2S_\tau[u,v]}&=0.
\end{aligned}
\]
Hence the quadratic cancellation results from two nonzero terms,
not from an absent operator variation.
\end{theorem}
\begin{proof}
The first displayed projector-sandwich identities are finite exact
matrix identities over the specified integer/Gaussian-integer
realization; see the independent integer replay retained in the
companion.  In the eigenbasis of $X_0$ the first jet has only
$4\leftrightarrow16$ blocks.  Substituting the exact identities
into the finite-dimensional Duhamel Hessian, \emph{including}
$X''_{uv}$, gives the two equal and opposite coefficients displayed.
Their common magnitude is strictly positive for $\tau>0$ and
$\langle u,u\rangle>0$.
\end{proof}

\begin{corollary}[Analytic scalar spectral profiles remain quadratic-blind]
\label{p13:cor:s3-analytic-profile-blind-v159}
On that same fixed massless spin stencil, for every scalar function
$f$ analytic in a neighborhood of the unperturbed spectrum
$\{-4,-2,2,4\}$,
\[
\boxed{D^2\operatorname{Tr}f(D_H)[u,v]=0
\qquad (u,v\in\operatorname{ran}P_{\mathrm b}).}
\]
In particular, changing only the scalar spectral profile or adding
a fixed commuting scalar mass cannot restore the missing quadratic
curvature response in this comparison model.
\end{corollary}
\begin{proof}
The explicit integer calculation certifies the complete projected
second link variations of the eleven moments
$\operatorname{Tr}D_H^n$ for $1\le n\le11$ as zero.
The second variation of $\operatorname{Tr}(z-D_H)^{-1}$ has poles only
at the four unperturbed eigenvalues, each of order at most three,
with a rational numerator of degree at most ten after collecting
a common degree-twelve denominator.  The eleven zero Laurent
coefficients force this rational function to vanish identically.
Cauchy's analytic functional calculus gives the claim.
No noncommuting physical mass or different coupled action sector is
excluded.
\end{proof}

\begin{proposition}[Strictly positive quartic response on a blind witness]
\label{p13:prop:s3-blind-quartic-v159}
For the archived integer link tangent $v$ satisfying
$P_{\mathrm b}v=v$ and $\|v\|^2=120$, and the SU(3)
Cartan generator $T=\operatorname{diag}(i/2,-i/2,0)$ with
$-2\operatorname{tr}T^2=1$, let
$U_e(s)=\exp(sv_eT)$.
For the fixed color-extended comparison heat action,
\[
\boxed{\left.\frac{d^4}{ds^4}S_\tau(U(s))\right|_{s=0}
=625\tau\bigl[(12\tau-1)e^{-4\tau}+e^{-16\tau}\bigr]>0
\quad(\tau>0).}
\]
Its small-$\tau$ leading coefficient is $45000\tau^3$.
\end{proposition}
\begin{proof}
The exact length-2 through length-18 closed-walk moments of the
finite stencil determine the fourth link variation of the trace
resolvent. Since its squared Dirac spectrum has only the poles
$4$ and $16$ and order at most five, nine successive integer
moments determine the rational function uniquely.  Residues of
that rational function give the displayed all-$\tau$ formula.
For positivity put $y=12\tau$: multiplying the bracket by
$e^{16\tau}$ yields $(y-1)e^y+1$, which vanishes at zero and has
positive derivative $ye^y$ for $y>0$.
The calculation does not select a physical QCD or Higgs coupling.
\end{proof}

\begin{proposition}[Conditional lower bound on same-carrier Hessian matching]
\label{p13:prop:s3-positive-face-matching-bound-v159}
\emph{If} the physical problem separately establishes that the full
same-carrier Hessian must equal the stipulated uniform face Hessian,
then any net correction $\Delta H=H_{\rm face}-H_{\rm spin}$
must obey on eight link modes and eight SU(3) color directions
\[
\begin{gathered}
(P_{\mathrm b}\otimes I_8)\Delta H(P_{\mathrm b}\otimes I_8)
=5\kappa(P_{\mathrm b}\otimes I_8),\\
\boxed{\operatorname{rank}\Delta H\ge64,\quad
\|\Delta H\|_{\rm op}\ge5\kappa,\quad
\|\Delta H\|_{\rm HS}\ge40\kappa.}
\end{gathered}
\]
These bounds are not a derived extra field, a new interaction term,
or evidence that the two tested actions are in fact the same owner.
\end{proposition}
\begin{proof}
The spin Hessian vanishes on $\operatorname{ran}P_{\mathrm b}$
and $LP_{\mathrm b}=5P_{\mathrm b}$, so the displayed compression
is forced.  Rank and operator- and Hilbert--Schmidt-norm inequalities
follow by orthogonal compression to the 64-dimensional tensor-product
subspace.
\end{proof}

\begin{definition}[Physical response target and source-typing firewall]
\label{p13:def:s3-physical-response-target-v159}
The \emph{physical} response of the complete generated action along
a Hurwitz curvature tangent $v$ is defined only if the source
provides either (i) a complete action-preserving configuration map
$T:\mathcal U_H\to\mathcal Q_x$ or (ii) a directly selected
physical action $S_{\rm phys}$ on the same finite carrier.
In case (i) the required response is
\[
\boxed{C_{\rm phys}(v)=D^2(S_{\rm Born}\circ T)[v,v]
=D^2S_{\rm Born}[DTv,DTv]
 +DS_{\rm Born}[D^2T[v,v]].}
\]
It must include every moving source, localizer, cutoff and trace
term, and it must use the actual Gauss-compatible physical carrier
and normalization for physical correlators or poles.
No theorem here asserts the existence of $T$, a numerical value for
$C_{\rm phys}$, or the equality of the comparison Hessians.
\end{definition}
\begin{remark}[Proof-chain stage]
\label{p13:rem:s3-physical-stage-v159}
The finite S$^3$ projectors and response identities above are
\emph{conditional mathematical results} on explicit finite
matrices.  The established position-sector Born action and the
momentum-shell QCD transfer retain their own physical roles.
Canonical G0 and GP1--GP3 are unchanged.  The graph may record
source/provenance or conditional mathematical interfaces but must
not turn these statements into unconditional physical mass,
transition, pole or decay results.  To test a physical observable,
evaluate the full same-owner action response or its properly
constructed physical correlator—not similarity of the two
comparison matrix spectra.
\end{remark}
% END SOURCE-INTEGRATED S3 FINITE RESPONSE v1_59

\printbibliography

% BEGIN FRAMEWORK v2.344: TWISTED BARYON RESPONSE EXECUTION FRONTIER
\section{The $C_3$-twisted baryon-source execution frontier}
\label{p12:sec:v344-twisted-baryon-frontier}
\begin{proposition}[Typed inputs still required for a mass-pole calculation]
\label{p12:prop:v344-baryon-character-execution-ledger}
The exact mathematical inputs are (i) the twisted $2T$ quintet
of Theorem~\ref{p1:thm:v344-induced-hurwitz-quintet},
(ii) the normalized rank-$25$ coefficient projector
of Theorem~\ref{p1:thm:v344-twisted-25-gram},
(iii) its Petersen intertwiner
of Proposition~\ref{p1:prop:v344-petersen-quintet-intertwiner},
and (iv) the nonzero gauge-invariant edge-fan character source
of Proposition~\ref{p10:prop:v344-threeface-character-witness}.
They suffice to formulate, but not execute, the physical test in
Proposition~\ref{p11:prop:v344-same-action-curvature-extraction}.
An action-level completion must supply one admissible $H_L$ and
physical mass source, the $qqQ/qsQ/ssQ$ source matrices at fixed
spin/charge, their $\chi,\bar\chi$ two-point Grams, the
flavor-dependent coupling amplitudes, and the physical pole
resolvents at the same owner and state.
\end{proposition}
\begin{proof}
The geometric transfer has eigenvalue $1/5$ on the quintet,
so its \emph{squared norm} on a normalized pure quintet source is
$1/25$.  But even a hypothetical scalar-gap model with
$V_\chi=\alpha T_u\mathcal J$ would give a pole-dependent Schur
term of schematic form
$- |\alpha|^2/[25(\Delta-z)]$ on any \emph{separately selected}
flavor subspace.  Neither $\alpha$, $\Delta$, the flavor selector,
nor its normalization is supplied by the finite character data.
The physical color-singlet determinant makes the $\chi$ modes
eligible, not dynamically selected.
\end{proof}
\begin{remark}[Status of the hadron-mass numerical screen]
Normal/inverse radical factors and assignments among
$E_{32},E_{36},E_{40}$ remain exploratory discrete mass
hypotheses, not derived physical eigenvalues.  The apparent
heavy-baryon finite-difference coincidence with $-E_{36}/25$
is a falsifiable candidate, not a theorem or input to the
Hamiltonian.  It must be checked across matched spin, charge,
flavor, and heavy-quark families without post-selection.
\end{remark}
% END FRAMEWORK v2.344


% BEGIN FRAMEWORK v2.345: F40 PHYSICAL EXECUTION LEDGER
\clearpage
\section{F40 quaternionic-source verification and action frontier}
\label{p12:sec:v345-f40-action-ledger}
\begin{proposition}[What the F40 certificate verifies and does not verify]
\label{p12:prop:v345-f40-execution-ledger}
The finite proof source
\texttt{Proofs/Paper00-PetersenHodgePlucker-v1\_11-verify.py}
verifies the $60\leftrightarrow12\times5$ arrow bijection,
complete native readout ranks $11,16,21$ and the strictly positive
cycle-frame polynomial.  The mathematical findings imply neither
an occupied $2I$ source in a physical QCD state nor the existence
of an admissible color-gauge-compatible $\iota_L$.
The necessary next physical inputs remain: the actual
source-to-$Z$ map, explicit F40 embedding into the constituted
singlet Hilbert space, an interaction/link-domain witness,
$B_L$ leakage, full Hamiltonian/source resolvents, and pole
or susceptibility extraction in one and the same action.
The group-shell walk and $24$-cell single-chart exclusion
must not be reinterpreted as a numerical hadron mass formula.
\end{proposition}
\begin{proof}
The first statements are enumerated and verified by the
referenced exact certificate.  The remainder are typed
requirements of the existing Paper~VII-A source/response
construction, not consequences of finite incidence alone.
\end{proof}
% END FRAMEWORK v2.345

% BEGIN FRAMEWORK v2.350: HURWITZ HAAR, FINITE SELECTOR AND CCR VERIFICATION
\clearpage
\section{Exact reference moments and the corrected finite-to-CCR test}
\label{p12:sec:v350-technical-repair}
The following computations separate Haar-reference identities,
finite commutants, and the much stronger convergence of actual
physical state sequences.

\begin{proposition}[Face-flux Haar moments and normalized 96-face variance]
\label{p12:prop:v350-flux-reference-witness}
For $U\in SU(3)$ define
$z(U)=\operatorname{tr}(U)/3$ and
$F(U)=|z(U)|^2-\operatorname{Re}z(U)$.
Then $-1/4\le F(U)\le3/2$ and, for normalized Haar measure,
\[
 \boxed{\mathbb E[F]=\frac19,\qquad
        \operatorname{Var}(F)=\frac{11}{162}.}
\]
For the $96$ triangular Hurwitz faces with independent Haar
link variables, set
$\mathcal F_H=96^{-1}\sum_{f}F(U_f)$. Each pair
of \emph{distinct} face holonomies is independent under this
reference measure. Consequently,
\[
 \boxed{\mathbb E[\mathcal F_H]=\frac19,\qquad
        \operatorname{Var}(\mathcal F_H)=\frac{11}{15552}.}
\]
The three-face edge fan obeys
$-3/4\le\sum_{f\supset e}F(U_f)\le9/2$.
These are Haar/reference bounds, not interacting energy shifts.
\end{proposition}
\begin{proof}
The identity $F=|z-1/2|^2-1/4$ yields the lower bound;
$|z|\le1$ and $\operatorname{Re}z\ge-1/2$ for
$SU(3)$ yield the upper bound. The character identities
$\chi_3\chi_{\bar3}=1+\chi_8$ and
$\operatorname{Re}\chi_3=(\chi_3+\chi_{\bar3})/2$ give
\[
 F=\frac{1+\chi_8}{9}-\frac{\chi_3+\chi_{\bar3}}6.
\]
Haar character orthogonality gives mean $1/9$ and variance
$1/81+2/36=11/162$. Distinct triangular faces have an
exclusive edge factor in each face holonomy after conditioning
on any common edge; independent Haar multiplication then makes
those two face holonomies jointly independent Haar variables.
Hence off-diagonal face covariances vanish, and division by
$96^2$ gives variance $(11/162)/96=11/15552$.
The edge bound is the sum of the three individual bounds.
\end{proof}

\begin{theorem}[Finite selector commutant and its actual logical domain]
\label{p12:thm:v350-wedderburn-exact-scope}
For a specified unital selector algebra
$\mathcal B_{\mathcal S}\subseteq\operatorname{End}(\mathcal H_C)$,
there is a finite unitary decomposition
\[
 \mathcal H_C=\bigoplus_\alpha V_\alpha\otimes M_\alpha,
 \quad
 \mathcal B_{\mathcal S}
 =\bigoplus_\alpha\operatorname{End}(V_\alpha)\otimes I_{M_\alpha},
\]
with commutant
\[
 \boxed{\mathcal B_{\mathcal S}'
 =\bigoplus_\alpha I_{V_\alpha}\otimes\operatorname{End}(M_\alpha).}
\]
This exhausts operators compatible with \emph{the supplied
finite selectors on the supplied carrier}; it does not
exclude new infrared collective observables or new
superselection sectors defined only on a refinement limit.
\end{theorem}
\begin{proof}
The finite-dimensional $C^*$-algebra structure theorem
provides the direct-sum decomposition; computing the
commutant block by block yields the stated multiplicity
algebra. A separate limit may possess observables not
represented by one uniformly fixed finite selector algebra,
so no general continuum-exhaustiveness claim follows.
\end{proof}

\begin{theorem}[Two precise levels of truncated bosonic CCR recovery]
\label{p12:thm:v350-ccr-expectation-vs-core}
On $\mathcal F_N=\operatorname{span}\{|0\rangle,\ldots,|N\rangle\}$
the truncated oscillator satisfies
\[
 [a_N,a_N^\dagger]
   =I-(N+1)|N\rangle\langle N|.
\]
For a state sequence $\omega_N$, the condition
$\epsilon_N:=(N+1)\omega_N(|N\rangle\langle N|)\to0$
is equivalent to \emph{expectation-level} recovery
$\omega_N([a_N,a_N^\dagger])\to1$, but is not alone
sufficient for regular Weyl-state convergence.
On the standard nested Fock spaces embedded into
$\ell^2(\mathbb N_0)$, $a_N,a_N^\dagger$ agree with the
canonical $a,a^\dagger$ on every finite-support vector
for all sufficiently large $N$, giving the exact CCR on the
inductive finite-particle core.
\end{theorem}
\begin{proof}
Taking the state expectation proves the equivalence.
The states $|N-1\rangle\langle N-1|$ have $\epsilon_N=0$
but divergent mean occupation, demonstrating insufficiency
of this test for a regular state limit. For every fixed
finite-support vector only finitely many oscillator levels
occur. Once $N$ exceeds those levels and their adjacent
raising images, the truncation does not affect the action
of $a,a^\dagger$ or their commutator on that vector.
The canonical core relation therefore holds exactly there.
\end{proof}

\begin{proposition}[Sufficient regular physical-state criterion]
\label{p12:prop:v350-ccr-state-tightness}
Embed density matrices $\rho_N$ on $\mathcal F_N$ by zero
extension into the standard Fock space. Suppose
$\sup_N\operatorname{Tr}(\rho_N\mathsf N)<\infty$ and
all fixed number-basis matrix elements converge.
Then $\rho_N$ converges in trace norm to a density matrix
$\rho$; hence the canonical Weyl expectations
$\operatorname{Tr}(\rho_N W(z))$ converge to those of a
regular Weyl state on the same representation.
If also $\epsilon_N\to0$, the truncated commutator
expectations tend to one. This criterion is sufficient,
not necessary for all possible CCR limits.
\end{proposition}
\begin{proof}
For the projection $P_{\le K}$ onto numbers at most $K$,
Markov's inequality gives
$\operatorname{Tr}(\rho_N(I-P_{\le K}))\le C/(K+1)$
uniformly. On each fixed finite block, convergence of all
matrix entries implies trace-norm convergence. The uniform
tail bound and the trace-class truncation estimate extend
this to convergence on the full space, yielding a positive
trace-one limit. Bounded Weyl operators then have convergent
expectations, and their canonical representation is strongly
continuous. The last statement follows from the exact finite
commutator identity.
\end{proof}

\begin{remark}[Proof provenance and no-unjustified-physical promotion]
\label{p12:rem:v350-technical-certificate-status}
The finite 24-cell moment/adjacency, incidence, polarized
transport, DGLAP, arithmetic shell and occupation-boundary
tests are reproduced by the following script:
\begin{quote}
\path{Proofs/Paper00-PaperVIIRepairProofs-v1_1.py}.
\end{quote}
These tests do not execute a physical nuclear spectral count,
interacting baryon source, spinorial exact intertwiner,
local Gauss-invariant truncation, interacting locality limit,
or pure-glue uniform mass-gap certificate.
\end{remark}
% END FRAMEWORK v2.350


% BEGIN FRAMEWORK v2.357: finite Clifford spectral witness and source no-go ledger
\clearpage
\section{Clifford comparison spectrum and the degree-three hadronic source frontier}
\label{p12:sec:v357-clifford-shell-certificate}
\begin{theorem}[Single comparison Hamiltonian with cumulative gap capacities]
\label{p12:thm:v357-clifford-shell-spectrum}
Under the explicitly \emph{mathematical} rotation identification of
Paper~I-D's spinor Fischer sequence, fix a maximum polynomial degree
$R\ge3$ and the finite space
$\mathscr P_{\le R}=\bigoplus_{n=0}^R P_n\otimes V_{1/2}$.
Let $\mathbf L=-i\mathbf q\times\nabla_{\mathbf q}$, and define
\[
 h_{\omega,\kappa}=\omega N+\kappa CD
 =\omega N+\kappa(N-\boldsymbol\sigma\cdot\mathbf L),
 \quad \omega>0,\quad 0<\kappa<\frac{\omega}{2R-1}.
\]
For a degree-$n$, orbital-$\ell$ state with total $j$, its eigenvalue is
\[
 E_{n\ell j}=\begin{cases}
 n\omega+\kappa(n-\ell),&j=\ell+\frac12,\\
 n\omega+\kappa(n+\ell+1),&j=\ell-\frac12.
 \end{cases}
\]
($\ell=n,n-2,\ldots\ge0$, and the second line requires $\ell\ge1$.)
There exist increasing thresholds in genuine gaps whose spectral
projectors have ranges $\mathscr F_0,\ldots,\mathscr F_R$ from the
Paper~I-D comparison filtration. For $R=7$ their cumulative ranks are
\[
 \boxed{2,8,20,28,50,82,126,184.}
\]
For every $3\le r\le R$ the gap immediately above the occupied
highest-spin component at $r\omega$ has size $2\kappa$.
\end{theorem}
\begin{proof}
The Pauli multiplication identity gives
$(\boldsymbol\sigma\cdot\mathbf q)
(\boldsymbol\sigma\cdot\nabla)=N-\boldsymbol\sigma\cdot\mathbf L$.
The angular-momentum Casimir identity gives
$\boldsymbol\sigma\cdot\mathbf L=\ell$ for
$j=\ell+1/2$ and $-(\ell+1)$ for $j=\ell-1/2$.
Thus $CD\ge0$, with zero precisely on the highest-spin module
$M_n$, and its smallest positive eigenvalue in degree $n\ge2$
is $2$. In every degree $n<r$, the energy is bounded above by
$n\omega+(2n+1)\kappa$, which is $<r\omega$ under the stated
uniform bound. At degree $r$ only $M_r$ has energy $r\omega$;
the remaining eigenvalues start at $r\omega+2\kappa$, and all
higher degrees start at $(r+1)\omega$. Taking a threshold in
the interval $(r\omega,r\omega+2\kappa)$ selects exactly
$\mathscr F_r$ for $r\ge3$. For $r\le2$, thresholds after all
states of the corresponding degrees and before the next degree
exist by the same bound. Direct dimension counting proves the
ranks. All statements concern the declared comparison operator.
\end{proof}
\begin{proposition}[Sufficient stability under a full finite perturbation]
\label{p12:prop:v357-shell-weyl-stability}
Let $U$ be a unitary identification of the \emph{entire} finite
comparison space with an actual finite one-body source Hilbert
space and suppose the latter has a self-adjoint Hamiltonian
$h_{\rm nuc}=Uh_{\omega,\kappa}U^\dagger+V$ there.
A comparison spectral gap of width $g>0$ at a cumulative rank
persists at the same rank whenever $\lVert V\rVert<g/2$.
This suffices to preserve its spectral counting, but does not
prove the required unitary identification or bound.
\end{proposition}
\begin{proof}
Each ordered eigenvalue moves by at most $\lVert V\rVert$
under a self-adjoint finite perturbation (Weyl's inequality).
The separation of the two adjacent eigenvalues is therefore
at least $g-2\lVert V\rVert>0$.
\end{proof}
\begin{proposition}[Two degree-three source-rank obstructions]
\label{p12:prop:v357-degree-three-source-no-go}
Let $X$ have complex dimension four and $\mathcal B_3=\Sym^3X$.
(i) If $S_3=\Sym^3T$ for one linear map $T:X\to Y$ with
rank $k\in\{0,1,2,3,4\}$, then
\[
 \operatorname{rank}S_3=\binom{k+2}{3}
 \in\{0,1,4,10,20\},
\]
so it cannot have the desired nuclear source rank $14$.
(ii) Under the full reference $\mathrm{Spin}(4)$ action,
\[
 \Sym^3(V_{1/2}\otimes V_{1/2})
 \cong(V_{3/2}\otimes V_{3/2})\oplus
 (V_{1/2}\otimes V_{1/2}).
\]
Every full-$\mathrm{Spin}(4)$-equivariant map from this source
has invariant kernel of dimension $0,4,16$ or $20$, never $6$.
These restrictions apply only to the stated factorized/equivariant
source classes and do not exclude a source using the already
selected oriented boundary or non-factorized baryonic dynamics.
\end{proposition}
\begin{proof}
The image of $\Sym^3T$ is $\Sym^3(\operatorname{im}T)$,
whose dimension is the first binomial.  The second identity is the
$\operatorname{Sym}^3$ Cauchy decomposition for a tensor product
of two two-dimensional modules.  Its inequivalent multiplicity-one
irreducible summands have dimensions $16$ and $4$; equivariant
kernels are sums of those summands.
\end{proof}
\begin{proposition}[Positive gauge averaging cannot create source nullity]
\label{p12:prop:v357-positive-gram-kernel}
For a measurable source family $S(U):\mathcal B_3\to\mathcal H$
and a finite positive gauge measure $\mu$ let
$G=\int S(U)^\dagger S(U)\,d\mu(U)$. Then
\[
 \boxed{\ker G=\{v:S(U)v=0\quad\mu\text{-almost everywhere}\}.}
\]
If $S(U)$ is continuous near a flat link configuration in the
support of $\mu$ and $S(I)$ is injective, then $G$ is strictly
positive. These conclusions are conditional on this particular
source-family and positive Gram construction.
\end{proposition}
\begin{proof}
For every $v$, $\langle v,Gv\rangle=
\int\lVert S(U)v\rVert^2d\mu(U)\ge0$. Vanishing is equivalent
to almost-everywhere vanishing of the nonnegative integrand.
Injectivity at $I$ gives a strictly positive least singular
value; continuity preserves a positive lower bound in a
neighborhood with positive $\mu$-measure.
\end{proof}
\begin{proposition}[Conditional heat-character sign change]
\label{p12:prop:v357-heat-spinflux-sign}
For the standalone central $SU(3)$ heat-kernel probability
measure with normalized representation moments
$\langle\chi_R\rangle_t=d_Re^{-tC_R/2}$,
consider one face of the spin--flux comparison coefficient.
With $C_{\mathbf3}=4/3$ and $C_{\mathbf8}=3$, its mean is
\[
 d(t)=\frac{1+8e^{-3t/2}}9-e^{-2t/3}.
\]
It is negative for sufficiently small $t>0$, positive for
sufficiently large $t$, and has exactly one positive zero,
$t_*=3.200226335068\ldots$ in this heat-time convention.
The sign is not selected by the Hurwitz incidence alone.
\end{proposition}
\begin{proof}
The fundamental--adjoint character identity gives $d(t)$.
One has $d(0)=0$, $d'(0)=-2/3$ and $d(t)\to1/9>0$.
With $x=e^{-t/6}$ the zero equation is
$1+8x^9-9x^4=0$ for $0<x<1$.
Its derivative $72x^8-36x^3=36x^3(2x^5-1)$ has a
single interior critical point. The function has opposite
signs near $x=0$ and $x=1$ and therefore precisely one
interior root; the numerical value follows by bisection.
\end{proof}
\begin{remark}[Unclosed physical identification]
\label{p12:rem:v357-hadronic-shell-frontier}
The comparison Hamiltonian is an existence witness, not a derived
QCD operator or a physical shell-gap prediction. The actual
$S_3\mathcal L_3^\sigma=0$ and complementary source positivity,
its orientation/rotation/parity intertwiner, and the same-action
physical nuclear spectral gaps remain separate unexecuted tests.
These results neither create GP4 nor change $m_p=1$,
$\mathcal H_p\cong\mathcal H_1$, or the established constitutional
proof chain. The diagnostic $E_{36}/E_{40}$ mass comparisons remain
in the technical material at their prior exploratory status.
\end{remark}
% END FRAMEWORK v2.357


% BEGIN FRAMEWORK v2.365: REFERENCE-NORMALIZATION WORKBENCH AND LIMITATIONS
\section{Wilson-fan reference calculations and physical-source limitations}
\label{p12:sec:v365-wilson-reference-tests}
For the three path transporters $A_1,A_2,A_3\in SU(3)$,
let $\epsilon/\sqrt6$ denote the normalized determinant tensor.
The identity
\[
  W_e=\langle\epsilon|A_1\otimes A_2\otimes A_3|\epsilon\rangle
\]
is independent of the order of the three fan faces and agrees
with Proposition~\ref{p10:prop:v2220-edge-fan-baryon}.
For a common normalized two-root Gram
$G=\left(\begin{smallmatrix}1&w\\\bar w&1\end{smallmatrix}\right)$,
$\det G=1-|w|^2$ and whitening produces $I$ on $\operatorname{supp}G$.
These identities are exact for all unitary path transporters,
with no Haar or physical vacuum assumption.

\begin{proposition}[Separated independent-Haar comparator]
\label{p12:prop:v365-haar-three-path-second-moment}
Let $A_1,A_2,A_3$ be three independent normalized-Haar $SU(3)$
matrices, each with its fundamental indices left uncorrelated by
other link insertions.  Then
\[
\boxed{\mathbb E W_e=0,\qquad \mathbb E|W_e|^2=\tfrac1{27}.}
\]
This is a \emph{single fan} reference moment; extending it to
192 mutually orthogonal \emph{oriented} fans requires the distinct
full-link overlap check and is not asserted by this proposition.
\end{proposition}
\begin{proof}
The first integral vanishes by the nontrivial fundamental
representation.  For the second, integrate each independent
$A_i$ using
$\int_{SU(3)}U_{ab}\overline U_{cd}\,dU
 =\delta_{ac}\delta_{bd}/3$.  There are three such $1/3$
factors; the two normalized epsilon contractions each contribute
one after their six terms are summed, giving $1/27$.
\end{proof}

\begin{remark}[Magnetic and jet comparisons are not action spectra]
\label{p12:rem:v365-haar-jet-quarantine}
The research proposal additionally reports the independent-Haar
192-fan Gram, the formal electric-Casimir value $4$, the
$96$-face equal-weight second-moment value $48$, a reversed-fan
third-/fourth-order selection with coefficient $3/16$, and the
factorized local-heat expression
$\prod_{i=1}^3 Z'(t_i)/(3Z(t_i))$,
$Z(t)=\int_{SU(3)}e^{t\operatorname{Re}\operatorname{tr}U}\,dU$.
Those belong to distinct \emph{declared comparison measures and
operator orderings}.  Without independent proof-route replay and
the complete same-action source normalization they are retained
as research certificates rather than promoted as the interacting
QCD action's pole masses or decay amplitudes.  In particular the
single-fan vacuum third-order insertion and the reversed-fan
first-nonzero fourth-order matrix element are different objects.
The precise confinement, baryon-number-one source and leakage
conditions remain in Papers VII-A and VII-B.
\end{remark}
% END FRAMEWORK v2.365



% BEGIN FRAMEWORK v3.025: FINITE QCD CHARACTER / POSITIVE-TYPE TECHNICAL CERTIFICATES
\section{Nonlinear QCD action sensitivity and $SU(3)$ character certificates}
\label{p13:sec:v3025-nonlinear-qcd-action}
The following finite compact-group calculations support
Propositions~\ref{p10:prop:v3025-threeface-action-sensitivity}
and~\ref{p10:prop:v3025-sextet-temporal-exclusion}.  They are
\emph{reference-action} results; a physical strong Hamiltonian,
Gauss-projected source, and time identification are not inferred.

\begin{proposition}[Independent-face Wilson character moments]
\label{p13:prop:v3025-character-moment-certificate}
For the $SU(3)$ heat kernel $K_t$ normalized with
$\int K_t(U)\,dU=1$, put
$x=(\chi_{\mathbf3}+\chi_{\overline{\mathbf3}})/6$ and $q=1-x$.
Then $\langle x\rangle_t=e^{-2t/3}$, and on the isolated three-face
reference $\langle W_e\rangle_{t,0}=e^{-2t}$.  With
$m_k=\langle x^k\rangle_t$, the deformation derivative is
\[
 \left.\partial_\lambda\langle W_e\rangle_{t,\lambda}
 \right|_{\lambda=0}
 =3e^{-4t/3}\{(1-3m_1+3m_2-m_3)
 -(1-m_1)(1-2m_1+m_2)\}.
\]
At the separately fixed isotropic exact-face heat time
$t_{\rm iso}\simeq 5.7532958466$, these quantities are
$1.0063538725\times10^{-5}$ and
$1.4817641219\times10^{-4}$, respectively.
\end{proposition}
\begin{proof}
The heat-character identity
$\langle\chi_{(p,q)}\rangle_t
=d_{p,q}e^{-tC_{p,q}/2}$ gives the first moment since
$C_{\mathbf3}=4/3$.  The second and third moments follow by
multiplication in the $SU(3)$ representation ring; the cubic
derivative follows by differentiating the normalized positive
integral and the relation $\partial_\lambda a=
\operatorname{Cov}_\lambda(q,q^2)$.  All numerical figures refer
only to the named isotropic face-time branch.
\end{proof}

\begin{proposition}[Uniform sextet negative-character certificate]
\label{p13:prop:v3025-sextet-character-certificate}
For $5.75\le t\le5.76$, put
\[
 N_{\mathbf6}(t)=\frac16\int_{SU(3)} K_t(U)
 \exp[-q(U)^2/10] \,\overline{\chi_{\mathbf6}(U)}\,dU.
\]
Then $N_{\mathbf6}(t)<-4.5\times10^{-5}$ throughout the interval.
This sign is a statement about one central Fourier coefficient, not
about positivity of the spatial magnetic measure.
\end{proposition}
\begin{proof}[Representation-ring and uniform-tail certificate]
Write an irreducible representation by Dynkin labels $(p,q)$.
The exact data are
\[
 d_{p,q}=\frac{(p+1)(q+1)(p+q+2)}2,\qquad
 C_{p,q}=\frac{p^2+q^2+pq+3p+3q}{3}.
\]
Multiplication by $\mathbf3$ uses the Pieri rule
\[
 (p,q)\otimes(1,0)
 =(p+1,q)\oplus(p-1,q+1)_{p\ge1}\oplus(p,q-1)_{q\ge1},
\]
and multiplication by $\overline{\mathbf3}$ uses its conjugate.
Consequently multiplication by $x$ is the mean of the two Pieri
operators divided by $6$, and
\[
 \boxed{q^2=\frac{19}{18}\mathbf1
 -\frac{11}{36}(\chi_{\mathbf3}+\chi_{\overline{\mathbf3}})
 +\frac1{36}(\chi_{\mathbf6}+\chi_{\overline{\mathbf6}})
 +\frac1{18}\chi_{\mathbf8}.}
\]
Starting with $\overline{\chi_{\mathbf6}}/6$, repeated
multiplication by this exact character polynomial evaluates the
Taylor terms $0\le n\le8$ of
$\exp[-q^2/10]$ through the heat-character identity
$\langle\chi_{(p,q)}\rangle_t=d_{p,q}e^{-tC_{p,q}/2}$.
On $SU(3)$, $0\le q\le3/2$, so the discarded Taylor remainder is
bounded in absolute value by
$(0.225)^9/9!$ after taking the normalized sextet character bound.
For $p+q\ge18$, the heat-character tail on the stated interval is
bounded by $5\times10^{-117}$ with the polynomial dimension/Casimir
envelope, and the change in the truncated sum from $t=5.75$ to
$t\le5.76$ is bounded by $0.01$ times the absolute differentiated
finite sum.  Directed interval evaluation gives
\[
 \boxed{N_{\mathbf6}(t)<-4.56\times10^{-5}
 \quad(5.75\le t\le5.76),}
\]
which implies the weaker bound in the statement.  The finite
representation-ring recurrence, term list, interval enclosures and
explicit tail test are retained as machine-checkable source-update
evidence; the result does not imply unique action selection.
\end{proof}
% END FRAMEWORK v3.025


% BEGIN FRAMEWORK v3.027: EXACT CARTAN GRAM AND CLIFFORD COMPOSITION CERTIFICATES
\section{Technical certificates for three-path Cartan and Clifford comparisons}
\label{p13:sec:v3027-three-path-certificates}
\begin{proposition}[Exact middle-profile Cartan Gram polynomial]
\label{p13:prop:v3027-cartan-gaussian-spectrum}
For $N$ the eight-neighbour Hurwitz adjacency, put
$D_\rho=\operatorname{diag}_{h\in2T}\rho(h)$ and
$\mathscr M=N\otimes I_2+D_\rho(N\otimes I_2)D_\rho^*$.
On the chosen exact Gaussian-integer two-color chart, the
$48\times48$ Hermitian matrix $2\mathscr M$ satisfies
\[
\boxed{(2\mathscr M+12I)(2\mathscr M+4I)
(2\mathscr M-8I)(2\mathscr M-24I)=0.}
\]
Moreover $\operatorname{tr}[(2\mathscr M)^k]$ for
$k=0,1,2,3$ equals $(48,0,4608,46080)$.
Its spectrum is therefore
$(-12)^{(8)},(-4)^{(24)},8^{(12)},24^{(4)}$.
At the middle Cartan profile the Gram
$12I+\frac{15+\sqrt5}{10}\mathscr M$
has spectrum
$(3-3/\sqrt5)^{(8)},(9-1/\sqrt5)^{(24)},
(18+2/\sqrt5)^{(12)},(30+6/\sqrt5)^{(4)}$.
\end{proposition}
\begin{proof}
Use the $24$ signed Hurwitz quaternions in the fixed $2T$ chart,
with their exact eight-neighbour dot-product relation.
Each entry of $2\mathscr M$ lies in $\mathbb Z[i]$, since the
fundamental quaternion representation has Gaussian-integer
entries after the factor $2$. Exact multiplication of this finite
matrix yields the displayed polynomial annihilator and four
integer traces. Solving the resulting Vandermonde system for the
four eigenvalue multiplicities gives $8,24,12,4$ (summing to $48$).
Finally substitution into the Gram affine expression proves the
four positive squared singular values, without floating rank
thresholds. This is a finite matrix certificate, not an independent
source-generated $600$-cell physical link.
\end{proof}

\begin{proposition}[Clifford-grade obstruction for the flat two-edge word]
\label{p13:prop:v3027-two-edge-clifford-obstruction}
In the actual Euclidean Clifford representation of Paper~VII-A let
$\Gamma_a=\sigma_1\otimes\sigma_a$ ($a=1,2,3$),
$\Gamma_4=\sigma_2\otimes I_2$,
$\Gamma_5=\Gamma_1\Gamma_2\Gamma_3\Gamma_4$, and
$S(n)=\Gamma_5(\Gamma\!\cdot n-I_4)$.
For arbitrary unit tangents $n,n_1,n_2$,
\[
\boxed{\operatorname{tr}_{\rm spin}
 [S(n)^*S(n_1)S(n_2)]=0.}
\]
For the three equal-weight, exactly enumerated intermediates
between a native Hurwitz pair, define
$K(n)=n_3\Gamma_1\Gamma_2+n_1\Gamma_2\Gamma_3+n_2\Gamma_3\Gamma_1$.
Their two natural ordered spin-word means are
\[
\begin{aligned}
T_0&=\tfrac13\sum_r S(n_{qr})S(n_{rq'})
=(1-\mu)I_4-cK(n),\\
T_5&=\tfrac13\sum_r S(n_{qr})\Gamma_5S(n_{rq'})
=\Gamma_5\big[(1+\mu)I_4-b(\Gamma\cdot n)+cK(n)\big],
\end{aligned}
\]
where
\[
 \mu=\frac1{\sqrt5},\quad
 b=\frac{\varphi\sqrt3}{3\sin(\pi/5)},\quad
 c=\frac{5\sqrt3-\sqrt{15}}{15}>0.
\]
If $a=1+\mu$, the nearest scalar multiple of the native word to
$T_5$ in Hilbert--Schmidt norm is
$k_*S(n)$ with $k_*=-(a+b)/2$, and
\[
\boxed{\|T_5-k_*S(n)\|_{\rm HS}^2
=2(a-b)^2+4c^2>0.}
\]
\end{proposition}
\begin{proof}
The target word lies in the Clifford grades $3$ and $4$, while the
product of two $S$ words with a spin-scalar propagator lies in
$0,1,2$; orthogonality of different grades gives the trace zero.
For the three intermediate nearest paths, exact quaternionic
inner products give $n_{qr}\cdot n_{rq'}=\mu$ and the unit tangent
sum fixed in the three-path theorem of Paper~VII-A.
Expand the products using
$\Gamma_a\Gamma_b+\Gamma_b\Gamma_a=2\delta_{ab}I_4$;
the scalar, vector and bivector components give the indicated
$T_0$ and $T_5$. The Clifford trace form makes these components
orthogonal, with $\|\Gamma_5\|_{\rm HS}^2=4$.
Minimizing the quadratic polynomial
$\|T_5-kS(n)\|_{\rm HS}^2$ over scalar $k$ gives
$k_*=-(a+b)/2$ and the strictly positive residue displayed.
This only excludes flat spin-scalar and single constant
$\Gamma_5$ intermediate insertions under equal-path assembly:
it does not exclude a complete action-owned spin/color resolvent.
\end{proof}

\begin{remark}[Companion exact verification and negative scope]
\label{p13:rem:v3027-three-path-audit-scope}
The finite $\mathbb Z[\varphi]$ checks of all five $600$-cell
coset returns, the $48\times48$ Gaussian-integer Cartan polynomial,
and the native $864\times864$ temporal-clutch matrix construction
are preserved in the associated geometric transition compiler
verification evidence. The core identities are algebraic; bounded
floating checks on the full regulator are only independent
consistency controls, not proofs of physical occurrence. Neither a
naive two-edge spin product nor the extended $600$-cell links have
been identified with the same-action Fock decay operator.
\end{remark}
% END FRAMEWORK v3.027 VII-D MODULE

\end{document}
